%%C020001 p

本章包含以下案例：


%%C020002 li

\begin{itemize}
\item 数值类型与字符串类型之间的转换


%%C020003 li

\item 数值类型的取值范围及其他属性


%%C020004 li

\item 生成伪随机数


%%C020005 li

\item 初始化伪随机数生成器内部状态的全部比特


%%C020006 li

\item 用原始字符串字面量避免转义字符


%%C020007 li

\item 创建cooked用户定义字面量


%%C020008 li

\item 创建raw用户定义字面量


%%C020009 li

\item 构建字符串工具库


%%C020010 li

\item 用正则表达式验证字符串格式


%%C020011 li

\item 用正则表达式解析字符串内容


%%C020012 li

\item 用正则表达式替换字符串内容


%%C020013 li

\item 用string\_view代替常量字符串引用

\end{itemize}


%%C020014 h2

\section{引言}


%%C020015 p

数字与字符串是任何编程语言的基础类型；所有其他类型都基于它们、由它们组合而成。开发者时刻都要面对诸如数字与字符串之间的转换、解析字符串或生成随机数之类的任务。本章着重用现代C++的语言与库特性，为这些常见任务提供实用的案例。


%%C020016 h2

\section{数值类型与字符串类型之间的转换}


%%C020017 p

在数值类型与字符串类型之间转换是一种无处不在的操作。在C++11之前，数字转字符串以及反向转换几乎没有支持，开发者大多只能求助于类型不安全的函数，并且通常不得不自己编写工具函数，以免一遍又一遍地重复同样的代码。从C++11起，标准库提供了用于数字与字符串之间相互转换的工具函数。本案例将介绍如何使用现代C++标准函数在数字与字符串之间双向转换。


%%C020018 h3

\subsection{准备工作}


%%C020019 p

本案例提到的所有工具函数都在\texttt{<\allowbreak{}string>}头文件中。


%%C020020 h3

\subsection{具体做法}


%%C020021 p

需要在数字与字符串之间转换时，使用下面这些标准转换函数：


%%C020022 li

\begin{itemize}
\item 要从整数或浮点类型转换为字符串类型，使用\texttt{std::\allowbreak{}to\_\allowbreak{}string()}或\texttt{std::\allowbreak{}to\_\allowbreak{}wstring()}，如下面的代码片段所示：


%%C020023 code

\begin{codeblockbg}
\begin{Verbatim}[commandchars=\\\{\},breaklines=true,breakanywhere=true,fontsize=\codefont,xleftmargin=17bp]
\PYG{k}{auto}\PYG{+w}{ }\PYG{n}{si}\PYG{+w}{ }\PYG{o}{=}\PYG{+w}{ }\PYG{n}{std}\PYG{o}{:}\PYG{o}{:}\PYG{n}{to\PYGZus{}string}\PYG{p}{(}\PYG{l+m+mi}{42}\PYG{p}{)}\PYG{p}{;}\PYG{+w}{      }\PYG{c+c1}{// si=\PYGZdq{}42\PYGZdq{}}
\PYG{k}{auto}\PYG{+w}{ }\PYG{n}{sl}\PYG{+w}{ }\PYG{o}{=}\PYG{+w}{ }\PYG{n}{std}\PYG{o}{:}\PYG{o}{:}\PYG{n}{to\PYGZus{}string}\PYG{p}{(}\PYG{l+m+mf}{42l}\PYG{p}{)}\PYG{p}{;}\PYG{+w}{     }\PYG{c+c1}{// sl=\PYGZdq{}42\PYGZdq{}}
\PYG{k}{auto}\PYG{+w}{ }\PYG{n}{su}\PYG{+w}{ }\PYG{o}{=}\PYG{+w}{ }\PYG{n}{std}\PYG{o}{:}\PYG{o}{:}\PYG{n}{to\PYGZus{}string}\PYG{p}{(}\PYG{l+m+mi}{42u}\PYG{p}{)}\PYG{p}{;}\PYG{+w}{     }\PYG{c+c1}{// su=\PYGZdq{}42\PYGZdq{}}
\PYG{k}{auto}\PYG{+w}{ }\PYG{n}{sd}\PYG{+w}{ }\PYG{o}{=}\PYG{+w}{ }\PYG{n}{std}\PYG{o}{:}\PYG{o}{:}\PYG{n}{to\PYGZus{}wstring}\PYG{p}{(}\PYG{l+m+mf}{42.0}\PYG{p}{)}\PYG{p}{;}\PYG{+w}{   }\PYG{c+c1}{// sd=L\PYGZdq{}42.000000\PYGZdq{}}
\PYG{k}{auto}\PYG{+w}{ }\PYG{n}{sld}\PYG{+w}{ }\PYG{o}{=}\PYG{+w}{ }\PYG{n}{std}\PYG{o}{:}\PYG{o}{:}\PYG{n}{to\PYGZus{}wstring}\PYG{p}{(}\PYG{l+m+mf}{42.0l}\PYG{p}{)}\PYG{p}{;}\PYG{+w}{ }\PYG{c+c1}{// sld=L\PYGZdq{}42.000000\PYGZdq{}}
\end{Verbatim}
\end{codeblockbg}


%%C020024 li

\item 要从字符串类型转换为整数类型，使用\texttt{std::\allowbreak{}stoi()}、\texttt{std::\allowbreak{}stol()}、\texttt{std::\allowbreak{}stoll()}、\texttt{std::\allowbreak{}stoul()}或\texttt{std::\allowbreak{}stoull()}；参见下面的代码片段：


%%C020025 code

\begin{codeblockbg}
\begin{Verbatim}[commandchars=\\\{\},breaklines=true,breakanywhere=true,fontsize=\codefont,xleftmargin=17bp]
\PYG{k}{auto}\PYG{+w}{ }\PYG{n}{i1}\PYG{+w}{ }\PYG{o}{=}\PYG{+w}{ }\PYG{n}{std}\PYG{o}{:}\PYG{o}{:}\PYG{n}{stoi}\PYG{p}{(}\PYG{l+s}{\PYGZdq{}}\PYG{l+s}{42}\PYG{l+s}{\PYGZdq{}}\PYG{p}{)}\PYG{p}{;}\PYG{+w}{                 }\PYG{c+c1}{// i1 = 42}
\PYG{k}{auto}\PYG{+w}{ }\PYG{n}{i2}\PYG{+w}{ }\PYG{o}{=}\PYG{+w}{ }\PYG{n}{std}\PYG{o}{:}\PYG{o}{:}\PYG{n}{stoi}\PYG{p}{(}\PYG{l+s}{\PYGZdq{}}\PYG{l+s}{101010}\PYG{l+s}{\PYGZdq{}}\PYG{p}{,}\PYG{+w}{ }\PYG{n+nb}{nullptr}\PYG{p}{,}\PYG{+w}{ }\PYG{l+m+mi}{2}\PYG{p}{)}\PYG{p}{;}\PYG{+w}{ }\PYG{c+c1}{// i2 = 42}
\PYG{k}{auto}\PYG{+w}{ }\PYG{n}{i3}\PYG{+w}{ }\PYG{o}{=}\PYG{+w}{ }\PYG{n}{std}\PYG{o}{:}\PYG{o}{:}\PYG{n}{stoi}\PYG{p}{(}\PYG{l+s}{\PYGZdq{}}\PYG{l+s}{052}\PYG{l+s}{\PYGZdq{}}\PYG{p}{,}\PYG{+w}{ }\PYG{n+nb}{nullptr}\PYG{p}{,}\PYG{+w}{ }\PYG{l+m+mi}{8}\PYG{p}{)}\PYG{p}{;}\PYG{+w}{    }\PYG{c+c1}{// i3 = 42}
\PYG{k}{auto}\PYG{+w}{ }\PYG{n}{i4}\PYG{+w}{ }\PYG{o}{=}\PYG{+w}{ }\PYG{n}{std}\PYG{o}{:}\PYG{o}{:}\PYG{n}{stoi}\PYG{p}{(}\PYG{l+s}{\PYGZdq{}}\PYG{l+s}{0x2A}\PYG{l+s}{\PYGZdq{}}\PYG{p}{,}\PYG{+w}{ }\PYG{n+nb}{nullptr}\PYG{p}{,}\PYG{+w}{ }\PYG{l+m+mi}{16}\PYG{p}{)}\PYG{p}{;}\PYG{+w}{  }\PYG{c+c1}{// i4 = 42}
\end{Verbatim}
\end{codeblockbg}


%%C020026 li

\item 要从字符串类型转换为浮点类型，使用\texttt{std::\allowbreak{}stof()}、\texttt{std::\allowbreak{}stod()}或\texttt{std::\allowbreak{}stold()}，如下面的代码片段所示：


%%C020027 code

\begin{codeblockbg}
\begin{Verbatim}[commandchars=\\\{\},breaklines=true,breakanywhere=true,fontsize=\codefont,xleftmargin=17bp]
\PYG{c+c1}{// d1 = 123.45000000000000}
\PYG{k}{auto}\PYG{+w}{ }\PYG{n}{d1}\PYG{+w}{ }\PYG{o}{=}\PYG{+w}{ }\PYG{n}{std}\PYG{o}{:}\PYG{o}{:}\PYG{n}{stod}\PYG{p}{(}\PYG{l+s}{\PYGZdq{}}\PYG{l+s}{123.45}\PYG{l+s}{\PYGZdq{}}\PYG{p}{)}\PYG{p}{;}
\PYG{c+c1}{// d2 = 123.45000000000000}
\PYG{k}{auto}\PYG{+w}{ }\PYG{n}{d2}\PYG{+w}{ }\PYG{o}{=}\PYG{+w}{ }\PYG{n}{std}\PYG{o}{:}\PYG{o}{:}\PYG{n}{stod}\PYG{p}{(}\PYG{l+s}{\PYGZdq{}}\PYG{l+s}{1.2345e+2}\PYG{l+s}{\PYGZdq{}}\PYG{p}{)}\PYG{p}{;}
\PYG{c+c1}{// d3 = 123.44999980926514}
\PYG{k}{auto}\PYG{+w}{ }\PYG{n}{d3}\PYG{+w}{ }\PYG{o}{=}\PYG{+w}{ }\PYG{n}{std}\PYG{o}{:}\PYG{o}{:}\PYG{n}{stod}\PYG{p}{(}\PYG{l+s}{\PYGZdq{}}\PYG{l+s}{0xF.6E6666p3}\PYG{l+s}{\PYGZdq{}}\PYG{p}{)}\PYG{p}{;}
\end{Verbatim}
\end{codeblockbg}

\end{itemize}


%%C020028 h3

\subsection{工作原理}


%%C020029 p

要在整型或浮点类型与字符串类型之间转换，可以使用\texttt{std::\allowbreak{}to\_\allowbreak{}string()}或\texttt{std::\allowbreak{}to\_\allowbreak{}wstring()}函数。这两个函数位于\texttt{<\allowbreak{}string>}头文件中，对有符号和无符号整数类型以及实数类型都有重载。它们产生的结果与用各类型合适的格式说明符调用\texttt{std::\allowbreak{}sprintf()}和\texttt{std::\allowbreak{}swprintf()}时产生的结果相同。下面的代码片段列出了这两个函数的全部重载。


%%C020030 code

\begin{codeblockbg}
\begin{Verbatim}[commandchars=\\\{\},breaklines=true,breakanywhere=true,fontsize=\codefont,xleftmargin=17bp]
\PYG{n}{std}\PYG{o}{:}\PYG{o}{:}\PYG{n}{string}\PYG{+w}{ }\PYG{n+nf}{to\PYGZus{}string}\PYG{p}{(}\PYG{k+kt}{int}\PYG{+w}{ }\PYG{n}{value}\PYG{p}{)}\PYG{p}{;}
\PYG{n}{std}\PYG{o}{:}\PYG{o}{:}\PYG{n}{string}\PYG{+w}{ }\PYG{n+nf}{to\PYGZus{}string}\PYG{p}{(}\PYG{k+kt}{long}\PYG{+w}{ }\PYG{n}{value}\PYG{p}{)}\PYG{p}{;}
\PYG{n}{std}\PYG{o}{:}\PYG{o}{:}\PYG{n}{string}\PYG{+w}{ }\PYG{n+nf}{to\PYGZus{}string}\PYG{p}{(}\PYG{k+kt}{long}\PYG{+w}{ }\PYG{k+kt}{long}\PYG{+w}{ }\PYG{n}{value}\PYG{p}{)}\PYG{p}{;}
\PYG{n}{std}\PYG{o}{:}\PYG{o}{:}\PYG{n}{string}\PYG{+w}{ }\PYG{n+nf}{to\PYGZus{}string}\PYG{p}{(}\PYG{k+kt}{unsigned}\PYG{+w}{ }\PYG{n}{value}\PYG{p}{)}\PYG{p}{;}
\PYG{n}{std}\PYG{o}{:}\PYG{o}{:}\PYG{n}{string}\PYG{+w}{ }\PYG{n+nf}{to\PYGZus{}string}\PYG{p}{(}\PYG{k+kt}{unsigned}\PYG{+w}{ }\PYG{k+kt}{long}\PYG{+w}{ }\PYG{n}{value}\PYG{p}{)}\PYG{p}{;}
\PYG{n}{std}\PYG{o}{:}\PYG{o}{:}\PYG{n}{string}\PYG{+w}{ }\PYG{n+nf}{to\PYGZus{}string}\PYG{p}{(}\PYG{k+kt}{unsigned}\PYG{+w}{ }\PYG{k+kt}{long}\PYG{+w}{ }\PYG{k+kt}{long}\PYG{+w}{ }\PYG{n}{value}\PYG{p}{)}\PYG{p}{;}
\PYG{n}{std}\PYG{o}{:}\PYG{o}{:}\PYG{n}{string}\PYG{+w}{ }\PYG{n+nf}{to\PYGZus{}string}\PYG{p}{(}\PYG{k+kt}{float}\PYG{+w}{ }\PYG{n}{value}\PYG{p}{)}\PYG{p}{;}
\PYG{n}{std}\PYG{o}{:}\PYG{o}{:}\PYG{n}{string}\PYG{+w}{ }\PYG{n+nf}{to\PYGZus{}string}\PYG{p}{(}\PYG{k+kt}{double}\PYG{+w}{ }\PYG{n}{value}\PYG{p}{)}\PYG{p}{;}
\PYG{n}{std}\PYG{o}{:}\PYG{o}{:}\PYG{n}{string}\PYG{+w}{ }\PYG{n+nf}{to\PYGZus{}string}\PYG{p}{(}\PYG{k+kt}{long}\PYG{+w}{ }\PYG{k+kt}{double}\PYG{+w}{ }\PYG{n}{value}\PYG{p}{)}\PYG{p}{;}
\PYG{n}{std}\PYG{o}{:}\PYG{o}{:}\PYG{n}{wstring}\PYG{+w}{ }\PYG{n+nf}{to\PYGZus{}wstring}\PYG{p}{(}\PYG{k+kt}{int}\PYG{+w}{ }\PYG{n}{value}\PYG{p}{)}\PYG{p}{;}
\PYG{n}{std}\PYG{o}{:}\PYG{o}{:}\PYG{n}{wstring}\PYG{+w}{ }\PYG{n+nf}{to\PYGZus{}wstring}\PYG{p}{(}\PYG{k+kt}{long}\PYG{+w}{ }\PYG{n}{value}\PYG{p}{)}\PYG{p}{;}
\PYG{n}{std}\PYG{o}{:}\PYG{o}{:}\PYG{n}{wstring}\PYG{+w}{ }\PYG{n+nf}{to\PYGZus{}wstring}\PYG{p}{(}\PYG{k+kt}{long}\PYG{+w}{ }\PYG{k+kt}{long}\PYG{+w}{ }\PYG{n}{value}\PYG{p}{)}\PYG{p}{;}
\PYG{n}{std}\PYG{o}{:}\PYG{o}{:}\PYG{n}{wstring}\PYG{+w}{ }\PYG{n+nf}{to\PYGZus{}wstring}\PYG{p}{(}\PYG{k+kt}{unsigned}\PYG{+w}{ }\PYG{n}{value}\PYG{p}{)}\PYG{p}{;}
\PYG{n}{std}\PYG{o}{:}\PYG{o}{:}\PYG{n}{wstring}\PYG{+w}{ }\PYG{n+nf}{to\PYGZus{}wstring}\PYG{p}{(}\PYG{k+kt}{unsigned}\PYG{+w}{ }\PYG{k+kt}{long}\PYG{+w}{ }\PYG{n}{value}\PYG{p}{)}\PYG{p}{;}
\PYG{n}{std}\PYG{o}{:}\PYG{o}{:}\PYG{n}{wstring}\PYG{+w}{ }\PYG{n+nf}{to\PYGZus{}wstring}\PYG{p}{(}\PYG{k+kt}{unsigned}\PYG{+w}{ }\PYG{k+kt}{long}\PYG{+w}{ }\PYG{k+kt}{long}\PYG{+w}{ }\PYG{n}{value}\PYG{p}{)}\PYG{p}{;}
\PYG{n}{std}\PYG{o}{:}\PYG{o}{:}\PYG{n}{wstring}\PYG{+w}{ }\PYG{n+nf}{to\PYGZus{}wstring}\PYG{p}{(}\PYG{k+kt}{float}\PYG{+w}{ }\PYG{n}{value}\PYG{p}{)}\PYG{p}{;}
\PYG{n}{std}\PYG{o}{:}\PYG{o}{:}\PYG{n}{wstring}\PYG{+w}{ }\PYG{n+nf}{to\PYGZus{}wstring}\PYG{p}{(}\PYG{k+kt}{double}\PYG{+w}{ }\PYG{n}{value}\PYG{p}{)}\PYG{p}{;}
\PYG{n}{std}\PYG{o}{:}\PYG{o}{:}\PYG{n}{wstring}\PYG{+w}{ }\PYG{n+nf}{to\PYGZus{}wstring}\PYG{p}{(}\PYG{k+kt}{long}\PYG{+w}{ }\PYG{k+kt}{double}\PYG{+w}{ }\PYG{n}{value}\PYG{p}{)}\PYG{p}{;}
\end{Verbatim}
\end{codeblockbg}


%%C020031 p

至于反方向的转换，则有一整套名字形如\textbf{ston}（\textbf{string to number}，即字符串转数字）的函数，其中\textbf{n}代表\textbf{i}（\texttt{integer}）、\textbf{l}（\texttt{long}）、\textbf{ll}（\texttt{long long}）、\textbf{ul}（\texttt{unsigned long}）或\textbf{ull}（\texttt{unsigned long long}）。下面的清单列出了这些函数，每个都有两个重载，一个以\texttt{std::\allowbreak{}string}作为第一个参数，另一个以\texttt{std::\allowbreak{}wstring}作为第一个参数：


%%C020032 code

\begin{codeblockbg}
\begin{Verbatim}[commandchars=\\\{\},breaklines=true,breakanywhere=true,fontsize=\codefont,xleftmargin=17bp]
\PYG{k+kt}{int}\PYG{+w}{ }\PYG{n+nf}{stoi}\PYG{p}{(}\PYG{k}{const}\PYG{+w}{ }\PYG{n}{std}\PYG{o}{:}\PYG{o}{:}\PYG{n}{string}\PYG{o}{\PYGZam{}}\PYG{+w}{ }\PYG{n}{str}\PYG{p}{,}\PYG{+w}{ }\PYG{n}{std}\PYG{o}{:}\PYG{o}{:}\PYG{k+kt}{size\PYGZus{}t}\PYG{o}{*}\PYG{+w}{ }\PYG{n}{pos}\PYG{+w}{ }\PYG{o}{=}\PYG{+w}{ }\PYG{l+m+mi}{0}\PYG{p}{,}
\PYG{+w}{         }\PYG{k+kt}{int}\PYG{+w}{ }\PYG{n}{base}\PYG{+w}{ }\PYG{o}{=}\PYG{+w}{ }\PYG{l+m+mi}{10}\PYG{p}{)}\PYG{p}{;}
\PYG{k+kt}{int}\PYG{+w}{ }\PYG{n+nf}{stoi}\PYG{p}{(}\PYG{k}{const}\PYG{+w}{ }\PYG{n}{std}\PYG{o}{:}\PYG{o}{:}\PYG{n}{wstring}\PYG{o}{\PYGZam{}}\PYG{+w}{ }\PYG{n}{str}\PYG{p}{,}\PYG{+w}{ }\PYG{n}{std}\PYG{o}{:}\PYG{o}{:}\PYG{k+kt}{size\PYGZus{}t}\PYG{o}{*}\PYG{+w}{ }\PYG{n}{pos}\PYG{+w}{ }\PYG{o}{=}\PYG{+w}{ }\PYG{l+m+mi}{0}\PYG{p}{,}
\PYG{+w}{         }\PYG{k+kt}{int}\PYG{+w}{ }\PYG{n}{base}\PYG{+w}{ }\PYG{o}{=}\PYG{+w}{ }\PYG{l+m+mi}{10}\PYG{p}{)}\PYG{p}{;}
\PYG{k+kt}{long}\PYG{+w}{ }\PYG{n+nf}{stol}\PYG{p}{(}\PYG{k}{const}\PYG{+w}{ }\PYG{n}{std}\PYG{o}{:}\PYG{o}{:}\PYG{n}{string}\PYG{o}{\PYGZam{}}\PYG{+w}{ }\PYG{n}{str}\PYG{p}{,}\PYG{+w}{ }\PYG{n}{std}\PYG{o}{:}\PYG{o}{:}\PYG{k+kt}{size\PYGZus{}t}\PYG{o}{*}\PYG{+w}{ }\PYG{n}{pos}\PYG{+w}{ }\PYG{o}{=}\PYG{+w}{ }\PYG{l+m+mi}{0}\PYG{p}{,}
\PYG{+w}{         }\PYG{k+kt}{int}\PYG{+w}{ }\PYG{n}{base}\PYG{+w}{ }\PYG{o}{=}\PYG{+w}{ }\PYG{l+m+mi}{10}\PYG{p}{)}\PYG{p}{;}
\PYG{k+kt}{long}\PYG{+w}{ }\PYG{n+nf}{stol}\PYG{p}{(}\PYG{k}{const}\PYG{+w}{ }\PYG{n}{std}\PYG{o}{:}\PYG{o}{:}\PYG{n}{wstring}\PYG{o}{\PYGZam{}}\PYG{+w}{ }\PYG{n}{str}\PYG{p}{,}\PYG{+w}{ }\PYG{n}{std}\PYG{o}{:}\PYG{o}{:}\PYG{k+kt}{size\PYGZus{}t}\PYG{o}{*}\PYG{+w}{ }\PYG{n}{pos}\PYG{+w}{ }\PYG{o}{=}\PYG{+w}{ }\PYG{l+m+mi}{0}\PYG{p}{,}
\PYG{+w}{         }\PYG{k+kt}{int}\PYG{+w}{ }\PYG{n}{base}\PYG{+w}{ }\PYG{o}{=}\PYG{+w}{ }\PYG{l+m+mi}{10}\PYG{p}{)}\PYG{p}{;}
\PYG{k+kt}{long}\PYG{+w}{ }\PYG{k+kt}{long}\PYG{+w}{ }\PYG{n+nf}{stoll}\PYG{p}{(}\PYG{k}{const}\PYG{+w}{ }\PYG{n}{std}\PYG{o}{:}\PYG{o}{:}\PYG{n}{string}\PYG{o}{\PYGZam{}}\PYG{+w}{ }\PYG{n}{str}\PYG{p}{,}\PYG{+w}{ }\PYG{n}{std}\PYG{o}{:}\PYG{o}{:}\PYG{k+kt}{size\PYGZus{}t}\PYG{o}{*}\PYG{+w}{ }\PYG{n}{pos}\PYG{+w}{ }\PYG{o}{=}\PYG{+w}{ }\PYG{l+m+mi}{0}\PYG{p}{,}
\PYG{+w}{                }\PYG{k+kt}{int}\PYG{+w}{ }\PYG{n}{base}\PYG{+w}{ }\PYG{o}{=}\PYG{+w}{ }\PYG{l+m+mi}{10}\PYG{p}{)}\PYG{p}{;}
\PYG{k+kt}{long}\PYG{+w}{ }\PYG{k+kt}{long}\PYG{+w}{ }\PYG{n+nf}{stoll}\PYG{p}{(}\PYG{k}{const}\PYG{+w}{ }\PYG{n}{std}\PYG{o}{:}\PYG{o}{:}\PYG{n}{wstring}\PYG{o}{\PYGZam{}}\PYG{+w}{ }\PYG{n}{str}\PYG{p}{,}\PYG{+w}{ }\PYG{n}{std}\PYG{o}{:}\PYG{o}{:}\PYG{k+kt}{size\PYGZus{}t}\PYG{o}{*}\PYG{+w}{ }\PYG{n}{pos}\PYG{+w}{ }\PYG{o}{=}\PYG{+w}{ }\PYG{l+m+mi}{0}\PYG{p}{,}
\PYG{+w}{                }\PYG{k+kt}{int}\PYG{+w}{ }\PYG{n}{base}\PYG{+w}{ }\PYG{o}{=}\PYG{+w}{ }\PYG{l+m+mi}{10}\PYG{p}{)}\PYG{p}{;}
\PYG{k+kt}{unsigned}\PYG{+w}{ }\PYG{k+kt}{long}\PYG{+w}{ }\PYG{n+nf}{stoul}\PYG{p}{(}\PYG{k}{const}\PYG{+w}{ }\PYG{n}{std}\PYG{o}{:}\PYG{o}{:}\PYG{n}{string}\PYG{o}{\PYGZam{}}\PYG{+w}{ }\PYG{n}{str}\PYG{p}{,}\PYG{+w}{ }\PYG{n}{std}\PYG{o}{:}\PYG{o}{:}\PYG{k+kt}{size\PYGZus{}t}\PYG{o}{*}\PYG{+w}{ }\PYG{n}{pos}\PYG{+w}{ }\PYG{o}{=}\PYG{+w}{ }\PYG{l+m+mi}{0}\PYG{p}{,}
\PYG{+w}{                    }\PYG{k+kt}{int}\PYG{+w}{ }\PYG{n}{base}\PYG{+w}{ }\PYG{o}{=}\PYG{+w}{ }\PYG{l+m+mi}{10}\PYG{p}{)}\PYG{p}{;}
\PYG{k+kt}{unsigned}\PYG{+w}{ }\PYG{k+kt}{long}\PYG{+w}{ }\PYG{n+nf}{stoul}\PYG{p}{(}\PYG{k}{const}\PYG{+w}{ }\PYG{n}{std}\PYG{o}{:}\PYG{o}{:}\PYG{n}{wstring}\PYG{o}{\PYGZam{}}\PYG{+w}{ }\PYG{n}{str}\PYG{p}{,}\PYG{+w}{ }\PYG{n}{std}\PYG{o}{:}\PYG{o}{:}\PYG{k+kt}{size\PYGZus{}t}\PYG{o}{*}\PYG{+w}{ }\PYG{n}{pos}\PYG{+w}{ }\PYG{o}{=}\PYG{+w}{ }\PYG{l+m+mi}{0}\PYG{p}{,}
\PYG{+w}{                    }\PYG{k+kt}{int}\PYG{+w}{ }\PYG{n}{base}\PYG{+w}{ }\PYG{o}{=}\PYG{+w}{ }\PYG{l+m+mi}{10}\PYG{p}{)}\PYG{p}{;}
\PYG{k+kt}{unsigned}\PYG{+w}{ }\PYG{k+kt}{long}\PYG{+w}{ }\PYG{k+kt}{long}\PYG{+w}{ }\PYG{n+nf}{stoull}\PYG{p}{(}\PYG{k}{const}\PYG{+w}{ }\PYG{n}{std}\PYG{o}{:}\PYG{o}{:}\PYG{n}{string}\PYG{o}{\PYGZam{}}\PYG{+w}{ }\PYG{n}{str}\PYG{p}{,}
\PYG{+w}{                          }\PYG{n}{std}\PYG{o}{:}\PYG{o}{:}\PYG{k+kt}{size\PYGZus{}t}\PYG{o}{*}\PYG{+w}{ }\PYG{n}{pos}\PYG{+w}{ }\PYG{o}{=}\PYG{+w}{ }\PYG{l+m+mi}{0}\PYG{p}{,}\PYG{+w}{ }\PYG{k+kt}{int}\PYG{+w}{ }\PYG{n}{base}\PYG{+w}{ }\PYG{o}{=}\PYG{+w}{ }\PYG{l+m+mi}{10}\PYG{p}{)}\PYG{p}{;}
\PYG{k+kt}{unsigned}\PYG{+w}{ }\PYG{k+kt}{long}\PYG{+w}{ }\PYG{k+kt}{long}\PYG{+w}{ }\PYG{n+nf}{stoull}\PYG{p}{(}\PYG{k}{const}\PYG{+w}{ }\PYG{n}{std}\PYG{o}{:}\PYG{o}{:}\PYG{n}{wstring}\PYG{o}{\PYGZam{}}\PYG{+w}{ }\PYG{n}{str}\PYG{p}{,}
\PYG{+w}{                          }\PYG{n}{std}\PYG{o}{:}\PYG{o}{:}\PYG{k+kt}{size\PYGZus{}t}\PYG{o}{*}\PYG{+w}{ }\PYG{n}{pos}\PYG{+w}{ }\PYG{o}{=}\PYG{+w}{ }\PYG{l+m+mi}{0}\PYG{p}{,}\PYG{+w}{ }\PYG{k+kt}{int}\PYG{+w}{ }\PYG{n}{base}\PYG{+w}{ }\PYG{o}{=}\PYG{+w}{ }\PYG{l+m+mi}{10}\PYG{p}{)}\PYG{p}{;}
\PYG{k+kt}{float}\PYG{+w}{       }\PYG{n+nf}{stof}\PYG{p}{(}\PYG{k}{const}\PYG{+w}{ }\PYG{n}{std}\PYG{o}{:}\PYG{o}{:}\PYG{n}{string}\PYG{o}{\PYGZam{}}\PYG{+w}{ }\PYG{n}{str}\PYG{p}{,}\PYG{+w}{ }\PYG{n}{std}\PYG{o}{:}\PYG{o}{:}\PYG{k+kt}{size\PYGZus{}t}\PYG{o}{*}\PYG{+w}{ }\PYG{n}{pos}\PYG{+w}{ }\PYG{o}{=}\PYG{+w}{ }\PYG{l+m+mi}{0}\PYG{p}{)}\PYG{p}{;}
\PYG{k+kt}{float}\PYG{+w}{       }\PYG{n+nf}{stof}\PYG{p}{(}\PYG{k}{const}\PYG{+w}{ }\PYG{n}{std}\PYG{o}{:}\PYG{o}{:}\PYG{n}{wstring}\PYG{o}{\PYGZam{}}\PYG{+w}{ }\PYG{n}{str}\PYG{p}{,}\PYG{+w}{ }\PYG{n}{std}\PYG{o}{:}\PYG{o}{:}\PYG{k+kt}{size\PYGZus{}t}\PYG{o}{*}\PYG{+w}{ }\PYG{n}{pos}\PYG{+w}{ }\PYG{o}{=}\PYG{+w}{ }\PYG{l+m+mi}{0}\PYG{p}{)}\PYG{p}{;}
\PYG{k+kt}{double}\PYG{+w}{      }\PYG{n+nf}{stod}\PYG{p}{(}\PYG{k}{const}\PYG{+w}{ }\PYG{n}{std}\PYG{o}{:}\PYG{o}{:}\PYG{n}{string}\PYG{o}{\PYGZam{}}\PYG{+w}{ }\PYG{n}{str}\PYG{p}{,}\PYG{+w}{ }\PYG{n}{std}\PYG{o}{:}\PYG{o}{:}\PYG{k+kt}{size\PYGZus{}t}\PYG{o}{*}\PYG{+w}{ }\PYG{n}{pos}\PYG{+w}{ }\PYG{o}{=}\PYG{+w}{ }\PYG{l+m+mi}{0}\PYG{p}{)}\PYG{p}{;}
\PYG{k+kt}{double}\PYG{+w}{      }\PYG{n+nf}{stod}\PYG{p}{(}\PYG{k}{const}\PYG{+w}{ }\PYG{n}{std}\PYG{o}{:}\PYG{o}{:}\PYG{n}{wstring}\PYG{o}{\PYGZam{}}\PYG{+w}{ }\PYG{n}{str}\PYG{p}{,}\PYG{+w}{ }\PYG{n}{std}\PYG{o}{:}\PYG{o}{:}\PYG{k+kt}{size\PYGZus{}t}\PYG{o}{*}\PYG{+w}{ }\PYG{n}{pos}\PYG{+w}{ }\PYG{o}{=}\PYG{+w}{ }\PYG{l+m+mi}{0}\PYG{p}{)}\PYG{p}{;}
\PYG{k+kt}{long}\PYG{+w}{ }\PYG{k+kt}{double}\PYG{+w}{ }\PYG{n+nf}{stold}\PYG{p}{(}\PYG{k}{const}\PYG{+w}{ }\PYG{n}{std}\PYG{o}{:}\PYG{o}{:}\PYG{n}{string}\PYG{o}{\PYGZam{}}\PYG{+w}{ }\PYG{n}{str}\PYG{p}{,}\PYG{+w}{ }\PYG{n}{std}\PYG{o}{:}\PYG{o}{:}\PYG{k+kt}{size\PYGZus{}t}\PYG{o}{*}\PYG{+w}{ }\PYG{n}{pos}\PYG{+w}{ }\PYG{o}{=}\PYG{+w}{ }\PYG{l+m+mi}{0}\PYG{p}{)}\PYG{p}{;}
\PYG{k+kt}{long}\PYG{+w}{ }\PYG{k+kt}{double}\PYG{+w}{ }\PYG{n+nf}{stold}\PYG{p}{(}\PYG{k}{const}\PYG{+w}{ }\PYG{n}{std}\PYG{o}{:}\PYG{o}{:}\PYG{n}{wstring}\PYG{o}{\PYGZam{}}\PYG{+w}{ }\PYG{n}{str}\PYG{p}{,}\PYG{+w}{ }\PYG{n}{std}\PYG{o}{:}\PYG{o}{:}\PYG{k+kt}{size\PYGZus{}t}\PYG{o}{*}\PYG{+w}{ }\PYG{n}{pos}\PYG{+w}{ }\PYG{o}{=}\PYG{+w}{ }\PYG{l+m+mi}{0}\PYG{p}{)}\PYG{p}{;}
\end{Verbatim}
\end{codeblockbg}


%%C020033 p

字符串转整型函数的工作方式是：先丢弃非空白字符之前的全部空白，然后尽可能多地取字符组成一个有符号或无符号数（视情况而定），再把它转换为所要求的整型（\texttt{stoi()}返回\texttt{integer}，\texttt{stoul()}返回\texttt{unsigned long}，以此类推）。在下面所有示例中，结果都是整数\texttt{42}，只有最后一个示例的结果是\texttt{-\allowbreak{}42}：


%%C020034 code

\begin{codeblockbg}
\begin{Verbatim}[commandchars=\\\{\},breaklines=true,breakanywhere=true,fontsize=\codefont,xleftmargin=17bp]
\PYG{k}{auto}\PYG{+w}{ }\PYG{n}{i1}\PYG{+w}{ }\PYG{o}{=}\PYG{+w}{ }\PYG{n}{std}\PYG{o}{:}\PYG{o}{:}\PYG{n}{stoi}\PYG{p}{(}\PYG{l+s}{\PYGZdq{}}\PYG{l+s}{42}\PYG{l+s}{\PYGZdq{}}\PYG{p}{)}\PYG{p}{;}\PYG{+w}{             }\PYG{c+c1}{// i1 = 42}
\PYG{k}{auto}\PYG{+w}{ }\PYG{n}{i2}\PYG{+w}{ }\PYG{o}{=}\PYG{+w}{ }\PYG{n}{std}\PYG{o}{:}\PYG{o}{:}\PYG{n}{stoi}\PYG{p}{(}\PYG{l+s}{\PYGZdq{}}\PYG{l+s}{   42}\PYG{l+s}{\PYGZdq{}}\PYG{p}{)}\PYG{p}{;}\PYG{+w}{          }\PYG{c+c1}{// i2 = 42}
\PYG{k}{auto}\PYG{+w}{ }\PYG{n}{i3}\PYG{+w}{ }\PYG{o}{=}\PYG{+w}{ }\PYG{n}{std}\PYG{o}{:}\PYG{o}{:}\PYG{n}{stoi}\PYG{p}{(}\PYG{l+s}{\PYGZdq{}}\PYG{l+s}{   42fortytwo}\PYG{l+s}{\PYGZdq{}}\PYG{p}{)}\PYG{p}{;}\PYG{+w}{  }\PYG{c+c1}{// i3 = 42}
\PYG{k}{auto}\PYG{+w}{ }\PYG{n}{i4}\PYG{+w}{ }\PYG{o}{=}\PYG{+w}{ }\PYG{n}{std}\PYG{o}{:}\PYG{o}{:}\PYG{n}{stoi}\PYG{p}{(}\PYG{l+s}{\PYGZdq{}}\PYG{l+s}{+42}\PYG{l+s}{\PYGZdq{}}\PYG{p}{)}\PYG{p}{;}\PYG{+w}{            }\PYG{c+c1}{// i4 = 42}
\PYG{k}{auto}\PYG{+w}{ }\PYG{n}{i5}\PYG{+w}{ }\PYG{o}{=}\PYG{+w}{ }\PYG{n}{std}\PYG{o}{:}\PYG{o}{:}\PYG{n}{stoi}\PYG{p}{(}\PYG{l+s}{\PYGZdq{}}\PYG{l+s}{\PYGZhy{}42}\PYG{l+s}{\PYGZdq{}}\PYG{p}{)}\PYG{p}{;}\PYG{+w}{            }\PYG{c+c1}{// i5 = \PYGZhy{}42}
\end{Verbatim}
\end{codeblockbg}


%%C020035 p

一个合法的整数可以由以下几部分组成：


%%C020036 li

\begin{itemize}
\item 一个符号，加号（\texttt{+}）或减号（\texttt{-}）（可选）。


%%C020037 li

\item 前缀\texttt{0}，表示八进制（可选）。


%%C020038 li

\item 前缀\texttt{0x}或\texttt{0X}，表示十六进制（可选）。


%%C020039 li

\item 一串数字。

\end{itemize}


%%C020040 p

可选前缀\texttt{0}（八进制）只在指定的进制为\texttt{8}或\texttt{0}时才生效。类似地，可选前缀\texttt{0x}或\texttt{0X}（十六进制）只在指定的进制为\texttt{16}或\texttt{0}时才生效。


%%C020041 p

把字符串转换为整数的函数有三个参数：


%%C020042 li

\begin{itemize}
\item 输入字符串。


%%C020043 li

\item 一个指针，它非空时接收已处理的字符数；这个数可以包含被丢弃的前导空白、符号以及进制前缀，因此不要把它与整数值的位数混淆。


%%C020044 li

\item 一个表示进制的数；默认为\texttt{10}。

\end{itemize}


%%C020045 p

输入字符串中的合法数字取决于进制。对二进制\texttt{2}，合法数字只有\texttt{0}和\texttt{1}；对\texttt{5}进制，是\texttt{01234}。对\texttt{11}进制，合法数字是\texttt{0-\allowbreak{}9}以及字符\texttt{A}和\texttt{a}。依此类推，直到\texttt{36}进制，其合法字符是\texttt{0-\allowbreak{}9}、\texttt{A-\allowbreak{}Z}和\texttt{a-\allowbreak{}z}。


%%C020046 p

下面是更多把各种进制下的数字字符串转换为十进制整数的示例。同样，所有情况下结果要么是\texttt{42}，要么是\texttt{-\allowbreak{}42}：


%%C020047 code

\begin{codeblockbg}
\begin{Verbatim}[commandchars=\\\{\},breaklines=true,breakanywhere=true,fontsize=\codefont,xleftmargin=17bp]
\PYG{k}{auto}\PYG{+w}{ }\PYG{n}{i6}\PYG{+w}{ }\PYG{o}{=}\PYG{+w}{ }\PYG{n}{std}\PYG{o}{:}\PYG{o}{:}\PYG{n}{stoi}\PYG{p}{(}\PYG{l+s}{\PYGZdq{}}\PYG{l+s}{052}\PYG{l+s}{\PYGZdq{}}\PYG{p}{,}\PYG{+w}{ }\PYG{n+nb}{nullptr}\PYG{p}{,}\PYG{+w}{ }\PYG{l+m+mi}{8}\PYG{p}{)}\PYG{p}{;}
\PYG{k}{auto}\PYG{+w}{ }\PYG{n}{i7}\PYG{+w}{ }\PYG{o}{=}\PYG{+w}{ }\PYG{n}{std}\PYG{o}{:}\PYG{o}{:}\PYG{n}{stoi}\PYG{p}{(}\PYG{l+s}{\PYGZdq{}}\PYG{l+s}{052}\PYG{l+s}{\PYGZdq{}}\PYG{p}{,}\PYG{+w}{ }\PYG{n+nb}{nullptr}\PYG{p}{,}\PYG{+w}{ }\PYG{l+m+mi}{0}\PYG{p}{)}\PYG{p}{;}
\PYG{k}{auto}\PYG{+w}{ }\PYG{n}{i8}\PYG{+w}{ }\PYG{o}{=}\PYG{+w}{ }\PYG{n}{std}\PYG{o}{:}\PYG{o}{:}\PYG{n}{stoi}\PYG{p}{(}\PYG{l+s}{\PYGZdq{}}\PYG{l+s}{0x2A}\PYG{l+s}{\PYGZdq{}}\PYG{p}{,}\PYG{+w}{ }\PYG{n+nb}{nullptr}\PYG{p}{,}\PYG{+w}{ }\PYG{l+m+mi}{16}\PYG{p}{)}\PYG{p}{;}
\PYG{k}{auto}\PYG{+w}{ }\PYG{n}{i9}\PYG{+w}{ }\PYG{o}{=}\PYG{+w}{ }\PYG{n}{std}\PYG{o}{:}\PYG{o}{:}\PYG{n}{stoi}\PYG{p}{(}\PYG{l+s}{\PYGZdq{}}\PYG{l+s}{0x2A}\PYG{l+s}{\PYGZdq{}}\PYG{p}{,}\PYG{+w}{ }\PYG{n+nb}{nullptr}\PYG{p}{,}\PYG{+w}{ }\PYG{l+m+mi}{0}\PYG{p}{)}\PYG{p}{;}
\PYG{k}{auto}\PYG{+w}{ }\PYG{n}{i10}\PYG{+w}{ }\PYG{o}{=}\PYG{+w}{ }\PYG{n}{std}\PYG{o}{:}\PYG{o}{:}\PYG{n}{stoi}\PYG{p}{(}\PYG{l+s}{\PYGZdq{}}\PYG{l+s}{101010}\PYG{l+s}{\PYGZdq{}}\PYG{p}{,}\PYG{+w}{ }\PYG{n+nb}{nullptr}\PYG{p}{,}\PYG{+w}{ }\PYG{l+m+mi}{2}\PYG{p}{)}\PYG{p}{;}
\PYG{k}{auto}\PYG{+w}{ }\PYG{n}{i11}\PYG{+w}{ }\PYG{o}{=}\PYG{+w}{ }\PYG{n}{std}\PYG{o}{:}\PYG{o}{:}\PYG{n}{stoi}\PYG{p}{(}\PYG{l+s}{\PYGZdq{}}\PYG{l+s}{22}\PYG{l+s}{\PYGZdq{}}\PYG{p}{,}\PYG{+w}{ }\PYG{n+nb}{nullptr}\PYG{p}{,}\PYG{+w}{ }\PYG{l+m+mi}{20}\PYG{p}{)}\PYG{p}{;}
\PYG{k}{auto}\PYG{+w}{ }\PYG{n}{i12}\PYG{+w}{ }\PYG{o}{=}\PYG{+w}{ }\PYG{n}{std}\PYG{o}{:}\PYG{o}{:}\PYG{n}{stoi}\PYG{p}{(}\PYG{l+s}{\PYGZdq{}}\PYG{l+s}{\PYGZhy{}22}\PYG{l+s}{\PYGZdq{}}\PYG{p}{,}\PYG{+w}{ }\PYG{n+nb}{nullptr}\PYG{p}{,}\PYG{+w}{ }\PYG{l+m+mi}{20}\PYG{p}{)}\PYG{p}{;}
\PYG{k}{auto}\PYG{+w}{ }\PYG{n}{pos}\PYG{+w}{ }\PYG{o}{=}\PYG{+w}{ }\PYG{k+kt}{size\PYGZus{}t}\PYG{p}{\PYGZob{}}\PYG{+w}{ }\PYG{l+m+mi}{0}\PYG{+w}{ }\PYG{p}{\PYGZcb{}}\PYG{p}{;}
\PYG{k}{auto}\PYG{+w}{ }\PYG{n}{i13}\PYG{+w}{ }\PYG{o}{=}\PYG{+w}{ }\PYG{n}{std}\PYG{o}{:}\PYG{o}{:}\PYG{n}{stoi}\PYG{p}{(}\PYG{l+s}{\PYGZdq{}}\PYG{l+s}{42}\PYG{l+s}{\PYGZdq{}}\PYG{p}{,}\PYG{+w}{ }\PYG{o}{\PYGZam{}}\PYG{n}{pos}\PYG{p}{)}\PYG{p}{;}\PYG{+w}{      }\PYG{c+c1}{// pos = 2}
\PYG{k}{auto}\PYG{+w}{ }\PYG{n}{i14}\PYG{+w}{ }\PYG{o}{=}\PYG{+w}{ }\PYG{n}{std}\PYG{o}{:}\PYG{o}{:}\PYG{n}{stoi}\PYG{p}{(}\PYG{l+s}{\PYGZdq{}}\PYG{l+s}{\PYGZhy{}42}\PYG{l+s}{\PYGZdq{}}\PYG{p}{,}\PYG{+w}{ }\PYG{o}{\PYGZam{}}\PYG{n}{pos}\PYG{p}{)}\PYG{p}{;}\PYG{+w}{     }\PYG{c+c1}{// pos = 3}
\PYG{k}{auto}\PYG{+w}{ }\PYG{n}{i15}\PYG{+w}{ }\PYG{o}{=}\PYG{+w}{ }\PYG{n}{std}\PYG{o}{:}\PYG{o}{:}\PYG{n}{stoi}\PYG{p}{(}\PYG{l+s}{\PYGZdq{}}\PYG{l+s}{  +42dec}\PYG{l+s}{\PYGZdq{}}\PYG{p}{,}\PYG{+w}{ }\PYG{o}{\PYGZam{}}\PYG{n}{pos}\PYG{p}{)}\PYG{p}{;}\PYG{c+c1}{// pos = 5}
\end{Verbatim}
\end{codeblockbg}


%%C020048 p

要注意的是，这些转换函数在转换失败时会抛出异常。可能抛出两种异常：


%%C020049 li

\begin{itemize}
\item \myttbr{std::\allowbreak{}invalid\_\allowbreak{}argument}：无法执行转换时：


%%C020050 code

\begin{codeblockbg}
\begin{Verbatim}[commandchars=\\\{\},breaklines=true,breakanywhere=true,fontsize=\codefont,xleftmargin=17bp]
\PYG{k}{try}
\PYG{p}{\PYGZob{}}
\PYG{+w}{   }\PYG{k}{auto}\PYG{+w}{ }\PYG{n}{i16}\PYG{+w}{ }\PYG{o}{=}\PYG{+w}{ }\PYG{n}{std}\PYG{o}{:}\PYG{o}{:}\PYG{n}{stoi}\PYG{p}{(}\PYG{l+s}{\PYGZdq{}}\PYG{l+s}{\PYGZdq{}}\PYG{p}{)}\PYG{p}{;}
\PYG{p}{\PYGZcb{}}
\PYG{k}{catch}\PYG{+w}{ }\PYG{p}{(}\PYG{n}{std}\PYG{o}{:}\PYG{o}{:}\PYG{n}{exception}\PYG{+w}{ }\PYG{k}{const}\PYG{+w}{ }\PYG{o}{\PYGZam{}}\PYG{+w}{ }\PYG{n}{e}\PYG{p}{)}
\PYG{p}{\PYGZob{}}
\PYG{+w}{   }\PYG{c+c1}{// 输出 \PYGZdq{}invalid stoi argument\PYGZdq{}}
\PYG{+w}{   }\PYG{n}{std}\PYG{o}{:}\PYG{o}{:}\PYG{n}{cout}\PYG{+w}{ }\PYG{o}{\PYGZlt{}}\PYG{o}{\PYGZlt{}}\PYG{+w}{ }\PYG{n}{e}\PYG{p}{.}\PYG{n}{what}\PYG{p}{(}\PYG{p}{)}\PYG{+w}{ }\PYG{o}{\PYGZlt{}}\PYG{o}{\PYGZlt{}}\PYG{+w}{ }\PYG{n}{std}\PYG{o}{:}\PYG{o}{:}\PYG{n}{endl}\PYG{p}{;}
\PYG{p}{\PYGZcb{}}
\end{Verbatim}
\end{codeblockbg}


%%C020051 li

\item \myttbr{std::\allowbreak{}out\_\allowbreak{}of\_\allowbreak{}range}：转换后的值超出结果类型的取值范围时（或者底层函数把\texttt{errno}置为\texttt{ERANGE}时）：


%%C020052 code

\begin{codeblockbg}
\begin{Verbatim}[commandchars=\\\{\},breaklines=true,breakanywhere=true,fontsize=\codefont,xleftmargin=17bp]
\PYG{k}{try}
\PYG{p}{\PYGZob{}}
\PYG{+w}{   }\PYG{c+c1}{// 正常}
\PYG{+w}{   }\PYG{k}{auto}\PYG{+w}{ }\PYG{n}{i17}\PYG{+w}{ }\PYG{o}{=}\PYG{+w}{ }\PYG{n}{std}\PYG{o}{:}\PYG{o}{:}\PYG{n}{stoll}\PYG{p}{(}\PYG{l+s}{\PYGZdq{}}\PYG{l+s}{12345678901234}\PYG{l+s}{\PYGZdq{}}\PYG{p}{)}\PYG{p}{;}
\PYG{+w}{   }\PYG{c+c1}{// 抛出 std::out\PYGZus{}of\PYGZus{}range}
\PYG{+w}{   }\PYG{k}{auto}\PYG{+w}{ }\PYG{n}{i18}\PYG{+w}{ }\PYG{o}{=}\PYG{+w}{ }\PYG{n}{std}\PYG{o}{:}\PYG{o}{:}\PYG{n}{stoi}\PYG{p}{(}\PYG{l+s}{\PYGZdq{}}\PYG{l+s}{12345678901234}\PYG{l+s}{\PYGZdq{}}\PYG{p}{)}\PYG{p}{;}
\PYG{p}{\PYGZcb{}}
\PYG{k}{catch}\PYG{+w}{ }\PYG{p}{(}\PYG{n}{std}\PYG{o}{:}\PYG{o}{:}\PYG{n}{exception}\PYG{+w}{ }\PYG{k}{const}\PYG{+w}{ }\PYG{o}{\PYGZam{}}\PYG{+w}{ }\PYG{n}{e}\PYG{p}{)}
\PYG{p}{\PYGZob{}}
\PYG{+w}{   }\PYG{c+c1}{// 输出 \PYGZdq{}stoi argument out of range\PYGZdq{}}
\PYG{+w}{   }\PYG{n}{std}\PYG{o}{:}\PYG{o}{:}\PYG{n}{cout}\PYG{+w}{ }\PYG{o}{\PYGZlt{}}\PYG{o}{\PYGZlt{}}\PYG{+w}{ }\PYG{n}{e}\PYG{p}{.}\PYG{n}{what}\PYG{p}{(}\PYG{p}{)}\PYG{+w}{ }\PYG{o}{\PYGZlt{}}\PYG{o}{\PYGZlt{}}\PYG{+w}{ }\PYG{n}{std}\PYG{o}{:}\PYG{o}{:}\PYG{n}{endl}\PYG{p}{;}
\PYG{p}{\PYGZcb{}}
\end{Verbatim}
\end{codeblockbg}

\end{itemize}


%%C020053 p

另一组把字符串转换为浮点类型的函数与之非常相似，只是没有表示进制的参数。合法的浮点值在输入字符串中可以有几种不同的表示形式：


%%C020054 li

\begin{itemize}
\item 十进制浮点表达式（可选符号、带可选小数点的十进制数字序列、可选的\texttt{e}或\texttt{E}后跟带可选符号的指数）。


%%C020055 li

\item 二进制浮点表达式（可选符号、\texttt{0x}或\texttt{0X}前缀、带可选小数点的十六进制数字序列、可选的\texttt{p}或\texttt{P}后跟带可选符号的指数）。


%%C020056 li

\item 无穷大表达式（可选符号后跟大小写不敏感的\texttt{INF}或\texttt{INFINITY}）。


%%C020057 li

\item 非数字表达式（可选符号后跟大小写不敏感的\texttt{NAN}，后面可能还跟其他字母数字字符）。

\end{itemize}


%%C020058 p

下面是把字符串转换为\texttt{double}的各种示例：


%%C020059 code

\begin{codeblockbg}
\begin{Verbatim}[commandchars=\\\{\},breaklines=true,breakanywhere=true,fontsize=\codefont,xleftmargin=17bp]
\PYG{k}{auto}\PYG{+w}{ }\PYG{n}{d1}\PYG{+w}{ }\PYG{o}{=}\PYG{+w}{ }\PYG{n}{std}\PYG{o}{:}\PYG{o}{:}\PYG{n}{stod}\PYG{p}{(}\PYG{l+s}{\PYGZdq{}}\PYG{l+s}{123.45}\PYG{l+s}{\PYGZdq{}}\PYG{p}{)}\PYG{p}{;}\PYG{+w}{         }\PYG{c+c1}{// d1 =  123.45000000000000}
\PYG{k}{auto}\PYG{+w}{ }\PYG{n}{d2}\PYG{+w}{ }\PYG{o}{=}\PYG{+w}{ }\PYG{n}{std}\PYG{o}{:}\PYG{o}{:}\PYG{n}{stod}\PYG{p}{(}\PYG{l+s}{\PYGZdq{}}\PYG{l+s}{+123.45}\PYG{l+s}{\PYGZdq{}}\PYG{p}{)}\PYG{p}{;}\PYG{+w}{        }\PYG{c+c1}{// d2 =  123.45000000000000}
\PYG{k}{auto}\PYG{+w}{ }\PYG{n}{d3}\PYG{+w}{ }\PYG{o}{=}\PYG{+w}{ }\PYG{n}{std}\PYG{o}{:}\PYG{o}{:}\PYG{n}{stod}\PYG{p}{(}\PYG{l+s}{\PYGZdq{}}\PYG{l+s}{\PYGZhy{}123.45}\PYG{l+s}{\PYGZdq{}}\PYG{p}{)}\PYG{p}{;}\PYG{+w}{        }\PYG{c+c1}{// d3 = \PYGZhy{}123.45000000000000}
\PYG{k}{auto}\PYG{+w}{ }\PYG{n}{d4}\PYG{+w}{ }\PYG{o}{=}\PYG{+w}{ }\PYG{n}{std}\PYG{o}{:}\PYG{o}{:}\PYG{n}{stod}\PYG{p}{(}\PYG{l+s}{\PYGZdq{}}\PYG{l+s}{  123.45}\PYG{l+s}{\PYGZdq{}}\PYG{p}{)}\PYG{p}{;}\PYG{+w}{       }\PYG{c+c1}{// d4 =  123.45000000000000}
\PYG{k}{auto}\PYG{+w}{ }\PYG{n}{d5}\PYG{+w}{ }\PYG{o}{=}\PYG{+w}{ }\PYG{n}{std}\PYG{o}{:}\PYG{o}{:}\PYG{n}{stod}\PYG{p}{(}\PYG{l+s}{\PYGZdq{}}\PYG{l+s}{  \PYGZhy{}123.45abc}\PYG{l+s}{\PYGZdq{}}\PYG{p}{)}\PYG{p}{;}\PYG{+w}{   }\PYG{c+c1}{// d5 = \PYGZhy{}123.45000000000000}
\PYG{k}{auto}\PYG{+w}{ }\PYG{n}{d6}\PYG{+w}{ }\PYG{o}{=}\PYG{+w}{ }\PYG{n}{std}\PYG{o}{:}\PYG{o}{:}\PYG{n}{stod}\PYG{p}{(}\PYG{l+s}{\PYGZdq{}}\PYG{l+s}{1.2345e+2}\PYG{l+s}{\PYGZdq{}}\PYG{p}{)}\PYG{p}{;}\PYG{+w}{      }\PYG{c+c1}{// d6 =  123.45000000000000}
\PYG{k}{auto}\PYG{+w}{ }\PYG{n}{d7}\PYG{+w}{ }\PYG{o}{=}\PYG{+w}{ }\PYG{n}{std}\PYG{o}{:}\PYG{o}{:}\PYG{n}{stod}\PYG{p}{(}\PYG{l+s}{\PYGZdq{}}\PYG{l+s}{0xF.6E6666p3}\PYG{l+s}{\PYGZdq{}}\PYG{p}{)}\PYG{p}{;}\PYG{+w}{   }\PYG{c+c1}{// d7 =  123.44999980926514}
\PYG{k}{auto}\PYG{+w}{ }\PYG{n}{d8}\PYG{+w}{ }\PYG{o}{=}\PYG{+w}{ }\PYG{n}{std}\PYG{o}{:}\PYG{o}{:}\PYG{n}{stod}\PYG{p}{(}\PYG{l+s}{\PYGZdq{}}\PYG{l+s}{INF}\PYG{l+s}{\PYGZdq{}}\PYG{p}{)}\PYG{p}{;}\PYG{+w}{            }\PYG{c+c1}{// d8 = inf}
\PYG{k}{auto}\PYG{+w}{ }\PYG{n}{d9}\PYG{+w}{ }\PYG{o}{=}\PYG{+w}{ }\PYG{n}{std}\PYG{o}{:}\PYG{o}{:}\PYG{n}{stod}\PYG{p}{(}\PYG{l+s}{\PYGZdq{}}\PYG{l+s}{\PYGZhy{}infinity}\PYG{l+s}{\PYGZdq{}}\PYG{p}{)}\PYG{p}{;}\PYG{+w}{      }\PYG{c+c1}{// d9 = \PYGZhy{}inf}
\PYG{k}{auto}\PYG{+w}{ }\PYG{n}{d10}\PYG{+w}{ }\PYG{o}{=}\PYG{+w}{ }\PYG{n}{std}\PYG{o}{:}\PYG{o}{:}\PYG{n}{stod}\PYG{p}{(}\PYG{l+s}{\PYGZdq{}}\PYG{l+s}{NAN}\PYG{l+s}{\PYGZdq{}}\PYG{p}{)}\PYG{p}{;}\PYG{+w}{           }\PYG{c+c1}{// d10 = nan}
\PYG{k}{auto}\PYG{+w}{ }\PYG{n}{d11}\PYG{+w}{ }\PYG{o}{=}\PYG{+w}{ }\PYG{n}{std}\PYG{o}{:}\PYG{o}{:}\PYG{n}{stod}\PYG{p}{(}\PYG{l+s}{\PYGZdq{}}\PYG{l+s}{\PYGZhy{}nanabc}\PYG{l+s}{\PYGZdq{}}\PYG{p}{)}\PYG{p}{;}\PYG{+w}{       }\PYG{c+c1}{// d11 = \PYGZhy{}nan}
\end{Verbatim}
\end{codeblockbg}


%%C020060 p

前面见到的形如\texttt{0xF.\allowbreak{}6E6666p3}的浮点二进制科学记数法不是本案例的主题。不过为了便于清晰理解，这里给出简要说明；若需细节，建议查阅其他参考资料。采用二进制科学记数法的浮点常量由几部分组成：


%%C020061 li

\begin{itemize}
\item 十六进制前缀\texttt{0x}。


%%C020062 li

\item 一个整数部分，本例中是\texttt{F}，其十进制值为15。


%%C020063 li

\item 一个小数部分，本例中是\texttt{6E6666}，二进制为\myttbr{011011100110011001100110}。要把它转换为十进制，需要把2的负幂相加：\myttbr{1/\allowbreak{}4 + 1/\allowbreak{}8 + 1/\allowbreak{}32 + 1/\allowbreak{}64 + 1/\allowbreak{}128 + ...}。


%%C020064 li

\item 一个后缀，表示2的幂；本例中\texttt{p3}表示2的3次幂。

\end{itemize}


%%C020065 p

十进制等价值由有效数字（整数部分与小数部分合成）乘以底数的指数次幂得到。对给定的十六进制二进制浮点字面量，有效数字为\texttt{15.\allowbreak{}4312499...}（注意第七位之后的数字未显示），底数为2，指数为3。因此结果是\texttt{15.\allowbreak{}4212499... * 8}，即\texttt{123.\allowbreak{}44999980926514}。


%%C020066 h3

\subsection{延伸阅读}


%%C020067 li

\begin{itemize}
\item \textit{数值类型的取值范围及其他属性}

\end{itemize}


%%C020068 h2

\section{数值类型的取值范围及其他属性}


%%C020069 p

有时需要了解并使用某种数值类型（例如\texttt{char}、\texttt{int}或\texttt{double}）所能表示的最小值与最大值。很多开发者为此使用标准C宏，例如\texttt{CHAR\_\allowbreak{}MIN}/\texttt{CHAR\_\allowbreak{}MAX}、\texttt{INT\_\allowbreak{}MIN}/\texttt{INT\_\allowbreak{}MAX}或\texttt{DBL\_\allowbreak{}MIN}/\texttt{DBL\_\allowbreak{}MAX}。C++提供了一个名为\texttt{numeric\_\allowbreak{}limits}的类模板，并为每种数值类型提供了特化，使你能够查询某个类型的最小值与最大值；它的能力还不止于此，另外还提供了用于查询类型属性的常量，例如类型是否有符号、表示其值需要多少比特、对浮点类型而言能否表示无穷大，等等。在C++11之前，\texttt{numeric\_\allowbreak{}limits<\allowbreak{}T>}的用途受到限制，因为它不能用在需要常量的地方（例如数组大小和switch分支）。正因如此，开发者更愿意在整个代码中使用C宏。到C++11，情况不再如此，因为\texttt{numeric\_\allowbreak{}limits<\allowbreak{}T>}的所有静态成员现在都是\texttt{constexpr}，这意味着它们可以用在任何需要常量表达式的地方。


%%C020070 h3

\subsection{准备工作}


%%C020071 p

\texttt{numeric\_\allowbreak{}limits<\allowbreak{}T>}类模板位于\texttt{<\allowbreak{}limits>}头文件的\texttt{std}命名空间中。


%%C020072 h3

\subsection{具体做法}


%%C020073 p

使用\myttbr{std::\allowbreak{}numeric\_\allowbreak{}limits<\allowbreak{}T>}查询数值类型\texttt{T}的各种属性：


%%C020074 li

\begin{itemize}
\item 用\texttt{min()}和\texttt{max()}静态方法获取某个类型的最小和最大有限值：


%%C020075 code

\begin{codeblockbg}
\begin{Verbatim}[commandchars=\\\{\},breaklines=true,breakanywhere=true,fontsize=\codefont,xleftmargin=17bp]
\PYG{k}{template}\PYG{o}{\PYGZlt{}}\PYG{k}{typename}\PYG{+w}{ }\PYG{n+nc}{T}\PYG{p}{,}\PYG{+w}{ }\PYG{k}{typename}\PYG{+w}{ }\PYG{n+nc}{I}\PYG{o}{\PYGZgt{}}
\PYG{n}{T}\PYG{+w}{ }\PYG{n}{minimum}\PYG{p}{(}\PYG{n}{I}\PYG{+w}{ }\PYG{k}{const}\PYG{+w}{ }\PYG{n}{start}\PYG{p}{,}\PYG{+w}{ }\PYG{n}{I}\PYG{+w}{ }\PYG{k}{const}\PYG{+w}{ }\PYG{n}{end}\PYG{p}{)}
\PYG{p}{\PYGZob{}}
\PYG{+w}{  }\PYG{n}{T}\PYG{+w}{ }\PYG{n}{minval}\PYG{+w}{ }\PYG{o}{=}\PYG{+w}{ }\PYG{n}{std}\PYG{o}{:}\PYG{o}{:}\PYG{n}{numeric\PYGZus{}limits}\PYG{o}{\PYGZlt{}}\PYG{n}{T}\PYG{o}{\PYGZgt{}}\PYG{o}{:}\PYG{o}{:}\PYG{n}{max}\PYG{p}{(}\PYG{p}{)}\PYG{p}{;}
\PYG{+w}{  }\PYG{k}{for}\PYG{+w}{ }\PYG{p}{(}\PYG{k}{auto}\PYG{+w}{ }\PYG{n}{i}\PYG{+w}{ }\PYG{o}{=}\PYG{+w}{ }\PYG{n}{start}\PYG{p}{;}\PYG{+w}{ }\PYG{n}{i}\PYG{+w}{ }\PYG{o}{\PYGZlt{}}\PYG{+w}{ }\PYG{n}{end}\PYG{p}{;}\PYG{+w}{ }\PYG{o}{+}\PYG{o}{+}\PYG{n}{i}\PYG{p}{)}
\PYG{+w}{  }\PYG{p}{\PYGZob{}}
\PYG{+w}{    }\PYG{k}{if}\PYG{+w}{ }\PYG{p}{(}\PYG{o}{*}\PYG{n}{i}\PYG{+w}{ }\PYG{o}{\PYGZlt{}}\PYG{+w}{ }\PYG{n}{minval}\PYG{p}{)}
\PYG{+w}{      }\PYG{n}{minval}\PYG{+w}{ }\PYG{o}{=}\PYG{+w}{ }\PYG{o}{*}\PYG{n}{i}\PYG{p}{;}
\PYG{+w}{  }\PYG{p}{\PYGZcb{}}
\PYG{+w}{  }\PYG{k}{return}\PYG{+w}{ }\PYG{n}{minval}\PYG{p}{;}
\PYG{p}{\PYGZcb{}}
\PYG{k+kt}{int}\PYG{+w}{ }\PYG{n}{range}\PYG{p}{[}\PYG{n}{std}\PYG{o}{:}\PYG{o}{:}\PYG{n}{numeric\PYGZus{}limits}\PYG{o}{\PYGZlt{}}\PYG{k+kt}{char}\PYG{o}{\PYGZgt{}}\PYG{o}{:}\PYG{o}{:}\PYG{n}{max}\PYG{p}{(}\PYG{p}{)}\PYG{+w}{ }\PYG{o}{+}\PYG{+w}{ }\PYG{l+m+mi}{1}\PYG{p}{]}\PYG{+w}{ }\PYG{o}{=}\PYG{+w}{ }\PYG{p}{\PYGZob{}}\PYG{+w}{ }\PYG{l+m+mi}{0}\PYG{+w}{ }\PYG{p}{\PYGZcb{}}\PYG{p}{;}
\PYG{k}{switch}\PYG{p}{(}\PYG{n}{get\PYGZus{}value}\PYG{p}{(}\PYG{p}{)}\PYG{p}{)}
\PYG{p}{\PYGZob{}}
\PYG{+w}{  }\PYG{k}{case}\PYG{+w}{ }\PYG{n+no}{std}\PYG{o}{:}\PYG{o}{:}\PYG{n+no}{numeric\PYGZus{}limits}\PYG{o}{\PYGZlt{}}\PYG{n+no}{int}\PYG{o}{\PYGZgt{}}\PYG{o}{:}\PYG{o}{:}\PYG{n+no}{min}\PYG{p}{(}\PYG{p}{)}\PYG{p}{:}
\PYG{+w}{  }\PYG{k}{break}\PYG{p}{;}
\PYG{p}{\PYGZcb{}}
\end{Verbatim}
\end{codeblockbg}


%%C020076 li

\item 用其他静态方法与静态常量获取数值类型的其他属性：


%%C020077 code

\begin{codeblockbg}
\begin{Verbatim}[commandchars=\\\{\},breaklines=true,breakanywhere=true,fontsize=\codefont,xleftmargin=17bp]
\PYG{k}{auto}\PYG{+w}{ }\PYG{n}{n}\PYG{+w}{ }\PYG{o}{=}\PYG{+w}{ }\PYG{l+m+mi}{42}\PYG{p}{;}
\PYG{n}{std}\PYG{o}{:}\PYG{o}{:}\PYG{n}{bitset}\PYG{o}{\PYGZlt{}}\PYG{n}{std}\PYG{o}{:}\PYG{o}{:}\PYG{n}{numeric\PYGZus{}limits}\PYG{o}{\PYGZlt{}}\PYG{k}{decltype}\PYG{p}{(}\PYG{n}{n}\PYG{p}{)}\PYG{o}{\PYGZgt{}}\PYG{o}{:}\PYG{o}{:}\PYG{n}{digits}\PYG{o}{\PYGZgt{}}
\PYG{+w}{  }\PYG{n}{bits}\PYG{+w}{ }\PYG{p}{\PYGZob{}}\PYG{+w}{ }\PYG{k}{static\PYGZus{}cast}\PYG{o}{\PYGZlt{}}\PYG{k+kt}{unsigned}\PYG{+w}{ }\PYG{k+kt}{long}\PYG{+w}{ }\PYG{k+kt}{long}\PYG{o}{\PYGZgt{}}\PYG{p}{(}\PYG{n}{n}\PYG{p}{)}\PYG{+w}{ }\PYG{p}{\PYGZcb{}}\PYG{p}{;}
\end{Verbatim}
\end{codeblockbg}

\end{itemize}


%%C020078 note

\begin{myNote}
在C++11中，\myttbr{std::\allowbreak{}numeric\_\allowbreak{}limits<\allowbreak{}T>}的使用位置不再有任何限制；因此，在现代C++代码中请优先使用它，而不是C宏。
\end{myNote}


%%C020079 h3

\subsection{工作原理}


%%C020080 p

\myttbr{std::\allowbreak{}numeric\_\allowbreak{}limits<\allowbreak{}T>}是一个类模板，使开发者能够查询数值类型的属性。实际的值由特化提供，标准库为所有内置数值类型（\texttt{char}、\texttt{short}、\texttt{int}、\texttt{long}、\texttt{float}、\texttt{double}等）都提供了特化。此外，第三方也可以为其他类型提供额外实现。例如，某个数值库实现了\texttt{bigint}整数类型和\texttt{decimal}类型，并为它们提供\texttt{numeric\_\allowbreak{}limits}的特化（如\myttbr{numeric\_\allowbreak{}limits<\allowbreak{}bigint>}和\myttbr{numeric\_\allowbreak{}limits<\allowbreak{}decimal>}）。


%%C020081 p

\texttt{<\allowbreak{}limits>}头文件中提供了下列数值类型的特化。注意\texttt{char16\_\allowbreak{}t}和\texttt{char32\_\allowbreak{}t}的特化是C++11新增的；其他特化此前就已存在。除上面列出的特化之外，库还为这些数值类型的每个\texttt{cv-\allowbreak{}qualified}版本提供了特化，它们与非限定版本的特化完全相同。例如考虑类型\texttt{int}，实际存在四个特化（而且完全相同）：\myttbr{numeric\_\allowbreak{}limits<\allowbreak{}int>}、\myttbr{numeric\_\allowbreak{}limits<\allowbreak{}const int>}、\myttbr{numeric\_\allowbreak{}limits<\allowbreak{}volatile int>}和\myttbr{numeric\_\allowbreak{}limits<\allowbreak{}const volatile int>}：


%%C020082 code

\begin{codeblockbg}
\begin{Verbatim}[commandchars=\\\{\},breaklines=true,breakanywhere=true,fontsize=\codefont,xleftmargin=17bp]
\PYG{k}{template}\PYG{o}{\PYGZlt{}}\PYG{o}{\PYGZgt{}}\PYG{+w}{ }\PYG{k}{class}\PYG{+w}{ }\PYG{n+nc}{numeric\PYGZus{}limits}\PYG{o}{\PYGZlt{}}\PYG{k+kt}{bool}\PYG{o}{\PYGZgt{}}\PYG{p}{;}
\PYG{k}{template}\PYG{o}{\PYGZlt{}}\PYG{o}{\PYGZgt{}}\PYG{+w}{ }\PYG{k}{class}\PYG{+w}{ }\PYG{n+nc}{numeric\PYGZus{}limits}\PYG{o}{\PYGZlt{}}\PYG{k+kt}{char}\PYG{o}{\PYGZgt{}}\PYG{p}{;}
\PYG{k}{template}\PYG{o}{\PYGZlt{}}\PYG{o}{\PYGZgt{}}\PYG{+w}{ }\PYG{k}{class}\PYG{+w}{ }\PYG{n+nc}{numeric\PYGZus{}limits}\PYG{o}{\PYGZlt{}}\PYG{k+kt}{signed}\PYG{+w}{ }\PYG{k+kt}{char}\PYG{o}{\PYGZgt{}}\PYG{p}{;}
\PYG{k}{template}\PYG{o}{\PYGZlt{}}\PYG{o}{\PYGZgt{}}\PYG{+w}{ }\PYG{k}{class}\PYG{+w}{ }\PYG{n+nc}{numeric\PYGZus{}limits}\PYG{o}{\PYGZlt{}}\PYG{k+kt}{unsigned}\PYG{+w}{ }\PYG{k+kt}{char}\PYG{o}{\PYGZgt{}}\PYG{p}{;}
\PYG{k}{template}\PYG{o}{\PYGZlt{}}\PYG{o}{\PYGZgt{}}\PYG{+w}{ }\PYG{k}{class}\PYG{+w}{ }\PYG{n+nc}{numeric\PYGZus{}limits}\PYG{o}{\PYGZlt{}}\PYG{k+kt}{wchar\PYGZus{}t}\PYG{o}{\PYGZgt{}}\PYG{p}{;}
\PYG{k}{template}\PYG{o}{\PYGZlt{}}\PYG{o}{\PYGZgt{}}\PYG{+w}{ }\PYG{k}{class}\PYG{+w}{ }\PYG{n+nc}{numeric\PYGZus{}limits}\PYG{o}{\PYGZlt{}}\PYG{k+kt}{char16\PYGZus{}t}\PYG{o}{\PYGZgt{}}\PYG{p}{;}
\PYG{k}{template}\PYG{o}{\PYGZlt{}}\PYG{o}{\PYGZgt{}}\PYG{+w}{ }\PYG{k}{class}\PYG{+w}{ }\PYG{n+nc}{numeric\PYGZus{}limits}\PYG{o}{\PYGZlt{}}\PYG{k+kt}{char32\PYGZus{}t}\PYG{o}{\PYGZgt{}}\PYG{p}{;}
\PYG{k}{template}\PYG{o}{\PYGZlt{}}\PYG{o}{\PYGZgt{}}\PYG{+w}{ }\PYG{k}{class}\PYG{+w}{ }\PYG{n+nc}{numeric\PYGZus{}limits}\PYG{o}{\PYGZlt{}}\PYG{k+kt}{short}\PYG{o}{\PYGZgt{}}\PYG{p}{;}
\PYG{k}{template}\PYG{o}{\PYGZlt{}}\PYG{o}{\PYGZgt{}}\PYG{+w}{ }\PYG{k}{class}\PYG{+w}{ }\PYG{n+nc}{numeric\PYGZus{}limits}\PYG{o}{\PYGZlt{}}\PYG{k+kt}{unsigned}\PYG{+w}{ }\PYG{k+kt}{short}\PYG{o}{\PYGZgt{}}\PYG{p}{;}
\PYG{k}{template}\PYG{o}{\PYGZlt{}}\PYG{o}{\PYGZgt{}}\PYG{+w}{ }\PYG{k}{class}\PYG{+w}{ }\PYG{n+nc}{numeric\PYGZus{}limits}\PYG{o}{\PYGZlt{}}\PYG{k+kt}{int}\PYG{o}{\PYGZgt{}}\PYG{p}{;}
\PYG{k}{template}\PYG{o}{\PYGZlt{}}\PYG{o}{\PYGZgt{}}\PYG{+w}{ }\PYG{k}{class}\PYG{+w}{ }\PYG{n+nc}{numeric\PYGZus{}limits}\PYG{o}{\PYGZlt{}}\PYG{k+kt}{unsigned}\PYG{+w}{ }\PYG{k+kt}{int}\PYG{o}{\PYGZgt{}}\PYG{p}{;}
\PYG{k}{template}\PYG{o}{\PYGZlt{}}\PYG{o}{\PYGZgt{}}\PYG{+w}{ }\PYG{k}{class}\PYG{+w}{ }\PYG{n+nc}{numeric\PYGZus{}limits}\PYG{o}{\PYGZlt{}}\PYG{k+kt}{long}\PYG{o}{\PYGZgt{}}\PYG{p}{;}
\PYG{k}{template}\PYG{o}{\PYGZlt{}}\PYG{o}{\PYGZgt{}}\PYG{+w}{ }\PYG{k}{class}\PYG{+w}{ }\PYG{n+nc}{numeric\PYGZus{}limits}\PYG{o}{\PYGZlt{}}\PYG{k+kt}{unsigned}\PYG{+w}{ }\PYG{k+kt}{long}\PYG{o}{\PYGZgt{}}\PYG{p}{;}
\PYG{k}{template}\PYG{o}{\PYGZlt{}}\PYG{o}{\PYGZgt{}}\PYG{+w}{ }\PYG{k}{class}\PYG{+w}{ }\PYG{n+nc}{numeric\PYGZus{}limits}\PYG{o}{\PYGZlt{}}\PYG{k+kt}{long}\PYG{+w}{ }\PYG{k+kt}{long}\PYG{o}{\PYGZgt{}}\PYG{p}{;}
\PYG{k}{template}\PYG{o}{\PYGZlt{}}\PYG{o}{\PYGZgt{}}\PYG{+w}{ }\PYG{k}{class}\PYG{+w}{ }\PYG{n+nc}{numeric\PYGZus{}limits}\PYG{o}{\PYGZlt{}}\PYG{k+kt}{unsigned}\PYG{+w}{ }\PYG{k+kt}{long}\PYG{+w}{ }\PYG{k+kt}{long}\PYG{o}{\PYGZgt{}}\PYG{p}{;}
\PYG{k}{template}\PYG{o}{\PYGZlt{}}\PYG{o}{\PYGZgt{}}\PYG{+w}{ }\PYG{k}{class}\PYG{+w}{ }\PYG{n+nc}{numeric\PYGZus{}limits}\PYG{o}{\PYGZlt{}}\PYG{k+kt}{float}\PYG{o}{\PYGZgt{}}\PYG{p}{;}
\PYG{k}{template}\PYG{o}{\PYGZlt{}}\PYG{o}{\PYGZgt{}}\PYG{+w}{ }\PYG{k}{class}\PYG{+w}{ }\PYG{n+nc}{numeric\PYGZus{}limits}\PYG{o}{\PYGZlt{}}\PYG{k+kt}{double}\PYG{o}{\PYGZgt{}}\PYG{p}{;}
\PYG{k}{template}\PYG{o}{\PYGZlt{}}\PYG{o}{\PYGZgt{}}\PYG{+w}{ }\PYG{k}{class}\PYG{+w}{ }\PYG{n+nc}{numeric\PYGZus{}limits}\PYG{o}{\PYGZlt{}}\PYG{k+kt}{long}\PYG{+w}{ }\PYG{k+kt}{double}\PYG{o}{\PYGZgt{}}\PYG{p}{;}
\end{Verbatim}
\end{codeblockbg}


%%C020083 p

如前所述，在C++11中\texttt{numeric\_\allowbreak{}limits}的所有静态成员都是\texttt{constexpr}，这意味着它们可以用在任何需要常量表达式的地方。相较C++宏，它们有几个主要优点：


%%C020084 li

\begin{itemize}
\item 它们更容易记住，因为你唯一需要知道的是类型名（这个你本来就知道），而不是数不清的宏名。

%%C020085 li

\item 它们支持C语言中不存在的类型，比如\texttt{char16\_\allowbreak{}t}和\texttt{char32\_\allowbreak{}t}。


%%C020086 li

\item 当模板中无法得知具体类型时，它们是唯一可行的方案。


%%C020087 li

\item 最小值和最大值只是它提供的众多类型属性中的两个，因此它的实际用途远超数值极限。顺带一提，正因如此，这个类也许本应叫作\myttbr{numeric\_\allowbreak{}properties}，而不是\texttt{numeric\_\allowbreak{}limits}。

\end{itemize}


%%C020088 p

下面的函数模板\myttbr{print\_\allowbreak{}type\_\allowbreak{}properties()}会打印该类型的最小和最大有限值以及其他信息：


%%C020089 code

\begin{codeblockbg}
\begin{Verbatim}[commandchars=\\\{\},breaklines=true,breakanywhere=true,fontsize=\codefont,xleftmargin=17bp]
\PYG{k}{template}\PYG{+w}{ }\PYG{o}{\PYGZlt{}}\PYG{k}{typename}\PYG{+w}{ }\PYG{n+nc}{T}\PYG{o}{\PYGZgt{}}
\PYG{k+kt}{void}\PYG{+w}{ }\PYG{n}{print\PYGZus{}type\PYGZus{}properties}\PYG{p}{(}\PYG{p}{)}
\PYG{p}{\PYGZob{}}
\PYG{+w}{  }\PYG{n}{std}\PYG{o}{:}\PYG{o}{:}\PYG{n}{cout}
\PYG{+w}{    }\PYG{o}{\PYGZlt{}}\PYG{o}{\PYGZlt{}}\PYG{+w}{ }\PYG{l+s}{\PYGZdq{}}\PYG{l+s}{min=}\PYG{l+s}{\PYGZdq{}}
\PYG{+w}{    }\PYG{o}{\PYGZlt{}}\PYG{o}{\PYGZlt{}}\PYG{+w}{ }\PYG{n}{std}\PYG{o}{:}\PYG{o}{:}\PYG{n}{numeric\PYGZus{}limits}\PYG{o}{\PYGZlt{}}\PYG{n}{T}\PYG{o}{\PYGZgt{}}\PYG{o}{:}\PYG{o}{:}\PYG{n}{min}\PYG{p}{(}\PYG{p}{)}\PYG{+w}{        }\PYG{o}{\PYGZlt{}}\PYG{o}{\PYGZlt{}}\PYG{+w}{ }\PYG{n}{std}\PYG{o}{:}\PYG{o}{:}\PYG{n}{endl}
\PYG{+w}{    }\PYG{o}{\PYGZlt{}}\PYG{o}{\PYGZlt{}}\PYG{+w}{ }\PYG{l+s}{\PYGZdq{}}\PYG{l+s}{max=}\PYG{l+s}{\PYGZdq{}}
\PYG{+w}{    }\PYG{o}{\PYGZlt{}}\PYG{o}{\PYGZlt{}}\PYG{+w}{ }\PYG{n}{std}\PYG{o}{:}\PYG{o}{:}\PYG{n}{numeric\PYGZus{}limits}\PYG{o}{\PYGZlt{}}\PYG{n}{T}\PYG{o}{\PYGZgt{}}\PYG{o}{:}\PYG{o}{:}\PYG{n}{max}\PYG{p}{(}\PYG{p}{)}\PYG{+w}{        }\PYG{o}{\PYGZlt{}}\PYG{o}{\PYGZlt{}}\PYG{+w}{ }\PYG{n}{std}\PYG{o}{:}\PYG{o}{:}\PYG{n}{endl}
\PYG{+w}{    }\PYG{o}{\PYGZlt{}}\PYG{o}{\PYGZlt{}}\PYG{+w}{ }\PYG{l+s}{\PYGZdq{}}\PYG{l+s}{bits=}\PYG{l+s}{\PYGZdq{}}
\PYG{+w}{    }\PYG{o}{\PYGZlt{}}\PYG{o}{\PYGZlt{}}\PYG{+w}{ }\PYG{n}{std}\PYG{o}{:}\PYG{o}{:}\PYG{n}{numeric\PYGZus{}limits}\PYG{o}{\PYGZlt{}}\PYG{n}{T}\PYG{o}{\PYGZgt{}}\PYG{o}{:}\PYG{o}{:}\PYG{n}{digits}\PYG{+w}{       }\PYG{o}{\PYGZlt{}}\PYG{o}{\PYGZlt{}}\PYG{+w}{ }\PYG{n}{std}\PYG{o}{:}\PYG{o}{:}\PYG{n}{endl}
\PYG{+w}{    }\PYG{o}{\PYGZlt{}}\PYG{o}{\PYGZlt{}}\PYG{+w}{ }\PYG{l+s}{\PYGZdq{}}\PYG{l+s}{decdigits=}\PYG{l+s}{\PYGZdq{}}
\PYG{+w}{    }\PYG{o}{\PYGZlt{}}\PYG{o}{\PYGZlt{}}\PYG{+w}{ }\PYG{n}{std}\PYG{o}{:}\PYG{o}{:}\PYG{n}{numeric\PYGZus{}limits}\PYG{o}{\PYGZlt{}}\PYG{n}{T}\PYG{o}{\PYGZgt{}}\PYG{o}{:}\PYG{o}{:}\PYG{n}{digits10}\PYG{+w}{     }\PYG{o}{\PYGZlt{}}\PYG{o}{\PYGZlt{}}\PYG{+w}{ }\PYG{n}{std}\PYG{o}{:}\PYG{o}{:}\PYG{n}{endl}
\PYG{+w}{    }\PYG{o}{\PYGZlt{}}\PYG{o}{\PYGZlt{}}\PYG{+w}{ }\PYG{l+s}{\PYGZdq{}}\PYG{l+s}{integral=}\PYG{l+s}{\PYGZdq{}}
\PYG{+w}{    }\PYG{o}{\PYGZlt{}}\PYG{o}{\PYGZlt{}}\PYG{+w}{ }\PYG{n}{std}\PYG{o}{:}\PYG{o}{:}\PYG{n}{numeric\PYGZus{}limits}\PYG{o}{\PYGZlt{}}\PYG{n}{T}\PYG{o}{\PYGZgt{}}\PYG{o}{:}\PYG{o}{:}\PYG{n}{is\PYGZus{}integer}\PYG{+w}{   }\PYG{o}{\PYGZlt{}}\PYG{o}{\PYGZlt{}}\PYG{+w}{ }\PYG{n}{std}\PYG{o}{:}\PYG{o}{:}\PYG{n}{endl}
\PYG{+w}{    }\PYG{o}{\PYGZlt{}}\PYG{o}{\PYGZlt{}}\PYG{+w}{ }\PYG{l+s}{\PYGZdq{}}\PYG{l+s}{signed=}\PYG{l+s}{\PYGZdq{}}
\PYG{+w}{    }\PYG{o}{\PYGZlt{}}\PYG{o}{\PYGZlt{}}\PYG{+w}{ }\PYG{n}{std}\PYG{o}{:}\PYG{o}{:}\PYG{n}{numeric\PYGZus{}limits}\PYG{o}{\PYGZlt{}}\PYG{n}{T}\PYG{o}{\PYGZgt{}}\PYG{o}{:}\PYG{o}{:}\PYG{n}{is\PYGZus{}signed}\PYG{+w}{    }\PYG{o}{\PYGZlt{}}\PYG{o}{\PYGZlt{}}\PYG{+w}{ }\PYG{n}{std}\PYG{o}{:}\PYG{o}{:}\PYG{n}{endl}
\PYG{+w}{    }\PYG{o}{\PYGZlt{}}\PYG{o}{\PYGZlt{}}\PYG{+w}{ }\PYG{l+s}{\PYGZdq{}}\PYG{l+s}{exact=}\PYG{l+s}{\PYGZdq{}}
\PYG{+w}{    }\PYG{o}{\PYGZlt{}}\PYG{o}{\PYGZlt{}}\PYG{+w}{ }\PYG{n}{std}\PYG{o}{:}\PYG{o}{:}\PYG{n}{numeric\PYGZus{}limits}\PYG{o}{\PYGZlt{}}\PYG{n}{T}\PYG{o}{\PYGZgt{}}\PYG{o}{:}\PYG{o}{:}\PYG{n}{is\PYGZus{}exact}\PYG{+w}{     }\PYG{o}{\PYGZlt{}}\PYG{o}{\PYGZlt{}}\PYG{+w}{ }\PYG{n}{std}\PYG{o}{:}\PYG{o}{:}\PYG{n}{endl}
\PYG{+w}{    }\PYG{o}{\PYGZlt{}}\PYG{o}{\PYGZlt{}}\PYG{+w}{ }\PYG{l+s}{\PYGZdq{}}\PYG{l+s}{infinity=}\PYG{l+s}{\PYGZdq{}}
\PYG{+w}{    }\PYG{o}{\PYGZlt{}}\PYG{o}{\PYGZlt{}}\PYG{+w}{ }\PYG{n}{std}\PYG{o}{:}\PYG{o}{:}\PYG{n}{numeric\PYGZus{}limits}\PYG{o}{\PYGZlt{}}\PYG{n}{T}\PYG{o}{\PYGZgt{}}\PYG{o}{:}\PYG{o}{:}\PYG{n}{has\PYGZus{}infinity}\PYG{+w}{ }\PYG{o}{\PYGZlt{}}\PYG{o}{\PYGZlt{}}\PYG{+w}{ }\PYG{n}{std}\PYG{o}{:}\PYG{o}{:}\PYG{n}{endl}\PYG{p}{;}
\PYG{p}{\PYGZcb{}}
\end{Verbatim}
\end{codeblockbg}


%%C020090 p

如果对无符号\texttt{short}、\texttt{int}和\texttt{double}调用\myttbr{print\_\allowbreak{}type\_\allowbreak{}properties()}函数，将得到如下输出：


%%C020091 table

\myTableBegin
\begin{xltabular}{\textwidth}{|>{\hsize=0.337\hsize\raggedright\arraybackslash}X|>{\hsize=0.321\hsize\raggedright\arraybackslash}X|>{\hsize=0.342\hsize\raggedright\arraybackslash}X|}
\hline
\rowcolor{tabhead} \textcolor{white}{\myttbr{unsigned short}} & \textcolor{white}{\myttbr{int}} & \textcolor{white}{\myttbr{double}} \\
\hline
\endfirsthead
\rowcolor{tabhead} \textcolor{white}{\myttbr{unsigned short}} & \textcolor{white}{\myttbr{int}} & \textcolor{white}{\myttbr{double}} \\
\hline
\endhead
min=0 max=65535 bits=16 decdigits=4 integral=1 signed=0 exact=1 infinity=0 & min=-2147483648 max=2147483647 bits=31 decdigits=9 integral=1 signed=1 exact=1 infinity=0 & min=2.22507e-308 max=1.79769e+308 bits=53 decdigits=15 integral=0 signed=1 exact=0 infinity=1 \\
\hline
\end{xltabular}
\myTableEnd


%%C020092 p

唯一需要注意的一点是\texttt{digits}与\texttt{digits10}两个常量之间的区别：


%%C020093 li

\begin{itemize}
\item 对整数类型，\texttt{digits}表示比特数（若存在符号位则不含符号位）再加上填充比特（如果有）；对浮点类型，它表示尾数的比特数。


%%C020094 li

\item \texttt{digits10}是该类型能够无损表示的十进制位数。为了更好地理解，我们来看\texttt{unsigned} \texttt{short}的情况。这是16位整数类型，可以表示0到


%%C020095 li

\item 它最多可以表示五位十进制数，即10,000到65,536，但并不能表示全部五位十进制数，因为65,537到99,999需要更多比特。因此，不需要更多比特就能表示的最大数字是四位十进制数（1,000到9,999）。这正是\texttt{digits10}所给出的值。对整数类型，它与常量\texttt{digits}有直接关系；对整数类型\texttt{T}，\texttt{digits10}的值为\myttbr{std::\allowbreak{}numeric\_\allowbreak{}limits<\allowbreak{}T>::\allowbreak{}digits * std::\allowbreak{}log10(\allowbreak{}2)}。

\end{itemize}


%%C020096 h2

\section{生成伪随机数}


%%C020097 p

生成随机数在许多应用中都必不可少，从游戏到密码学，从抽样到预测。不过，\textit{随机数}这一说法其实并不准确，因为用数学公式生成数字的过程是确定性的，产生的并不是真正的随机数，而是看起来随机的数，被称为\textit{伪随机数}。真正的随机性只能依靠基于物理过程的硬件设备来实现，而且即便如此也值得商榷，因为有人甚至认为宇宙本身也是确定性的。现代C++通过伪随机数库提供生成伪随机数的支持，该库包含数字生成器和分布。理论上它也能产生真正的随机数，但在实践中，这些也可能只是伪随机的。


%%C020098 h3

\subsection{准备工作}


%%C020099 p

本案例讨论标准库对生成伪随机数的支持。关键在于理解随机数与伪随机数的区别。另一方面，熟悉各种统计分布会有所帮助。但有一点是必须的：你要知道什么是均匀分布，因为库中所有引擎产生的数字都服从均匀分布。


%%C020100 h3

\subsection{具体做法}


%%C020101 p

要在应用中生成伪随机数，应当执行以下步骤：


%%C020102 p

包含头文件\texttt{<\allowbreak{}random>}：


%%C020103 code

\begin{codeblockbg}
\begin{Verbatim}[commandchars=\\\{\},breaklines=true,breakanywhere=true,fontsize=\codefont,xleftmargin=17bp]
\PYG{c+cp}{\PYGZsh{}}\PYG{c+cp}{include}\PYG{+w}{ }\PYG{c+cpf}{\PYGZlt{}random\PYGZgt{}}
\end{Verbatim}
\end{codeblockbg}


%%C020104 p

用\myttbr{std::\allowbreak{}random\_\allowbreak{}device}生成器为伪随机引擎播种：


%%C020105 code

\begin{codeblockbg}
\begin{Verbatim}[commandchars=\\\{\},breaklines=true,breakanywhere=true,fontsize=\codefont,xleftmargin=17bp]
\PYG{n}{std}\PYG{o}{:}\PYG{o}{:}\PYG{n}{random\PYGZus{}device}\PYG{+w}{ }\PYG{n}{rd}\PYG{p}{\PYGZob{}}\PYG{p}{\PYGZcb{}}\PYG{p}{;}
\end{Verbatim}
\end{codeblockbg}


%%C020106 p

从可用的引擎中选一个用来生成数字，并用一个


%%C020107 p

随机种子初始化它：


%%C020108 code

\begin{codeblockbg}
\begin{Verbatim}[commandchars=\\\{\},breaklines=true,breakanywhere=true,fontsize=\codefont,xleftmargin=17bp]
\PYG{k}{auto}\PYG{+w}{ }\PYG{n}{mtgen}\PYG{+w}{ }\PYG{o}{=}\PYG{+w}{ }\PYG{n}{std}\PYG{o}{:}\PYG{o}{:}\PYG{n}{mt19937}\PYG{p}{\PYGZob{}}\PYG{+w}{ }\PYG{n}{rd}\PYG{p}{(}\PYG{p}{)}\PYG{+w}{ }\PYG{p}{\PYGZcb{}}\PYG{p}{;}
\end{Verbatim}
\end{codeblockbg}


%%C020109 p

从可用的分布中选一个，把引擎的输出转换为


%%C020110 p

所需的某种统计分布：


%%C020111 code

\begin{codeblockbg}
\begin{Verbatim}[commandchars=\\\{\},breaklines=true,breakanywhere=true,fontsize=\codefont,xleftmargin=17bp]
\PYG{k}{auto}\PYG{+w}{ }\PYG{n}{ud}\PYG{+w}{ }\PYG{o}{=}\PYG{+w}{ }\PYG{n}{std}\PYG{o}{:}\PYG{o}{:}\PYG{n}{uniform\PYGZus{}int\PYGZus{}distribution}\PYG{o}{\PYGZlt{}}\PYG{o}{\PYGZgt{}}\PYG{p}{\PYGZob{}}\PYG{+w}{ }\PYG{l+m+mi}{1}\PYG{p}{,}\PYG{+w}{ }\PYG{l+m+mi}{6}\PYG{+w}{ }\PYG{p}{\PYGZcb{}}\PYG{p}{;}
\end{Verbatim}
\end{codeblockbg}


%%C020112 p

生成伪随机数：


%%C020113 code

\begin{codeblockbg}
\begin{Verbatim}[commandchars=\\\{\},breaklines=true,breakanywhere=true,fontsize=\codefont,xleftmargin=17bp]
\PYG{k}{for}\PYG{p}{(}\PYG{k}{auto}\PYG{+w}{ }\PYG{n}{i}\PYG{+w}{ }\PYG{o}{=}\PYG{+w}{ }\PYG{l+m+mi}{0}\PYG{p}{;}\PYG{+w}{ }\PYG{n}{i}\PYG{+w}{ }\PYG{o}{\PYGZlt{}}\PYG{+w}{ }\PYG{l+m+mi}{20}\PYG{p}{;}\PYG{+w}{ }\PYG{o}{+}\PYG{o}{+}\PYG{n}{i}\PYG{p}{)}
\PYG{+w}{  }\PYG{k}{auto}\PYG{+w}{ }\PYG{n}{number}\PYG{+w}{ }\PYG{o}{=}\PYG{+w}{ }\PYG{n}{ud}\PYG{p}{(}\PYG{n}{mtgen}\PYG{p}{)}\PYG{p}{;}
\end{Verbatim}
\end{codeblockbg}


%%C020114 h3

\subsection{工作原理}


%%C020115 p

伪随机数库包含两类组件：


%%C020116 li

\begin{itemize}
\item \textit{引擎}，即随机数生成器；它们可以产生服从均匀分布的伪随机数，或者在可用时产生真正的随机数。


%%C020117 li

\item \textit{分布}，把引擎的输出转换为某种统计分布。

\end{itemize}


%%C020118 p

所有引擎（\texttt{random\_\allowbreak{}device}除外）都产生服从均匀分布的整数，并且所有引擎都实现了以下方法：


%%C020119 li

\begin{itemize}
\item \texttt{min()}：静态方法，返回生成器能够产生的最小值。


%%C020120 li

\item \texttt{max()}：静态方法，返回生成器能够产生的最大值。


%%C020121 li

\item \texttt{seed()}：用一个起始值初始化算法（\texttt{random\_\allowbreak{}device}除外，它无法播种）。


%%C020122 li

\item \texttt{operator()}：生成一个在\texttt{min()}与\texttt{max()}之间均匀分布的新数字。


%%C020123 li

\item \texttt{discard()}：生成并丢弃指定数量的伪随机数。

\end{itemize}


%%C020124 p

可用的引擎如下：


%%C020125 li

\begin{itemize}
\item \myttbr{linear\_\allowbreak{}congruential\_\allowbreak{}engine}：线性同余生成器，使用以下公式产生数字：

\end{itemize}


%%C020126 p

\textit{x(i) = (A * x(i-1) + C) mod M}


%%C020127 li

\begin{itemize}
\item \myttbr{mersenne\_\allowbreak{}twister\_\allowbreak{}engine}：梅森旋转（Mersenne twister）生成器，它维护一个\textit{W * (N-1) * R}位的值；每次需要生成数字时，从中取出\textit{W}位。当所有比特都用完后，它通过移位和混合比特来扭转这个大值，从而得到一组新的可取出比特。


%%C020128 li

\item \myttbr{subtract\_\allowbreak{}with\_\allowbreak{}carry\_\allowbreak{}engine}：实现带进位减法（subtract with carry）算法的生成器，基于以下公式：

\end{itemize}


%%C020129 p

x(i) = (x(i - R) - x(i - S) - cy(i - 1)) mod M


%%C020130 p

上式中，\textit{cy}的定义为：


%%C020131 p

cy(i) = x(i - S) - x(i - R) - cy(i - 1)<0?1:0


%%C020132 p

此外，库还提供引擎适配器，它们本身也是引擎，包装另一个引擎，并基于基础引擎的输出产生数字。引擎适配器实现前面提到的与基础引擎相同的方法。可用的引擎适配器如下：


%%C020133 li

\begin{itemize}
\item \myttbr{discard\_\allowbreak{}block\_\allowbreak{}engine}：对基础引擎产生的每P个数字为一块，只保留其中R个，其余丢弃。


%%C020134 li

\item \myttbr{independent\_\allowbreak{}bits\_\allowbreak{}engine}：产生的数字位数与基础引擎不同的生成器。


%%C020135 li

\item \myttbr{shuffle\_\allowbreak{}order\_\allowbreak{}engine}：维护一张由基础引擎产生的K个数字组成的乱序表，返回表中的数字，并用基础引擎新生成的数字替换它们。

\end{itemize}


%%C020136 p

所有这些引擎和引擎适配器产生的都是伪随机数。不过，库还提供另一个名为\texttt{random\_\allowbreak{}device}的引擎，它理应产生非确定性数字，但这并不是硬性约束，因为随机的物理熵源未必可用。因此，\texttt{random\_\allowbreak{}device}的实现实际上可能基于伪随机引擎。\texttt{random\_\allowbreak{}device}类不能像其他引擎那样播种，并且有一个额外的方法\texttt{entropy()}，返回随机设备的熵：对确定性生成器为0，对非确定性生成器不为0。然而，这并不是判断设备究竟确定还是不确定的可靠方法。例如，GNU \texttt{libstdc++}和LLVM \texttt{libc++}都实现的是非确定性设备，但返回的熵为\texttt{0}。另一方面，\texttt{VC++}和\texttt{boost.\allowbreak{}random}返回的熵分别为\texttt{32}和\texttt{10}。


%%C020137 p

所有这些生成器产生的都是服从均匀分布的整数。然而，这只是大多数应用中所需的众多统计分布之一。为了能按其他分布产生数字（整数或实数），库提供了若干称为\textit{分布}的类，按各自实现的统计分布转换引擎的输出。可用的分布如下：


%%C020138 table

\myTableBegin
\begin{xltabular}{\textwidth}{|>{\hsize=0.114\hsize\raggedright\arraybackslash}X|>{\hsize=0.451\hsize\raggedright\arraybackslash}X|>{\hsize=0.115\hsize\raggedright\arraybackslash}X|>{\hsize=0.320\hsize\raggedright\arraybackslash}X|}
\hline
\rowcolor{tabhead} \textcolor{white}{\textbf{类型}} & \textcolor{white}{\textbf{类名}} & \textcolor{white}{\textbf{数字}} & \textcolor{white}{\textbf{统计分布}} \\
\hline
\endfirsthead
\rowcolor{tabhead} \textcolor{white}{\textbf{类型}} & \textcolor{white}{\textbf{类名}} & \textcolor{white}{\textbf{数字}} & \textcolor{white}{\textbf{统计分布}} \\
\hline
\endhead
均匀 & \myttbr{uniform\_\allowbreak{}int\_\allowbreak{}distribution} & 整数 & 均匀\\
\hline
 & \myttbr{uniform\_\allowbreak{}real\_\allowbreak{}distribution} & 实数 & 均匀\\
\hline
伯努利 & \myttbr{bernoulli\_\allowbreak{}distribution} & 布尔 & 伯努利\\
\hline
 & \myttbr{binomial\_\allowbreak{}distribution} & 整数 & 二项\\
\hline
 & \myttbr{negative\_\allowbreak{}binomial\_\allowbreak{}distribution} & 整数 & 负二项\\
\hline
 & \myttbr{geometric\_\allowbreak{}distribution} & 整数 & 几何\\
\hline
泊松 & \myttbr{poisson\_\allowbreak{}distribution} & 整数 & 泊松\\
\hline
 & \myttbr{exponential\_\allowbreak{}distribution} & 实数 & 指数\\
\hline
 & \myttbr{gamma\_\allowbreak{}distribution} & 实数 & 伽马\\
\hline
 & \myttbr{weibull\_\allowbreak{}distribution} & 实数 & 威布尔\\
\hline
 & \myttbr{extreme\_\allowbreak{}value\_\allowbreak{}distribution} & 实数 & 极值\\
\hline
正态 & \myttbr{normal\_\allowbreak{}distribution} & 实数 & 标准正态（高斯）\\
\hline
 & \myttbr{lognormal\_\allowbreak{}distribution} & 实数 & 对数正态\\
\hline
 & \myttbr{chi\_\allowbreak{}squared\_\allowbreak{}distribution} & 实数 & 卡方\\
\hline
 & \myttbr{cauchy\_\allowbreak{}distribution} & 实数 & 柯西\\
\hline
 & \myttbr{fisher\_\allowbreak{}f\_\allowbreak{}distribution} & 实数 & Fisher F分布\\
\hline
 & \myttbr{student\_\allowbreak{}t\_\allowbreak{}distribution} & 实数 & Student t分布\\
\hline
抽样 & \myttbr{discrete\_\allowbreak{}distribution} & 整数 & 离散\\
\hline
 & \myttbr{piecewise\_\allowbreak{}constant\_\allowbreak{}distribution} & 实数 & 在常量子区间上分布的值\\
\hline
 & \myttbr{piecewise\_\allowbreak{}linear\_\allowbreak{}distribution} & 实数 & 在给定子区间上分布的值\\
\hline
\end{xltabular}
\myTableEnd


%%C020139 p

库提供的每种引擎都各有优缺点。线性同余引擎内部状态小，但速度不算快。另一方面，带进位减法引擎非常快，但内部状态需要更多内存。梅森旋转是其中最慢的，也是内部状态最大的，不过若初始化得当，它能产生最长的非重复数字序列。下面的示例中我们使用\texttt{std::\allowbreak{}mt19937}，这是一个32位的梅森旋转引擎，内部状态为19,937比特。


%%C020140 p

生成随机数最简单的方式如下：


%%C020141 code

\begin{codeblockbg}
\begin{Verbatim}[commandchars=\\\{\},breaklines=true,breakanywhere=true,fontsize=\codefont,xleftmargin=17bp]
\PYG{k}{auto}\PYG{+w}{ }\PYG{n}{mtgen}\PYG{+w}{ }\PYG{o}{=}\PYG{+w}{ }\PYG{n}{std}\PYG{o}{:}\PYG{o}{:}\PYG{n}{mt19937}\PYG{+w}{ }\PYG{p}{\PYGZob{}}\PYG{p}{\PYGZcb{}}\PYG{p}{;}
\PYG{k}{for}\PYG{+w}{ }\PYG{p}{(}\PYG{k}{auto}\PYG{+w}{ }\PYG{n}{i}\PYG{+w}{ }\PYG{o}{=}\PYG{+w}{ }\PYG{l+m+mi}{0}\PYG{p}{;}\PYG{+w}{ }\PYG{n}{i}\PYG{+w}{ }\PYG{o}{\PYGZlt{}}\PYG{+w}{ }\PYG{l+m+mi}{10}\PYG{p}{;}\PYG{+w}{ }\PYG{o}{+}\PYG{o}{+}\PYG{n}{i}\PYG{p}{)}
\PYG{+w}{  }\PYG{n}{std}\PYG{o}{:}\PYG{o}{:}\PYG{n}{cout}\PYG{+w}{ }\PYG{o}{\PYGZlt{}}\PYG{o}{\PYGZlt{}}\PYG{+w}{ }\PYG{n}{mtgen}\PYG{p}{(}\PYG{p}{)}\PYG{+w}{ }\PYG{o}{\PYGZlt{}}\PYG{o}{\PYGZlt{}}\PYG{+w}{ }\PYG{n}{std}\PYG{o}{:}\PYG{o}{:}\PYG{n}{endl}\PYG{p}{;}
\end{Verbatim}
\end{codeblockbg}


%%C020142 p

在这个例子中，\texttt{mtgen}是一个\texttt{std::\allowbreak{}mt19937}梅森旋转引擎。要生成数字，只需使用调用运算符，它会推进内部状态并返回下一个伪随机数。然而这段代码有缺陷，因为引擎没有播种。结果是它总是产生相同的数字序列，在大多数情况下这大概并不是你想要的结果。


%%C020143 p

初始化引擎有多种做法。一种常见于C的rand库的做法是使用当前时间。在现代C++中，它应当写成这样：


%%C020144 code

\begin{codeblockbg}
\begin{Verbatim}[commandchars=\\\{\},breaklines=true,breakanywhere=true,fontsize=\codefont,xleftmargin=17bp]
\PYG{k}{auto}\PYG{+w}{ }\PYG{n}{seed}\PYG{+w}{ }\PYG{o}{=}\PYG{+w}{ }\PYG{n}{std}\PYG{o}{:}\PYG{o}{:}\PYG{n}{chrono}\PYG{o}{:}\PYG{o}{:}\PYG{n}{high\PYGZus{}resolution\PYGZus{}clock}\PYG{o}{:}\PYG{o}{:}\PYG{n}{now}\PYG{p}{(}\PYG{p}{)}
\PYG{+w}{            }\PYG{p}{.}\PYG{n}{time\PYGZus{}since\PYGZus{}epoch}\PYG{p}{(}\PYG{p}{)}
\PYG{+w}{            }\PYG{p}{.}\PYG{n}{count}\PYG{p}{(}\PYG{p}{)}\PYG{p}{;}
\PYG{k}{auto}\PYG{+w}{ }\PYG{n}{mtgen}\PYG{+w}{ }\PYG{o}{=}\PYG{+w}{ }\PYG{n}{std}\PYG{o}{:}\PYG{o}{:}\PYG{n}{mt19937}\PYG{p}{\PYGZob{}}\PYG{+w}{ }\PYG{k}{static\PYGZus{}cast}\PYG{o}{\PYGZlt{}}\PYG{k+kt}{unsigned}\PYG{+w}{ }\PYG{k+kt}{int}\PYG{o}{\PYGZgt{}}\PYG{p}{(}\PYG{n}{seed}\PYG{p}{)}\PYG{+w}{ }\PYG{p}{\PYGZcb{}}\PYG{p}{;}
\end{Verbatim}
\end{codeblockbg}


%%C020145 p

在这个例子中，\texttt{seed}是一个数字，表示从时钟纪元到当前时刻的计时周期数。然后用这个数字为引擎播种。这种做法的问题在于，\texttt{seed}的值实际上是确定性的，在某些类型的应用中可能容易被攻击。更可靠的做法是用真正的随机数为生成器播种。\myttbr{std::\allowbreak{}random\_\allowbreak{}device}类是一个理应返回真正随机数的引擎，不过其实现实际上可能基于伪随机生成器：


%%C020146 code

\begin{codeblockbg}
\begin{Verbatim}[commandchars=\\\{\},breaklines=true,breakanywhere=true,fontsize=\codefont,xleftmargin=17bp]
\PYG{n}{std}\PYG{o}{:}\PYG{o}{:}\PYG{n}{random\PYGZus{}device}\PYG{+w}{ }\PYG{n}{rd}\PYG{p}{;}
\PYG{k}{auto}\PYG{+w}{ }\PYG{n}{mtgen}\PYG{+w}{ }\PYG{o}{=}\PYG{+w}{ }\PYG{n}{std}\PYG{o}{:}\PYG{o}{:}\PYG{n}{mt19937}\PYG{+w}{ }\PYG{p}{\PYGZob{}}\PYG{n}{rd}\PYG{p}{(}\PYG{p}{)}\PYG{p}{\PYGZcb{}}\PYG{p}{;}
\end{Verbatim}
\end{codeblockbg}


%%C020147 p

所有引擎产生的数字都服从均匀分布。要把结果转换为其他统计分布，必须使用分布类。为了展示生成的数字如何按所选分布分布，我们使用下面的函数。该函数生成指定数量的伪随机数，并在map中统计它们的重复次数。然后用map中的值绘制类似条形图的图形，显示每个数字出现的频率：


%%C020148 code

\begin{codeblockbg}
\begin{Verbatim}[commandchars=\\\{\},breaklines=true,breakanywhere=true,fontsize=\codefont,xleftmargin=17bp]
\PYG{k+kt}{void}\PYG{+w}{ }\PYG{n+nf}{generate\PYGZus{}and\PYGZus{}print}\PYG{p}{(}
\PYG{+w}{  }\PYG{n}{std}\PYG{o}{:}\PYG{o}{:}\PYG{n}{function}\PYG{o}{\PYGZlt{}}\PYG{k+kt}{int}\PYG{p}{(}\PYG{k+kt}{void}\PYG{p}{)}\PYG{o}{\PYGZgt{}}\PYG{+w}{ }\PYG{n}{gen}\PYG{p}{,}
\PYG{+w}{  }\PYG{k+kt}{int}\PYG{+w}{ }\PYG{k}{const}\PYG{+w}{ }\PYG{n}{iterations}\PYG{+w}{ }\PYG{o}{=}\PYG{+w}{ }\PYG{l+m+mi}{10000}\PYG{p}{)}
\PYG{p}{\PYGZob{}}
\PYG{+w}{  }\PYG{c+c1}{// 用 map 保存数字及其重复次数}
\PYG{+w}{  }\PYG{k}{auto}\PYG{+w}{ }\PYG{n}{data}\PYG{+w}{ }\PYG{o}{=}\PYG{+w}{ }\PYG{n}{std}\PYG{o}{:}\PYG{o}{:}\PYG{n}{map}\PYG{o}{\PYGZlt{}}\PYG{k+kt}{int}\PYG{p}{,}\PYG{+w}{ }\PYG{k+kt}{int}\PYG{o}{\PYGZgt{}}\PYG{p}{\PYGZob{}}\PYG{p}{\PYGZcb{}}\PYG{p}{;}
\PYG{c+c1}{// 生成随机数}
\PYG{k}{for}\PYG{+w}{ }\PYG{p}{(}\PYG{k}{auto}\PYG{+w}{ }\PYG{n}{n}\PYG{+w}{ }\PYG{o}{=}\PYG{+w}{ }\PYG{l+m+mi}{0}\PYG{p}{;}\PYG{+w}{ }\PYG{n}{n}\PYG{+w}{ }\PYG{o}{\PYGZlt{}}\PYG{+w}{ }\PYG{n}{iterations}\PYG{p}{;}\PYG{+w}{ }\PYG{o}{+}\PYG{o}{+}\PYG{n}{n}\PYG{p}{)}
\PYG{+w}{  }\PYG{o}{+}\PYG{o}{+}\PYG{n}{data}\PYG{p}{[}\PYG{n}{gen}\PYG{p}{(}\PYG{p}{)}\PYG{p}{]}\PYG{p}{;}
\PYG{c+c1}{// 找出重复次数最多的元素}
\PYG{k}{auto}\PYG{+w}{ }\PYG{n}{max}\PYG{+w}{ }\PYG{o}{=}\PYG{+w}{ }\PYG{n}{std}\PYG{o}{:}\PYG{o}{:}\PYG{n}{max\PYGZus{}element}\PYG{p}{(}
\PYG{+w}{           }\PYG{n}{std}\PYG{o}{:}\PYG{o}{:}\PYG{n}{begin}\PYG{p}{(}\PYG{n}{data}\PYG{p}{)}\PYG{p}{,}\PYG{+w}{ }\PYG{n}{std}\PYG{o}{:}\PYG{o}{:}\PYG{n}{end}\PYG{p}{(}\PYG{n}{data}\PYG{p}{)}\PYG{p}{,}
\PYG{+w}{           }\PYG{p}{[}\PYG{p}{]}\PYG{p}{(}\PYG{k}{auto}\PYG{+w}{ }\PYG{n}{kvp1}\PYG{p}{,}\PYG{+w}{ }\PYG{k}{auto}\PYG{+w}{ }\PYG{n}{kvp2}\PYG{p}{)}\PYG{+w}{ }\PYG{p}{\PYGZob{}}
\PYG{+w}{  }\PYG{k}{return}\PYG{+w}{ }\PYG{n}{kvp1}\PYG{p}{.}\PYG{n}{second}\PYG{+w}{ }\PYG{o}{\PYGZlt{}}\PYG{+w}{ }\PYG{n}{kvp2}\PYG{p}{.}\PYG{n}{second}\PYG{p}{;}\PYG{+w}{ }\PYG{p}{\PYGZcb{}}\PYG{p}{)}\PYG{p}{;}
\PYG{c+c1}{// 打印条形}
\PYG{k}{for}\PYG{+w}{ }\PYG{p}{(}\PYG{k}{auto}\PYG{+w}{ }\PYG{n}{i}\PYG{+w}{ }\PYG{o}{=}\PYG{+w}{ }\PYG{n}{max}\PYG{o}{\PYGZhy{}}\PYG{o}{\PYGZgt{}}\PYG{n}{second}\PYG{+w}{ }\PYG{o}{/}\PYG{+w}{ }\PYG{l+m+mi}{200}\PYG{p}{;}\PYG{+w}{ }\PYG{n}{i}\PYG{+w}{ }\PYG{o}{\PYGZgt{}}\PYG{+w}{ }\PYG{l+m+mi}{0}\PYG{p}{;}\PYG{+w}{ }\PYG{o}{\PYGZhy{}}\PYG{o}{\PYGZhy{}}\PYG{n}{i}\PYG{p}{)}
\PYG{p}{\PYGZob{}}
\PYG{+w}{  }\PYG{k}{for}\PYG{+w}{ }\PYG{p}{(}\PYG{k}{auto}\PYG{+w}{ }\PYG{n}{kvp}\PYG{+w}{ }\PYG{o}{:}\PYG{+w}{ }\PYG{n}{data}\PYG{p}{)}
\PYG{+w}{  }\PYG{p}{\PYGZob{}}
\PYG{+w}{    }\PYG{n}{std}\PYG{o}{:}\PYG{o}{:}\PYG{n}{cout}
\PYG{+w}{      }\PYG{o}{\PYGZlt{}}\PYG{o}{\PYGZlt{}}\PYG{+w}{ }\PYG{n}{std}\PYG{o}{:}\PYG{o}{:}\PYG{n}{fixed}\PYG{+w}{ }\PYG{o}{\PYGZlt{}}\PYG{o}{\PYGZlt{}}\PYG{+w}{ }\PYG{n}{std}\PYG{o}{:}\PYG{o}{:}\PYG{n}{setprecision}\PYG{p}{(}\PYG{l+m+mi}{1}\PYG{p}{)}\PYG{+w}{ }\PYG{o}{\PYGZlt{}}\PYG{o}{\PYGZlt{}}\PYG{+w}{ }\PYG{n}{std}\PYG{o}{:}\PYG{o}{:}\PYG{n}{setw}\PYG{p}{(}\PYG{l+m+mi}{3}\PYG{p}{)}
\PYG{+w}{      }\PYG{o}{\PYGZlt{}}\PYG{o}{\PYGZlt{}}\PYG{+w}{ }\PYG{p}{(}\PYG{n}{kvp}\PYG{p}{.}\PYG{n}{second}\PYG{+w}{ }\PYG{o}{/}\PYG{+w}{ }\PYG{l+m+mi}{200}\PYG{+w}{ }\PYG{o}{\PYGZgt{}}\PYG{o}{=}\PYG{+w}{ }\PYG{n}{i}\PYG{+w}{ }\PYG{o}{?}\PYG{+w}{ }\PYG{p}{(}\PYG{k+kt}{char}\PYG{p}{)}\PYG{l+m+mi}{219}\PYG{+w}{ }\PYG{o}{:}\PYG{+w}{ }\PYG{l+s+sc}{\PYGZsq{}}\PYG{l+s+sc}{ }\PYG{l+s+sc}{\PYGZsq{}}\PYG{p}{)}\PYG{p}{;}
\PYG{+w}{  }\PYG{p}{\PYGZcb{}}
\PYG{+w}{  }\PYG{n}{std}\PYG{o}{:}\PYG{o}{:}\PYG{n}{cout}\PYG{+w}{ }\PYG{o}{\PYGZlt{}}\PYG{o}{\PYGZlt{}}\PYG{+w}{ }\PYG{n}{std}\PYG{o}{:}\PYG{o}{:}\PYG{n}{endl}\PYG{p}{;}
\PYG{p}{\PYGZcb{}}
\PYG{c+c1}{// 打印数字}
\PYG{k}{for}\PYG{+w}{ }\PYG{p}{(}\PYG{k}{auto}\PYG{+w}{ }\PYG{n}{kvp}\PYG{+w}{ }\PYG{o}{:}\PYG{+w}{ }\PYG{n}{data}\PYG{p}{)}
\PYG{p}{\PYGZob{}}
\PYG{+w}{  }\PYG{n}{std}\PYG{o}{:}\PYG{o}{:}\PYG{n}{cout}
\PYG{+w}{    }\PYG{o}{\PYGZlt{}}\PYG{o}{\PYGZlt{}}\PYG{+w}{ }\PYG{n}{std}\PYG{o}{:}\PYG{o}{:}\PYG{n}{fixed}\PYG{+w}{ }\PYG{o}{\PYGZlt{}}\PYG{o}{\PYGZlt{}}\PYG{+w}{ }\PYG{n}{std}\PYG{o}{:}\PYG{o}{:}\PYG{n}{setprecision}\PYG{p}{(}\PYG{l+m+mi}{1}\PYG{p}{)}\PYG{+w}{ }\PYG{o}{\PYGZlt{}}\PYG{o}{\PYGZlt{}}\PYG{+w}{ }\PYG{n}{std}\PYG{o}{:}\PYG{o}{:}\PYG{n}{setw}\PYG{p}{(}\PYG{l+m+mi}{3}\PYG{p}{)}
\PYG{+w}{    }\PYG{o}{\PYGZlt{}}\PYG{o}{\PYGZlt{}}\PYG{+w}{ }\PYG{n}{kvp}\PYG{p}{.}\PYG{n}{first}\PYG{p}{;}
\PYG{p}{\PYGZcb{}}
\PYG{+w}{  }\PYG{n}{std}\PYG{o}{:}\PYG{o}{:}\PYG{n}{cout}\PYG{+w}{ }\PYG{o}{\PYGZlt{}}\PYG{o}{\PYGZlt{}}\PYG{+w}{ }\PYG{n}{std}\PYG{o}{:}\PYG{o}{:}\PYG{n}{endl}\PYG{p}{;}
\PYG{p}{\PYGZcb{}}
\end{Verbatim}
\end{codeblockbg}


%%C020149 p

下面的代码使用\texttt{std::\allowbreak{}mt19937}引擎生成\texttt{[1, 6]}区间内均匀分布的随机数，这基本上就是掷骰子所得的结果：


%%C020150 code

\begin{codeblockbg}
\begin{Verbatim}[commandchars=\\\{\},breaklines=true,breakanywhere=true,fontsize=\codefont,xleftmargin=17bp]
\PYG{n}{std}\PYG{o}{:}\PYG{o}{:}\PYG{n}{random\PYGZus{}device}\PYG{+w}{ }\PYG{n}{rd}\PYG{p}{\PYGZob{}}\PYG{p}{\PYGZcb{}}\PYG{p}{;}
\PYG{k}{auto}\PYG{+w}{ }\PYG{n}{mtgen}\PYG{+w}{ }\PYG{o}{=}\PYG{+w}{ }\PYG{n}{std}\PYG{o}{:}\PYG{o}{:}\PYG{n}{mt19937}\PYG{p}{\PYGZob{}}\PYG{+w}{ }\PYG{n}{rd}\PYG{p}{(}\PYG{p}{)}\PYG{+w}{ }\PYG{p}{\PYGZcb{}}\PYG{p}{;}
\PYG{k}{auto}\PYG{+w}{ }\PYG{n}{ud}\PYG{+w}{ }\PYG{o}{=}\PYG{+w}{ }\PYG{n}{std}\PYG{o}{:}\PYG{o}{:}\PYG{n}{uniform\PYGZus{}int\PYGZus{}distribution}\PYG{o}{\PYGZlt{}}\PYG{o}{\PYGZgt{}}\PYG{p}{\PYGZob{}}\PYG{+w}{ }\PYG{l+m+mi}{1}\PYG{p}{,}\PYG{+w}{ }\PYG{l+m+mi}{6}\PYG{+w}{ }\PYG{p}{\PYGZcb{}}\PYG{p}{;}
\PYG{n}{generate\PYGZus{}and\PYGZus{}print}\PYG{p}{(}\PYG{p}{[}\PYG{o}{\PYGZam{}}\PYG{n}{mtgen}\PYG{p}{,}\PYG{+w}{ }\PYG{o}{\PYGZam{}}\PYG{n}{ud}\PYG{p}{]}\PYG{p}{(}\PYG{p}{)}\PYG{+w}{ }\PYG{p}{\PYGZob{}}\PYG{k}{return}\PYG{+w}{ }\PYG{n}{ud}\PYG{p}{(}\PYG{n}{mtgen}\PYG{p}{)}\PYG{p}{;}\PYG{+w}{ }\PYG{p}{\PYGZcb{}}\PYG{p}{)}\PYG{p}{;}
\end{Verbatim}
\end{codeblockbg}


%%C020151 p

程序的输出如下：


%%C020152 img

\myFigure{images/ch02/fig-2-1.png}{}


%%C020153 p

在下一个也是最后一个例子中，我们把分布改为均值为\texttt{5}、标准差为\texttt{2}的正态分布。该分布产生实数，因此为了使用前面的\myttbr{generate\_\allowbreak{}and\_\allowbreak{}print()}函数，必须把数字四舍五入为整数：


%%C020154 code

\begin{codeblockbg}
\begin{Verbatim}[commandchars=\\\{\},breaklines=true,breakanywhere=true,fontsize=\codefont,xleftmargin=17bp]
\PYG{n}{std}\PYG{o}{:}\PYG{o}{:}\PYG{n}{random\PYGZus{}device}\PYG{+w}{ }\PYG{n}{rd}\PYG{p}{\PYGZob{}}\PYG{p}{\PYGZcb{}}\PYG{p}{;}
\PYG{k}{auto}\PYG{+w}{ }\PYG{n}{mtgen}\PYG{+w}{ }\PYG{o}{=}\PYG{+w}{ }\PYG{n}{std}\PYG{o}{:}\PYG{o}{:}\PYG{n}{mt19937}\PYG{p}{\PYGZob{}}\PYG{+w}{ }\PYG{n}{rd}\PYG{p}{(}\PYG{p}{)}\PYG{+w}{ }\PYG{p}{\PYGZcb{}}\PYG{p}{;}
\PYG{k}{auto}\PYG{+w}{ }\PYG{n}{nd}\PYG{+w}{ }\PYG{o}{=}\PYG{+w}{ }\PYG{n}{std}\PYG{o}{:}\PYG{o}{:}\PYG{n}{normal\PYGZus{}distribution}\PYG{o}{\PYGZlt{}}\PYG{o}{\PYGZgt{}}\PYG{p}{\PYGZob{}}\PYG{+w}{ }\PYG{l+m+mi}{5}\PYG{p}{,}\PYG{+w}{ }\PYG{l+m+mi}{2}\PYG{+w}{ }\PYG{p}{\PYGZcb{}}\PYG{p}{;}
\PYG{n}{generate\PYGZus{}and\PYGZus{}print}\PYG{p}{(}
\PYG{+w}{  }\PYG{p}{[}\PYG{o}{\PYGZam{}}\PYG{n}{mtgen}\PYG{p}{,}\PYG{+w}{ }\PYG{o}{\PYGZam{}}\PYG{n}{nd}\PYG{p}{]}\PYG{p}{(}\PYG{p}{)}\PYG{+w}{ }\PYG{p}{\PYGZob{}}
\PYG{+w}{    }\PYG{k}{return}\PYG{+w}{ }\PYG{k}{static\PYGZus{}cast}\PYG{o}{\PYGZlt{}}\PYG{k+kt}{int}\PYG{o}{\PYGZgt{}}\PYG{p}{(}\PYG{n}{std}\PYG{o}{:}\PYG{o}{:}\PYG{n}{round}\PYG{p}{(}\PYG{n}{nd}\PYG{p}{(}\PYG{n}{mtgen}\PYG{p}{)}\PYG{p}{)}\PYG{p}{)}\PYG{p}{;}\PYG{+w}{ }\PYG{p}{\PYGZcb{}}\PYG{p}{)}\PYG{p}{;}
\end{Verbatim}
\end{codeblockbg}


%%C020155 p

前面代码的输出将如下：


%%C020156 img

\myFigure{images/ch02/fig-2-2.png}{}


%%C020157 h3

\subsection{延伸阅读}


%%C020158 li

\begin{itemize}
\item \textit{初始化伪随机数生成器内部状态的全部比特}

\end{itemize}


%%C020159 h2

\section{初始化伪随机数生成器内部状态的全部比特}


%%C020160 p

上一个案例中，我们了解了伪随机数库及其组件，以及如何用它按不同统计分布产生数字。那个案例忽略了一个重要因素：伪随机数生成器的正确初始化。本案例将学习如何初始化生成器，以产生最佳的伪随机数序列。


%%C020161 h3

\subsection{准备工作}


%%C020162 p

你应当先阅读上一个案例\textit{生成伪随机数}，以了解伪随机数库提供了哪些内容。


%%C020163 h3

\subsection{具体做法}


%%C020164 p

要正确初始化伪随机数生成器以产生最佳的伪随机数序列，请执行以下步骤：


%%C020165 p

使用\myttbr{std::\allowbreak{}random\_\allowbreak{}device}产生用作播种


%%C020166 p

值的随机数：


%%C020167 code

\begin{codeblockbg}
\begin{Verbatim}[commandchars=\\\{\},breaklines=true,breakanywhere=true,fontsize=\codefont,xleftmargin=17bp]
\PYG{n}{std}\PYG{o}{:}\PYG{o}{:}\PYG{n}{random\PYGZus{}device}\PYG{+w}{ }\PYG{n}{rd}\PYG{p}{;}
\end{Verbatim}
\end{codeblockbg}


%%C020168 p

为引擎的全部内部比特生成随机数据：


%%C020169 code

\begin{codeblockbg}
\begin{Verbatim}[commandchars=\\\{\},breaklines=true,breakanywhere=true,fontsize=\codefont,xleftmargin=17bp]
\PYG{n}{std}\PYG{o}{:}\PYG{o}{:}\PYG{n}{array}\PYG{o}{\PYGZlt{}}\PYG{k+kt}{int}\PYG{p}{,}\PYG{+w}{ }\PYG{n}{std}\PYG{o}{:}\PYG{o}{:}\PYG{n}{mt19937}\PYG{o}{:}\PYG{o}{:}\PYG{n}{state\PYGZus{}size}\PYG{o}{\PYGZgt{}}\PYG{+w}{ }\PYG{n}{seed\PYGZus{}data}\PYG{+w}{ }\PYG{p}{\PYGZob{}}\PYG{p}{\PYGZcb{}}\PYG{p}{;}
\PYG{n}{std}\PYG{o}{:}\PYG{o}{:}\PYG{n}{generate}\PYG{p}{(}\PYG{n}{std}\PYG{o}{:}\PYG{o}{:}\PYG{n}{begin}\PYG{p}{(}\PYG{n}{seed\PYGZus{}data}\PYG{p}{)}\PYG{p}{,}\PYG{+w}{ }\PYG{n}{std}\PYG{o}{:}\PYG{o}{:}\PYG{n}{end}\PYG{p}{(}\PYG{n}{seed\PYGZus{}data}\PYG{p}{)}\PYG{p}{,}
\PYG{+w}{              }\PYG{n}{std}\PYG{o}{:}\PYG{o}{:}\PYG{n}{ref}\PYG{p}{(}\PYG{n}{rd}\PYG{p}{)}\PYG{p}{)}\PYG{p}{;}
\end{Verbatim}
\end{codeblockbg}


%%C020170 p

用前面生成的伪随机


%%C020171 p

数据创建\texttt{std::\allowbreak{}seed\_\allowbreak{}seq}对象：


%%C020172 code

\begin{codeblockbg}
\begin{Verbatim}[commandchars=\\\{\},breaklines=true,breakanywhere=true,fontsize=\codefont,xleftmargin=17bp]
\PYG{n}{std}\PYG{o}{:}\PYG{o}{:}\PYG{n}{seed\PYGZus{}seq}\PYG{+w}{ }\PYG{n+nf}{seq}\PYG{p}{(}\PYG{n}{std}\PYG{o}{:}\PYG{o}{:}\PYG{n}{begin}\PYG{p}{(}\PYG{n}{seed\PYGZus{}data}\PYG{p}{)}\PYG{p}{,}\PYG{+w}{ }\PYG{n}{std}\PYG{o}{:}\PYG{o}{:}\PYG{n}{end}\PYG{p}{(}\PYG{n}{seed\PYGZus{}data}\PYG{p}{)}\PYG{p}{)}\PYG{p}{;}
\end{Verbatim}
\end{codeblockbg}


%%C020173 p

创建引擎对象并初始化表示引擎内部状态的


%%C020174 p

全部比特；例如\texttt{mt19937}有19,937比特的内部状态：


%%C020175 code

\begin{codeblockbg}
\begin{Verbatim}[commandchars=\\\{\},breaklines=true,breakanywhere=true,fontsize=\codefont,xleftmargin=17bp]
\PYG{k}{auto}\PYG{+w}{ }\PYG{n}{eng}\PYG{+w}{ }\PYG{o}{=}\PYG{+w}{ }\PYG{n}{std}\PYG{o}{:}\PYG{o}{:}\PYG{n}{mt19937}\PYG{p}{\PYGZob{}}\PYG{+w}{ }\PYG{n}{seq}\PYG{+w}{ }\PYG{p}{\PYGZcb{}}\PYG{p}{;}
\end{Verbatim}
\end{codeblockbg}

%%C020176 p

根据应用的需求选用合适的分布：


%%C020177 code

\begin{codeblockbg}
\begin{Verbatim}[commandchars=\\\{\},breaklines=true,breakanywhere=true,fontsize=\codefont,xleftmargin=17bp]
\PYG{k}{auto}\PYG{+w}{ }\PYG{n}{dist}\PYG{+w}{ }\PYG{o}{=}\PYG{+w}{ }\PYG{n}{std}\PYG{o}{:}\PYG{o}{:}\PYG{n}{uniform\PYGZus{}real\PYGZus{}distribution}\PYG{o}{\PYGZlt{}}\PYG{o}{\PYGZgt{}}\PYG{p}{\PYGZob{}}\PYG{+w}{ }\PYG{l+m+mi}{0}\PYG{p}{,}\PYG{+w}{ }\PYG{l+m+mi}{1}\PYG{+w}{ }\PYG{p}{\PYGZcb{}}\PYG{p}{;}
\end{Verbatim}
\end{codeblockbg}


%%C020178 h3

\subsection{工作原理}


%%C020179 p

上一个案例中的所有示例都用\texttt{std::\allowbreak{}mt19937}引擎生成伪随机数。虽然梅森旋转算法（Mersenne twister）比其他引擎慢，但它能产生最长的非重复序列，且频谱特性最好。不过，按上一个案例中的方式初始化引擎并不能获得这种效果。经过仔细分析（这超出了本案例乃至本书的范围）可以证明，该引擎倾向于反复生成某些值而遗漏另一些值，于是产生的数不服从均匀分布，而更接近二项分布或泊松分布。问题在于\texttt{mt19937}的内部状态有624个32位整数，而上一案例的示例只初始化了其中一个。


%%C020180 p

使用伪随机数库时，请记住下面这条经验法则（见信息框）：


%%C020181 note

\begin{myNote}
为了获得最佳结果，引擎在生成数字之前必须把全部内部状态正确初始化。
\end{myNote}


%%C020182 p

伪随机数库为此专门提供了一个类，叫作\texttt{std::\allowbreak{}seed\_\allowbreak{}seq}。它是一个生成器，可以用任意多个32位整数作种子，并产生指定数量的、均匀分布在32位空间中的整数。


%%C020183 p

在\textit{具体做法}一节的代码中，我们定义了一个名为\texttt{seed\_\allowbreak{}data}的数组，其32位整数个数与\texttt{mt19937}生成器的内部状态相同，即624个整数。然后用\myttbr{std::\allowbreak{}random\_\allowbreak{}device}产生的随机数初始化该数组。之后用这个数组为\texttt{std::\allowbreak{}seed\_\allowbreak{}seq}提供种子，再用它为\texttt{mt19937}生成器提供种子。


%%C020184 h2

\section{创建cooked用户定义字面量}


%%C020185 p

字面量是内置类型（数值、布尔、字符、字符串和指针）的常量，在程序中不能修改。语言定义了一系列前缀和后缀来修饰字面量（前缀和后缀实际上是字面量的一部分）。C++11允许通过定义称为\textit{字面量运算符}（literal operator）的函数来创建用户定义字面量，这些函数为字面量引入后缀。它们只适用于数值、字符和字符串类型。这既为将来版本定义标准字面量提供了可能，也让开发者可以创建自己的字面量。本案例将介绍如何创建自己的cooked字面量。


%%C020186 h3

\subsection{准备工作}


%%C020187 p

用户定义字面量有两种形式：\textit{raw}和\textit{cooked}。raw字面量不经过编译器处理，而cooked字面量是经编译器处理后的值（例如处理字符串中的转义序列，或者从字面量0xBAD识别出整数2898）。raw字面量只适用于整型和浮点类型，而cooked字面量还适用于字符和字符串字面量。


%%C020188 h3

\subsection{具体做法}


%%C020189 p

创建cooked用户定义字面量应遵循以下步骤：


%%C020190 p

把字面量定义在独立的命名空间中，以避免名称冲突。


%%C020191 p

用户定义后缀一律以下划线（\texttt{\_}）开头。


%%C020192 p

为cooked字面量定义如下形式的字面量运算符：


%%C020193 code

\begin{codeblockbg}
\begin{Verbatim}[commandchars=\\\{\},breaklines=true,breakanywhere=true,fontsize=\codefont,xleftmargin=17bp]
\PYG{n}{T}\PYG{+w}{ }\PYG{k}{operator}\PYG{+w}{ }\PYG{l+s}{\PYGZdq{}}\PYG{l+s}{\PYGZdq{}}\PYG{+w}{ }\PYG{n}{\PYGZus{}suffix}\PYG{p}{(}\PYG{k+kt}{unsigned}\PYG{+w}{ }\PYG{k+kt}{long}\PYG{+w}{ }\PYG{k+kt}{long}\PYG{+w}{ }\PYG{k+kt}{int}\PYG{p}{)}\PYG{p}{;}
\PYG{n}{T}\PYG{+w}{ }\PYG{k}{operator}\PYG{+w}{ }\PYG{l+s}{\PYGZdq{}}\PYG{l+s}{\PYGZdq{}}\PYG{+w}{ }\PYG{n}{\PYGZus{}suffix}\PYG{p}{(}\PYG{k+kt}{long}\PYG{+w}{ }\PYG{k+kt}{double}\PYG{p}{)}\PYG{p}{;}
\PYG{n}{T}\PYG{+w}{ }\PYG{k}{operator}\PYG{+w}{ }\PYG{l+s}{\PYGZdq{}}\PYG{l+s}{\PYGZdq{}}\PYG{+w}{ }\PYG{n}{\PYGZus{}suffix}\PYG{p}{(}\PYG{k+kt}{char}\PYG{p}{)}\PYG{p}{;}
\PYG{n}{T}\PYG{+w}{ }\PYG{k}{operator}\PYG{+w}{ }\PYG{l+s}{\PYGZdq{}}\PYG{l+s}{\PYGZdq{}}\PYG{+w}{ }\PYG{n}{\PYGZus{}suffix}\PYG{p}{(}\PYG{k+kt}{wchar\PYGZus{}t}\PYG{p}{)}\PYG{p}{;}
\PYG{n}{T}\PYG{+w}{ }\PYG{k}{operator}\PYG{+w}{ }\PYG{l+s}{\PYGZdq{}}\PYG{l+s}{\PYGZdq{}}\PYG{+w}{ }\PYG{n}{\PYGZus{}suffix}\PYG{p}{(}\PYG{k+kt}{char16\PYGZus{}t}\PYG{p}{)}\PYG{p}{;}
\PYG{n}{T}\PYG{+w}{ }\PYG{k}{operator}\PYG{+w}{ }\PYG{l+s}{\PYGZdq{}}\PYG{l+s}{\PYGZdq{}}\PYG{+w}{ }\PYG{n}{\PYGZus{}suffix}\PYG{p}{(}\PYG{k+kt}{char32\PYGZus{}t}\PYG{p}{)}\PYG{p}{;}
\PYG{n}{T}\PYG{+w}{ }\PYG{k}{operator}\PYG{+w}{ }\PYG{l+s}{\PYGZdq{}}\PYG{l+s}{\PYGZdq{}}\PYG{+w}{ }\PYG{n}{\PYGZus{}suffix}\PYG{p}{(}\PYG{k+kt}{char}\PYG{+w}{ }\PYG{k}{const}\PYG{+w}{ }\PYG{o}{*}\PYG{p}{,}\PYG{+w}{ }\PYG{n}{std}\PYG{o}{:}\PYG{o}{:}\PYG{k+kt}{size\PYGZus{}t}\PYG{p}{)}\PYG{p}{;}
\PYG{n}{T}\PYG{+w}{ }\PYG{k}{operator}\PYG{+w}{ }\PYG{l+s}{\PYGZdq{}}\PYG{l+s}{\PYGZdq{}}\PYG{+w}{ }\PYG{n}{\PYGZus{}suffix}\PYG{p}{(}\PYG{k+kt}{wchar\PYGZus{}t}\PYG{+w}{ }\PYG{k}{const}\PYG{+w}{ }\PYG{o}{*}\PYG{p}{,}\PYG{+w}{ }\PYG{n}{std}\PYG{o}{:}\PYG{o}{:}\PYG{k+kt}{size\PYGZus{}t}\PYG{p}{)}\PYG{p}{;}
\PYG{n}{T}\PYG{+w}{ }\PYG{k}{operator}\PYG{+w}{ }\PYG{l+s}{\PYGZdq{}}\PYG{l+s}{\PYGZdq{}}\PYG{+w}{ }\PYG{n}{\PYGZus{}suffix}\PYG{p}{(}\PYG{k+kt}{char16\PYGZus{}t}\PYG{+w}{ }\PYG{k}{const}\PYG{+w}{ }\PYG{o}{*}\PYG{p}{,}\PYG{+w}{ }\PYG{n}{std}\PYG{o}{:}\PYG{o}{:}\PYG{k+kt}{size\PYGZus{}t}\PYG{p}{)}\PYG{p}{;}
\PYG{n}{T}\PYG{+w}{ }\PYG{k}{operator}\PYG{+w}{ }\PYG{l+s}{\PYGZdq{}}\PYG{l+s}{\PYGZdq{}}\PYG{+w}{ }\PYG{n}{\PYGZus{}suffix}\PYG{p}{(}\PYG{k+kt}{char32\PYGZus{}t}\PYG{+w}{ }\PYG{k}{const}\PYG{+w}{ }\PYG{o}{*}\PYG{p}{,}\PYG{+w}{ }\PYG{n}{std}\PYG{o}{:}\PYG{o}{:}\PYG{k+kt}{size\PYGZus{}t}\PYG{p}{)}\PYG{p}{;}
\end{Verbatim}
\end{codeblockbg}


%%C020194 p

下面的示例创建一个用于表示千字节的用户定义字面量：


%%C020195 code

\begin{codeblockbg}
\begin{Verbatim}[commandchars=\\\{\},breaklines=true,breakanywhere=true,fontsize=\codefont,xleftmargin=17bp]
\PYG{k}{namespace}\PYG{+w}{ }\PYG{n+nn}{compunits}
\PYG{p}{\PYGZob{}}
\PYG{+w}{  }\PYG{k}{constexpr}\PYG{+w}{ }\PYG{k+kt}{size\PYGZus{}t}\PYG{+w}{ }\PYG{k}{operator}\PYG{+w}{ }\PYG{l+s}{\PYGZdq{}}\PYG{l+s}{\PYGZdq{}}\PYG{+w}{ }\PYG{n}{\PYGZus{}KB}\PYG{p}{(}\PYG{k+kt}{unsigned}\PYG{+w}{ }\PYG{k+kt}{long}\PYG{+w}{ }\PYG{k+kt}{long}\PYG{+w}{ }\PYG{k}{const}\PYG{+w}{ }\PYG{n}{size}\PYG{p}{)}
\PYG{+w}{  }\PYG{p}{\PYGZob{}}
\PYG{+w}{    }\PYG{k}{return}\PYG{+w}{ }\PYG{k}{static\PYGZus{}cast}\PYG{o}{\PYGZlt{}}\PYG{k+kt}{size\PYGZus{}t}\PYG{o}{\PYGZgt{}}\PYG{p}{(}\PYG{n}{size}\PYG{+w}{ }\PYG{o}{*}\PYG{+w}{ }\PYG{l+m+mi}{1024}\PYG{p}{)}\PYG{p}{;}
\PYG{+w}{  }\PYG{p}{\PYGZcb{}}
\PYG{p}{\PYGZcb{}}
\PYG{k}{auto}\PYG{+w}{ }\PYG{n}{size}\PYG{p}{\PYGZob{}}\PYG{+w}{ }\PYG{l+m+mi}{4}\PYG{n}{\PYGZus{}KB}\PYG{+w}{ }\PYG{p}{\PYGZcb{}}\PYG{p}{;}\PYG{+w}{         }\PYG{c+c1}{// size\PYGZus{}t size = 4096;}
\PYG{k}{using}\PYG{+w}{ }\PYG{k+kt}{byte}\PYG{+w}{ }\PYG{o}{=}\PYG{+w}{ }\PYG{k+kt}{unsigned}\PYG{+w}{ }\PYG{k+kt}{char}\PYG{p}{;}
\PYG{k}{auto}\PYG{+w}{ }\PYG{n}{buffer}\PYG{+w}{ }\PYG{o}{=}\PYG{+w}{ }\PYG{n}{std}\PYG{o}{:}\PYG{o}{:}\PYG{n}{array}\PYG{o}{\PYGZlt{}}\PYG{k+kt}{byte}\PYG{p}{,}\PYG{+w}{ }\PYG{l+m+mi}{1}\PYG{n}{\PYGZus{}KB}\PYG{o}{\PYGZgt{}}\PYG{p}{\PYGZob{}}\PYG{p}{\PYGZcb{}}\PYG{p}{;}
\end{Verbatim}
\end{codeblockbg}


%%C020196 h3

\subsection{工作原理}


%%C020197 p

当编译器遇到带有用户定义后缀\texttt{S}的用户定义字面量时（第三方后缀总是以下划线开头，因为不带下划线的后缀保留给标准库），它会做一次非限定名称查找，以找出名为\texttt{“operator ”“ S}的函数。如果找到了，就根据字面量的类型和字面量运算符的类型调用它；否则编译器将报错。


%%C020198 p

在\textit{具体做法}一节的示例中，字面量运算符名为\texttt{operator “”} \texttt{\_\allowbreak{}KB}，参数类型为\myttbr{unsigned long long int}。这是处理整型时字面量运算符唯一可用的整型参数类型。同样，对于浮点用户定义字面量，参数类型必须是\texttt{long double}，因为对于数值类型，字面量运算符必须能够处理尽可能大的值。该字面量运算符返回一个\texttt{constexpr}值，因此可以用在需要编译期值的地方，例如上面示例中指定数组大小。


%%C020199 p

当编译器识别出用户定义字面量并需要调用相应的字面量运算符时，它会按以下规则从重载集中挑选重载：


%%C020200 li

\begin{itemize}
\item \textbf{整型字面量}：按以下顺序调用：接受\texttt{unsigned long long}的运算符、接受\texttt{const char*}的raw字面量运算符，或字面量运算符模板。


%%C020201 li

\item \textbf{浮点字面量}：按以下顺序调用：接受\texttt{long double}的运算符、接受\texttt{const char*}的raw字面量运算符，或字面量运算符模板。


%%C020202 li

\item \textbf{字符字面量}：根据字符类型（\texttt{char}、\texttt{wchar\_\allowbreak{}t}、\texttt{char16\_\allowbreak{}t}和\texttt{char32\_\allowbreak{}t}）调用相应的运算符。


%%C020203 li

\item \textbf{字符串字面量}：根据字符串类型调用相应的运算符，该运算符接受指向字符序列的指针和长度。

\end{itemize}


%%C020204 p

下面的示例定义了一套单位与数量的系统。我们要处理千克、件、升以及其他类型的单位。这在处理订单的系统中很有用，因为需要为每件商品指定数量和单位。在命名空间\texttt{units}中定义了以下内容：


%%C020205 li

\begin{itemize}
\item 表示可能的单位类型（千克、米、升和件）的限定作用域枚举：


%%C020206 code

\begin{codeblockbg}
\begin{Verbatim}[commandchars=\\\{\},breaklines=true,breakanywhere=true,fontsize=\codefont,xleftmargin=17bp]
\PYG{k}{enum}\PYG{+w}{ }\PYG{k}{class}\PYG{+w}{ }\PYG{n+nc}{unit}\PYG{+w}{ }\PYG{p}{\PYGZob{}}\PYG{+w}{ }\PYG{n}{kilogram}\PYG{p}{,}\PYG{+w}{ }\PYG{n}{liter}\PYG{p}{,}\PYG{+w}{ }\PYG{n}{meter}\PYG{p}{,}\PYG{+w}{ }\PYG{n}{piece}\PYG{p}{,}\PYG{+w}{ }\PYG{p}{\PYGZcb{}}\PYG{p}{;}
\end{Verbatim}
\end{codeblockbg}


%%C020207 li

\item 用于表示特定单位数量的类模板（例如3.5千克或42件）：


%%C020208 code

\begin{codeblockbg}
\begin{Verbatim}[commandchars=\\\{\},breaklines=true,breakanywhere=true,fontsize=\codefont,xleftmargin=17bp]
\PYG{k}{template}\PYG{+w}{ }\PYG{o}{\PYGZlt{}}\PYG{n}{unit}\PYG{+w}{ }\PYG{n}{U}\PYG{o}{\PYGZgt{}}
\PYG{k}{class}\PYG{+w}{ }\PYG{n+nc}{quantity}
\PYG{p}{\PYGZob{}}
\PYG{+w}{  }\PYG{k}{const}\PYG{+w}{ }\PYG{k+kt}{double}\PYG{+w}{ }\PYG{n}{amount}\PYG{p}{;}
\PYG{+w}{  }\PYG{k}{public}\PYG{o}{:}
\PYG{+w}{    }\PYG{k}{constexpr}\PYG{+w}{ }\PYG{k}{explicit}\PYG{+w}{ }\PYG{n}{quantity}\PYG{p}{(}\PYG{k+kt}{double}\PYG{+w}{ }\PYG{k}{const}\PYG{+w}{ }\PYG{n}{a}\PYG{p}{)}\PYG{+w}{ }\PYG{o}{:}
\PYG{+w}{      }\PYG{n}{amount}\PYG{p}{(}\PYG{n}{a}\PYG{p}{)}\PYG{+w}{ }\PYG{p}{\PYGZob{}}\PYG{p}{\PYGZcb{}}
\PYG{+w}{  }\PYG{k}{explicit}\PYG{+w}{ }\PYG{k}{operator}\PYG{+w}{ }\PYG{k+kt}{double}\PYG{p}{(}\PYG{p}{)}\PYG{+w}{ }\PYG{k}{const}\PYG{+w}{ }\PYG{p}{\PYGZob{}}\PYG{+w}{ }\PYG{k}{return}\PYG{+w}{ }\PYG{n}{amount}\PYG{p}{;}\PYG{+w}{ }\PYG{p}{\PYGZcb{}}
\PYG{p}{\PYGZcb{}}\PYG{p}{;}
\end{Verbatim}
\end{codeblockbg}


%%C020209 li

\item 为\texttt{quantity}类模板定义\texttt{operator+}与\texttt{operator-}函数，以便对数量做加法和减法：


%%C020210 code

\begin{codeblockbg}
\begin{Verbatim}[commandchars=\\\{\},breaklines=true,breakanywhere=true,fontsize=\codefont,xleftmargin=17bp]
\PYG{k}{template}\PYG{+w}{ }\PYG{o}{\PYGZlt{}}\PYG{n}{unit}\PYG{+w}{ }\PYG{n}{U}\PYG{o}{\PYGZgt{}}
\PYG{k}{constexpr}\PYG{+w}{ }\PYG{n}{quantity}\PYG{o}{\PYGZlt{}}\PYG{n}{U}\PYG{o}{\PYGZgt{}}\PYG{+w}{ }\PYG{k}{operator}\PYG{o}{+}\PYG{p}{(}\PYG{n}{quantity}\PYG{o}{\PYGZlt{}}\PYG{n}{U}\PYG{o}{\PYGZgt{}}\PYG{+w}{ }\PYG{k}{const}\PYG{+w}{ }\PYG{o}{\PYGZam{}}\PYG{n}{q1}\PYG{p}{,}
\PYG{+w}{                                }\PYG{n}{quantity}\PYG{o}{\PYGZlt{}}\PYG{n}{U}\PYG{o}{\PYGZgt{}}\PYG{+w}{ }\PYG{k}{const}\PYG{+w}{ }\PYG{o}{\PYGZam{}}\PYG{n}{q2}\PYG{p}{)}
\PYG{p}{\PYGZob{}}
\PYG{+w}{  }\PYG{k}{return}\PYG{+w}{ }\PYG{n}{quantity}\PYG{o}{\PYGZlt{}}\PYG{n}{U}\PYG{o}{\PYGZgt{}}\PYG{p}{(}\PYG{k}{static\PYGZus{}cast}\PYG{o}{\PYGZlt{}}\PYG{k+kt}{double}\PYG{o}{\PYGZgt{}}\PYG{p}{(}\PYG{n}{q1}\PYG{p}{)}\PYG{+w}{ }\PYG{o}{+}
\PYG{+w}{                     }\PYG{k}{static\PYGZus{}cast}\PYG{o}{\PYGZlt{}}\PYG{k+kt}{double}\PYG{o}{\PYGZgt{}}\PYG{p}{(}\PYG{n}{q2}\PYG{p}{)}\PYG{p}{)}\PYG{p}{;}
\PYG{p}{\PYGZcb{}}
\PYG{k}{template}\PYG{+w}{ }\PYG{o}{\PYGZlt{}}\PYG{n}{unit}\PYG{+w}{ }\PYG{n}{U}\PYG{o}{\PYGZgt{}}
\PYG{k}{constexpr}\PYG{+w}{ }\PYG{n}{quantity}\PYG{o}{\PYGZlt{}}\PYG{n}{U}\PYG{o}{\PYGZgt{}}\PYG{+w}{ }\PYG{k}{operator}\PYG{o}{\PYGZhy{}}\PYG{p}{(}\PYG{n}{quantity}\PYG{o}{\PYGZlt{}}\PYG{n}{U}\PYG{o}{\PYGZgt{}}\PYG{+w}{ }\PYG{k}{const}\PYG{+w}{ }\PYG{o}{\PYGZam{}}\PYG{n}{q1}\PYG{p}{,}
\PYG{+w}{                                }\PYG{n}{quantity}\PYG{o}{\PYGZlt{}}\PYG{n}{U}\PYG{o}{\PYGZgt{}}\PYG{+w}{ }\PYG{k}{const}\PYG{+w}{ }\PYG{o}{\PYGZam{}}\PYG{n}{q2}\PYG{p}{)}
\PYG{p}{\PYGZob{}}
\PYG{+w}{  }\PYG{k}{return}\PYG{+w}{ }\PYG{n}{quantity}\PYG{o}{\PYGZlt{}}\PYG{n}{U}\PYG{o}{\PYGZgt{}}\PYG{p}{(}\PYG{k}{static\PYGZus{}cast}\PYG{o}{\PYGZlt{}}\PYG{k+kt}{double}\PYG{o}{\PYGZgt{}}\PYG{p}{(}\PYG{n}{q1}\PYG{p}{)}\PYG{+w}{ }\PYG{o}{\PYGZhy{}}
\PYG{+w}{                     }\PYG{k}{static\PYGZus{}cast}\PYG{o}{\PYGZlt{}}\PYG{k+kt}{double}\PYG{o}{\PYGZgt{}}\PYG{p}{(}\PYG{n}{q2}\PYG{p}{)}\PYG{p}{)}\PYG{p}{;}
\PYG{p}{\PYGZcb{}}
\end{Verbatim}
\end{codeblockbg}


%%C020211 li

\item 用于创建\texttt{quantity}字面量的字面量运算符，定义在名为\texttt{unit\_\allowbreak{}literals}的内层命名空间中。这样做的目的是避免与其他命名空间的字面量发生名称冲突。如果真的发生冲突，开发者可以在需要使用字面量的作用域中引入合适的命名空间，从而选择要用的一组字面量：


%%C020212 code

\begin{codeblockbg}
\begin{Verbatim}[commandchars=\\\{\},breaklines=true,breakanywhere=true,fontsize=\codefont,xleftmargin=17bp]
\PYG{k}{namespace}\PYG{+w}{ }\PYG{n+nn}{unit\PYGZus{}literals}
\PYG{p}{\PYGZob{}}
\PYG{+w}{  }\PYG{k}{constexpr}\PYG{+w}{ }\PYG{n}{quantity}\PYG{o}{\PYGZlt{}}\PYG{n}{unit}\PYG{o}{:}\PYG{o}{:}\PYG{n}{kilogram}\PYG{o}{\PYGZgt{}}\PYG{+w}{ }\PYG{k}{operator}\PYG{+w}{ }\PYG{l+s}{\PYGZdq{}}\PYG{l+s}{\PYGZdq{}}\PYG{+w}{ }\PYG{n}{\PYGZus{}kg}\PYG{p}{(}
\PYG{+w}{      }\PYG{k+kt}{long}\PYG{+w}{ }\PYG{k+kt}{double}\PYG{+w}{ }\PYG{k}{const}\PYG{+w}{ }\PYG{n}{amount}\PYG{p}{)}
\PYG{+w}{  }\PYG{p}{\PYGZob{}}
\PYG{+w}{    }\PYG{k}{return}\PYG{+w}{ }\PYG{n}{quantity}\PYG{o}{\PYGZlt{}}\PYG{n}{unit}\PYG{o}{:}\PYG{o}{:}\PYG{n}{kilogram}\PYG{o}{\PYGZgt{}}
\PYG{+w}{      }\PYG{p}{\PYGZob{}}\PYG{+w}{ }\PYG{k}{static\PYGZus{}cast}\PYG{o}{\PYGZlt{}}\PYG{k+kt}{double}\PYG{o}{\PYGZgt{}}\PYG{p}{(}\PYG{n}{amount}\PYG{p}{)}\PYG{+w}{ }\PYG{p}{\PYGZcb{}}\PYG{p}{;}
\PYG{+w}{  }\PYG{p}{\PYGZcb{}}
\PYG{k}{constexpr}\PYG{+w}{ }\PYG{n}{quantity}\PYG{o}{\PYGZlt{}}\PYG{n}{unit}\PYG{o}{:}\PYG{o}{:}\PYG{n}{kilogram}\PYG{o}{\PYGZgt{}}\PYG{+w}{ }\PYG{k}{operator}\PYG{+w}{ }\PYG{l+s}{\PYGZdq{}}\PYG{l+s}{\PYGZdq{}}\PYG{+w}{ }\PYG{n}{\PYGZus{}kg}\PYG{p}{(}
\PYG{+w}{   }\PYG{k+kt}{unsigned}\PYG{+w}{ }\PYG{k+kt}{long}\PYG{+w}{ }\PYG{k+kt}{long}\PYG{+w}{ }\PYG{k}{const}\PYG{+w}{ }\PYG{n}{amount}\PYG{p}{)}
\PYG{p}{\PYGZob{}}
\PYG{+w}{  }\PYG{k}{return}\PYG{+w}{ }\PYG{n}{quantity}\PYG{o}{\PYGZlt{}}\PYG{n}{unit}\PYG{o}{:}\PYG{o}{:}\PYG{n}{kilogram}\PYG{o}{\PYGZgt{}}
\PYG{+w}{    }\PYG{p}{\PYGZob{}}\PYG{+w}{ }\PYG{k}{static\PYGZus{}cast}\PYG{o}{\PYGZlt{}}\PYG{k+kt}{double}\PYG{o}{\PYGZgt{}}\PYG{p}{(}\PYG{n}{amount}\PYG{p}{)}\PYG{+w}{ }\PYG{p}{\PYGZcb{}}\PYG{p}{;}
\PYG{p}{\PYGZcb{}}
\PYG{k}{constexpr}\PYG{+w}{ }\PYG{n}{quantity}\PYG{o}{\PYGZlt{}}\PYG{n}{unit}\PYG{o}{:}\PYG{o}{:}\PYG{n}{liter}\PYG{o}{\PYGZgt{}}\PYG{+w}{ }\PYG{k}{operator}\PYG{+w}{ }\PYG{l+s}{\PYGZdq{}}\PYG{l+s}{\PYGZdq{}}\PYG{+w}{ }\PYG{n}{\PYGZus{}l}\PYG{p}{(}
\PYG{+w}{   }\PYG{k+kt}{long}\PYG{+w}{ }\PYG{k+kt}{double}\PYG{+w}{ }\PYG{k}{const}\PYG{+w}{ }\PYG{n}{amount}\PYG{p}{)}
\PYG{p}{\PYGZob{}}
\PYG{+w}{   }\PYG{k}{return}\PYG{+w}{ }\PYG{n}{quantity}\PYG{o}{\PYGZlt{}}\PYG{n}{unit}\PYG{o}{:}\PYG{o}{:}\PYG{n}{liter}\PYG{o}{\PYGZgt{}}
\PYG{+w}{     }\PYG{p}{\PYGZob{}}\PYG{+w}{ }\PYG{k}{static\PYGZus{}cast}\PYG{o}{\PYGZlt{}}\PYG{k+kt}{double}\PYG{o}{\PYGZgt{}}\PYG{p}{(}\PYG{n}{amount}\PYG{p}{)}\PYG{+w}{ }\PYG{p}{\PYGZcb{}}\PYG{p}{;}
\PYG{p}{\PYGZcb{}}
\PYG{k}{constexpr}\PYG{+w}{ }\PYG{n}{quantity}\PYG{o}{\PYGZlt{}}\PYG{n}{unit}\PYG{o}{:}\PYG{o}{:}\PYG{n}{meter}\PYG{o}{\PYGZgt{}}\PYG{+w}{ }\PYG{k}{operator}\PYG{+w}{ }\PYG{l+s}{\PYGZdq{}}\PYG{l+s}{\PYGZdq{}}\PYG{+w}{ }\PYG{n}{\PYGZus{}m}\PYG{p}{(}
\PYG{+w}{   }\PYG{k+kt}{long}\PYG{+w}{ }\PYG{k+kt}{double}\PYG{+w}{ }\PYG{k}{const}\PYG{+w}{ }\PYG{n}{amount}\PYG{p}{)}
\PYG{p}{\PYGZob{}}
\PYG{+w}{  }\PYG{k}{return}\PYG{+w}{ }\PYG{n}{quantity}\PYG{o}{\PYGZlt{}}\PYG{n}{unit}\PYG{o}{:}\PYG{o}{:}\PYG{n}{meter}\PYG{o}{\PYGZgt{}}
\PYG{+w}{    }\PYG{p}{\PYGZob{}}\PYG{+w}{ }\PYG{k}{static\PYGZus{}cast}\PYG{o}{\PYGZlt{}}\PYG{k+kt}{double}\PYG{o}{\PYGZgt{}}\PYG{p}{(}\PYG{n}{amount}\PYG{p}{)}\PYG{+w}{ }\PYG{p}{\PYGZcb{}}\PYG{p}{;}
\PYG{p}{\PYGZcb{}}
\PYG{+w}{  }\PYG{k}{constexpr}\PYG{+w}{ }\PYG{n}{quantity}\PYG{o}{\PYGZlt{}}\PYG{n}{unit}\PYG{o}{:}\PYG{o}{:}\PYG{n}{piece}\PYG{o}{\PYGZgt{}}\PYG{+w}{ }\PYG{k}{operator}\PYG{+w}{ }\PYG{l+s}{\PYGZdq{}}\PYG{l+s}{\PYGZdq{}}\PYG{+w}{ }\PYG{n}{\PYGZus{}pcs}\PYG{p}{(}
\PYG{+w}{     }\PYG{k+kt}{unsigned}\PYG{+w}{ }\PYG{k+kt}{long}\PYG{+w}{ }\PYG{k+kt}{long}\PYG{+w}{ }\PYG{k}{const}\PYG{+w}{ }\PYG{n}{amount}\PYG{p}{)}
\PYG{+w}{  }\PYG{p}{\PYGZob{}}
\PYG{+w}{    }\PYG{k}{return}\PYG{+w}{ }\PYG{n}{quantity}\PYG{o}{\PYGZlt{}}\PYG{n}{unit}\PYG{o}{:}\PYG{o}{:}\PYG{n}{piece}\PYG{o}{\PYGZgt{}}
\PYG{+w}{      }\PYG{p}{\PYGZob{}}\PYG{+w}{ }\PYG{k}{static\PYGZus{}cast}\PYG{o}{\PYGZlt{}}\PYG{k+kt}{double}\PYG{o}{\PYGZgt{}}\PYG{p}{(}\PYG{n}{amount}\PYG{p}{)}\PYG{+w}{ }\PYG{p}{\PYGZcb{}}\PYG{p}{;}
\PYG{+w}{  }\PYG{p}{\PYGZcb{}}
\PYG{p}{\PYGZcb{}}
\end{Verbatim}
\end{codeblockbg}

\end{itemize}


%%C020213 p

仔细观察可以发现，前面定义的字面量运算符并不相同：


%%C020214 li

\begin{itemize}
\item \texttt{\_\allowbreak{}kg}同时为整型和浮点字面量定义，因此可以创建整型和浮点值，如\texttt{1\_\allowbreak{}kg}和\texttt{1.\allowbreak{}0\_\allowbreak{}kg}。


%%C020215 li

\item \texttt{\_\allowbreak{}l}和\texttt{\_\allowbreak{}m}只为浮点字面量定义，这意味着只能用浮点数表示这些单位的数量字面量，如\texttt{4.\allowbreak{}5\_\allowbreak{}l}和\texttt{10.\allowbreak{}0\_\allowbreak{}m}。


%%C020216 li

\item \texttt{\_\allowbreak{}pcs}只为整型字面量定义，这意味着只能表示整数件的数量，如\texttt{42\_\allowbreak{}pcs}。

\end{itemize}


%%C020217 p

有了这些字面量运算符，就可以对各种数量进行运算。下面的示例既包含合法的运算，也包含非法的运算：


%%C020218 code

\begin{codeblockbg}
\begin{Verbatim}[commandchars=\\\{\},breaklines=true,breakanywhere=true,fontsize=\codefont,xleftmargin=17bp]
\PYG{k}{using}\PYG{+w}{ }\PYG{k}{namespace}\PYG{+w}{ }\PYG{n+nn}{units}\PYG{p}{;}
\PYG{k}{using}\PYG{+w}{ }\PYG{k}{namespace}\PYG{+w}{ }\PYG{n+nn}{unit\PYGZus{}literals}\PYG{p}{;}
\PYG{k}{auto}\PYG{+w}{ }\PYG{n}{q1}\PYG{p}{\PYGZob{}}\PYG{+w}{ }\PYG{l+m+mi}{1}\PYG{n}{\PYGZus{}kg}\PYG{+w}{ }\PYG{p}{\PYGZcb{}}\PYG{p}{;}\PYG{+w}{    }\PYG{c+c1}{// 正确}
\PYG{k}{auto}\PYG{+w}{ }\PYG{n}{q2}\PYG{p}{\PYGZob{}}\PYG{+w}{ }\PYG{l+m+mf}{4.5}\PYG{n}{\PYGZus{}kg}\PYG{+w}{ }\PYG{p}{\PYGZcb{}}\PYG{p}{;}\PYG{+w}{  }\PYG{c+c1}{// 正确}
\PYG{k}{auto}\PYG{+w}{ }\PYG{n}{q3}\PYG{p}{\PYGZob{}}\PYG{+w}{ }\PYG{n}{q1}\PYG{+w}{ }\PYG{o}{+}\PYG{+w}{ }\PYG{n}{q2}\PYG{+w}{ }\PYG{p}{\PYGZcb{}}\PYG{p}{;}\PYG{+w}{ }\PYG{c+c1}{// 正确}
\PYG{k}{auto}\PYG{+w}{ }\PYG{n}{q4}\PYG{p}{\PYGZob{}}\PYG{+w}{ }\PYG{n}{q2}\PYG{+w}{ }\PYG{o}{\PYGZhy{}}\PYG{+w}{ }\PYG{n}{q1}\PYG{+w}{ }\PYG{p}{\PYGZcb{}}\PYG{p}{;}\PYG{+w}{ }\PYG{c+c1}{// 正确}
\PYG{c+c1}{// 错误：不能把米和件相加}
\PYG{k}{auto}\PYG{+w}{ }\PYG{n}{q5}\PYG{p}{\PYGZob{}}\PYG{+w}{ }\PYG{l+m+mf}{1.0}\PYG{n}{\PYGZus{}m}\PYG{+w}{ }\PYG{o}{+}\PYG{+w}{ }\PYG{l+m+mi}{1}\PYG{n}{\PYGZus{}pcs}\PYG{+w}{ }\PYG{p}{\PYGZcb{}}\PYG{p}{;}
\PYG{c+c1}{// 错误：升的数量不能是整数}
\PYG{k}{auto}\PYG{+w}{ }\PYG{n}{q6}\PYG{p}{\PYGZob{}}\PYG{+w}{ }\PYG{l+m+mi}{1}\PYG{n}{\PYGZus{}l}\PYG{+w}{ }\PYG{p}{\PYGZcb{}}\PYG{p}{;}
\PYG{c+c1}{// 错误：件数只能是整数}
\PYG{k}{auto}\PYG{+w}{ }\PYG{n}{q7}\PYG{p}{\PYGZob{}}\PYG{+w}{ }\PYG{l+m+mf}{2.0}\PYG{n}{\PYGZus{}pcs}\PYG{p}{\PYGZcb{}}
\end{Verbatim}
\end{codeblockbg}


%%C020219 p

\texttt{q1}是1千克的数量，即一个整数值。因为存在重载的\texttt{operator “”} \myttbr{\_\allowbreak{}kg(\allowbreak{}unsigned long long const)}，所以可以从整数1正确创建该字面量。类似地，\texttt{q2}是4.5千克的数量，即一个实数值。因为存在重载版本\texttt{overload} \myttbr{operator “” \_\allowbreak{}kg(\allowbreak{}long double)}，所以可以从双精度浮点值4.5创建该字面量。


%%C020220 p

另一方面，\texttt{q6}是1升的数量。由于不存在重载的\texttt{operator “”} \myttbr{\_\allowbreak{}l(\allowbreak{}unsigned long long)}，该字面量无法创建。它需要一个接受\texttt{unsigned long long}的重载，但这样的重载并不存在。类似地，\texttt{q7}是2.0件的数量，但件字面量只能由整数值创建，因此这里又会产生一个编译错误。


%%C020221 h3

\subsection{更多内容}


%%C020222 p

虽然从C++11起就有了用户定义字面量，但标准字面量运算符直到C++14才提供。下面列出这些标准字面量运算符：


%%C020223 li

\begin{itemize}
\item 用于定义\texttt{std::\allowbreak{}basic\_\allowbreak{}string}字面量的\texttt{operator“”s}：


%%C020224 code

\begin{codeblockbg}
\begin{Verbatim}[commandchars=\\\{\},breaklines=true,breakanywhere=true,fontsize=\codefont,xleftmargin=17bp]
\PYG{k}{using}\PYG{+w}{ }\PYG{k}{namespace}\PYG{+w}{ }\PYG{n+nn}{std}\PYG{o}{:}\PYG{o}{:}\PYG{n+nn}{string\PYGZus{}literals}\PYG{p}{;}
\PYG{k}{auto}\PYG{+w}{ }\PYG{n}{s1}\PYG{p}{\PYGZob{}}\PYG{+w}{  }\PYG{l+s}{\PYGZdq{}}\PYG{l+s}{text}\PYG{l+s}{\PYGZdq{}}\PYG{n}{s}\PYG{+w}{ }\PYG{p}{\PYGZcb{}}\PYG{p}{;}\PYG{+w}{ }\PYG{c+c1}{// std::string}
\PYG{k}{auto}\PYG{+w}{ }\PYG{n}{s2}\PYG{p}{\PYGZob{}}\PYG{+w}{ }\PYG{l+s+sa}{L}\PYG{l+s}{\PYGZdq{}}\PYG{l+s}{text}\PYG{l+s}{\PYGZdq{}}\PYG{n}{s}\PYG{+w}{ }\PYG{p}{\PYGZcb{}}\PYG{p}{;}\PYG{+w}{ }\PYG{c+c1}{// std::wstring}
\PYG{k}{auto}\PYG{+w}{ }\PYG{n}{s3}\PYG{p}{\PYGZob{}}\PYG{+w}{ }\PYG{l+s+sa}{u}\PYG{l+s}{\PYGZdq{}}\PYG{l+s}{text}\PYG{l+s}{\PYGZdq{}}\PYG{n}{s}\PYG{+w}{ }\PYG{p}{\PYGZcb{}}\PYG{p}{;}\PYG{+w}{ }\PYG{c+c1}{// std::u16string}
\PYG{k}{auto}\PYG{+w}{ }\PYG{n}{s4}\PYG{p}{\PYGZob{}}\PYG{+w}{ }\PYG{l+s+sa}{U}\PYG{l+s}{\PYGZdq{}}\PYG{l+s}{text}\PYG{l+s}{\PYGZdq{}}\PYG{n}{s}\PYG{+w}{ }\PYG{p}{\PYGZcb{}}\PYG{p}{;}\PYG{+w}{ }\PYG{c+c1}{// std::u32string}
\end{Verbatim}
\end{codeblockbg}


%%C020225 li

\item \texttt{operator“”h}、\texttt{operator“”min}、\texttt{operator“”s}、\texttt{operator“”ms}、\texttt{operator“”us}和\texttt{operator“”ns}，用于创建\myttbr{std::\allowbreak{}chrono::\allowbreak{}duration}值：


%%C020226 code

\begin{codeblockbg}
\begin{Verbatim}[commandchars=\\\{\},breaklines=true,breakanywhere=true,fontsize=\codefont,xleftmargin=17bp]
\PYG{k}{using}\PYG{+w}{ }\PYG{k}{namespace}\PYG{+w}{ }\PYG{n+nn}{std}\PYG{o}{:}\PYG{o}{:}\PYG{n+nn}{literals}\PYG{o}{:}\PYG{o}{:}\PYG{n+nn}{chrono\PYGZus{}literals}\PYG{p}{;}
\PYG{c+c1}{// std::chrono::duration\PYGZlt{}long long\PYGZgt{}}
\PYG{k}{auto}\PYG{+w}{ }\PYG{n}{timer}\PYG{+w}{ }\PYG{p}{\PYGZob{}}\PYG{l+m+mi}{2}\PYG{n}{h}\PYG{+w}{ }\PYG{o}{+}\PYG{+w}{ }\PYG{l+m+mi}{42}\PYG{n}{min}\PYG{+w}{ }\PYG{o}{+}\PYG{+w}{ }\PYG{l+m+mi}{15}\PYG{n}{s}\PYG{p}{\PYGZcb{}}\PYG{p}{;}
\end{Verbatim}
\end{codeblockbg}


%%C020227 li

\item \texttt{operator“”if}、\texttt{operator“”i}和\texttt{operator“”il}，用于创建\texttt{std::\allowbreak{}complex}值：


%%C020228 code

\begin{codeblockbg}
\begin{Verbatim}[commandchars=\\\{\},breaklines=true,breakanywhere=true,fontsize=\codefont,xleftmargin=17bp]
\PYG{k}{using}\PYG{+w}{ }\PYG{k}{namespace}\PYG{+w}{ }\PYG{n+nn}{std}\PYG{o}{:}\PYG{o}{:}\PYG{n+nn}{literals}\PYG{o}{:}\PYG{o}{:}\PYG{n+nn}{complex\PYGZus{}literals}\PYG{p}{;}
\PYG{k}{auto}\PYG{+w}{ }\PYG{n}{c}\PYG{p}{\PYGZob{}}\PYG{+w}{ }\PYG{l+m+mf}{12.0}\PYG{+w}{ }\PYG{o}{+}\PYG{+w}{ }\PYG{l+m+mf}{4.5}\PYG{n}{i}\PYG{+w}{ }\PYG{p}{\PYGZcb{}}\PYG{p}{;}\PYG{+w}{ }\PYG{c+c1}{// std::complex\PYGZlt{}double\PYGZgt{}}
\end{Verbatim}
\end{codeblockbg}

\end{itemize}


%%C020229 h3

\subsection{延伸阅读}


%%C020230 li

\begin{itemize}
\item \textit{用原始字符串字面量避免转义字符}


%%C020231 li

\item \textit{创建raw用户定义字面量}

\end{itemize}


%%C020232 h2

\section{创建raw用户定义字面量}


%%C020233 p

上一个案例介绍了C++11如何让库实现者和开发者创建用户定义字面量，以及C++14标准中可用的用户定义字面量。不过，用户定义字面量有两种形式：cooked形式，字面量值在传给字面量运算符之前已由编译器处理；raw形式，字面量不经过编译器解析。后者只适用于整型和浮点类型。本案例将介绍如何创建raw用户定义字面量。


%%C020234 h3

\subsection{准备工作}


%%C020235 p

在继续本案例之前，强烈建议你先阅读上一个案例\textit{创建cooked用户定义字面量}，因为这里不再重复用户定义字面量的一般细节。


%%C020236 p

为了演示raw用户定义字面量的创建方式，我们将定义二进制字面量。这些二进制字面量可以是8位、16位和32位（无符号）类型。这些类型称为\texttt{byte8}、\texttt{byte16}和\texttt{byte32}，我们要创建的字面量称为\texttt{\_\allowbreak{}b8}、\texttt{\_\allowbreak{}b16}和\texttt{\_\allowbreak{}b32}。


%%C020237 h3

\subsection{具体做法}


%%C020238 p

创建raw用户定义字面量应遵循以下步骤：


%%C020239 p

把字面量定义在独立的命名空间中，以避免名称冲突。


%%C020240 p

用户定义后缀一律以下划线（\texttt{\_}）开头。


%%C020241 p

定义如下形式的字面量运算符或字面量运算符模板：


%%C020242 code

\begin{codeblockbg}
\begin{Verbatim}[commandchars=\\\{\},breaklines=true,breakanywhere=true,fontsize=\codefont,xleftmargin=17bp]
\PYG{n}{T}\PYG{+w}{ }\PYG{k}{operator}\PYG{+w}{ }\PYG{l+s}{\PYGZdq{}}\PYG{l+s}{\PYGZdq{}}\PYG{+w}{ }\PYG{n}{\PYGZus{}suffix}\PYG{p}{(}\PYG{k}{const}\PYG{+w}{ }\PYG{k+kt}{char}\PYG{o}{*}\PYG{p}{)}\PYG{p}{;}
\PYG{k}{template}\PYG{o}{\PYGZlt{}}\PYG{k+kt}{char}\PYG{p}{.}\PYG{p}{.}\PYG{p}{.}\PYG{o}{\PYGZgt{}}\PYG{+w}{ }\PYG{n}{T}\PYG{+w}{ }\PYG{k}{operator}\PYG{+w}{ }\PYG{l+s}{\PYGZdq{}}\PYG{l+s}{\PYGZdq{}}\PYG{+w}{ }\PYG{n}{\PYGZus{}suffix}\PYG{p}{(}\PYG{p}{)}\PYG{p}{;}
\end{Verbatim}
\end{codeblockbg}


%%C020243 p

下面的示例给出8位、16位和32位二进制字面量的一种可能实现：


%%C020244 code

\begin{codeblockbg}
\begin{Verbatim}[commandchars=\\\{\},breaklines=true,breakanywhere=true,fontsize=\codefont,xleftmargin=17bp]
\PYG{k}{namespace}\PYG{+w}{ }\PYG{n+nn}{binary}
\PYG{p}{\PYGZob{}}
\PYG{+w}{  }\PYG{k}{using}\PYG{+w}{ }\PYG{n}{byte8}\PYG{+w}{  }\PYG{o}{=}\PYG{+w}{ }\PYG{k+kt}{unsigned}\PYG{+w}{ }\PYG{k+kt}{char}\PYG{p}{;}
\PYG{+w}{  }\PYG{k}{using}\PYG{+w}{ }\PYG{n}{byte16}\PYG{+w}{ }\PYG{o}{=}\PYG{+w}{ }\PYG{k+kt}{unsigned}\PYG{+w}{ }\PYG{k+kt}{short}\PYG{p}{;}
\PYG{+w}{  }\PYG{k}{using}\PYG{+w}{ }\PYG{n}{byte32}\PYG{+w}{ }\PYG{o}{=}\PYG{+w}{ }\PYG{k+kt}{unsigned}\PYG{+w}{ }\PYG{k+kt}{int}\PYG{p}{;}
\PYG{k}{namespace}\PYG{+w}{ }\PYG{n+nn}{binary\PYGZus{}literals}
\PYG{p}{\PYGZob{}}
\PYG{+w}{  }\PYG{k}{namespace}\PYG{+w}{ }\PYG{n+nn}{binary\PYGZus{}literals\PYGZus{}internals}
\PYG{+w}{  }\PYG{p}{\PYGZob{}}
\PYG{+w}{    }\PYG{k}{template}\PYG{+w}{ }\PYG{o}{\PYGZlt{}}\PYG{k}{typename}\PYG{+w}{ }\PYG{n+nc}{CharT}\PYG{p}{,}\PYG{+w}{ }\PYG{k+kt}{char}\PYG{p}{.}\PYG{p}{.}\PYG{p}{.}\PYG{+w}{ }\PYG{n}{bits}\PYG{o}{\PYGZgt{}}
\PYG{+w}{    }\PYG{k}{struct}\PYG{+w}{ }\PYG{n+nc}{binary\PYGZus{}struct}\PYG{p}{;}
\PYG{k}{template}\PYG{+w}{ }\PYG{o}{\PYGZlt{}}\PYG{k}{typename}\PYG{+w}{ }\PYG{n+nc}{CharT}\PYG{p}{,}\PYG{+w}{ }\PYG{k+kt}{char}\PYG{p}{.}\PYG{p}{.}\PYG{p}{.}\PYG{+w}{ }\PYG{n}{bits}\PYG{o}{\PYGZgt{}}
\PYG{k}{struct}\PYG{+w}{ }\PYG{n+nc}{binary\PYGZus{}struct}\PYG{o}{\PYGZlt{}}\PYG{n}{CharT}\PYG{p}{,}\PYG{+w}{ }\PYG{l+s+sc}{\PYGZsq{}}\PYG{l+s+sc}{0}\PYG{l+s+sc}{\PYGZsq{}}\PYG{p}{,}\PYG{+w}{ }\PYG{n}{bits}\PYG{p}{.}\PYG{p}{.}\PYG{p}{.}\PYG{o}{\PYGZgt{}}
\PYG{p}{\PYGZob{}}
\PYG{+w}{  }\PYG{k}{static}\PYG{+w}{ }\PYG{k}{constexpr}\PYG{+w}{ }\PYG{n}{CharT}\PYG{+w}{ }\PYG{n}{value}\PYG{p}{\PYGZob{}}
\PYG{+w}{    }\PYG{n}{binary\PYGZus{}struct}\PYG{o}{\PYGZlt{}}\PYG{n}{CharT}\PYG{p}{,}\PYG{+w}{ }\PYG{n}{bits}\PYG{p}{.}\PYG{p}{.}\PYG{p}{.}\PYG{o}{\PYGZgt{}}\PYG{o}{:}\PYG{o}{:}\PYG{n}{value}\PYG{+w}{ }\PYG{p}{\PYGZcb{}}\PYG{p}{;}
\PYG{p}{\PYGZcb{}}\PYG{p}{;}
\PYG{k}{template}\PYG{+w}{ }\PYG{o}{\PYGZlt{}}\PYG{k}{typename}\PYG{+w}{ }\PYG{n+nc}{CharT}\PYG{p}{,}\PYG{+w}{ }\PYG{k+kt}{char}\PYG{p}{.}\PYG{p}{.}\PYG{p}{.}\PYG{+w}{ }\PYG{n}{bits}\PYG{o}{\PYGZgt{}}
\PYG{k}{struct}\PYG{+w}{ }\PYG{n+nc}{binary\PYGZus{}struct}\PYG{o}{\PYGZlt{}}\PYG{n}{CharT}\PYG{p}{,}\PYG{+w}{ }\PYG{l+s+sc}{\PYGZsq{}}\PYG{l+s+sc}{1}\PYG{l+s+sc}{\PYGZsq{}}\PYG{p}{,}\PYG{+w}{ }\PYG{n}{bits}\PYG{p}{.}\PYG{p}{.}\PYG{p}{.}\PYG{o}{\PYGZgt{}}
\PYG{p}{\PYGZob{}}
\PYG{+w}{  }\PYG{k}{static}\PYG{+w}{ }\PYG{k}{constexpr}\PYG{+w}{ }\PYG{n}{CharT}\PYG{+w}{ }\PYG{n}{value}\PYG{p}{\PYGZob{}}
\PYG{+w}{    }\PYG{k}{static\PYGZus{}cast}\PYG{o}{\PYGZlt{}}\PYG{n}{CharT}\PYG{o}{\PYGZgt{}}\PYG{p}{(}\PYG{l+m+mi}{1}\PYG{+w}{ }\PYG{o}{\PYGZlt{}}\PYG{o}{\PYGZlt{}}\PYG{+w}{ }\PYG{k}{sizeof}\PYG{p}{.}\PYG{p}{.}\PYG{p}{.}\PYG{p}{(}\PYG{n}{bits}\PYG{p}{)}\PYG{p}{)}\PYG{+w}{ }\PYG{o}{|}
\PYG{+w}{    }\PYG{n}{binary\PYGZus{}struct}\PYG{o}{\PYGZlt{}}\PYG{n}{CharT}\PYG{p}{,}\PYG{+w}{ }\PYG{n}{bits}\PYG{p}{.}\PYG{p}{.}\PYG{p}{.}\PYG{o}{\PYGZgt{}}\PYG{o}{:}\PYG{o}{:}\PYG{n}{value}\PYG{+w}{ }\PYG{p}{\PYGZcb{}}\PYG{p}{;}
\PYG{p}{\PYGZcb{}}\PYG{p}{;}
\PYG{+w}{  }\PYG{k}{template}\PYG{+w}{ }\PYG{o}{\PYGZlt{}}\PYG{k}{typename}\PYG{+w}{ }\PYG{n+nc}{CharT}\PYG{o}{\PYGZgt{}}
\PYG{+w}{  }\PYG{k}{struct}\PYG{+w}{ }\PYG{n+nc}{binary\PYGZus{}struct}\PYG{o}{\PYGZlt{}}\PYG{n}{CharT}\PYG{o}{\PYGZgt{}}
\PYG{+w}{  }\PYG{p}{\PYGZob{}}
\PYG{+w}{    }\PYG{k}{static}\PYG{+w}{ }\PYG{k}{constexpr}\PYG{+w}{ }\PYG{n}{CharT}\PYG{+w}{ }\PYG{n}{value}\PYG{p}{\PYGZob{}}\PYG{+w}{ }\PYG{l+m+mi}{0}\PYG{+w}{ }\PYG{p}{\PYGZcb{}}\PYG{p}{;}
\PYG{+w}{  }\PYG{p}{\PYGZcb{}}\PYG{p}{;}
\PYG{p}{\PYGZcb{}}
\PYG{k}{template}\PYG{o}{\PYGZlt{}}\PYG{k+kt}{char}\PYG{p}{.}\PYG{p}{.}\PYG{p}{.}\PYG{+w}{ }\PYG{n}{bits}\PYG{o}{\PYGZgt{}}
\PYG{k}{constexpr}\PYG{+w}{ }\PYG{n}{byte8}\PYG{+w}{ }\PYG{k}{operator}\PYG{l+s}{\PYGZdq{}}\PYG{l+s}{\PYGZdq{}}\PYG{n}{\PYGZus{}b8}\PYG{p}{(}\PYG{p}{)}
\PYG{p}{\PYGZob{}}
\PYG{+w}{  }\PYG{k}{static\PYGZus{}assert}\PYG{p}{(}
\PYG{+w}{    }\PYG{k}{sizeof}\PYG{p}{.}\PYG{p}{.}\PYG{p}{.}\PYG{p}{(}\PYG{n}{bits}\PYG{p}{)}\PYG{+w}{ }\PYG{o}{\PYGZlt{}}\PYG{o}{=}\PYG{+w}{ }\PYG{l+m+mi}{8}\PYG{p}{,}
\PYG{+w}{    }\PYG{l+s}{\PYGZdq{}}\PYG{l+s}{binary literal b8 must be up to 8 digits long}\PYG{l+s}{\PYGZdq{}}\PYG{p}{)}\PYG{p}{;}
\PYG{+w}{  }\PYG{k}{return}\PYG{+w}{ }\PYG{n}{binary\PYGZus{}literals\PYGZus{}internals}\PYG{o}{:}\PYG{o}{:}
\PYG{+w}{            }\PYG{n}{binary\PYGZus{}struct}\PYG{o}{\PYGZlt{}}\PYG{n}{byte8}\PYG{p}{,}\PYG{+w}{ }\PYG{n}{bits}\PYG{p}{.}\PYG{p}{.}\PYG{p}{.}\PYG{o}{\PYGZgt{}}\PYG{o}{:}\PYG{o}{:}\PYG{n}{value}\PYG{p}{;}
\PYG{p}{\PYGZcb{}}
\PYG{k}{template}\PYG{o}{\PYGZlt{}}\PYG{k+kt}{char}\PYG{p}{.}\PYG{p}{.}\PYG{p}{.}\PYG{+w}{ }\PYG{n}{bits}\PYG{o}{\PYGZgt{}}
\PYG{k}{constexpr}\PYG{+w}{ }\PYG{n}{byte16}\PYG{+w}{ }\PYG{k}{operator}\PYG{l+s}{\PYGZdq{}}\PYG{l+s}{\PYGZdq{}}\PYG{n}{\PYGZus{}b16}\PYG{p}{(}\PYG{p}{)}
\PYG{p}{\PYGZob{}}
\PYG{+w}{  }\PYG{k}{static\PYGZus{}assert}\PYG{p}{(}
\PYG{+w}{    }\PYG{k}{sizeof}\PYG{p}{.}\PYG{p}{.}\PYG{p}{.}\PYG{p}{(}\PYG{n}{bits}\PYG{p}{)}\PYG{+w}{ }\PYG{o}{\PYGZlt{}}\PYG{o}{=}\PYG{+w}{ }\PYG{l+m+mi}{16}\PYG{p}{,}
\PYG{+w}{    }\PYG{l+s}{\PYGZdq{}}\PYG{l+s}{binary literal b16 must be up to 16 digits long}\PYG{l+s}{\PYGZdq{}}\PYG{p}{)}\PYG{p}{;}
\PYG{+w}{  }\PYG{k}{return}\PYG{+w}{ }\PYG{n}{binary\PYGZus{}literals\PYGZus{}internals}\PYG{o}{:}\PYG{o}{:}
\PYG{+w}{            }\PYG{n}{binary\PYGZus{}struct}\PYG{o}{\PYGZlt{}}\PYG{n}{byte16}\PYG{p}{,}\PYG{+w}{ }\PYG{n}{bits}\PYG{p}{.}\PYG{p}{.}\PYG{p}{.}\PYG{o}{\PYGZgt{}}\PYG{o}{:}\PYG{o}{:}\PYG{n}{value}\PYG{p}{;}
\PYG{p}{\PYGZcb{}}
\PYG{k}{template}\PYG{o}{\PYGZlt{}}\PYG{k+kt}{char}\PYG{p}{.}\PYG{p}{.}\PYG{p}{.}\PYG{+w}{ }\PYG{n}{bits}\PYG{o}{\PYGZgt{}}
\PYG{k}{constexpr}\PYG{+w}{ }\PYG{n}{byte32}\PYG{+w}{ }\PYG{k}{operator}\PYG{l+s}{\PYGZdq{}}\PYG{l+s}{\PYGZdq{}}\PYG{n}{\PYGZus{}b32}\PYG{p}{(}\PYG{p}{)}
\PYG{p}{\PYGZob{}}
\PYG{+w}{  }\PYG{k}{static\PYGZus{}assert}\PYG{p}{(}
\PYG{+w}{     }\PYG{k}{sizeof}\PYG{p}{.}\PYG{p}{.}\PYG{p}{.}\PYG{p}{(}\PYG{n}{bits}\PYG{p}{)}\PYG{+w}{ }\PYG{o}{\PYGZlt{}}\PYG{o}{=}\PYG{+w}{ }\PYG{l+m+mi}{32}\PYG{p}{,}
\PYG{+w}{     }\PYG{l+s}{\PYGZdq{}}\PYG{l+s}{binary literal b32 must be up to 32 digits long}\PYG{l+s}{\PYGZdq{}}\PYG{p}{)}\PYG{p}{;}
\PYG{+w}{  }\PYG{k}{return}\PYG{+w}{ }\PYG{n}{binary\PYGZus{}literals\PYGZus{}internals}\PYG{o}{:}\PYG{o}{:}
\PYG{+w}{            }\PYG{n}{binary\PYGZus{}struct}\PYG{o}{\PYGZlt{}}\PYG{n}{byte32}\PYG{p}{,}\PYG{+w}{ }\PYG{n}{bits}\PYG{p}{.}\PYG{p}{.}\PYG{p}{.}\PYG{o}{\PYGZgt{}}\PYG{o}{:}\PYG{o}{:}\PYG{n}{value}\PYG{p}{;}
\PYG{p}{\PYGZcb{}}
\PYG{+w}{  }\PYG{p}{\PYGZcb{}}
\PYG{p}{\PYGZcb{}}
\end{Verbatim}
\end{codeblockbg}


%%C020245 h3

\subsection{工作原理}


%%C020246 p

上一节的实现使我们能够定义形如1010\_b8（十进制10的\texttt{byte8}值）或000010101100\_b16（十进制2130496的\texttt{byte16}值）的二进制字面量。不过，我们要确保不超过每种类型的位数。换言之，像111100001\_b8这样的值应当非法，编译器应当报错。


%%C020247 p

首先，我们把所有内容定义在名为\texttt{binary}的命名空间中，并先引入几个类型别名（\texttt{byte8}、\texttt{byte16}和\texttt{byte32}）。


%%C020248 p

字面量运算符模板定义在名为\myttbr{binary\_\allowbreak{}literal\_\allowbreak{}internals}的嵌套命名空间中。这是避免与其他命名空间的字面量运算符发生名称冲突的好做法。万一真的发生冲突，你可以在恰当的作用域中选用合适的命名空间（例如在一个函数或代码块中使用一个命名空间，在另一个函数或代码块中使用另一个命名空间）。


%%C020249 p

三个字面量运算符模板非常相似，不同的只有名称（\texttt{\_\allowbreak{}b8}、\texttt{\_\allowbreak{}16}和\texttt{\_\allowbreak{}b32}）、返回类型（\texttt{byte8}、\texttt{byte16}和\texttt{byte32}），以及静态断言中检查位数的条件。


%%C020250 p

我们会在后面的案例中详细探讨可变参数模板与模板递归；不过为了便于理解，这里先说明这一实现的工作原理：\texttt{bits}是一个模板参数包，它不是单个值，而是该模板可能被实例化时用到的所有值。例如，对于字面量\texttt{1010\_\allowbreak{}b8}，字面量运算符模板会实例化为\myttbr{operator“” \_\allowbreak{}b8<\allowbreak{}‘1’, ‘0’, ‘1’, ‘0’>()}。在计算二进制值之前，先检查字面量的位数。对于\texttt{\_\allowbreak{}b8}，位数不能超过八位（包括末尾的零）。类似地，\texttt{\_\allowbreak{}b16}最多16位，\texttt{\_\allowbreak{}b32}最多32位。为此我们使用\texttt{sizeof...}运算符，它返回参数包中元素的数量（这里是\texttt{bits}）。


%%C020251 p

如果位数正确，就可以展开参数包，递归计算该二进制字面量所表示的十进制值。这借助一个额外的类模板及其特化完成。这些模板定义在另一个嵌套命名空间\myttbr{binary\_\allowbreak{}literals\_\allowbreak{}internals}中。这同样是好的做法，因为它把实现细节向使用者隐藏（除非使用者做了恰当的限定），除非有显式的using namespace指令将它们引入当前命名空间。

%%C020252 note

\begin{myNote}
尽管这看起来像递归，但它并不是真正的运行期递归，因为编译器展开并由模板生成代码之后，最终得到的实际上是对参数个数不同的重载函数的调用。这一点在后面\textit{编写参数个数可变的函数模板}案例中还会进一步说明。
\end{myNote}


%%C020253 p

\texttt{binary\_\allowbreak{}struct}类模板有一个表示函数返回类型的模板类型\texttt{CharT}（之所以需要它，是因为字面量运算符模板应当返回\texttt{byte8}、\texttt{byte16}或\texttt{byte32}），以及一个参数包：


%%C020254 code

\begin{codeblockbg}
\begin{Verbatim}[commandchars=\\\{\},breaklines=true,breakanywhere=true,fontsize=\codefont,xleftmargin=17bp]
\PYG{k}{template}\PYG{+w}{ }\PYG{o}{\PYGZlt{}}\PYG{k}{typename}\PYG{+w}{ }\PYG{n+nc}{CharT}\PYG{p}{,}\PYG{+w}{ }\PYG{k+kt}{char}\PYG{p}{.}\PYG{p}{.}\PYG{p}{.}\PYG{+w}{ }\PYG{n}{bits}\PYG{o}{\PYGZgt{}}
\PYG{k}{struct}\PYG{+w}{ }\PYG{n+nc}{binary\PYGZus{}struct}\PYG{p}{;}
\end{Verbatim}
\end{codeblockbg}


%%C020255 p

这个类模板提供了若干带参数包分解的特化。当参数包的第一位是‘0’时，计算出的值保持不变，我们继续展开参数包的其余部分。如果参数包的第一位是‘1’，则新值是1左移参数包剩余部分的位数、再与参数包其余部分的值相或的结果：


%%C020256 code

\begin{codeblockbg}
\begin{Verbatim}[commandchars=\\\{\},breaklines=true,breakanywhere=true,fontsize=\codefont,xleftmargin=17bp]
\PYG{k}{template}\PYG{+w}{ }\PYG{o}{\PYGZlt{}}\PYG{k}{typename}\PYG{+w}{ }\PYG{n+nc}{CharT}\PYG{p}{,}\PYG{+w}{ }\PYG{k+kt}{char}\PYG{p}{.}\PYG{p}{.}\PYG{p}{.}\PYG{+w}{ }\PYG{n}{bits}\PYG{o}{\PYGZgt{}}
\PYG{k}{struct}\PYG{+w}{ }\PYG{n+nc}{binary\PYGZus{}struct}\PYG{o}{\PYGZlt{}}\PYG{n}{CharT}\PYG{p}{,}\PYG{+w}{ }\PYG{l+s+sc}{\PYGZsq{}}\PYG{l+s+sc}{0}\PYG{l+s+sc}{\PYGZsq{}}\PYG{p}{,}\PYG{+w}{ }\PYG{n}{bits}\PYG{p}{.}\PYG{p}{.}\PYG{p}{.}\PYG{o}{\PYGZgt{}}
\PYG{p}{\PYGZob{}}
\PYG{+w}{  }\PYG{k}{static}\PYG{+w}{ }\PYG{k}{constexpr}\PYG{+w}{ }\PYG{n}{CharT}\PYG{+w}{ }\PYG{n}{value}\PYG{p}{\PYGZob{}}
\PYG{+w}{    }\PYG{n}{binary\PYGZus{}struct}\PYG{o}{\PYGZlt{}}\PYG{n}{CharT}\PYG{p}{,}\PYG{+w}{ }\PYG{n}{bits}\PYG{p}{.}\PYG{p}{.}\PYG{p}{.}\PYG{o}{\PYGZgt{}}\PYG{o}{:}\PYG{o}{:}\PYG{n}{value}\PYG{+w}{ }\PYG{p}{\PYGZcb{}}\PYG{p}{;}
\PYG{p}{\PYGZcb{}}\PYG{p}{;}
\PYG{k}{template}\PYG{+w}{ }\PYG{o}{\PYGZlt{}}\PYG{k}{typename}\PYG{+w}{ }\PYG{n+nc}{CharT}\PYG{p}{,}\PYG{+w}{ }\PYG{k+kt}{char}\PYG{p}{.}\PYG{p}{.}\PYG{p}{.}\PYG{+w}{ }\PYG{n}{bits}\PYG{o}{\PYGZgt{}}
\PYG{k}{struct}\PYG{+w}{ }\PYG{n+nc}{binary\PYGZus{}struct}\PYG{o}{\PYGZlt{}}\PYG{n}{CharT}\PYG{p}{,}\PYG{+w}{ }\PYG{l+s+sc}{\PYGZsq{}}\PYG{l+s+sc}{1}\PYG{l+s+sc}{\PYGZsq{}}\PYG{p}{,}\PYG{+w}{ }\PYG{n}{bits}\PYG{p}{.}\PYG{p}{.}\PYG{p}{.}\PYG{o}{\PYGZgt{}}
\PYG{p}{\PYGZob{}}
\PYG{+w}{  }\PYG{k}{static}\PYG{+w}{ }\PYG{k}{constexpr}\PYG{+w}{ }\PYG{n}{CharT}\PYG{+w}{ }\PYG{n}{value}\PYG{p}{\PYGZob{}}
\PYG{+w}{    }\PYG{k}{static\PYGZus{}cast}\PYG{o}{\PYGZlt{}}\PYG{n}{CharT}\PYG{o}{\PYGZgt{}}\PYG{p}{(}\PYG{l+m+mi}{1}\PYG{+w}{ }\PYG{o}{\PYGZlt{}}\PYG{o}{\PYGZlt{}}\PYG{+w}{ }\PYG{k}{sizeof}\PYG{p}{.}\PYG{p}{.}\PYG{p}{.}\PYG{p}{(}\PYG{n}{bits}\PYG{p}{)}\PYG{p}{)}\PYG{+w}{ }\PYG{o}{|}
\PYG{+w}{    }\PYG{n}{binary\PYGZus{}struct}\PYG{o}{\PYGZlt{}}\PYG{n}{CharT}\PYG{p}{,}\PYG{+w}{ }\PYG{n}{bits}\PYG{p}{.}\PYG{p}{.}\PYG{p}{.}\PYG{o}{\PYGZgt{}}\PYG{o}{:}\PYG{o}{:}\PYG{n}{value}\PYG{+w}{ }\PYG{p}{\PYGZcb{}}\PYG{p}{;}
\PYG{p}{\PYGZcb{}}\PYG{p}{;}
\end{Verbatim}
\end{codeblockbg}


%%C020257 p

最后一个特化处理参数包为空的情况；此时返回0：


%%C020258 code

\begin{codeblockbg}
\begin{Verbatim}[commandchars=\\\{\},breaklines=true,breakanywhere=true,fontsize=\codefont,xleftmargin=17bp]
\PYG{k}{template}\PYG{+w}{ }\PYG{o}{\PYGZlt{}}\PYG{k}{typename}\PYG{+w}{ }\PYG{n+nc}{CharT}\PYG{o}{\PYGZgt{}}
\PYG{k}{struct}\PYG{+w}{ }\PYG{n+nc}{binary\PYGZus{}struct}\PYG{o}{\PYGZlt{}}\PYG{n}{CharT}\PYG{o}{\PYGZgt{}}
\PYG{p}{\PYGZob{}}
\PYG{+w}{  }\PYG{k}{static}\PYG{+w}{ }\PYG{k}{constexpr}\PYG{+w}{ }\PYG{n}{CharT}\PYG{+w}{ }\PYG{n}{value}\PYG{p}{\PYGZob{}}\PYG{+w}{ }\PYG{l+m+mi}{0}\PYG{+w}{ }\PYG{p}{\PYGZcb{}}\PYG{p}{;}
\PYG{p}{\PYGZcb{}}\PYG{p}{;}
\end{Verbatim}
\end{codeblockbg}


%%C020259 p

定义好这些辅助类之后，就可以按预期实现\texttt{byte8}、\texttt{byte16}和\texttt{byte32}二进制字面量了。注意，为了使用字面量运算符模板，需要把命名空间\texttt{binary\_\allowbreak{}literals}的内容引入当前命名空间：


%%C020260 code

\begin{codeblockbg}
\begin{Verbatim}[commandchars=\\\{\},breaklines=true,breakanywhere=true,fontsize=\codefont,xleftmargin=17bp]
\PYG{k}{using}\PYG{+w}{ }\PYG{k}{namespace}\PYG{+w}{ }\PYG{n+nn}{binary}\PYG{p}{;}
\PYG{k}{using}\PYG{+w}{ }\PYG{k}{namespace}\PYG{+w}{ }\PYG{n+nn}{binary\PYGZus{}literals}\PYG{p}{;}
\PYG{k}{auto}\PYG{+w}{ }\PYG{n}{b1}\PYG{+w}{ }\PYG{o}{=}\PYG{+w}{ }\PYG{l+m+mi}{1010}\PYG{n}{\PYGZus{}b8}\PYG{p}{;}
\PYG{k}{auto}\PYG{+w}{ }\PYG{n}{b2}\PYG{+w}{ }\PYG{o}{=}\PYG{+w}{ }\PYG{l+m+mi}{101010101010}\PYG{n}{\PYGZus{}b16}\PYG{p}{;}
\PYG{k}{auto}\PYG{+w}{ }\PYG{n}{b3}\PYG{+w}{ }\PYG{o}{=}\PYG{+w}{ }\PYG{l+m+mi}{101010101010101010101010}\PYG{n}{\PYGZus{}b32}\PYG{p}{;}
\end{Verbatim}
\end{codeblockbg}


%%C020261 p

下面的定义会触发编译器错误，因为\texttt{static\_\allowbreak{}assert}中的条件不满足：


%%C020262 code

\begin{codeblockbg}
\begin{Verbatim}[commandchars=\\\{\},breaklines=true,breakanywhere=true,fontsize=\codefont,xleftmargin=17bp]
\PYG{c+c1}{// b8 二进制字面量最多 8 位数字}
\PYG{k}{auto}\PYG{+w}{ }\PYG{n}{b4}\PYG{+w}{ }\PYG{o}{=}\PYG{+w}{ }\PYG{l+m+mo}{0011111111}\PYG{n}{\PYGZus{}b8}\PYG{p}{;}
\PYG{c+c1}{// b16 二进制字面量最多 16 位数字}
\PYG{k}{auto}\PYG{+w}{ }\PYG{n}{b5}\PYG{+w}{ }\PYG{o}{=}\PYG{+w}{ }\PYG{l+m+mo}{001111111111111111}\PYG{n}{\PYGZus{}b16}\PYG{p}{;}
\PYG{c+c1}{// b32 二进制字面量最多 32 位数字}
\PYG{k}{auto}\PYG{+w}{ }\PYG{n}{b6}\PYG{+w}{ }\PYG{o}{=}\PYG{+w}{ }\PYG{l+m+mo}{0011111111111111111111111111111111}\PYG{n}{\PYGZus{}b32}\PYG{p}{;}
\end{Verbatim}
\end{codeblockbg}


%%C020263 h3

\subsection{延伸阅读}


%%C020264 li

\begin{itemize}
\item \textit{用原始字符串字面量避免转义字符}s


%%C020265 li

\item \textit{创建cooked用户定义字面量}


%%C020266 li

\item \textit{编写参数个数可变的函数模板}案例，见\texttt{第3章}，\textit{函数探究}


%%C020267 li

\item \textit{创建类型别名与别名模板}案例，见\texttt{第1章}，\textit{现代C++语言} \textit{核心特性}


\end{itemize}


%%C020268 h2

\section{用原始字符串字面量避免转义字符}


%%C020269 p

字符串中可能包含特殊字符，例如不可打印字符（换行符、水平制表符和垂直制表符等）、字符串和字符的分隔符（双引号和单引号），或者任意的八进制、十六进制或Unicode值。这些特殊字符用转义序列引入，转义序列以反斜杠开头，后面跟字符本身（例如\texttt{‘}和\texttt{“}）、其指定字母（例如表示换行的\texttt{n}、表示水平制表的\texttt{t}），或者其值（例如八进制050、十六进制XF7或Unicode U16F0）。因此，反斜杠字符本身必须用另一个反斜杠转义。这使得字面量字符串更加复杂，也更难阅读。


%%C020270 p

为了避免转义字符，C++11引入了不处理转义序列的原始字符串字面量。本案例将介绍如何使用各种形式的原始字符串字面量。


%%C020271 h3

\subsection{准备工作}


%%C020272 p

在本案例以及本书其余部分中，我将使用\texttt{s}后缀来定义\texttt{basic\_\allowbreak{}string}字面量。这一点已在\textit{创建cooked用户定义} \textit{字面量}案例中介绍过。


%%C020273 h3

\subsection{具体做法}


%%C020274 p

为了避免转义字符，可以按下面的方式定义字符串字面量：


%%C020275 p

\texttt{R“( literal )\allowbreak{}”}为默认形式：


%%C020276 code

\begin{codeblockbg}
\begin{Verbatim}[commandchars=\\\{\},breaklines=true,breakanywhere=true,fontsize=\codefont,xleftmargin=17bp]
\PYG{k}{auto}\PYG{+w}{ }\PYG{n}{filename}\PYG{+w}{ }\PYG{p}{\PYGZob{}}\PYG{l+s+sa}{R}\PYG{l+s}{\PYGZdq{}}\PYG{l+s+dl}{(}\PYG{l+s}{C:\PYGZbs{}Users\PYGZbs{}Marius\PYGZbs{}Documents\PYGZbs{}}\PYG{l+s+dl}{)}\PYG{l+s}{\PYGZdq{}}\PYG{n}{s}\PYG{p}{\PYGZcb{}}\PYG{p}{;}
\PYG{k}{auto}\PYG{+w}{ }\PYG{n}{pattern}\PYG{+w}{ }\PYG{p}{\PYGZob{}}\PYG{l+s+sa}{R}\PYG{l+s}{\PYGZdq{}}\PYG{l+s+dl}{(}\PYG{l+s}{(\PYGZbs{}w+)=(\PYGZbs{}d+)\PYGZdl{}}\PYG{l+s+dl}{)}\PYG{l+s}{\PYGZdq{}}\PYG{n}{s}\PYG{p}{\PYGZcb{}}\PYG{p}{;}
\PYG{k}{auto}\PYG{+w}{ }\PYG{n}{sqlselect}\PYG{+w}{ }\PYG{p}{\PYGZob{}}
\PYG{+w}{  }\PYG{l+s+sa}{R}\PYG{l+s}{\PYGZdq{}}\PYG{l+s+dl}{(}\PYG{l+s}{SELECT *}
\PYG{l+s}{  FROM Books}
\PYG{l+s}{  WHERE Publisher=\PYGZsq{}Paktpub\PYGZsq{}}
\PYG{l+s}{  ORDER BY PubDate DESC}\PYG{l+s+dl}{)}\PYG{l+s}{\PYGZdq{}}\PYG{n}{s}\PYG{p}{\PYGZcb{}}\PYG{p}{;}
\end{Verbatim}
\end{codeblockbg}


%%C020277 p

\myttbr{R“delimiter( literal )\allowbreak{}delimiter”}，其中\texttt{delimiter}是任意字符


%%C020278 p

序列；当序列\texttt{)\allowbreak{}“}确实应当是字符串的一部分时，该字符序列不能出现在实际字符串中。下面是一个以\texttt{!!}作为分隔符的例子：


%%C020279 code

\begin{codeblockbg}
\begin{Verbatim}[commandchars=\\\{\},breaklines=true,breakanywhere=true,fontsize=\codefont,xleftmargin=17bp]
\PYG{k}{auto}\PYG{+w}{ }\PYG{n}{text}\PYG{p}{\PYGZob{}}\PYG{+w}{ }\PYG{l+s+sa}{R}\PYG{l+s}{\PYGZdq{}}\PYG{l+s+dl}{!!}\PYG{l+s+dl}{(}\PYG{l+s}{This text contains both \PYGZdq{}( and )\PYGZdq{}.}\PYG{l+s+dl}{)!!}\PYG{l+s}{\PYGZdq{}}\PYG{n}{s}\PYG{+w}{ }\PYG{p}{\PYGZcb{}}\PYG{p}{;}
\PYG{n}{std}\PYG{o}{:}\PYG{o}{:}\PYG{n}{cout}\PYG{+w}{ }\PYG{o}{\PYGZlt{}}\PYG{o}{\PYGZlt{}}\PYG{+w}{ }\PYG{n}{text}\PYG{+w}{ }\PYG{o}{\PYGZlt{}}\PYG{o}{\PYGZlt{}}\PYG{+w}{ }\PYG{n}{std}\PYG{o}{:}\PYG{o}{:}\PYG{n}{endl}\PYG{p}{;}
\end{Verbatim}
\end{codeblockbg}


%%C020280 h3

\subsection{工作原理}


%%C020281 p

使用原始字符串字面量时，转义序列不会被处理，字符串的实际内容就写在分隔符之间（也就是说，所见即所得）。下面的例子给出的原始字面量字符串看上去相同；不过第二个字符串仍然含有转义字符。由于使用原始字符串字面量时这些字符不会被处理，它们会原样输出到结果中：


%%C020282 code

\begin{codeblockbg}
\begin{Verbatim}[commandchars=\\\{\},breaklines=true,breakanywhere=true,fontsize=\codefont,xleftmargin=17bp]
\PYG{k}{auto}\PYG{+w}{ }\PYG{n}{filename1}\PYG{+w}{ }\PYG{p}{\PYGZob{}}\PYG{l+s+sa}{R}\PYG{l+s}{\PYGZdq{}}\PYG{l+s+dl}{(}\PYG{l+s}{C:\PYGZbs{}Users\PYGZbs{}Marius\PYGZbs{}Documents\PYGZbs{}}\PYG{l+s+dl}{)}\PYG{l+s}{\PYGZdq{}}\PYG{n}{s}\PYG{p}{\PYGZcb{}}\PYG{p}{;}
\PYG{k}{auto}\PYG{+w}{ }\PYG{n}{filename2}\PYG{+w}{ }\PYG{p}{\PYGZob{}}\PYG{l+s+sa}{R}\PYG{l+s}{\PYGZdq{}}\PYG{l+s+dl}{(}\PYG{l+s}{C:\PYGZbs{}\PYGZbs{}Users\PYGZbs{}\PYGZbs{}Marius\PYGZbs{}\PYGZbs{}Documents\PYGZbs{}\PYGZbs{}}\PYG{l+s+dl}{)}\PYG{l+s}{\PYGZdq{}}\PYG{n}{s}\PYG{p}{\PYGZcb{}}\PYG{p}{;}
\PYG{c+c1}{// 输出 C:\PYGZbs{}Users\PYGZbs{}Marius\PYGZbs{}Documents\PYGZbs{}}
\PYG{c+c1}{std::cout \PYGZlt{}\PYGZlt{} filename1 \PYGZlt{}\PYGZlt{} std::endl;}
\PYG{c+c1}{// 输出 C:\PYGZbs{}\PYGZbs{}Users\PYGZbs{}\PYGZbs{}Marius\PYGZbs{}\PYGZbs{}Documents\PYGZbs{}\PYGZbs{}}
\PYG{c+c1}{std::cout \PYGZlt{}\PYGZlt{} filename2 \PYGZlt{}\PYGZlt{} std::endl;}
\end{Verbatim}
\end{codeblockbg}


%%C020283 p

如果文本中必须包含\texttt{)\allowbreak{}“}序列，就必须使用不同的分隔符，即\myttbr{R“delimiter( literal )\allowbreak{}delimiter”}形式。按照标准，分隔符中允许的字符是：


%%C020284 p

\textit{基本源字符集中的任意成员，但空格、左圆括号（右} \textit{圆括号即“）”、反斜杠\textbackslash{}，以及表示水平制表、} \textit{垂直制表、换页和换行的控制字符除外。}


%%C020285 p

原始字符串字面量可以带\texttt{L}、\texttt{u8}、\texttt{u}、\texttt{U}之一作为前缀，以表示宽字符、UTF-8、UTF-16或UTF-32字符串字面量。下面是这类字符串字面量的例子。注意，字符串末尾出现字符串字面量\texttt{operator “”s}会让编译器把类型推导为各种字符串类，而不是字符数组：


%%C020286 code

\begin{codeblockbg}
\begin{Verbatim}[commandchars=\\\{\},breaklines=true,breakanywhere=true,fontsize=\codefont,xleftmargin=17bp]
\PYG{k}{auto}\PYG{+w}{ }\PYG{n}{t1}\PYG{p}{\PYGZob{}}\PYG{+w}{ }\PYG{l+s+sa}{LR}\PYG{l+s}{\PYGZdq{}}\PYG{l+s+dl}{(}\PYG{l+s}{text}\PYG{l+s+dl}{)}\PYG{l+s}{\PYGZdq{}}\PYG{+w}{  }\PYG{p}{\PYGZcb{}}\PYG{p}{;}\PYG{+w}{  }\PYG{c+c1}{// const wchar\PYGZus{}t*}
\PYG{k}{auto}\PYG{+w}{ }\PYG{n}{t2}\PYG{p}{\PYGZob{}}\PYG{+w}{ }\PYG{l+s+sa}{u8R}\PYG{l+s}{\PYGZdq{}}\PYG{l+s+dl}{(}\PYG{l+s}{text}\PYG{l+s+dl}{)}\PYG{l+s}{\PYGZdq{}}\PYG{+w}{ }\PYG{p}{\PYGZcb{}}\PYG{p}{;}\PYG{+w}{  }\PYG{c+c1}{// const char*}
\PYG{k}{auto}\PYG{+w}{ }\PYG{n}{t3}\PYG{p}{\PYGZob{}}\PYG{+w}{ }\PYG{l+s+sa}{uR}\PYG{l+s}{\PYGZdq{}}\PYG{l+s+dl}{(}\PYG{l+s}{text}\PYG{l+s+dl}{)}\PYG{l+s}{\PYGZdq{}}\PYG{+w}{  }\PYG{p}{\PYGZcb{}}\PYG{p}{;}\PYG{+w}{  }\PYG{c+c1}{// const char16\PYGZus{}t*}
\PYG{k}{auto}\PYG{+w}{ }\PYG{n}{t4}\PYG{p}{\PYGZob{}}\PYG{+w}{ }\PYG{l+s+sa}{UR}\PYG{l+s}{\PYGZdq{}}\PYG{l+s+dl}{(}\PYG{l+s}{text}\PYG{l+s+dl}{)}\PYG{l+s}{\PYGZdq{}}\PYG{+w}{  }\PYG{p}{\PYGZcb{}}\PYG{p}{;}\PYG{+w}{  }\PYG{c+c1}{// const char32\PYGZus{}t*}
\PYG{k}{auto}\PYG{+w}{ }\PYG{n}{t5}\PYG{p}{\PYGZob{}}\PYG{+w}{ }\PYG{l+s+sa}{LR}\PYG{l+s}{\PYGZdq{}}\PYG{l+s+dl}{(}\PYG{l+s}{text}\PYG{l+s+dl}{)}\PYG{l+s}{\PYGZdq{}}\PYG{n}{s}\PYG{+w}{  }\PYG{p}{\PYGZcb{}}\PYG{p}{;}\PYG{+w}{ }\PYG{c+c1}{// wstring}
\PYG{k}{auto}\PYG{+w}{ }\PYG{n}{t6}\PYG{p}{\PYGZob{}}\PYG{+w}{ }\PYG{l+s+sa}{u8R}\PYG{l+s}{\PYGZdq{}}\PYG{l+s+dl}{(}\PYG{l+s}{text}\PYG{l+s+dl}{)}\PYG{l+s}{\PYGZdq{}}\PYG{n}{s}\PYG{+w}{ }\PYG{p}{\PYGZcb{}}\PYG{p}{;}\PYG{+w}{ }\PYG{c+c1}{// string}
\PYG{k}{auto}\PYG{+w}{ }\PYG{n}{t7}\PYG{p}{\PYGZob{}}\PYG{+w}{ }\PYG{l+s+sa}{uR}\PYG{l+s}{\PYGZdq{}}\PYG{l+s+dl}{(}\PYG{l+s}{text}\PYG{l+s+dl}{)}\PYG{l+s}{\PYGZdq{}}\PYG{n}{s}\PYG{+w}{  }\PYG{p}{\PYGZcb{}}\PYG{p}{;}\PYG{+w}{ }\PYG{c+c1}{// u16string}
\PYG{k}{auto}\PYG{+w}{ }\PYG{n}{t8}\PYG{p}{\PYGZob{}}\PYG{+w}{ }\PYG{l+s+sa}{UR}\PYG{l+s}{\PYGZdq{}}\PYG{l+s+dl}{(}\PYG{l+s}{text}\PYG{l+s+dl}{)}\PYG{l+s}{\PYGZdq{}}\PYG{n}{s}\PYG{+w}{  }\PYG{p}{\PYGZcb{}}\PYG{p}{;}\PYG{+w}{ }\PYG{c+c1}{// u32string}
\end{Verbatim}
\end{codeblockbg}


%%C020287 h3

\subsection{延伸阅读}


%%C020288 li

\begin{itemize}
\item \textit{创建cooked用户定义字面量}

\end{itemize}


%%C020289 h2

\section{构建字符串工具库}


%%C020290 p

标准库中的字符串类型是通用实现，缺少许多有用的方法，例如改变大小写、去除首尾字符、拆分，以及其他满足不同开发者需求的操作。确实存在提供丰富字符串功能的第三方库。不过，本案例将实现几个简单却很实用的方法，你在实践中经常需要用到它们。目的主要是看看字符串方法和标准通用算法如何用于操作字符串，同时也留下一份可在应用中复用的代码参考。


%%C020291 p

本案例将实现一个小型字符串工具库，提供以下函数：


%%C020292 li

\begin{itemize}
\item 把字符串转换为小写或大写。


%%C020293 li

\item 反转字符串。


%%C020294 li

\item 去除字符串开头和/或结尾的空白字符。


%%C020295 li

\item 去除字符串开头和/或结尾的一组指定字符。


%%C020296 li

\item 删除字符串中任意位置出现的某个字符。


%%C020297 li

\item 用指定分隔符把字符串切分为若干记号。

\end{itemize}


%%C020298 h3

\subsection{准备工作}


%%C020299 p

我们要实现的字符串库应当适用于所有标准字符串类型：\texttt{std::\allowbreak{}string}、\texttt{std::\allowbreak{}wstring}、\texttt{std::\allowbreak{}u16string}和\texttt{std::\allowbreak{}u32string}。为了避免写出\myttbr{std::\allowbreak{}basic\_\allowbreak{}string<\allowbreak{}CharT,} \myttbr{std::\allowbreak{}char\_\allowbreak{}traits<\allowbreak{}CharT>\allowbreak{}, std::\allowbreak{}allocator<\allowbreak{}CharT>>}这样冗长的名字，我们为字符串和字符串流使用下面的别名模板：


%%C020300 code

\begin{codeblockbg}
\begin{Verbatim}[commandchars=\\\{\},breaklines=true,breakanywhere=true,fontsize=\codefont,xleftmargin=17bp]
\PYG{k}{template}\PYG{+w}{ }\PYG{o}{\PYGZlt{}}\PYG{k}{typename}\PYG{+w}{ }\PYG{n+nc}{CharT}\PYG{o}{\PYGZgt{}}
\PYG{k}{using}\PYG{+w}{ }\PYG{n}{tstring}\PYG{+w}{ }\PYG{o}{=}
\PYG{+w}{   }\PYG{n}{std}\PYG{o}{:}\PYG{o}{:}\PYG{n}{basic\PYGZus{}string}\PYG{o}{\PYGZlt{}}\PYG{n}{CharT}\PYG{p}{,}\PYG{+w}{ }\PYG{n}{std}\PYG{o}{:}\PYG{o}{:}\PYG{n}{char\PYGZus{}traits}\PYG{o}{\PYGZlt{}}\PYG{n}{CharT}\PYG{o}{\PYGZgt{}}\PYG{p}{,}
\PYG{+w}{                     }\PYG{n}{std}\PYG{o}{:}\PYG{o}{:}\PYG{n}{allocator}\PYG{o}{\PYGZlt{}}\PYG{n}{CharT}\PYG{o}{\PYGZgt{}}\PYG{o}{\PYGZgt{}}\PYG{p}{;}
\PYG{k}{template}\PYG{+w}{ }\PYG{o}{\PYGZlt{}}\PYG{k}{typename}\PYG{+w}{ }\PYG{n+nc}{CharT}\PYG{o}{\PYGZgt{}}
\PYG{k}{using}\PYG{+w}{ }\PYG{n}{tstringstream}\PYG{+w}{ }\PYG{o}{=}
\PYG{+w}{   }\PYG{n}{std}\PYG{o}{:}\PYG{o}{:}\PYG{n}{basic\PYGZus{}stringstream}\PYG{o}{\PYGZlt{}}\PYG{n}{CharT}\PYG{p}{,}\PYG{+w}{ }\PYG{n}{std}\PYG{o}{:}\PYG{o}{:}\PYG{n}{char\PYGZus{}traits}\PYG{o}{\PYGZlt{}}\PYG{n}{CharT}\PYG{o}{\PYGZgt{}}\PYG{p}{,}
\PYG{+w}{                           }\PYG{n}{std}\PYG{o}{:}\PYG{o}{:}\PYG{n}{allocator}\PYG{o}{\PYGZlt{}}\PYG{n}{CharT}\PYG{o}{\PYGZgt{}}\PYG{o}{\PYGZgt{}}\PYG{p}{;}
\end{Verbatim}
\end{codeblockbg}


%%C020301 p

要实现这些字符串辅助函数，需要包含头文件\texttt{<\allowbreak{}string>}以获得字符串支持，包含\texttt{<\allowbreak{}algorithm>}以获得将要使用的通用标准算法。


%%C020302 p

本案例的所有示例都使用C++14中为字符串提供的标准用户定义字面量运算符，为此需要显式使用\myttbr{std::\allowbreak{}string\_\allowbreak{}literals}命名空间。


%%C020303 h3

\subsection{具体做法}


%%C020304 p

要把字符串转换为小写或大写，可以对字符串中的字符应用\texttt{tolower()}或


%%C020305 p

\texttt{toupper()}函数，使用通用算法\texttt{std::\allowbreak{}transform()}：


%%C020306 code

\begin{codeblockbg}
\begin{Verbatim}[commandchars=\\\{\},breaklines=true,breakanywhere=true,fontsize=\codefont,xleftmargin=17bp]
\PYG{k}{template}\PYG{o}{\PYGZlt{}}\PYG{k}{typename}\PYG{+w}{ }\PYG{n+nc}{CharT}\PYG{o}{\PYGZgt{}}
\PYG{k+kr}{inline}\PYG{+w}{ }\PYG{n}{tstring}\PYG{o}{\PYGZlt{}}\PYG{n}{CharT}\PYG{o}{\PYGZgt{}}\PYG{+w}{ }\PYG{n}{to\PYGZus{}upper}\PYG{p}{(}\PYG{n}{tstring}\PYG{o}{\PYGZlt{}}\PYG{n}{CharT}\PYG{o}{\PYGZgt{}}\PYG{+w}{ }\PYG{n}{text}\PYG{p}{)}
\PYG{p}{\PYGZob{}}
\PYG{+w}{  }\PYG{n}{std}\PYG{o}{:}\PYG{o}{:}\PYG{n}{transform}\PYG{p}{(}\PYG{n}{std}\PYG{o}{:}\PYG{o}{:}\PYG{n}{begin}\PYG{p}{(}\PYG{n}{text}\PYG{p}{)}\PYG{p}{,}\PYG{+w}{ }\PYG{n}{std}\PYG{o}{:}\PYG{o}{:}\PYG{n}{end}\PYG{p}{(}\PYG{n}{text}\PYG{p}{)}\PYG{p}{,}
\PYG{+w}{                 }\PYG{n}{std}\PYG{o}{:}\PYG{o}{:}\PYG{n}{begin}\PYG{p}{(}\PYG{n}{text}\PYG{p}{)}\PYG{p}{,}\PYG{+w}{ }\PYG{n}{toupper}\PYG{p}{)}\PYG{p}{;}
\PYG{+w}{  }\PYG{k}{return}\PYG{+w}{ }\PYG{n}{text}\PYG{p}{;}
\PYG{p}{\PYGZcb{}}
\PYG{k}{template}\PYG{o}{\PYGZlt{}}\PYG{k}{typename}\PYG{+w}{ }\PYG{n+nc}{CharT}\PYG{o}{\PYGZgt{}}
\PYG{k+kr}{inline}\PYG{+w}{ }\PYG{n}{tstring}\PYG{o}{\PYGZlt{}}\PYG{n}{CharT}\PYG{o}{\PYGZgt{}}\PYG{+w}{ }\PYG{n}{to\PYGZus{}lower}\PYG{p}{(}\PYG{n}{tstring}\PYG{o}{\PYGZlt{}}\PYG{n}{CharT}\PYG{o}{\PYGZgt{}}\PYG{+w}{ }\PYG{n}{text}\PYG{p}{)}
\PYG{p}{\PYGZob{}}
\PYG{+w}{  }\PYG{n}{std}\PYG{o}{:}\PYG{o}{:}\PYG{n}{transform}\PYG{p}{(}\PYG{n}{std}\PYG{o}{:}\PYG{o}{:}\PYG{n}{begin}\PYG{p}{(}\PYG{n}{text}\PYG{p}{)}\PYG{p}{,}\PYG{+w}{ }\PYG{n}{std}\PYG{o}{:}\PYG{o}{:}\PYG{n}{end}\PYG{p}{(}\PYG{n}{text}\PYG{p}{)}\PYG{p}{,}
\PYG{+w}{                 }\PYG{n}{std}\PYG{o}{:}\PYG{o}{:}\PYG{n}{begin}\PYG{p}{(}\PYG{n}{text}\PYG{p}{)}\PYG{p}{,}\PYG{+w}{ }\PYG{n}{tolower}\PYG{p}{)}\PYG{p}{;}
\PYG{+w}{  }\PYG{k}{return}\PYG{+w}{ }\PYG{n}{text}\PYG{p}{;}
\PYG{p}{\PYGZcb{}}
\end{Verbatim}
\end{codeblockbg}


%%C020307 p

要反转字符串，可以使用通用算法\texttt{std::\allowbreak{}reverse()}：


%%C020308 code

\begin{codeblockbg}
\begin{Verbatim}[commandchars=\\\{\},breaklines=true,breakanywhere=true,fontsize=\codefont,xleftmargin=17bp]
\PYG{k}{template}\PYG{o}{\PYGZlt{}}\PYG{k}{typename}\PYG{+w}{ }\PYG{n+nc}{CharT}\PYG{o}{\PYGZgt{}}
\PYG{k+kr}{inline}\PYG{+w}{ }\PYG{n}{tstring}\PYG{o}{\PYGZlt{}}\PYG{n}{CharT}\PYG{o}{\PYGZgt{}}\PYG{+w}{ }\PYG{n}{reverse}\PYG{p}{(}\PYG{n}{tstring}\PYG{o}{\PYGZlt{}}\PYG{n}{CharT}\PYG{o}{\PYGZgt{}}\PYG{+w}{ }\PYG{n}{text}\PYG{p}{)}
\PYG{p}{\PYGZob{}}
\PYG{+w}{  }\PYG{n}{std}\PYG{o}{:}\PYG{o}{:}\PYG{n}{reverse}\PYG{p}{(}\PYG{n}{std}\PYG{o}{:}\PYG{o}{:}\PYG{n}{begin}\PYG{p}{(}\PYG{n}{text}\PYG{p}{)}\PYG{p}{,}\PYG{+w}{ }\PYG{n}{std}\PYG{o}{:}\PYG{o}{:}\PYG{n}{end}\PYG{p}{(}\PYG{n}{text}\PYG{p}{)}\PYG{p}{)}\PYG{p}{;}
\PYG{+w}{  }\PYG{k}{return}\PYG{+w}{ }\PYG{n}{text}\PYG{p}{;}
\PYG{p}{\PYGZcb{}}
\end{Verbatim}
\end{codeblockbg}


%%C020309 p

要去除字符串开头、结尾或两端的字符，可以使用\texttt{std::\allowbreak{}basic\_\allowbreak{}string}的


%%C020310 p

方法\myttbr{find\_\allowbreak{}first\_\allowbreak{}not\_\allowbreak{}of()}和\myttbr{find\_\allowbreak{}last\_\allowbreak{}not\_\allowbreak{}of()}：


%%C020311 code

\begin{codeblockbg}
\begin{Verbatim}[commandchars=\\\{\},breaklines=true,breakanywhere=true,fontsize=\codefont,xleftmargin=17bp]
\PYG{k}{template}\PYG{o}{\PYGZlt{}}\PYG{k}{typename}\PYG{+w}{ }\PYG{n+nc}{CharT}\PYG{o}{\PYGZgt{}}
\PYG{k+kr}{inline}\PYG{+w}{ }\PYG{n}{tstring}\PYG{o}{\PYGZlt{}}\PYG{n}{CharT}\PYG{o}{\PYGZgt{}}\PYG{+w}{ }\PYG{n}{trim}\PYG{p}{(}\PYG{n}{tstring}\PYG{o}{\PYGZlt{}}\PYG{n}{CharT}\PYG{o}{\PYGZgt{}}\PYG{+w}{ }\PYG{k}{const}\PYG{+w}{ }\PYG{o}{\PYGZam{}}\PYG{+w}{ }\PYG{n}{text}\PYG{p}{)}
\PYG{p}{\PYGZob{}}
\PYG{+w}{  }\PYG{k}{auto}\PYG{+w}{ }\PYG{n}{first}\PYG{p}{\PYGZob{}}\PYG{+w}{ }\PYG{n}{text}\PYG{p}{.}\PYG{n}{find\PYGZus{}first\PYGZus{}not\PYGZus{}of}\PYG{p}{(}\PYG{l+s+sc}{\PYGZsq{}}\PYG{l+s+sc}{ }\PYG{l+s+sc}{\PYGZsq{}}\PYG{p}{)}\PYG{+w}{ }\PYG{p}{\PYGZcb{}}\PYG{p}{;}
\PYG{+w}{  }\PYG{k}{auto}\PYG{+w}{ }\PYG{n}{last}\PYG{p}{\PYGZob{}}\PYG{+w}{ }\PYG{n}{text}\PYG{p}{.}\PYG{n}{find\PYGZus{}last\PYGZus{}not\PYGZus{}of}\PYG{p}{(}\PYG{l+s+sc}{\PYGZsq{}}\PYG{l+s+sc}{ }\PYG{l+s+sc}{\PYGZsq{}}\PYG{p}{)}\PYG{+w}{ }\PYG{p}{\PYGZcb{}}\PYG{p}{;}
\PYG{+w}{  }\PYG{k}{return}\PYG{+w}{ }\PYG{n}{text}\PYG{p}{.}\PYG{n}{substr}\PYG{p}{(}\PYG{n}{first}\PYG{p}{,}\PYG{+w}{ }\PYG{p}{(}\PYG{n}{last}\PYG{+w}{ }\PYG{o}{\PYGZhy{}}\PYG{+w}{ }\PYG{n}{first}\PYG{+w}{ }\PYG{o}{+}\PYG{+w}{ }\PYG{l+m+mi}{1}\PYG{p}{)}\PYG{p}{)}\PYG{p}{;}
\PYG{p}{\PYGZcb{}}
\PYG{k}{template}\PYG{o}{\PYGZlt{}}\PYG{k}{typename}\PYG{+w}{ }\PYG{n+nc}{CharT}\PYG{o}{\PYGZgt{}}
\PYG{k+kr}{inline}\PYG{+w}{ }\PYG{n}{tstring}\PYG{o}{\PYGZlt{}}\PYG{n}{CharT}\PYG{o}{\PYGZgt{}}\PYG{+w}{ }\PYG{n}{trimleft}\PYG{p}{(}\PYG{n}{tstring}\PYG{o}{\PYGZlt{}}\PYG{n}{CharT}\PYG{o}{\PYGZgt{}}\PYG{+w}{ }\PYG{k}{const}\PYG{+w}{ }\PYG{o}{\PYGZam{}}\PYG{+w}{ }\PYG{n}{text}\PYG{p}{)}
\PYG{p}{\PYGZob{}}
\PYG{+w}{  }\PYG{k}{auto}\PYG{+w}{ }\PYG{n}{first}\PYG{p}{\PYGZob{}}\PYG{+w}{ }\PYG{n}{text}\PYG{p}{.}\PYG{n}{find\PYGZus{}first\PYGZus{}not\PYGZus{}of}\PYG{p}{(}\PYG{l+s+sc}{\PYGZsq{}}\PYG{l+s+sc}{ }\PYG{l+s+sc}{\PYGZsq{}}\PYG{p}{)}\PYG{+w}{ }\PYG{p}{\PYGZcb{}}\PYG{p}{;}
\PYG{+w}{  }\PYG{k}{return}\PYG{+w}{ }\PYG{n}{text}\PYG{p}{.}\PYG{n}{substr}\PYG{p}{(}\PYG{n}{first}\PYG{p}{,}\PYG{+w}{ }\PYG{n}{text}\PYG{p}{.}\PYG{n}{size}\PYG{p}{(}\PYG{p}{)}\PYG{+w}{ }\PYG{o}{\PYGZhy{}}\PYG{+w}{ }\PYG{n}{first}\PYG{p}{)}\PYG{p}{;}
\PYG{p}{\PYGZcb{}}
\PYG{k}{template}\PYG{o}{\PYGZlt{}}\PYG{k}{typename}\PYG{+w}{ }\PYG{n+nc}{CharT}\PYG{o}{\PYGZgt{}}
\PYG{k+kr}{inline}\PYG{+w}{ }\PYG{n}{tstring}\PYG{o}{\PYGZlt{}}\PYG{n}{CharT}\PYG{o}{\PYGZgt{}}\PYG{+w}{ }\PYG{n}{trimright}\PYG{p}{(}\PYG{n}{tstring}\PYG{o}{\PYGZlt{}}\PYG{n}{CharT}\PYG{o}{\PYGZgt{}}\PYG{+w}{ }\PYG{k}{const}\PYG{+w}{ }\PYG{o}{\PYGZam{}}\PYG{+w}{ }\PYG{n}{text}\PYG{p}{)}
\PYG{p}{\PYGZob{}}
\PYG{+w}{  }\PYG{k}{auto}\PYG{+w}{ }\PYG{n}{last}\PYG{p}{\PYGZob{}}\PYG{+w}{ }\PYG{n}{text}\PYG{p}{.}\PYG{n}{find\PYGZus{}last\PYGZus{}not\PYGZus{}of}\PYG{p}{(}\PYG{l+s+sc}{\PYGZsq{}}\PYG{l+s+sc}{ }\PYG{l+s+sc}{\PYGZsq{}}\PYG{p}{)}\PYG{+w}{ }\PYG{p}{\PYGZcb{}}\PYG{p}{;}
\PYG{+w}{  }\PYG{k}{return}\PYG{+w}{ }\PYG{n}{text}\PYG{p}{.}\PYG{n}{substr}\PYG{p}{(}\PYG{l+m+mi}{0}\PYG{p}{,}\PYG{+w}{ }\PYG{n}{last}\PYG{+w}{ }\PYG{o}{+}\PYG{+w}{ }\PYG{l+m+mi}{1}\PYG{p}{)}\PYG{p}{;}
\PYG{p}{\PYGZcb{}}
\end{Verbatim}
\end{codeblockbg}


%%C020312 p

要去除字符串中某个给定字符集合里的字符，可以使用


%%C020313 p

\texttt{std::\allowbreak{}basic\_\allowbreak{}string}的方法\myttbr{find\_\allowbreak{}first\_\allowbreak{}not\_\allowbreak{}of()}和\myttbr{find\_\allowbreak{}last\_\allowbreak{}not\_\allowbreak{}of()}的重载版本，它们接受一个字符串参数，用于指定要查找的字符集合：


%%C020314 code

\begin{codeblockbg}
\begin{Verbatim}[commandchars=\\\{\},breaklines=true,breakanywhere=true,fontsize=\codefont,xleftmargin=17bp]
\PYG{k}{template}\PYG{o}{\PYGZlt{}}\PYG{k}{typename}\PYG{+w}{ }\PYG{n+nc}{CharT}\PYG{o}{\PYGZgt{}}
\PYG{k+kr}{inline}\PYG{+w}{ }\PYG{n}{tstring}\PYG{o}{\PYGZlt{}}\PYG{n}{CharT}\PYG{o}{\PYGZgt{}}\PYG{+w}{ }\PYG{n}{trim}\PYG{p}{(}\PYG{n}{tstring}\PYG{o}{\PYGZlt{}}\PYG{n}{CharT}\PYG{o}{\PYGZgt{}}\PYG{+w}{ }\PYG{k}{const}\PYG{+w}{ }\PYG{o}{\PYGZam{}}\PYG{+w}{ }\PYG{n}{text}\PYG{p}{,}
\PYG{+w}{                           }\PYG{n}{tstring}\PYG{o}{\PYGZlt{}}\PYG{n}{CharT}\PYG{o}{\PYGZgt{}}\PYG{+w}{ }\PYG{k}{const}\PYG{+w}{ }\PYG{o}{\PYGZam{}}\PYG{+w}{ }\PYG{n}{chars}\PYG{p}{)}
\PYG{p}{\PYGZob{}}
\PYG{+w}{  }\PYG{k}{auto}\PYG{+w}{ }\PYG{n}{first}\PYG{p}{\PYGZob{}}\PYG{+w}{ }\PYG{n}{text}\PYG{p}{.}\PYG{n}{find\PYGZus{}first\PYGZus{}not\PYGZus{}of}\PYG{p}{(}\PYG{n}{chars}\PYG{p}{)}\PYG{+w}{ }\PYG{p}{\PYGZcb{}}\PYG{p}{;}
\PYG{+w}{  }\PYG{k}{auto}\PYG{+w}{ }\PYG{n}{last}\PYG{p}{\PYGZob{}}\PYG{+w}{ }\PYG{n}{text}\PYG{p}{.}\PYG{n}{find\PYGZus{}last\PYGZus{}not\PYGZus{}of}\PYG{p}{(}\PYG{n}{chars}\PYG{p}{)}\PYG{+w}{ }\PYG{p}{\PYGZcb{}}\PYG{p}{;}
\PYG{+w}{  }\PYG{k}{return}\PYG{+w}{ }\PYG{n}{text}\PYG{p}{.}\PYG{n}{substr}\PYG{p}{(}\PYG{n}{first}\PYG{p}{,}\PYG{+w}{ }\PYG{p}{(}\PYG{n}{last}\PYG{+w}{ }\PYG{o}{\PYGZhy{}}\PYG{+w}{ }\PYG{n}{first}\PYG{+w}{ }\PYG{o}{+}\PYG{+w}{ }\PYG{l+m+mi}{1}\PYG{p}{)}\PYG{p}{)}\PYG{p}{;}
\PYG{p}{\PYGZcb{}}
\PYG{k}{template}\PYG{o}{\PYGZlt{}}\PYG{k}{typename}\PYG{+w}{ }\PYG{n+nc}{CharT}\PYG{o}{\PYGZgt{}}
\PYG{k+kr}{inline}\PYG{+w}{ }\PYG{n}{tstring}\PYG{o}{\PYGZlt{}}\PYG{n}{CharT}\PYG{o}{\PYGZgt{}}\PYG{+w}{ }\PYG{n}{trimleft}\PYG{p}{(}\PYG{n}{tstring}\PYG{o}{\PYGZlt{}}\PYG{n}{CharT}\PYG{o}{\PYGZgt{}}\PYG{+w}{ }\PYG{k}{const}\PYG{+w}{ }\PYG{o}{\PYGZam{}}\PYG{+w}{ }\PYG{n}{text}\PYG{p}{,}
\PYG{+w}{                               }\PYG{n}{tstring}\PYG{o}{\PYGZlt{}}\PYG{n}{CharT}\PYG{o}{\PYGZgt{}}\PYG{+w}{ }\PYG{k}{const}\PYG{+w}{ }\PYG{o}{\PYGZam{}}\PYG{+w}{ }\PYG{n}{chars}\PYG{p}{)}
\PYG{p}{\PYGZob{}}
\PYG{+w}{  }\PYG{k}{auto}\PYG{+w}{ }\PYG{n}{first}\PYG{p}{\PYGZob{}}\PYG{+w}{ }\PYG{n}{text}\PYG{p}{.}\PYG{n}{find\PYGZus{}first\PYGZus{}not\PYGZus{}of}\PYG{p}{(}\PYG{n}{chars}\PYG{p}{)}\PYG{+w}{ }\PYG{p}{\PYGZcb{}}\PYG{p}{;}
\PYG{+w}{  }\PYG{k}{return}\PYG{+w}{ }\PYG{n}{text}\PYG{p}{.}\PYG{n}{substr}\PYG{p}{(}\PYG{n}{first}\PYG{p}{,}\PYG{+w}{ }\PYG{n}{text}\PYG{p}{.}\PYG{n}{size}\PYG{p}{(}\PYG{p}{)}\PYG{+w}{ }\PYG{o}{\PYGZhy{}}\PYG{+w}{ }\PYG{n}{first}\PYG{p}{)}\PYG{p}{;}
\PYG{p}{\PYGZcb{}}
\PYG{k}{template}\PYG{o}{\PYGZlt{}}\PYG{k}{typename}\PYG{+w}{ }\PYG{n+nc}{CharT}\PYG{o}{\PYGZgt{}}
\PYG{k+kr}{inline}\PYG{+w}{ }\PYG{n}{tstring}\PYG{o}{\PYGZlt{}}\PYG{n}{CharT}\PYG{o}{\PYGZgt{}}\PYG{+w}{ }\PYG{n}{trimright}\PYG{p}{(}\PYG{n}{tstring}\PYG{o}{\PYGZlt{}}\PYG{n}{CharT}\PYG{o}{\PYGZgt{}}\PYG{+w}{ }\PYG{k}{const}\PYG{+w}{ }\PYG{o}{\PYGZam{}}\PYG{n}{text}\PYG{p}{,}
\PYG{+w}{                                }\PYG{n}{tstring}\PYG{o}{\PYGZlt{}}\PYG{n}{CharT}\PYG{o}{\PYGZgt{}}\PYG{+w}{ }\PYG{k}{const}\PYG{+w}{ }\PYG{o}{\PYGZam{}}\PYG{n}{chars}\PYG{p}{)}
\PYG{p}{\PYGZob{}}
\PYG{+w}{  }\PYG{k}{auto}\PYG{+w}{ }\PYG{n}{last}\PYG{p}{\PYGZob{}}\PYG{+w}{ }\PYG{n}{text}\PYG{p}{.}\PYG{n}{find\PYGZus{}last\PYGZus{}not\PYGZus{}of}\PYG{p}{(}\PYG{n}{chars}\PYG{p}{)}\PYG{+w}{ }\PYG{p}{\PYGZcb{}}\PYG{p}{;}
\PYG{+w}{  }\PYG{k}{return}\PYG{+w}{ }\PYG{n}{text}\PYG{p}{.}\PYG{n}{substr}\PYG{p}{(}\PYG{l+m+mi}{0}\PYG{p}{,}\PYG{+w}{ }\PYG{n}{last}\PYG{+w}{ }\PYG{o}{+}\PYG{+w}{ }\PYG{l+m+mi}{1}\PYG{p}{)}\PYG{p}{;}
\PYG{p}{\PYGZcb{}}
\end{Verbatim}
\end{codeblockbg}


%%C020315 p

要从字符串中删除字符，可以使用\texttt{std::\allowbreak{}remove\_\allowbreak{}if()}和


%%C020316 p

\myttbr{std::\allowbreak{}basic\_\allowbreak{}string::\allowbreak{}erase()}：


%%C020317 code

\begin{codeblockbg}
\begin{Verbatim}[commandchars=\\\{\},breaklines=true,breakanywhere=true,fontsize=\codefont,xleftmargin=17bp]
\PYG{k}{template}\PYG{o}{\PYGZlt{}}\PYG{k}{typename}\PYG{+w}{ }\PYG{n+nc}{CharT}\PYG{o}{\PYGZgt{}}
\PYG{k+kr}{inline}\PYG{+w}{ }\PYG{n}{tstring}\PYG{o}{\PYGZlt{}}\PYG{n}{CharT}\PYG{o}{\PYGZgt{}}\PYG{+w}{ }\PYG{n}{remove}\PYG{p}{(}\PYG{n}{tstring}\PYG{o}{\PYGZlt{}}\PYG{n}{CharT}\PYG{o}{\PYGZgt{}}\PYG{+w}{ }\PYG{n}{text}\PYG{p}{,}
\PYG{+w}{                             }\PYG{n}{CharT}\PYG{+w}{ }\PYG{k}{const}\PYG{+w}{ }\PYG{n}{ch}\PYG{p}{)}
\PYG{p}{\PYGZob{}}
\PYG{+w}{  }\PYG{k}{auto}\PYG{+w}{ }\PYG{n}{start}\PYG{+w}{ }\PYG{o}{=}\PYG{+w}{ }\PYG{n}{std}\PYG{o}{:}\PYG{o}{:}\PYG{n}{remove\PYGZus{}if}\PYG{p}{(}
\PYG{+w}{                  }\PYG{n}{std}\PYG{o}{:}\PYG{o}{:}\PYG{n}{begin}\PYG{p}{(}\PYG{n}{text}\PYG{p}{)}\PYG{p}{,}\PYG{+w}{ }\PYG{n}{std}\PYG{o}{:}\PYG{o}{:}\PYG{n}{end}\PYG{p}{(}\PYG{n}{text}\PYG{p}{)}\PYG{p}{,}
\PYG{+w}{                  }\PYG{p}{[}\PYG{o}{=}\PYG{p}{]}\PYG{p}{(}\PYG{n}{CharT}\PYG{+w}{ }\PYG{k}{const}\PYG{+w}{ }\PYG{n}{c}\PYG{p}{)}\PYG{+w}{ }\PYG{p}{\PYGZob{}}\PYG{k}{return}\PYG{+w}{ }\PYG{n}{c}\PYG{+w}{ }\PYG{o}{=}\PYG{o}{=}\PYG{+w}{  }\PYG{n}{ch}\PYG{p}{;}\PYG{+w}{ }\PYG{p}{\PYGZcb{}}\PYG{p}{)}\PYG{p}{;}
\PYG{+w}{  }\PYG{n}{text}\PYG{p}{.}\PYG{n}{erase}\PYG{p}{(}\PYG{n}{start}\PYG{p}{,}\PYG{+w}{ }\PYG{n}{std}\PYG{o}{:}\PYG{o}{:}\PYG{n}{end}\PYG{p}{(}\PYG{n}{text}\PYG{p}{)}\PYG{p}{)}\PYG{p}{;}
\PYG{+w}{  }\PYG{k}{return}\PYG{+w}{ }\PYG{n}{text}\PYG{p}{;}
\PYG{p}{\PYGZcb{}}
\end{Verbatim}
\end{codeblockbg}


%%C020318 p

要按指定的分隔符拆分字符串，可以使用\texttt{std::\allowbreak{}getline()}从


%%C020319 p

一个用字符串内容初始化的\myttbr{std::\allowbreak{}basic\_\allowbreak{}stringstream}中读取。从流中提取出的记号被依次放入字符串向量：


%%C020320 code

\begin{codeblockbg}
\begin{Verbatim}[commandchars=\\\{\},breaklines=true,breakanywhere=true,fontsize=\codefont,xleftmargin=17bp]
\PYG{k}{template}\PYG{o}{\PYGZlt{}}\PYG{k}{typename}\PYG{+w}{ }\PYG{n+nc}{CharT}\PYG{o}{\PYGZgt{}}
\PYG{k+kr}{inline}\PYG{+w}{ }\PYG{n}{std}\PYG{o}{:}\PYG{o}{:}\PYG{n}{vector}\PYG{o}{\PYGZlt{}}\PYG{n}{tstring}\PYG{o}{\PYGZlt{}}\PYG{n}{CharT}\PYG{o}{\PYGZgt{}}\PYG{o}{\PYGZgt{}}\PYG{+w}{ }\PYG{n}{split}
\PYG{+w}{   }\PYG{p}{(}\PYG{n}{tstring}\PYG{o}{\PYGZlt{}}\PYG{n}{CharT}\PYG{o}{\PYGZgt{}}\PYG{+w}{ }\PYG{n}{text}\PYG{p}{,}\PYG{+w}{ }\PYG{n}{CharT}\PYG{+w}{ }\PYG{k}{const}\PYG{+w}{ }\PYG{n}{delimiter}\PYG{p}{)}
\PYG{p}{\PYGZob{}}
\PYG{+w}{  }\PYG{k}{auto}\PYG{+w}{ }\PYG{n}{sstr}\PYG{+w}{ }\PYG{o}{=}\PYG{+w}{ }\PYG{n}{tstringstream}\PYG{o}{\PYGZlt{}}\PYG{n}{CharT}\PYG{o}{\PYGZgt{}}\PYG{p}{\PYGZob{}}\PYG{+w}{ }\PYG{n}{text}\PYG{+w}{ }\PYG{p}{\PYGZcb{}}\PYG{p}{;}
\PYG{+w}{  }\PYG{k}{auto}\PYG{+w}{ }\PYG{n}{tokens}\PYG{+w}{ }\PYG{o}{=}\PYG{+w}{ }\PYG{n}{std}\PYG{o}{:}\PYG{o}{:}\PYG{n}{vector}\PYG{o}{\PYGZlt{}}\PYG{n}{tstring}\PYG{o}{\PYGZlt{}}\PYG{n}{CharT}\PYG{o}{\PYGZgt{}}\PYG{o}{\PYGZgt{}}\PYG{p}{\PYGZob{}}\PYG{p}{\PYGZcb{}}\PYG{p}{;}
\PYG{+w}{  }\PYG{k}{auto}\PYG{+w}{ }\PYG{n}{token}\PYG{+w}{ }\PYG{o}{=}\PYG{+w}{ }\PYG{n}{tstring}\PYG{o}{\PYGZlt{}}\PYG{n}{CharT}\PYG{o}{\PYGZgt{}}\PYG{p}{\PYGZob{}}\PYG{p}{\PYGZcb{}}\PYG{p}{;}
\PYG{+w}{  }\PYG{k}{while}\PYG{+w}{ }\PYG{p}{(}\PYG{n}{std}\PYG{o}{:}\PYG{o}{:}\PYG{n}{getline}\PYG{p}{(}\PYG{n}{sstr}\PYG{p}{,}\PYG{+w}{ }\PYG{n}{token}\PYG{p}{,}\PYG{+w}{ }\PYG{n}{delimiter}\PYG{p}{)}\PYG{p}{)}
\PYG{+w}{  }\PYG{p}{\PYGZob{}}
\PYG{+w}{    }\PYG{k}{if}\PYG{+w}{ }\PYG{p}{(}\PYG{o}{!}\PYG{n}{token}\PYG{p}{.}\PYG{n}{empty}\PYG{p}{(}\PYG{p}{)}\PYG{p}{)}\PYG{+w}{ }\PYG{n}{tokens}\PYG{p}{.}\PYG{n}{push\PYGZus{}back}\PYG{p}{(}\PYG{n}{token}\PYG{p}{)}\PYG{p}{;}
\PYG{+w}{  }\PYG{p}{\PYGZcb{}}
\PYG{+w}{  }\PYG{k}{return}\PYG{+w}{ }\PYG{n}{tokens}\PYG{p}{;}
\PYG{p}{\PYGZcb{}}
\end{Verbatim}
\end{codeblockbg}


%%C020321 h3

\subsection{工作原理}


%%C020322 p

实现这个库中的工具函数时，我们有两种选择：


%%C020323 li

\begin{itemize}
\item 函数通过引用修改传入的字符串。


%%C020324 li

\item 函数不改变原字符串，而是返回一个新字符串。

\end{itemize}


%%C020325 p

第二种选择的好处是保留原字符串，这在很多情况下很有帮助。否则，在这些情况下就得先复制一份字符串，再修改副本。本案例给出的实现采用第二种做法。


%%C020326 p

我们在\textit{具体做法}一节中实现的头两个函数是\texttt{to\_\allowbreak{}upper()}和\texttt{to\_\allowbreak{}lower()}。这两个函数把字符串内容改为大写或小写。最简单的实现方式是使用\texttt{std::\allowbreak{}transform()}标准算法。这是一个通用算法，它把一个函数应用到某个范围（由起始和结束迭代器定义）的每个元素上，并把结果存入另一个只需指定起始迭代器的范围。输出范围可以与输入范围相同，我们正是这样转换字符串的。所应用的函数是\texttt{toupper()}或\texttt{tolower()}：


%%C020327 code

\begin{codeblockbg}
\begin{Verbatim}[commandchars=\\\{\},breaklines=true,breakanywhere=true,fontsize=\codefont,xleftmargin=17bp]
\PYG{k}{auto}\PYG{+w}{ }\PYG{n}{ut}\PYG{p}{\PYGZob{}}\PYG{+w}{ }\PYG{n}{string\PYGZus{}library}\PYG{o}{:}\PYG{o}{:}\PYG{n}{to\PYGZus{}upper}\PYG{p}{(}\PYG{l+s}{\PYGZdq{}}\PYG{l+s}{this is not UPPERCASE}\PYG{l+s}{\PYGZdq{}}\PYG{n}{s}\PYG{p}{)}\PYG{+w}{ }\PYG{p}{\PYGZcb{}}\PYG{p}{;}
\PYG{c+c1}{// ut = \PYGZdq{}THIS IS NOT UPPERCASE\PYGZdq{}}
\PYG{k}{auto}\PYG{+w}{ }\PYG{n}{lt}\PYG{p}{\PYGZob{}}\PYG{+w}{ }\PYG{n}{string\PYGZus{}library}\PYG{o}{:}\PYG{o}{:}\PYG{n}{to\PYGZus{}lower}\PYG{p}{(}\PYG{l+s}{\PYGZdq{}}\PYG{l+s}{THIS IS NOT lowercase}\PYG{l+s}{\PYGZdq{}}\PYG{n}{s}\PYG{p}{)}\PYG{+w}{ }\PYG{p}{\PYGZcb{}}\PYG{p}{;}
\PYG{c+c1}{// lt = \PYGZdq{}this is not lowercase\PYGZdq{}}
\end{Verbatim}
\end{codeblockbg}


%%C020328 p

接下来考虑的函数是\texttt{reverse()}，顾名思义，它反转字符串的内容。为此我们使用了\texttt{std::\allowbreak{}reverse()}标准算法。这个通用算法反转由起始和结束迭代器定义的范围内各元素的顺序：


%%C020329 code

\begin{codeblockbg}
\begin{Verbatim}[commandchars=\\\{\},breaklines=true,breakanywhere=true,fontsize=\codefont,xleftmargin=17bp]
\PYG{k}{auto}\PYG{+w}{ }\PYG{n}{rt}\PYG{p}{\PYGZob{}}\PYG{n}{string\PYGZus{}library}\PYG{o}{:}\PYG{o}{:}\PYG{n}{reverse}\PYG{p}{(}\PYG{l+s}{\PYGZdq{}}\PYG{l+s}{cookbook}\PYG{l+s}{\PYGZdq{}}\PYG{n}{s}\PYG{p}{)}\PYG{p}{\PYGZcb{}}\PYG{p}{;}\PYG{+w}{ }\PYG{c+c1}{// rt = \PYGZdq{}koobkooc\PYGZdq{}}
\end{Verbatim}
\end{codeblockbg}


%%C020330 p

说到去除字符，字符串可以在开头、结尾或两端进行修剪。因此我们实现了三个不同的函数：\texttt{trim()}删除两端的字符，\texttt{trimleft()}删除字符串开头的字符，\texttt{trimright()}删除字符串结尾的字符。这些函数的第一个版本只删除空格。为了找到需要删除的部分，我们使用\texttt{std::\allowbreak{}basic\_\allowbreak{}string}的\myttbr{find\_\allowbreak{}first\_\allowbreak{}not\_\allowbreak{}of()}和\myttbr{find\_\allowbreak{}last\_\allowbreak{}not\_\allowbreak{}of()}方法。它们返回字符串中第一个和最后一个不是指定字符的字符。随后调用\texttt{std::\allowbreak{}basic\_\allowbreak{}string}的\texttt{substr()}方法返回一个新字符串。\texttt{substr()}方法接受字符串中的一个索引和要复制到新字符串中的元素个数：


%%C020331 code

\begin{codeblockbg}
\begin{Verbatim}[commandchars=\\\{\},breaklines=true,breakanywhere=true,fontsize=\codefont,xleftmargin=17bp]
\PYG{k}{auto}\PYG{+w}{ }\PYG{n}{text1}\PYG{p}{\PYGZob{}}\PYG{l+s}{\PYGZdq{}}\PYG{l+s}{   this is an example   }\PYG{l+s}{\PYGZdq{}}\PYG{n}{s}\PYG{p}{\PYGZcb{}}\PYG{p}{;}
\PYG{c+c1}{// t1 = \PYGZdq{}this is an example\PYGZdq{}}
\PYG{k}{auto}\PYG{+w}{ }\PYG{n}{t1}\PYG{p}{\PYGZob{}}\PYG{+w}{ }\PYG{n}{string\PYGZus{}library}\PYG{o}{:}\PYG{o}{:}\PYG{n}{trim}\PYG{p}{(}\PYG{n}{text1}\PYG{p}{)}\PYG{+w}{ }\PYG{p}{\PYGZcb{}}\PYG{p}{;}
\PYG{c+c1}{// t2 = \PYGZdq{}this is an example   \PYGZdq{}}
\PYG{k}{auto}\PYG{+w}{ }\PYG{n}{t2}\PYG{p}{\PYGZob{}}\PYG{+w}{ }\PYG{n}{string\PYGZus{}library}\PYG{o}{:}\PYG{o}{:}\PYG{n}{trimleft}\PYG{p}{(}\PYG{n}{text1}\PYG{p}{)}\PYG{+w}{ }\PYG{p}{\PYGZcb{}}\PYG{p}{;}
\PYG{c+c1}{// t3 = \PYGZdq{}   this is an example\PYGZdq{}}
\PYG{k}{auto}\PYG{+w}{ }\PYG{n}{t3}\PYG{p}{\PYGZob{}}\PYG{+w}{ }\PYG{n}{string\PYGZus{}library}\PYG{o}{:}\PYG{o}{:}\PYG{n}{trimright}\PYG{p}{(}\PYG{n}{text1}\PYG{p}{)}\PYG{+w}{ }\PYG{p}{\PYGZcb{}}\PYG{p}{;}
\end{Verbatim}
\end{codeblockbg}


%%C020332 p

有时可能需要先删除字符串中的其他字符，再删除空格。为此，我们为修剪函数提供了指定待删除字符集合的重载版本。该集合同样以字符串的形式给出。其实现与前一个版本非常相似，因为\myttbr{find\_\allowbreak{}first\_\allowbreak{}not\_\allowbreak{}of()}和\myttbr{find\_\allowbreak{}last\_\allowbreak{}not\_\allowbreak{}of()}都有接受字符串参数的重载，该字符串包含要从查找中排除的字符：


%%C020333 code

\begin{codeblockbg}
\begin{Verbatim}[commandchars=\\\{\},breaklines=true,breakanywhere=true,fontsize=\codefont,xleftmargin=17bp]
\PYG{k}{auto}\PYG{+w}{ }\PYG{n}{chars1}\PYG{p}{\PYGZob{}}\PYG{l+s}{\PYGZdq{}}\PYG{l+s}{ !\PYGZpc{}}\PYG{l+s+se}{\PYGZbs{}n}\PYG{l+s+se}{\PYGZbs{}r}\PYG{l+s}{\PYGZdq{}}\PYG{n}{s}\PYG{p}{\PYGZcb{}}\PYG{p}{;}
\PYG{k}{auto}\PYG{+w}{ }\PYG{n}{text3}\PYG{p}{\PYGZob{}}\PYG{l+s}{\PYGZdq{}}\PYG{l+s}{!!  this \PYGZpc{} needs a lot}\PYG{l+s+se}{\PYGZbs{}r}\PYG{l+s}{of trimming  !}\PYG{l+s+se}{\PYGZbs{}n}\PYG{l+s}{\PYGZdq{}}\PYG{n}{s}\PYG{p}{\PYGZcb{}}\PYG{p}{;}
\PYG{k}{auto}\PYG{+w}{ }\PYG{n}{t7}\PYG{p}{\PYGZob{}}\PYG{+w}{ }\PYG{n}{string\PYGZus{}library}\PYG{o}{:}\PYG{o}{:}\PYG{n}{trim}\PYG{p}{(}\PYG{n}{text3}\PYG{p}{,}\PYG{+w}{ }\PYG{n}{chars1}\PYG{p}{)}\PYG{+w}{ }\PYG{p}{\PYGZcb{}}\PYG{p}{;}
\PYG{c+c1}{// t7 = \PYGZdq{}this \PYGZpc{} needs a lot\PYGZbs{}rof trimming\PYGZdq{}}
\PYG{k}{auto}\PYG{+w}{ }\PYG{n}{t8}\PYG{p}{\PYGZob{}}\PYG{+w}{ }\PYG{n}{string\PYGZus{}library}\PYG{o}{:}\PYG{o}{:}\PYG{n}{trimleft}\PYG{p}{(}\PYG{n}{text3}\PYG{p}{,}\PYG{+w}{ }\PYG{n}{chars1}\PYG{p}{)}\PYG{+w}{ }\PYG{p}{\PYGZcb{}}\PYG{p}{;}
\PYG{c+c1}{// t8 = \PYGZdq{}this \PYGZpc{} needs a lot\PYGZbs{}rof trimming  !\PYGZbs{}n\PYGZdq{}}
\PYG{k}{auto}\PYG{+w}{ }\PYG{n}{t9}\PYG{p}{\PYGZob{}}\PYG{+w}{ }\PYG{n}{string\PYGZus{}library}\PYG{o}{:}\PYG{o}{:}\PYG{n}{trimright}\PYG{p}{(}\PYG{n}{text3}\PYG{p}{,}\PYG{+w}{ }\PYG{n}{chars1}\PYG{p}{)}\PYG{+w}{ }\PYG{p}{\PYGZcb{}}\PYG{p}{;}
\PYG{c+c1}{// t9 = \PYGZdq{}!!  this \PYGZpc{} needs a lot\PYGZbs{}rof trimming\PYGZdq{}}
\end{Verbatim}
\end{codeblockbg}


%%C020334 p

如果需要删除字符串任意位置上的字符，修剪函数就帮不上忙了，因为它们只处理字符串开头和结尾处的连续字符序列。为此，我们实现了一个简单的\texttt{remove()}方法。它使用\texttt{std:\allowbreak{}remove\_\allowbreak{}if()}标准算法。\texttt{std::\allowbreak{}remove()}和\texttt{std::\allowbreak{}remove\_\allowbreak{}if()}的工作方式初看可能不太直观。它们通过重排范围内容（使用移动赋值）来从由首尾迭代器定义的范围内移除满足条件的元素。需要移除的元素被放到范围末尾，函数返回指向范围中第一个代表被移除元素的迭代器。这个迭代器实际上定义了修改后范围的新末尾。如果没有元素被移除，返回的迭代器就是原范围的末尾迭代器。然后用这个返回的迭代器调用\myttbr{std::\allowbreak{}basic\_\allowbreak{}string::\allowbreak{}erase()}方法，真正删除由两个迭代器所定义的那部分字符串内容。本例中的两个迭代器分别是\texttt{std::\allowbreak{}remove\_\allowbreak{}if()}返回的迭代器和字符串末尾：


%%C020335 code

\begin{codeblockbg}
\begin{Verbatim}[commandchars=\\\{\},breaklines=true,breakanywhere=true,fontsize=\codefont,xleftmargin=17bp]
\PYG{k}{auto}\PYG{+w}{ }\PYG{n}{text4}\PYG{p}{\PYGZob{}}\PYG{l+s}{\PYGZdq{}}\PYG{l+s}{must remove all * from text**}\PYG{l+s}{\PYGZdq{}}\PYG{n}{s}\PYG{p}{\PYGZcb{}}\PYG{p}{;}
\PYG{k}{auto}\PYG{+w}{ }\PYG{n}{t10}\PYG{p}{\PYGZob{}}\PYG{+w}{ }\PYG{n}{string\PYGZus{}library}\PYG{o}{:}\PYG{o}{:}\PYG{n}{remove}\PYG{p}{(}\PYG{n}{text4}\PYG{p}{,}\PYG{+w}{ }\PYG{l+s+sc}{\PYGZsq{}}\PYG{l+s+sc}{*}\PYG{l+s+sc}{\PYGZsq{}}\PYG{p}{)}\PYG{+w}{ }\PYG{p}{\PYGZcb{}}\PYG{p}{;}
\PYG{c+c1}{// t10 = \PYGZdq{}must remove all  from text\PYGZdq{}}
\PYG{k}{auto}\PYG{+w}{ }\PYG{n}{t11}\PYG{p}{\PYGZob{}}\PYG{+w}{ }\PYG{n}{string\PYGZus{}library}\PYG{o}{:}\PYG{o}{:}\PYG{n}{remove}\PYG{p}{(}\PYG{n}{text4}\PYG{p}{,}\PYG{+w}{ }\PYG{l+s+sc}{\PYGZsq{}}\PYG{l+s+sc}{!}\PYG{l+s+sc}{\PYGZsq{}}\PYG{p}{)}\PYG{+w}{ }\PYG{p}{\PYGZcb{}}\PYG{p}{;}
\PYG{c+c1}{// t11 = \PYGZdq{}must remove all * from text**\PYGZdq{}}
\end{Verbatim}
\end{codeblockbg}

%%C020336 p

我们实现的最后一个方法按指定分隔符拆分字符串内容。实现方式有多种，这里用的是\texttt{std::\allowbreak{}getline()}。该函数从输入流中读取字符，直到遇到指定分隔符为止，并把读到的字符放进一个字符串。在开始从输入缓冲区读取之前，它会对输出字符串调用\texttt{erase()}清空内容。在循环中调用这个方法就能产生若干词元，放进一个vector。在我们的实现中，空词元会从结果集中跳过：


%%C020337 code

\begin{codeblockbg}
\begin{Verbatim}[commandchars=\\\{\},breaklines=true,breakanywhere=true,fontsize=\codefont,xleftmargin=17bp]
\PYG{k}{auto}\PYG{+w}{ }\PYG{n}{text5}\PYG{p}{\PYGZob{}}\PYG{l+s}{\PYGZdq{}}\PYG{l+s}{this text will be split   }\PYG{l+s}{\PYGZdq{}}\PYG{n}{s}\PYG{p}{\PYGZcb{}}\PYG{p}{;}
\PYG{k}{auto}\PYG{+w}{ }\PYG{n}{tokens1}\PYG{p}{\PYGZob{}}\PYG{+w}{ }\PYG{n}{string\PYGZus{}library}\PYG{o}{:}\PYG{o}{:}\PYG{n}{split}\PYG{p}{(}\PYG{n}{text5}\PYG{p}{,}\PYG{+w}{ }\PYG{l+s+sc}{\PYGZsq{}}\PYG{l+s+sc}{ }\PYG{l+s+sc}{\PYGZsq{}}\PYG{p}{)}\PYG{+w}{ }\PYG{p}{\PYGZcb{}}\PYG{p}{;}
\PYG{c+c1}{// tokens1 = \PYGZob{}\PYGZdq{}this\PYGZdq{}, \PYGZdq{}text\PYGZdq{}, \PYGZdq{}will\PYGZdq{}, \PYGZdq{}be\PYGZdq{}, \PYGZdq{}split\PYGZdq{}\PYGZcb{}}
\PYG{k}{auto}\PYG{+w}{ }\PYG{n}{tokens2}\PYG{p}{\PYGZob{}}\PYG{+w}{ }\PYG{n}{string\PYGZus{}library}\PYG{o}{:}\PYG{o}{:}\PYG{n}{split}\PYG{p}{(}\PYG{l+s}{\PYGZdq{}}\PYG{l+s}{\PYGZdq{}}\PYG{n}{s}\PYG{p}{,}\PYG{+w}{ }\PYG{l+s+sc}{\PYGZsq{}}\PYG{l+s+sc}{ }\PYG{l+s+sc}{\PYGZsq{}}\PYG{p}{)}\PYG{+w}{ }\PYG{p}{\PYGZcb{}}\PYG{p}{;}
\PYG{c+c1}{// tokens2 = \PYGZob{}\PYGZcb{}}
\end{Verbatim}
\end{codeblockbg}


%%C020338 h3

\subsection{延伸阅读}


%%C020339 li

\begin{itemize}
\item \textit{创建cooked用户定义字面量}


%%C020340 li

\item \textit{创建类型别名与别名模板}案例，见\texttt{Chapter 1}\textit{现代C++语言核心特性}

\end{itemize}


%%C020341 h2

\section{用正则表达式验证字符串格式}


%%C020342 p

正则表达式是一种用于在文本中做模式匹配与替换的语言。C++11通过\texttt{<\allowbreak{}regex>}头文件中提供的一组类、算法和迭代器，在标准库中支持正则表达式。本案例将说明如何用正则表达式验证字符串是否匹配某个模式（例如验证电子邮件地址或IP地址的格式）。


%%C020343 h3

\subsection{准备工作}


%%C020344 p

本案例中，我们会在必要时解释所用正则表达式的细节。不过，要想使用C++标准库的正则表达式功能，你至少应当对正则表达式有一些基本了解。正则表达式语法与标准的完整说明超出了本书的范围；如果你不熟悉正则表达式，建议先多读一些相关资料，再继续学习与正则表达式相关的案例。


%%C020345 h3

\subsection{具体做法}


%%C020346 p

要验证字符串是否匹配某个正则表达式，请按以下步骤操作：


%%C020347 p

包含头文件\texttt{<\allowbreak{}regex>}与\texttt{<\allowbreak{}string>}，以及命名空间


%%C020348 p

\myttbr{std::\allowbreak{}string\_\allowbreak{}literals}，以便使用C++14标准中针对字符串的用户定义字面量：


%%C020349 code

\begin{codeblockbg}
\begin{Verbatim}[commandchars=\\\{\},breaklines=true,breakanywhere=true,fontsize=\codefont,xleftmargin=17bp]
\PYG{c+cp}{\PYGZsh{}}\PYG{c+cp}{include}\PYG{+w}{ }\PYG{c+cpf}{\PYGZlt{}regex\PYGZgt{}}
\PYG{c+cp}{\PYGZsh{}}\PYG{c+cp}{include}\PYG{+w}{ }\PYG{c+cpf}{\PYGZlt{}string\PYGZgt{}}
\PYG{k}{using}\PYG{+w}{ }\PYG{k}{namespace}\PYG{+w}{ }\PYG{n+nn}{std}\PYG{o}{:}\PYG{o}{:}\PYG{n+nn}{string\PYGZus{}literals}\PYG{p}{;}
\end{Verbatim}
\end{codeblockbg}


%%C020350 p

用原始字符串字面量来书写正则表达式，以避免转义


%%C020351 p

反斜杠（它出现的频率可能很高）。下面这个正则表达式可以校验大多数电子邮件格式：


%%C020352 code

\begin{codeblockbg}
\begin{Verbatim}[commandchars=\\\{\},breaklines=true,breakanywhere=true,fontsize=\codefont,xleftmargin=17bp]
\PYG{k}{auto}\PYG{+w}{ }\PYG{n}{pattern}\PYG{+w}{ }\PYG{p}{\PYGZob{}}\PYG{l+s+sa}{R}\PYG{l+s}{\PYGZdq{}}\PYG{l+s+dl}{(}\PYG{l+s}{\PYGZca{}[A\PYGZhy{}Z0\PYGZhy{}9.\PYGZus{}\PYGZpc{}+\PYGZhy{}]+@[A\PYGZhy{}Z0\PYGZhy{}9.\PYGZhy{}]+\PYGZbs{}.[A\PYGZhy{}Z]\PYGZob{}2,\PYGZcb{}\PYGZdl{}}\PYG{l+s+dl}{)}\PYG{l+s}{\PYGZdq{}}\PYG{n}{s}\PYG{p}{\PYGZcb{}}\PYG{p}{;}
\end{Verbatim}
\end{codeblockbg}


%%C020353 p

创建一个\texttt{std::\allowbreak{}regex}/\texttt{std::\allowbreak{}wregex}对象（取决于所使用的字符集）


%%C020354 p

来封装该正则表达式：


%%C020355 code

\begin{codeblockbg}
\begin{Verbatim}[commandchars=\\\{\},breaklines=true,breakanywhere=true,fontsize=\codefont,xleftmargin=17bp]
\PYG{k}{auto}\PYG{+w}{ }\PYG{n}{rx}\PYG{+w}{ }\PYG{o}{=}\PYG{+w}{ }\PYG{n}{std}\PYG{o}{:}\PYG{o}{:}\PYG{n}{regex}\PYG{p}{\PYGZob{}}\PYG{n}{pattern}\PYG{p}{\PYGZcb{}}\PYG{p}{;}
\end{Verbatim}
\end{codeblockbg}


%%C020356 p

若要忽略大小写或指定其他解析选项，请使用带有


%%C020357 p

正则表达式标志这一额外参数的重载构造函数：


%%C020358 code

\begin{codeblockbg}
\begin{Verbatim}[commandchars=\\\{\},breaklines=true,breakanywhere=true,fontsize=\codefont,xleftmargin=17bp]
\PYG{k}{auto}\PYG{+w}{ }\PYG{n}{rx}\PYG{+w}{ }\PYG{o}{=}\PYG{+w}{ }\PYG{n}{std}\PYG{o}{:}\PYG{o}{:}\PYG{n}{regex}\PYG{p}{\PYGZob{}}\PYG{n}{pattern}\PYG{p}{,}\PYG{+w}{ }\PYG{n}{std}\PYG{o}{:}\PYG{o}{:}\PYG{n}{regex\PYGZus{}constants}\PYG{o}{:}\PYG{o}{:}\PYG{n}{icase}\PYG{p}{\PYGZcb{}}\PYG{p}{;}
\end{Verbatim}
\end{codeblockbg}


%%C020359 p

用\myttbr{std::\allowbreak{}regex\_\allowbreak{}match()}把正则表达式与整个字符串进行匹配：


%%C020360 code

\begin{codeblockbg}
\begin{Verbatim}[commandchars=\\\{\},breaklines=true,breakanywhere=true,fontsize=\codefont,xleftmargin=17bp]
\PYG{k}{auto}\PYG{+w}{ }\PYG{n}{valid}\PYG{+w}{ }\PYG{o}{=}\PYG{+w}{ }\PYG{n}{std}\PYG{o}{:}\PYG{o}{:}\PYG{n}{regex\PYGZus{}match}\PYG{p}{(}\PYG{l+s}{\PYGZdq{}}\PYG{l+s}{marius@domain.com}\PYG{l+s}{\PYGZdq{}}\PYG{n}{s}\PYG{p}{,}\PYG{+w}{ }\PYG{n}{rx}\PYG{p}{)}\PYG{p}{;}
\end{Verbatim}
\end{codeblockbg}


%%C020361 h3

\subsection{工作原理}


%%C020362 p

以校验电子邮件地址格式这个问题为例，尽管它看起来很简单，实践中却很难找到一个能覆盖所有合法电子邮件格式的简单正则表达式。本案例并不打算找出那个终极正则表达式，而是使用一个对大多数情况都足够好的正则表达式。为此我们使用的正则表达式如下：


%%C020363 code

\begin{codeblockbg}
\begin{Verbatim}[commandchars=\\\{\},breaklines=true,breakanywhere=true,fontsize=\codefont,xleftmargin=17bp]
\PYG{o}{\PYGZca{}}\PYG{p}{[}\PYG{n}{A}\PYG{o}{\PYGZhy{}}\PYG{n}{Z0}\PYG{l+m+mf}{\PYGZhy{}9.}\PYG{n}{\PYGZus{}}\PYG{o}{\PYGZpc{}}\PYG{o}{+}\PYG{o}{\PYGZhy{}}\PYG{p}{]}\PYG{o}{+}\PYG{err}{@}\PYG{p}{[}\PYG{n}{A}\PYG{o}{\PYGZhy{}}\PYG{n}{Z0}\PYG{l+m+mf}{\PYGZhy{}9.}\PYG{o}{\PYGZhy{}}\PYG{p}{]}\PYG{o}{+}\PYG{err}{\PYGZbs{}}\PYG{p}{.}\PYG{p}{[}\PYG{n}{A}\PYG{o}{\PYGZhy{}}\PYG{n}{Z}\PYG{p}{]}\PYG{p}{\PYGZob{}}\PYG{l+m+mi}{2}\PYG{p}{,}\PYG{p}{\PYGZcb{}}\PYG{n}{\PYGZdl{}}
\end{Verbatim}
\end{codeblockbg}


%%C020364 p

下表解释了该正则表达式的结构：


%%C020365 table

\myTableBegin
\begin{xltabular}{\textwidth}{|>{\hsize=0.212\hsize\raggedright\arraybackslash}X|>{\hsize=0.788\hsize\raggedright\arraybackslash}X|}
\hline
\rowcolor{tabhead} \textcolor{white}{\textbf{组成部分}} & \textcolor{white}{\textbf{说明}} \\
\hline
\endfirsthead
\rowcolor{tabhead} \textcolor{white}{\textbf{组成部分}} & \textcolor{white}{\textbf{说明}} \\
\hline
\endhead
\myttbr{\textasciicircum{}} & 字符串开头\\
\hline
\myttbr{[A-\allowbreak{}Z0-\allowbreak{}9.\_\%+-\allowbreak{}]+} & 至少一个A-Z、0-9范围内的字符，或 -、\%，+、- 之一，表示电子邮件地址的本地部分\\
\hline
\myttbr{@} & 字符@ \\
\hline
\myttbr{[A-\allowbreak{}Z0-\allowbreak{}9.-\allowbreak{}]+} & 至少一个A-Z、0-9范围内的字符，或 -、\%，+、- 之一，表示域名部分的主机名\\
\hline
\myttbr{\textbackslash{}.} & 分隔域主机名与标签的点\\
\hline
\myttbr{[A-Z]\{2,\}} & 域的DNS标签，长度为2到63个字符\\
\hline
\myttbr{\$} & 字符串结尾\\
\hline
\end{xltabular}
\myTableEnd


%%C020366 p

要记住，实践中域名由主机名加上一串以点分隔的DNS标签组成，例如\texttt{localhost}、\texttt{gmail.\allowbreak{}com}或\texttt{yahoo.\allowbreak{}co.\allowbreak{}uk}。我们使用的这个正则表达式无法匹配没有DNS标签的域名，比如localhost（像\texttt{root@localhost}这样的电子邮件地址是合法的）。域名也可以是用方括号括起来的IP地址，例如\texttt{[192.\allowbreak{}168.\allowbreak{}100.\allowbreak{}11]}（如\myttbr{john.\allowbreak{}doe@[192.\allowbreak{}168.\allowbreak{}100.\allowbreak{}11]}）。包含这类域名的电子邮件地址不会匹配上面定义的正则表达式。尽管这些比较少见的形式无法匹配，该正则表达式仍能覆盖大多数电子邮件格式。


%%C020367 note

\begin{myNote}[tip]
本章示例中的正则表达式仅用于教学目的，不应原样用于生产代码。如前所述，这个示例并未覆盖所有可能的电子邮件格式。
\end{myNote}


%%C020368 p

我们首先包含必要的头文件：\texttt{<\allowbreak{}regex>}用于正则表达式，\texttt{<\allowbreak{}string>}用于字符串。下面给出的\texttt{is\_\allowbreak{}valid\_\allowbreak{}email()}函数（内容基本上就是\textit{具体做法}一节的示例）接收一个表示电子邮件地址的字符串，返回一个布尔值，指示该电子邮件的格式是否合法。我们首先构造一个\texttt{std::\allowbreak{}regex}对象，封装由原始字符串字面量给出的正则表达式。使用原始字符串字面量很有帮助，因为这样可以避免转义反斜杠——反斜杠在正则表达式中同样用于转义字符。随后函数调用\myttbr{std::\allowbreak{}regex\_\allowbreak{}match()}，传入输入文本和正则表达式：


%%C020369 code

\begin{codeblockbg}
\begin{Verbatim}[commandchars=\\\{\},breaklines=true,breakanywhere=true,fontsize=\codefont,xleftmargin=17bp]
\PYG{k+kt}{bool}\PYG{+w}{ }\PYG{n+nf}{is\PYGZus{}valid\PYGZus{}email\PYGZus{}format}\PYG{p}{(}\PYG{n}{std}\PYG{o}{:}\PYG{o}{:}\PYG{n}{string}\PYG{+w}{ }\PYG{k}{const}\PYG{+w}{ }\PYG{o}{\PYGZam{}}\PYG{+w}{ }\PYG{n}{email}\PYG{p}{)}
\PYG{p}{\PYGZob{}}
\PYG{+w}{  }\PYG{k}{auto}\PYG{+w}{ }\PYG{n}{pattern}\PYG{+w}{ }\PYG{p}{\PYGZob{}}\PYG{l+s+sa}{R}\PYG{l+s}{\PYGZdq{}}\PYG{l+s+dl}{(}\PYG{l+s}{\PYGZca{}[A\PYGZhy{}Z0\PYGZhy{}9.\PYGZus{}\PYGZpc{}+\PYGZhy{}]+@[A\PYGZhy{}Z0\PYGZhy{}9.\PYGZhy{}]+\PYGZbs{}.[A\PYGZhy{}Z]\PYGZob{}2,\PYGZcb{}\PYGZdl{}}\PYG{l+s+dl}{)}\PYG{l+s}{\PYGZdq{}}\PYG{n}{s}\PYG{p}{\PYGZcb{}}\PYG{p}{;}
\PYG{k}{auto}\PYG{+w}{ }\PYG{n}{rx}\PYG{+w}{ }\PYG{o}{=}\PYG{+w}{ }\PYG{n}{std}\PYG{o}{:}\PYG{o}{:}\PYG{n}{regex}\PYG{p}{\PYGZob{}}\PYG{n}{pattern}\PYG{p}{\PYGZcb{}}\PYG{p}{;}
\PYG{+w}{  }\PYG{k}{return}\PYG{+w}{ }\PYG{n}{std}\PYG{o}{:}\PYG{o}{:}\PYG{n}{regex\PYGZus{}match}\PYG{p}{(}\PYG{n}{email}\PYG{p}{,}\PYG{+w}{ }\PYG{n}{rx}\PYG{p}{)}\PYG{p}{;}
\PYG{p}{\PYGZcb{}}
\end{Verbatim}
\end{codeblockbg}


%%C020370 p

\myttbr{std::\allowbreak{}regex\_\allowbreak{}match()}方法尝试把正则表达式与整个字符串匹配。匹配成功则返回\texttt{true}，否则返回\texttt{false}：


%%C020371 code

\begin{codeblockbg}
\begin{Verbatim}[commandchars=\\\{\},breaklines=true,breakanywhere=true,fontsize=\codefont,xleftmargin=17bp]
\PYG{k}{auto}\PYG{+w}{ }\PYG{n}{ltest}\PYG{+w}{ }\PYG{o}{=}\PYG{+w}{ }\PYG{p}{[}\PYG{p}{]}\PYG{p}{(}\PYG{n}{std}\PYG{o}{:}\PYG{o}{:}\PYG{n}{string}\PYG{+w}{ }\PYG{k}{const}\PYG{+w}{ }\PYG{o}{\PYGZam{}}\PYG{+w}{ }\PYG{n}{email}\PYG{p}{)}
\PYG{p}{\PYGZob{}}
\PYG{+w}{  }\PYG{n}{std}\PYG{o}{:}\PYG{o}{:}\PYG{n}{cout}\PYG{+w}{ }\PYG{o}{\PYGZlt{}}\PYG{o}{\PYGZlt{}}\PYG{+w}{ }\PYG{n}{std}\PYG{o}{:}\PYG{o}{:}\PYG{n}{setw}\PYG{p}{(}\PYG{l+m+mi}{30}\PYG{p}{)}\PYG{+w}{ }\PYG{o}{\PYGZlt{}}\PYG{o}{\PYGZlt{}}\PYG{+w}{ }\PYG{n}{std}\PYG{o}{:}\PYG{o}{:}\PYG{n}{left}
\PYG{+w}{            }\PYG{o}{\PYGZlt{}}\PYG{o}{\PYGZlt{}}\PYG{+w}{ }\PYG{n}{email}\PYG{+w}{ }\PYG{o}{\PYGZlt{}}\PYG{o}{\PYGZlt{}}\PYG{+w}{ }\PYG{l+s}{\PYGZdq{}}\PYG{l+s}{ : }\PYG{l+s}{\PYGZdq{}}
\PYG{+w}{            }\PYG{o}{\PYGZlt{}}\PYG{o}{\PYGZlt{}}\PYG{+w}{ }\PYG{p}{(}\PYG{n}{is\PYGZus{}valid\PYGZus{}email\PYGZus{}format}\PYG{p}{(}\PYG{n}{email}\PYG{p}{)}\PYG{+w}{ }\PYG{o}{?}
\PYG{+w}{               }\PYG{l+s}{\PYGZdq{}}\PYG{l+s}{valid format}\PYG{l+s}{\PYGZdq{}}\PYG{+w}{ }\PYG{o}{:}\PYG{+w}{ }\PYG{l+s}{\PYGZdq{}}\PYG{l+s}{invalid format}\PYG{l+s}{\PYGZdq{}}\PYG{p}{)}
\PYG{+w}{            }\PYG{o}{\PYGZlt{}}\PYG{o}{\PYGZlt{}}\PYG{+w}{ }\PYG{n}{std}\PYG{o}{:}\PYG{o}{:}\PYG{n}{endl}\PYG{p}{;}
\PYG{p}{\PYGZcb{}}\PYG{p}{;}
\PYG{n}{ltest}\PYG{p}{(}\PYG{l+s}{\PYGZdq{}}\PYG{l+s}{JOHN.DOE@DOMAIN.COM}\PYG{l+s}{\PYGZdq{}}\PYG{n}{s}\PYG{p}{)}\PYG{p}{;}\PYG{+w}{         }\PYG{c+c1}{// valid format}
\PYG{n}{ltest}\PYG{p}{(}\PYG{l+s}{\PYGZdq{}}\PYG{l+s}{JOHNDOE@DOMAIL.CO.UK}\PYG{l+s}{\PYGZdq{}}\PYG{n}{s}\PYG{p}{)}\PYG{p}{;}\PYG{+w}{        }\PYG{c+c1}{// valid format}
\PYG{n}{ltest}\PYG{p}{(}\PYG{l+s}{\PYGZdq{}}\PYG{l+s}{JOHNDOE@DOMAIL.INFO}\PYG{l+s}{\PYGZdq{}}\PYG{n}{s}\PYG{p}{)}\PYG{p}{;}\PYG{+w}{         }\PYG{c+c1}{// valid format}
\PYG{n}{ltest}\PYG{p}{(}\PYG{l+s}{\PYGZdq{}}\PYG{l+s}{J.O.H.N\PYGZus{}D.O.E@DOMAIN.INFO}\PYG{l+s}{\PYGZdq{}}\PYG{n}{s}\PYG{p}{)}\PYG{p}{;}\PYG{+w}{   }\PYG{c+c1}{// valid format}
\PYG{n}{ltest}\PYG{p}{(}\PYG{l+s}{\PYGZdq{}}\PYG{l+s}{ROOT@LOCALHOST}\PYG{l+s}{\PYGZdq{}}\PYG{n}{s}\PYG{p}{)}\PYG{p}{;}\PYG{+w}{              }\PYG{c+c1}{// invalid format}
\PYG{n}{ltest}\PYG{p}{(}\PYG{l+s}{\PYGZdq{}}\PYG{l+s}{john.doe@domain.com}\PYG{l+s}{\PYGZdq{}}\PYG{n}{s}\PYG{p}{)}\PYG{p}{;}\PYG{+w}{         }\PYG{c+c1}{// invalid format}
\end{Verbatim}
\end{codeblockbg}


%%C020372 p

在这个简单的测试中，唯一不匹配正则表达式的电子邮件是\texttt{ROOT@LOCALHOST}和\myttbr{john.\allowbreak{}doe@domain.\allowbreak{}com}。前者包含一个没有带点DNS标签的域名，这种情况未被该正则表达式覆盖。后者只含小写字母，而正则表达式中本地部分与域名允许的字符集是大写字母A到Z。


%%C020373 p

与其为正则表达式增加更多合法字符（例如\texttt{[A-\allowbreak{}Za-\allowbreak{}z0-\allowbreak{}9.\_\%+-\allowbreak{}]}）而使它变复杂，不如直接声明匹配时忽略大小写。这可以通过给\texttt{std::\allowbreak{}basic\_\allowbreak{}regex}类的构造函数多传一个参数来实现。为此可用的常量定义在\texttt{regex\_\allowbreak{}constants}命名空间中。对\myttbr{is\_\allowbreak{}valid\_\allowbreak{}email\_\allowbreak{}format()}做如下细微改动，就能忽略大小写，让同时含小写和大写字母的电子邮件正确匹配该正则表达式：


%%C020374 code

\begin{codeblockbg}
\begin{Verbatim}[commandchars=\\\{\},breaklines=true,breakanywhere=true,fontsize=\codefont,xleftmargin=17bp]
\PYG{k+kt}{bool}\PYG{+w}{ }\PYG{n+nf}{is\PYGZus{}valid\PYGZus{}email\PYGZus{}format}\PYG{p}{(}\PYG{n}{std}\PYG{o}{:}\PYG{o}{:}\PYG{n}{string}\PYG{+w}{ }\PYG{k}{const}\PYG{+w}{ }\PYG{o}{\PYGZam{}}\PYG{+w}{ }\PYG{n}{email}\PYG{p}{)}
\PYG{p}{\PYGZob{}}
\PYG{+w}{  }\PYG{k}{auto}\PYG{+w}{ }\PYG{n}{rx}\PYG{+w}{ }\PYG{o}{=}\PYG{+w}{ }\PYG{n}{std}\PYG{o}{:}\PYG{o}{:}\PYG{n}{regex}\PYG{p}{\PYGZob{}}
\PYG{+w}{    }\PYG{l+s+sa}{R}\PYG{l+s}{\PYGZdq{}}\PYG{l+s+dl}{(}\PYG{l+s}{\PYGZca{}[A\PYGZhy{}Z0\PYGZhy{}9.\PYGZus{}\PYGZpc{}+\PYGZhy{}]+@[A\PYGZhy{}Z0\PYGZhy{}9.\PYGZhy{}]+\PYGZbs{}.[A\PYGZhy{}Z]\PYGZob{}2,\PYGZcb{}\PYGZdl{}}\PYG{l+s+dl}{)}\PYG{l+s}{\PYGZdq{}}\PYG{n}{s}\PYG{p}{,}
\PYG{+w}{    }\PYG{n}{std}\PYG{o}{:}\PYG{o}{:}\PYG{n}{regex\PYGZus{}constants}\PYG{o}{:}\PYG{o}{:}\PYG{n}{icase}\PYG{p}{\PYGZcb{}}\PYG{p}{;}
\PYG{+w}{  }\PYG{k}{return}\PYG{+w}{ }\PYG{n}{std}\PYG{o}{:}\PYG{o}{:}\PYG{n}{regex\PYGZus{}match}\PYG{p}{(}\PYG{n}{email}\PYG{p}{,}\PYG{+w}{ }\PYG{n}{rx}\PYG{p}{)}\PYG{p}{;}
\PYG{p}{\PYGZcb{}}
\end{Verbatim}
\end{codeblockbg}


%%C020375 p

这个\myttbr{is\_\allowbreak{}valid\_\allowbreak{}email\_\allowbreak{}format()}函数相当简单；如果把正则表达式与待匹配文本一起作为参数传入，它就能用于任何匹配任务。不过，若能用一个函数同时处理多字节字符串（\texttt{std::\allowbreak{}string}）和宽字符串（\texttt{std::\allowbreak{}wstring}）就更好了。只要创建一个以字符类型为模板参数的函数模板，就能做到这一点：


%%C020376 code

\begin{codeblockbg}
\begin{Verbatim}[commandchars=\\\{\},breaklines=true,breakanywhere=true,fontsize=\codefont,xleftmargin=17bp]
\PYG{k}{template}\PYG{+w}{ }\PYG{o}{\PYGZlt{}}\PYG{k}{typename}\PYG{+w}{ }\PYG{n+nc}{CharT}\PYG{o}{\PYGZgt{}}
\PYG{k}{using}\PYG{+w}{ }\PYG{n}{tstring}\PYG{+w}{ }\PYG{o}{=}\PYG{+w}{ }\PYG{n}{std}\PYG{o}{:}\PYG{o}{:}\PYG{n}{basic\PYGZus{}string}\PYG{o}{\PYGZlt{}}\PYG{n}{CharT}\PYG{p}{,}\PYG{+w}{ }\PYG{n}{std}\PYG{o}{:}\PYG{o}{:}\PYG{n}{char\PYGZus{}traits}\PYG{o}{\PYGZlt{}}\PYG{n}{CharT}\PYG{o}{\PYGZgt{}}\PYG{p}{,}
\PYG{+w}{                                  }\PYG{n}{std}\PYG{o}{:}\PYG{o}{:}\PYG{n}{allocator}\PYG{o}{\PYGZlt{}}\PYG{n}{CharT}\PYG{o}{\PYGZgt{}}\PYG{o}{\PYGZgt{}}\PYG{p}{;}
\PYG{k}{template}\PYG{+w}{ }\PYG{o}{\PYGZlt{}}\PYG{k}{typename}\PYG{+w}{ }\PYG{n+nc}{CharT}\PYG{o}{\PYGZgt{}}
\PYG{k+kt}{bool}\PYG{+w}{ }\PYG{n}{is\PYGZus{}valid\PYGZus{}format}\PYG{p}{(}\PYG{n}{tstring}\PYG{o}{\PYGZlt{}}\PYG{n}{CharT}\PYG{o}{\PYGZgt{}}\PYG{+w}{ }\PYG{k}{const}\PYG{+w}{ }\PYG{o}{\PYGZam{}}\PYG{+w}{ }\PYG{n}{pattern}\PYG{p}{,}
\PYG{+w}{                     }\PYG{n}{tstring}\PYG{o}{\PYGZlt{}}\PYG{n}{CharT}\PYG{o}{\PYGZgt{}}\PYG{+w}{ }\PYG{k}{const}\PYG{+w}{ }\PYG{o}{\PYGZam{}}\PYG{+w}{ }\PYG{n}{text}\PYG{p}{)}
\PYG{p}{\PYGZob{}}
\PYG{+w}{  }\PYG{k}{auto}\PYG{+w}{ }\PYG{n}{rx}\PYG{+w}{ }\PYG{o}{=}\PYG{+w}{ }\PYG{n}{std}\PYG{o}{:}\PYG{o}{:}\PYG{n}{basic\PYGZus{}regex}\PYG{o}{\PYGZlt{}}\PYG{n}{CharT}\PYG{o}{\PYGZgt{}}\PYG{p}{\PYGZob{}}
\PYG{+w}{    }\PYG{n}{pattern}\PYG{p}{,}\PYG{+w}{ }\PYG{n}{std}\PYG{o}{:}\PYG{o}{:}\PYG{n}{regex\PYGZus{}constants}\PYG{o}{:}\PYG{o}{:}\PYG{n}{icase}\PYG{+w}{ }\PYG{p}{\PYGZcb{}}\PYG{p}{;}
\PYG{+w}{  }\PYG{k}{return}\PYG{+w}{ }\PYG{n}{std}\PYG{o}{:}\PYG{o}{:}\PYG{n}{regex\PYGZus{}match}\PYG{p}{(}\PYG{n}{text}\PYG{p}{,}\PYG{+w}{ }\PYG{n}{rx}\PYG{p}{)}\PYG{p}{;}
\PYG{p}{\PYGZcb{}}
\end{Verbatim}
\end{codeblockbg}


%%C020377 p

我们首先为\texttt{std::\allowbreak{}basic\_\allowbreak{}string}创建一个别名模板，以简化它的用法。新的\myttbr{is\_\allowbreak{}valid\_\allowbreak{}format()}函数是一个函数模板，与我们的\texttt{is\_\allowbreak{}valid\_\allowbreak{}email()}实现非常相似\textbf{. }不过现在我们使用\myttbr{std::\allowbreak{}basic\_\allowbreak{}regex<\allowbreak{}CharT>}，而不是\texttt{std::\allowbreak{}regex}这个typedef（它等价于\myttbr{std::\allowbreak{}basic\_\allowbreak{}regex<\allowbreak{}char>}），并且模式作为第一个参数传入。接着我们为宽字符串实现一个新函数\myttbr{is\_\allowbreak{}valid\_\allowbreak{}email\_\allowbreak{}format\_\allowbreak{}w()}，它依赖这个函数模板。这个函数模板也可以复用来实现其他校验，例如车牌号是否符合特定格式：


%%C020378 code

\begin{codeblockbg}
\begin{Verbatim}[commandchars=\\\{\},breaklines=true,breakanywhere=true,fontsize=\codefont,xleftmargin=17bp]
\PYG{k+kt}{bool}\PYG{+w}{ }\PYG{n+nf}{is\PYGZus{}valid\PYGZus{}email\PYGZus{}format\PYGZus{}w}\PYG{p}{(}\PYG{n}{std}\PYG{o}{:}\PYG{o}{:}\PYG{n}{wstring}\PYG{+w}{ }\PYG{k}{const}\PYG{+w}{ }\PYG{o}{\PYGZam{}}\PYG{+w}{ }\PYG{n}{text}\PYG{p}{)}
\PYG{p}{\PYGZob{}}
\PYG{+w}{  }\PYG{k}{return}\PYG{+w}{ }\PYG{n}{is\PYGZus{}valid\PYGZus{}format}\PYG{p}{(}
\PYG{+w}{    }\PYG{l+s+sa}{LR}\PYG{l+s}{\PYGZdq{}}\PYG{l+s+dl}{(}\PYG{l+s}{\PYGZca{}[A\PYGZhy{}Z0\PYGZhy{}9.\PYGZus{}\PYGZpc{}+\PYGZhy{}]+@[A\PYGZhy{}Z0\PYGZhy{}9.\PYGZhy{}]+\PYGZbs{}.[A\PYGZhy{}Z]\PYGZob{}2,\PYGZcb{}\PYGZdl{}}\PYG{l+s+dl}{)}\PYG{l+s}{\PYGZdq{}}\PYG{n}{s}\PYG{p}{,}
\PYG{+w}{    }\PYG{n}{text}\PYG{p}{)}\PYG{p}{;}
\PYG{p}{\PYGZcb{}}
\PYG{k}{auto}\PYG{+w}{ }\PYG{n}{ltest2}\PYG{+w}{ }\PYG{o}{=}\PYG{+w}{ }\PYG{p}{[}\PYG{p}{]}\PYG{p}{(}\PYG{k}{auto}\PYG{+w}{ }\PYG{k}{const}\PYG{+w}{ }\PYG{o}{\PYGZam{}}\PYG{+w}{ }\PYG{n}{email}\PYG{p}{)}
\PYG{p}{\PYGZob{}}
\PYG{+w}{  }\PYG{n}{std}\PYG{o}{:}\PYG{o}{:}\PYG{n}{wcout}\PYG{+w}{ }\PYG{o}{\PYGZlt{}}\PYG{o}{\PYGZlt{}}\PYG{+w}{ }\PYG{n}{std}\PYG{o}{:}\PYG{o}{:}\PYG{n}{setw}\PYG{p}{(}\PYG{l+m+mi}{30}\PYG{p}{)}\PYG{+w}{ }\PYG{o}{\PYGZlt{}}\PYG{o}{\PYGZlt{}}\PYG{+w}{ }\PYG{n}{std}\PYG{o}{:}\PYG{o}{:}\PYG{n}{left}
\PYG{+w}{     }\PYG{o}{\PYGZlt{}}\PYG{o}{\PYGZlt{}}\PYG{+w}{ }\PYG{n}{email}\PYG{+w}{ }\PYG{o}{\PYGZlt{}}\PYG{o}{\PYGZlt{}}\PYG{+w}{ }\PYG{l+s+sa}{L}\PYG{l+s}{\PYGZdq{}}\PYG{l+s}{ : }\PYG{l+s}{\PYGZdq{}}
\PYG{+w}{     }\PYG{o}{\PYGZlt{}}\PYG{o}{\PYGZlt{}}\PYG{+w}{ }\PYG{p}{(}\PYG{n}{is\PYGZus{}valid\PYGZus{}email\PYGZus{}format\PYGZus{}w}\PYG{p}{(}\PYG{n}{email}\PYG{p}{)}\PYG{+w}{ }\PYG{o}{?}\PYG{+w}{ }\PYG{l+s+sa}{L}\PYG{l+s}{\PYGZdq{}}\PYG{l+s}{valid}\PYG{l+s}{\PYGZdq{}}\PYG{+w}{ }\PYG{o}{:}\PYG{+w}{ }\PYG{l+s+sa}{L}\PYG{l+s}{\PYGZdq{}}\PYG{l+s}{invalid}\PYG{l+s}{\PYGZdq{}}\PYG{p}{)}
\PYG{+w}{     }\PYG{o}{\PYGZlt{}}\PYG{o}{\PYGZlt{}}\PYG{+w}{ }\PYG{n}{std}\PYG{o}{:}\PYG{o}{:}\PYG{n}{endl}\PYG{p}{;}
\PYG{p}{\PYGZcb{}}\PYG{p}{;}
\PYG{n}{ltest2}\PYG{p}{(}\PYG{l+s+sa}{L}\PYG{l+s}{\PYGZdq{}}\PYG{l+s}{JOHN.DOE@DOMAIN.COM}\PYG{l+s}{\PYGZdq{}}\PYG{n}{s}\PYG{p}{)}\PYG{p}{;}\PYG{+w}{       }\PYG{c+c1}{// valid}
\PYG{n}{ltest2}\PYG{p}{(}\PYG{l+s+sa}{L}\PYG{l+s}{\PYGZdq{}}\PYG{l+s}{JOHNDOE@DOMAIL.CO.UK}\PYG{l+s}{\PYGZdq{}}\PYG{n}{s}\PYG{p}{)}\PYG{p}{;}\PYG{+w}{      }\PYG{c+c1}{// valid}
\PYG{n}{ltest2}\PYG{p}{(}\PYG{l+s+sa}{L}\PYG{l+s}{\PYGZdq{}}\PYG{l+s}{JOHNDOE@DOMAIL.INFO}\PYG{l+s}{\PYGZdq{}}\PYG{n}{s}\PYG{p}{)}\PYG{p}{;}\PYG{+w}{       }\PYG{c+c1}{// valid}
\PYG{n}{ltest2}\PYG{p}{(}\PYG{l+s+sa}{L}\PYG{l+s}{\PYGZdq{}}\PYG{l+s}{J.O.H.N\PYGZus{}D.O.E@DOMAIN.INFO}\PYG{l+s}{\PYGZdq{}}\PYG{n}{s}\PYG{p}{)}\PYG{p}{;}\PYG{+w}{ }\PYG{c+c1}{// valid}
\PYG{n}{ltest2}\PYG{p}{(}\PYG{l+s+sa}{L}\PYG{l+s}{\PYGZdq{}}\PYG{l+s}{ROOT@LOCALHOST}\PYG{l+s}{\PYGZdq{}}\PYG{n}{s}\PYG{p}{)}\PYG{p}{;}\PYG{+w}{            }\PYG{c+c1}{// invalid}
\PYG{n}{ltest2}\PYG{p}{(}\PYG{l+s+sa}{L}\PYG{l+s}{\PYGZdq{}}\PYG{l+s}{john.doe@domain.com}\PYG{l+s}{\PYGZdq{}}\PYG{n}{s}\PYG{p}{)}\PYG{p}{;}\PYG{+w}{       }\PYG{c+c1}{// valid}
\end{Verbatim}
\end{codeblockbg}


%%C020379 p

上面所有示例中，唯一不匹配的还是\texttt{ROOT@LOCAHOST}，这与预期一致。


%%C020380 p

\myttbr{std::\allowbreak{}regex\_\allowbreak{}match()}方法实际上有多个重载版本，其中一些带一个指向\myttbr{std::\allowbreak{}match\_\allowbreak{}results}对象的引用参数，用来保存匹配结果。如果没有匹配，\myttbr{std::\allowbreak{}match\_\allowbreak{}results}为空，其大小为0。反之，若有匹配，\myttbr{std::\allowbreak{}match\_\allowbreak{}results}对象非空，其大小为1加上匹配到的子表达式个数。


%%C020381 p

下面这个版本的函数使用了上述重载，并返回一个\texttt{std::\allowbreak{}smatch}对象，其中存放匹配到的子表达式。注意正则表达式有所改变，定义了三个捕获组——一个对应本地部分，一个对应域名的主机名部分，还有一个对应DNS标签。若匹配成功，\texttt{std::\allowbreak{}smatch}对象将包含四个子匹配对象：第一个匹配整个字符串，第二个对应第一个捕获组（本地部分），第三个对应第二个捕获组（主机名），第四个对应第三个也是最后一个捕获组（DNS标签）。结果以tuple返回，其中第一项表示成功还是失败：


%%C020382 code

\begin{codeblockbg}
\begin{Verbatim}[commandchars=\\\{\},breaklines=true,breakanywhere=true,fontsize=\codefont,xleftmargin=17bp]
\PYG{n}{std}\PYG{o}{:}\PYG{o}{:}\PYG{n}{tuple}\PYG{o}{\PYGZlt{}}\PYG{k+kt}{bool}\PYG{p}{,}\PYG{+w}{ }\PYG{n}{std}\PYG{o}{:}\PYG{o}{:}\PYG{n}{string}\PYG{p}{,}\PYG{+w}{ }\PYG{n}{std}\PYG{o}{:}\PYG{o}{:}\PYG{n}{string}\PYG{p}{,}\PYG{+w}{ }\PYG{n}{std}\PYG{o}{:}\PYG{o}{:}\PYG{n}{string}\PYG{o}{\PYGZgt{}}
\PYG{n}{is\PYGZus{}valid\PYGZus{}email\PYGZus{}format\PYGZus{}with\PYGZus{}result}\PYG{p}{(}\PYG{n}{std}\PYG{o}{:}\PYG{o}{:}\PYG{n}{string}\PYG{+w}{ }\PYG{k}{const}\PYG{+w}{ }\PYG{o}{\PYGZam{}}\PYG{+w}{ }\PYG{n}{email}\PYG{p}{)}
\PYG{p}{\PYGZob{}}
\PYG{+w}{  }\PYG{k}{auto}\PYG{+w}{ }\PYG{n}{rx}\PYG{+w}{ }\PYG{o}{=}\PYG{+w}{ }\PYG{n}{std}\PYG{o}{:}\PYG{o}{:}\PYG{n}{regex}\PYG{p}{\PYGZob{}}
\PYG{+w}{    }\PYG{l+s+sa}{R}\PYG{l+s}{\PYGZdq{}}\PYG{l+s+dl}{(}\PYG{l+s}{\PYGZca{}([A\PYGZhy{}Z0\PYGZhy{}9.\PYGZus{}\PYGZpc{}+\PYGZhy{}]+)@([A\PYGZhy{}Z0\PYGZhy{}9.\PYGZhy{}]+)\PYGZbs{}.([A\PYGZhy{}Z]\PYGZob{}2,\PYGZcb{})\PYGZdl{}}\PYG{l+s+dl}{)}\PYG{l+s}{\PYGZdq{}}\PYG{n}{s}\PYG{p}{,}
\PYG{+w}{    }\PYG{n}{std}\PYG{o}{:}\PYG{o}{:}\PYG{n}{regex\PYGZus{}constants}\PYG{o}{:}\PYG{o}{:}\PYG{n}{icase}\PYG{+w}{ }\PYG{p}{\PYGZcb{}}\PYG{p}{;}
\PYG{+w}{  }\PYG{k}{auto}\PYG{+w}{ }\PYG{n}{result}\PYG{+w}{ }\PYG{o}{=}\PYG{+w}{ }\PYG{n}{std}\PYG{o}{:}\PYG{o}{:}\PYG{n}{smatch}\PYG{p}{\PYGZob{}}\PYG{p}{\PYGZcb{}}\PYG{p}{;}
\PYG{+w}{  }\PYG{k}{auto}\PYG{+w}{ }\PYG{n}{success}\PYG{+w}{ }\PYG{o}{=}\PYG{+w}{ }\PYG{n}{std}\PYG{o}{:}\PYG{o}{:}\PYG{n}{regex\PYGZus{}match}\PYG{p}{(}\PYG{n}{email}\PYG{p}{,}\PYG{+w}{ }\PYG{n}{result}\PYG{p}{,}\PYG{+w}{ }\PYG{n}{rx}\PYG{p}{)}\PYG{p}{;}
\PYG{+w}{  }\PYG{k}{return}\PYG{+w}{ }\PYG{n}{std}\PYG{o}{:}\PYG{o}{:}\PYG{n}{make\PYGZus{}tuple}\PYG{p}{(}
\PYG{+w}{    }\PYG{n}{success}\PYG{p}{,}
\PYG{+w}{    }\PYG{n}{success}\PYG{+w}{ }\PYG{o}{?}\PYG{+w}{ }\PYG{n}{result}\PYG{p}{[}\PYG{l+m+mi}{1}\PYG{p}{]}\PYG{p}{.}\PYG{n}{str}\PYG{p}{(}\PYG{p}{)}\PYG{+w}{ }\PYG{o}{:}\PYG{+w}{ }\PYG{l+s}{\PYGZdq{}}\PYG{l+s}{\PYGZdq{}}\PYG{n}{s}\PYG{p}{,}
\PYG{+w}{    }\PYG{n}{success}\PYG{+w}{ }\PYG{o}{?}\PYG{+w}{ }\PYG{n}{result}\PYG{p}{[}\PYG{l+m+mi}{2}\PYG{p}{]}\PYG{p}{.}\PYG{n}{str}\PYG{p}{(}\PYG{p}{)}\PYG{+w}{ }\PYG{o}{:}\PYG{+w}{ }\PYG{l+s}{\PYGZdq{}}\PYG{l+s}{\PYGZdq{}}\PYG{n}{s}\PYG{p}{,}
\PYG{+w}{    }\PYG{n}{success}\PYG{+w}{ }\PYG{o}{?}\PYG{+w}{ }\PYG{n}{result}\PYG{p}{[}\PYG{l+m+mi}{3}\PYG{p}{]}\PYG{p}{.}\PYG{n}{str}\PYG{p}{(}\PYG{p}{)}\PYG{+w}{ }\PYG{o}{:}\PYG{+w}{ }\PYG{l+s}{\PYGZdq{}}\PYG{l+s}{\PYGZdq{}}\PYG{n}{s}\PYG{p}{)}\PYG{p}{;}
\PYG{p}{\PYGZcb{}}
\end{Verbatim}
\end{codeblockbg}


%%C020383 p

沿用上面的代码，我们用C++17的结构化绑定把tuple的内容解包到具名变量中：


%%C020384 code

\begin{codeblockbg}
\begin{Verbatim}[commandchars=\\\{\},breaklines=true,breakanywhere=true,fontsize=\codefont,xleftmargin=17bp]
\PYG{k}{auto}\PYG{+w}{ }\PYG{n}{ltest3}\PYG{+w}{ }\PYG{o}{=}\PYG{+w}{ }\PYG{p}{[}\PYG{p}{]}\PYG{p}{(}\PYG{n}{std}\PYG{o}{:}\PYG{o}{:}\PYG{n}{string}\PYG{+w}{ }\PYG{k}{const}\PYG{+w}{ }\PYG{o}{\PYGZam{}}\PYG{+w}{ }\PYG{n}{email}\PYG{p}{)}
\PYG{p}{\PYGZob{}}
\PYG{+w}{  }\PYG{k}{auto}\PYG{+w}{ }\PYG{p}{[}\PYG{n}{valid}\PYG{p}{,}\PYG{+w}{ }\PYG{n}{localpart}\PYG{p}{,}\PYG{+w}{ }\PYG{n}{hostname}\PYG{p}{,}\PYG{+w}{ }\PYG{n}{dnslabel}\PYG{p}{]}\PYG{+w}{ }\PYG{o}{=}
\PYG{+w}{   }\PYG{n}{is\PYGZus{}valid\PYGZus{}email\PYGZus{}format\PYGZus{}with\PYGZus{}result}\PYG{p}{(}\PYG{n}{email}\PYG{p}{)}\PYG{p}{;}
\PYG{n}{std}\PYG{o}{:}\PYG{o}{:}\PYG{n}{cout}\PYG{+w}{ }\PYG{o}{\PYGZlt{}}\PYG{o}{\PYGZlt{}}\PYG{+w}{ }\PYG{n}{std}\PYG{o}{:}\PYG{o}{:}\PYG{n}{setw}\PYG{p}{(}\PYG{l+m+mi}{30}\PYG{p}{)}\PYG{+w}{ }\PYG{o}{\PYGZlt{}}\PYG{o}{\PYGZlt{}}\PYG{+w}{ }\PYG{n}{std}\PYG{o}{:}\PYG{o}{:}\PYG{n}{left}
\PYG{+w}{   }\PYG{o}{\PYGZlt{}}\PYG{o}{\PYGZlt{}}\PYG{+w}{ }\PYG{n}{email}\PYG{+w}{ }\PYG{o}{\PYGZlt{}}\PYG{o}{\PYGZlt{}}\PYG{+w}{ }\PYG{l+s}{\PYGZdq{}}\PYG{l+s}{ : }\PYG{l+s}{\PYGZdq{}}
\PYG{+w}{   }\PYG{o}{\PYGZlt{}}\PYG{o}{\PYGZlt{}}\PYG{+w}{ }\PYG{n}{std}\PYG{o}{:}\PYG{o}{:}\PYG{n}{setw}\PYG{p}{(}\PYG{l+m+mi}{10}\PYG{p}{)}\PYG{+w}{ }\PYG{o}{\PYGZlt{}}\PYG{o}{\PYGZlt{}}\PYG{+w}{ }\PYG{p}{(}\PYG{n}{valid}\PYG{+w}{ }\PYG{o}{?}\PYG{+w}{ }\PYG{l+s}{\PYGZdq{}}\PYG{l+s}{valid}\PYG{l+s}{\PYGZdq{}}\PYG{+w}{ }\PYG{o}{:}\PYG{+w}{ }\PYG{l+s}{\PYGZdq{}}\PYG{l+s}{invalid}\PYG{l+s}{\PYGZdq{}}\PYG{p}{)}
\PYG{+w}{   }\PYG{o}{\PYGZlt{}}\PYG{o}{\PYGZlt{}}\PYG{+w}{ }\PYG{l+s}{\PYGZdq{}}\PYG{l+s}{local=}\PYG{l+s}{\PYGZdq{}}\PYG{+w}{ }\PYG{o}{\PYGZlt{}}\PYG{o}{\PYGZlt{}}\PYG{+w}{ }\PYG{n}{localpart}
\PYG{+w}{   }\PYG{o}{\PYGZlt{}}\PYG{o}{\PYGZlt{}}\PYG{+w}{ }\PYG{l+s}{\PYGZdq{}}\PYG{l+s}{;domain=}\PYG{l+s}{\PYGZdq{}}\PYG{+w}{ }\PYG{o}{\PYGZlt{}}\PYG{o}{\PYGZlt{}}\PYG{+w}{ }\PYG{n}{hostname}
\PYG{+w}{   }\PYG{o}{\PYGZlt{}}\PYG{o}{\PYGZlt{}}\PYG{+w}{ }\PYG{l+s}{\PYGZdq{}}\PYG{l+s}{;dns=}\PYG{l+s}{\PYGZdq{}}\PYG{+w}{ }\PYG{o}{\PYGZlt{}}\PYG{o}{\PYGZlt{}}\PYG{+w}{ }\PYG{n}{dnslabel}
\PYG{+w}{   }\PYG{o}{\PYGZlt{}}\PYG{o}{\PYGZlt{}}\PYG{+w}{ }\PYG{n}{std}\PYG{o}{:}\PYG{o}{:}\PYG{n}{endl}\PYG{p}{;}
\PYG{p}{\PYGZcb{}}\PYG{p}{;}
\PYG{n}{ltest3}\PYG{p}{(}\PYG{l+s}{\PYGZdq{}}\PYG{l+s}{JOHN.DOE@DOMAIN.COM}\PYG{l+s}{\PYGZdq{}}\PYG{n}{s}\PYG{p}{)}\PYG{p}{;}
\PYG{n}{ltest3}\PYG{p}{(}\PYG{l+s}{\PYGZdq{}}\PYG{l+s}{JOHNDOE@DOMAIL.CO.UK}\PYG{l+s}{\PYGZdq{}}\PYG{n}{s}\PYG{p}{)}\PYG{p}{;}
\PYG{n}{ltest3}\PYG{p}{(}\PYG{l+s}{\PYGZdq{}}\PYG{l+s}{JOHNDOE@DOMAIL.INFO}\PYG{l+s}{\PYGZdq{}}\PYG{n}{s}\PYG{p}{)}\PYG{p}{;}
\PYG{n}{ltest3}\PYG{p}{(}\PYG{l+s}{\PYGZdq{}}\PYG{l+s}{J.O.H.N\PYGZus{}D.O.E@DOMAIN.INFO}\PYG{l+s}{\PYGZdq{}}\PYG{n}{s}\PYG{p}{)}\PYG{p}{;}
\PYG{n}{ltest3}\PYG{p}{(}\PYG{l+s}{\PYGZdq{}}\PYG{l+s}{ROOT@LOCALHOST}\PYG{l+s}{\PYGZdq{}}\PYG{n}{s}\PYG{p}{)}\PYG{p}{;}
\PYG{n}{ltest3}\PYG{p}{(}\PYG{l+s}{\PYGZdq{}}\PYG{l+s}{john.doe@domain.com}\PYG{l+s}{\PYGZdq{}}\PYG{n}{s}\PYG{p}{)}\PYG{p}{;}
\end{Verbatim}
\end{codeblockbg}


%%C020385 p

程序的输出如下：


%%C020386 code

\begin{codeblockbg}
\begin{Verbatim}[commandchars=\\\{\},breaklines=true,breakanywhere=true,fontsize=\codefont,xleftmargin=17bp]
\PYG{n}{JOHN}\PYG{p}{.}\PYG{n}{DOE}\PYG{err}{@}\PYG{n}{DOMAIN}\PYG{p}{.}\PYG{n}{COM}\PYG{+w}{            }\PYG{o}{:}\PYG{+w}{ }\PYG{n}{valid}
\PYG{+w}{   }\PYG{n}{local}\PYG{o}{=}\PYG{n}{JOHN}\PYG{p}{.}\PYG{n}{DOE}\PYG{p}{;}\PYG{n}{domain}\PYG{o}{=}\PYG{n}{DOMAIN}\PYG{p}{;}\PYG{n}{dns}\PYG{o}{=}\PYG{n}{COM}
\PYG{n}{JOHNDOE}\PYG{err}{@}\PYG{n}{DOMAIL}\PYG{p}{.}\PYG{n}{CO}\PYG{p}{.}\PYG{n}{UK}\PYG{+w}{           }\PYG{o}{:}\PYG{+w}{ }\PYG{n}{valid}
\PYG{+w}{   }\PYG{n}{local}\PYG{o}{=}\PYG{n}{JOHNDOE}\PYG{p}{;}\PYG{n}{domain}\PYG{o}{=}\PYG{n}{DOMAIL}\PYG{p}{.}\PYG{n}{CO}\PYG{p}{;}\PYG{n}{dns}\PYG{o}{=}\PYG{n}{UK}
\PYG{n}{JOHNDOE}\PYG{err}{@}\PYG{n}{DOMAIL}\PYG{p}{.}\PYG{n}{INFO}\PYG{+w}{            }\PYG{o}{:}\PYG{+w}{ }\PYG{n}{valid}
\PYG{+w}{   }\PYG{n}{local}\PYG{o}{=}\PYG{n}{JOHNDOE}\PYG{p}{;}\PYG{n}{domain}\PYG{o}{=}\PYG{n}{DOMAIL}\PYG{p}{;}\PYG{n}{dns}\PYG{o}{=}\PYG{n}{INFO}
\PYG{n}{J}\PYG{p}{.}\PYG{n}{O}\PYG{p}{.}\PYG{n}{H}\PYG{p}{.}\PYG{n}{N\PYGZus{}D}\PYG{p}{.}\PYG{n}{O}\PYG{p}{.}\PYG{n}{E}\PYG{err}{@}\PYG{n}{DOMAIN}\PYG{p}{.}\PYG{n}{INFO}\PYG{+w}{      }\PYG{o}{:}\PYG{+w}{ }\PYG{n}{valid}
\PYG{+w}{   }\PYG{n}{local}\PYG{o}{=}\PYG{n}{J}\PYG{p}{.}\PYG{n}{O}\PYG{p}{.}\PYG{n}{H}\PYG{p}{.}\PYG{n}{N\PYGZus{}D}\PYG{p}{.}\PYG{n}{O}\PYG{p}{.}\PYG{n}{E}\PYG{p}{;}\PYG{n}{domain}\PYG{o}{=}\PYG{n}{DOMAIN}\PYG{p}{;}\PYG{n}{dns}\PYG{o}{=}\PYG{n}{INFO}
\PYG{n}{ROOT}\PYG{err}{@}\PYG{n}{LOCALHOST}\PYG{+w}{                 }\PYG{o}{:}\PYG{+w}{ }\PYG{n}{invalid}
\PYG{+w}{   }\PYG{n}{local}\PYG{o}{=}\PYG{p}{;}\PYG{n}{domain}\PYG{o}{=}\PYG{p}{;}\PYG{n}{dns}\PYG{o}{=}
\PYG{n}{john}\PYG{p}{.}\PYG{n}{doe}\PYG{err}{@}\PYG{n}{domain}\PYG{p}{.}\PYG{n}{com}\PYG{+w}{            }\PYG{o}{:}\PYG{+w}{ }\PYG{n}{valid}
\PYG{+w}{   }\PYG{n}{local}\PYG{o}{=}\PYG{n}{john}\PYG{p}{.}\PYG{n}{doe}\PYG{p}{;}\PYG{n}{domain}\PYG{o}{=}\PYG{n}{domain}\PYG{p}{;}\PYG{n}{dns}\PYG{o}{=}\PYG{n}{com}
\end{Verbatim}
\end{codeblockbg}


%%C020387 h3

\subsection{更多内容}


%%C020388 p

正则表达式有多种版本，C++标准库支持其中六种：ECMAScript、basic POSIX、extended POSIX、awk、grep和egrep（带\texttt{-\allowbreak{}E}选项的grep）。默认使用的是ECMAScript语法；若要使用其他语法，必须在定义正则表达式时显式指定。除了指定语法，还可以指定解析选项，例如匹配时忽略大小写。


%%C020389 p

标准库提供的类和算法比我们目前看到的更多。库中主要的类如下（它们都是类模板，为方便使用，标准库针对不同字符类型提供了\texttt{typedef}）：


%%C020390 li

\begin{itemize}
\item 类模板\texttt{std::\allowbreak{}basic\_\allowbreak{}regex}定义正则表达式对象：


%%C020391 code

\begin{codeblockbg}
\begin{Verbatim}[commandchars=\\\{\},breaklines=true,breakanywhere=true,fontsize=\codefont,xleftmargin=17bp]
\PYG{k}{typedef}\PYG{+w}{ }\PYG{n}{basic\PYGZus{}regex}\PYG{o}{\PYGZlt{}}\PYG{k+kt}{char}\PYG{o}{\PYGZgt{}}\PYG{+w}{    }\PYG{n}{regex}\PYG{p}{;}
\PYG{k}{typedef}\PYG{+w}{ }\PYG{n}{basic\PYGZus{}regex}\PYG{o}{\PYGZlt{}}\PYG{k+kt}{wchar\PYGZus{}t}\PYG{o}{\PYGZgt{}}\PYG{+w}{ }\PYG{n}{wregex}\PYG{p}{;}
\end{Verbatim}
\end{codeblockbg}


%%C020392 li

\item 类模板\texttt{std::\allowbreak{}sub\_\allowbreak{}match}表示与某个捕获组匹配的字符序列；该类实际上派生自\texttt{std::\allowbreak{}pair}，其\texttt{first}和\texttt{second}成员分别是指向匹配序列中第一个字符和末尾之后位置的迭代器；若没有匹配序列，两个迭代器相等：


%%C020393 code

\begin{codeblockbg}
\begin{Verbatim}[commandchars=\\\{\},breaklines=true,breakanywhere=true,fontsize=\codefont,xleftmargin=17bp]
\PYG{k}{typedef}\PYG{+w}{ }\PYG{n}{sub\PYGZus{}match}\PYG{o}{\PYGZlt{}}\PYG{k}{const}\PYG{+w}{ }\PYG{k+kt}{char}\PYG{+w}{ }\PYG{o}{*}\PYG{o}{\PYGZgt{}}\PYG{+w}{            }\PYG{n}{csub\PYGZus{}match}\PYG{p}{;}
\PYG{k}{typedef}\PYG{+w}{ }\PYG{n}{sub\PYGZus{}match}\PYG{o}{\PYGZlt{}}\PYG{k}{const}\PYG{+w}{ }\PYG{k+kt}{wchar\PYGZus{}t}\PYG{+w}{ }\PYG{o}{*}\PYG{o}{\PYGZgt{}}\PYG{+w}{         }\PYG{n}{wcsub\PYGZus{}match}\PYG{p}{;}
\PYG{k}{typedef}\PYG{+w}{ }\PYG{n}{sub\PYGZus{}match}\PYG{o}{\PYGZlt{}}\PYG{n}{string}\PYG{o}{:}\PYG{o}{:}\PYG{n}{const\PYGZus{}iterator}\PYG{o}{\PYGZgt{}}\PYG{+w}{  }\PYG{n}{ssub\PYGZus{}match}\PYG{p}{;}
\PYG{k}{typedef}\PYG{+w}{ }\PYG{n}{sub\PYGZus{}match}\PYG{o}{\PYGZlt{}}\PYG{n}{wstring}\PYG{o}{:}\PYG{o}{:}\PYG{n}{const\PYGZus{}iterator}\PYG{o}{\PYGZgt{}}\PYG{+w}{ }\PYG{n}{wssub\PYGZus{}match}\PYG{p}{;}
\end{Verbatim}
\end{codeblockbg}


%%C020394 li

\item 类模板\myttbr{std::\allowbreak{}match\_\allowbreak{}results}是匹配结果的集合；第一个元素始终是目标中的完整匹配，其余元素是各子表达式的匹配：


%%C020395 code

\begin{codeblockbg}
\begin{Verbatim}[commandchars=\\\{\},breaklines=true,breakanywhere=true,fontsize=\codefont,xleftmargin=17bp]
\PYG{k}{typedef}\PYG{+w}{ }\PYG{n}{match\PYGZus{}results}\PYG{o}{\PYGZlt{}}\PYG{k}{const}\PYG{+w}{ }\PYG{k+kt}{char}\PYG{+w}{ }\PYG{o}{*}\PYG{o}{\PYGZgt{}}\PYG{+w}{            }\PYG{n}{cmatch}\PYG{p}{;}
\PYG{k}{typedef}\PYG{+w}{ }\PYG{n}{match\PYGZus{}results}\PYG{o}{\PYGZlt{}}\PYG{k}{const}\PYG{+w}{ }\PYG{k+kt}{wchar\PYGZus{}t}\PYG{+w}{ }\PYG{o}{*}\PYG{o}{\PYGZgt{}}\PYG{+w}{         }\PYG{n}{wcmatch}\PYG{p}{;}
\PYG{k}{typedef}\PYG{+w}{ }\PYG{n}{match\PYGZus{}results}\PYG{o}{\PYGZlt{}}\PYG{n}{string}\PYG{o}{:}\PYG{o}{:}\PYG{n}{const\PYGZus{}iterator}\PYG{o}{\PYGZgt{}}\PYG{+w}{  }\PYG{n}{smatch}\PYG{p}{;}
\PYG{k}{typedef}\PYG{+w}{ }\PYG{n}{match\PYGZus{}results}\PYG{o}{\PYGZlt{}}\PYG{n}{wstring}\PYG{o}{:}\PYG{o}{:}\PYG{n}{const\PYGZus{}iterator}\PYG{o}{\PYGZgt{}}\PYG{+w}{ }\PYG{n}{wsmatch}\PYG{p}{;}
\end{Verbatim}
\end{codeblockbg}

\end{itemize}


%%C020396 p

正则表达式标准库中可用的算法如下：


%%C020397 li

\begin{itemize}
\item \myttbr{std::\allowbreak{}regex\_\allowbreak{}match()}：尝试把正则表达式（由\texttt{std::\allowbreak{}basic\_\allowbreak{}regex}实例表示）与整个字符串匹配。


%%C020398 li

\item \myttbr{std::\allowbreak{}regex\_\allowbreak{}search()}：尝试把正则表达式（由\texttt{std::\allowbreak{}basic\_\allowbreak{}regex}实例表示）与字符串的一部分（也可以是整个字符串）匹配。


%%C020399 li

\item \myttbr{std::\allowbreak{}regex\_\allowbreak{}replace()}：按指定格式替换正则表达式的匹配结果。

\end{itemize}


%%C020400 p

正则表达式标准库中可用的迭代器如下：


%%C020401 li

\begin{itemize}
\item \myttbr{std::\allowbreak{}regex\_\allowbreak{}interator}：常量前向迭代器，用于遍历字符串中某个模式的所有出现位置。它持有指向\texttt{std::\allowbreak{}basic\_\allowbreak{}regex}的指针，该对象必须存活到迭代器销毁为止。在创建时以及每次递增时，迭代器会调用\myttbr{std::\allowbreak{}regex\_\allowbreak{}search()}，并保存一份该算法返回的\myttbr{std::\allowbreak{}match\_\allowbreak{}results}对象的副本。


%%C020402 li

\item \myttbr{std::\allowbreak{}regex\_\allowbreak{}token\_\allowbreak{}iterator}：常量前向迭代器，用于遍历字符串中正则表达式每次匹配的各个子匹配。它内部使用\myttbr{std::\allowbreak{}regex\_\allowbreak{}iterator}逐个走过子匹配。由于它保存指向\texttt{std::\allowbreak{}basic\_\allowbreak{}regex}实例的指针，该正则表达式对象必须存活到迭代器销毁为止。

\end{itemize}


%%C020403 h3

\subsection{延伸阅读}


%%C020404 li

\begin{itemize}
\item \textit{用正则表达式解析字符串内容}


%%C020405 li

\item \textit{用正则表达式替换字符串内容}


%%C020406 li

\item \textit{用结构化绑定处理多返回值}案例，见\texttt{Chapter 1}\textit{现代C++语言核心特性}

\end{itemize}


%%C020407 h2

\section{用正则表达式解析字符串内容}


%%C020408 p

上一个案例中，我们了解了如何用\myttbr{std::\allowbreak{}regex\_\allowbreak{}match()}验证字符串内容是否匹配某种格式。标准库还提供了另一个算法\myttbr{std::\allowbreak{}regex\_\allowbreak{}search()}，它把正则表达式与字符串的任意部分匹配，而不像\texttt{regex\_\allowbreak{}match()}那样只匹配整个字符串。不过，这个函数无法遍历输入字符串中某个正则表达式的所有出现位置。为此，我们需要使用库中提供的某个迭代器类。

%%C020409 p

本案例将介绍如何用正则表达式解析字符串的内容。为此，我们考虑解析一个包含名值对的文本文件这一问题。每个这样的名值对单独占一行，格式为\texttt{name = value}，但以\texttt{\#}开头的行表示注释，必须忽略。下面是一个例子：


%%C020410 code

\begin{codeblockbg}
\begin{Verbatim}[commandchars=\\\{\},breaklines=true,breakanywhere=true,fontsize=\codefont,xleftmargin=17bp]
\PYG{c+cp}{\PYGZsh{}}\PYG{c+cp}{remove \PYGZsh{} to uncomment the following lines}
\PYG{n}{timeout}\PYG{o}{=}\PYG{l+m+mi}{120}
\PYG{n}{server}\PYG{+w}{ }\PYG{o}{=}\PYG{+w}{ }\PYG{l+m+mf}{127.0}\PYG{l+m+mf}{.0}\PYG{l+m+mf}{.1}

\PYG{c+cp}{\PYGZsh{}}\PYG{c+cp}{retrycount=3}
\end{Verbatim}
\end{codeblockbg}


%%C020411 h3

\subsection{准备工作}


%%C020412 p

关于C++11中正则表达式支持的一般信息，请参考\textit{用正则表达式}\textit{验证字符串格式}案例。继续本案例需要具备正则表达式的基础知识。


%%C020413 p

在下面的示例中，\texttt{text}是按下述方式定义的变量：


%%C020414 code

\begin{codeblockbg}
\begin{Verbatim}[commandchars=\\\{\},breaklines=true,breakanywhere=true,fontsize=\codefont,xleftmargin=17bp]
\PYG{k}{auto}\PYG{+w}{ }\PYG{n}{text}\PYG{+w}{ }\PYG{p}{\PYGZob{}}
\PYG{+w}{  }\PYG{l+s+sa}{R}\PYG{l+s}{\PYGZdq{}}\PYG{l+s+dl}{(}
\PYG{l+s}{    \PYGZsh{}remove \PYGZsh{} to uncomment the following lines}
\PYG{l+s}{    timeout=120}
\PYG{l+s}{    server = 127.0.0.1}
\PYG{l+s}{  \PYGZsh{}retrycount=3}
\PYG{l+s+dl}{)}\PYG{l+s}{\PYGZdq{}}\PYG{n}{s}\PYG{p}{\PYGZcb{}}\PYG{p}{;}
\end{Verbatim}
\end{codeblockbg}


%%C020415 h3

\subsection{具体做法}


%%C020416 p

要在字符串中搜索正则表达式的出现位置，应执行以下步骤：


%%C020417 p

包含头文件\texttt{<\allowbreak{}regex>}和\texttt{<\allowbreak{}string>}，以及命名空间


%%C020418 p

\myttbr{std::\allowbreak{}string\_\allowbreak{}literals}，它提供C++14标准中用于字符串的用户定义字面量：


%%C020419 code

\begin{codeblockbg}
\begin{Verbatim}[commandchars=\\\{\},breaklines=true,breakanywhere=true,fontsize=\codefont,xleftmargin=17bp]
\PYG{c+cp}{\PYGZsh{}}\PYG{c+cp}{include}\PYG{+w}{ }\PYG{c+cpf}{\PYGZlt{}regex\PYGZgt{}}
\PYG{c+cp}{\PYGZsh{}}\PYG{c+cp}{include}\PYG{+w}{ }\PYG{c+cpf}{\PYGZlt{}string\PYGZgt{}}
\PYG{k}{using}\PYG{+w}{ }\PYG{k}{namespace}\PYG{+w}{ }\PYG{n+nn}{std}\PYG{o}{:}\PYG{o}{:}\PYG{n+nn}{string\PYGZus{}literals}\PYG{p}{;}
\end{Verbatim}
\end{codeblockbg}


%%C020420 p

用raw字符串字面量来书写正则表达式，以避免转义


%%C020421 p

反斜杠（反斜杠出现得相当频繁）。下面的正则表达式校验前面提出的文件格式：


%%C020422 code

\begin{codeblockbg}
\begin{Verbatim}[commandchars=\\\{\},breaklines=true,breakanywhere=true,fontsize=\codefont,xleftmargin=17bp]
\PYG{k}{auto}\PYG{+w}{ }\PYG{n}{pattern}\PYG{+w}{ }\PYG{p}{\PYGZob{}}\PYG{l+s+sa}{R}\PYG{l+s}{\PYGZdq{}}\PYG{l+s+dl}{(}\PYG{l+s}{\PYGZca{}(?!\PYGZsh{})(\PYGZbs{}w+)\PYGZbs{}s*=\PYGZbs{}s*([\PYGZbs{}w\PYGZbs{}d]+[\PYGZbs{}w\PYGZbs{}d.\PYGZus{},\PYGZbs{}\PYGZhy{}:]*)\PYGZdl{}}\PYG{l+s+dl}{)}\PYG{l+s}{\PYGZdq{}}\PYG{n}{s}\PYG{p}{\PYGZcb{}}\PYG{p}{;}
\end{Verbatim}
\end{codeblockbg}


%%C020423 p

创建\texttt{std::\allowbreak{}regex}/\texttt{std::\allowbreak{}wregex}对象（取决于所用的字符集），


%%C020424 p

用它封装该正则表达式：


%%C020425 code

\begin{codeblockbg}
\begin{Verbatim}[commandchars=\\\{\},breaklines=true,breakanywhere=true,fontsize=\codefont,xleftmargin=17bp]
\PYG{k}{auto}\PYG{+w}{ }\PYG{n}{rx}\PYG{+w}{ }\PYG{o}{=}\PYG{+w}{ }\PYG{n}{std}\PYG{o}{:}\PYG{o}{:}\PYG{n}{regex}\PYG{p}{\PYGZob{}}\PYG{n}{pattern}\PYG{p}{\PYGZcb{}}\PYG{p}{;}
\end{Verbatim}
\end{codeblockbg}


%%C020426 p

要在给定文本中搜索正则表达式的第一个匹配，可使用


%%C020427 p

通用算法\myttbr{std::\allowbreak{}regex\_\allowbreak{}search()}（示例1）：


%%C020428 code

\begin{codeblockbg}
\begin{Verbatim}[commandchars=\\\{\},breaklines=true,breakanywhere=true,fontsize=\codefont,xleftmargin=17bp]
\PYG{k}{auto}\PYG{+w}{ }\PYG{n}{match}\PYG{+w}{ }\PYG{o}{=}\PYG{+w}{ }\PYG{n}{std}\PYG{o}{:}\PYG{o}{:}\PYG{n}{smatch}\PYG{p}{\PYGZob{}}\PYG{p}{\PYGZcb{}}\PYG{p}{;}
\PYG{k}{if}\PYG{+w}{ }\PYG{p}{(}\PYG{n}{std}\PYG{o}{:}\PYG{o}{:}\PYG{n}{regex\PYGZus{}search}\PYG{p}{(}\PYG{n}{text}\PYG{p}{,}\PYG{+w}{ }\PYG{n}{match}\PYG{p}{,}\PYG{+w}{ }\PYG{n}{rx}\PYG{p}{)}\PYG{p}{)}
\PYG{p}{\PYGZob{}}
\PYG{+w}{  }\PYG{n}{std}\PYG{o}{:}\PYG{o}{:}\PYG{n}{cout}\PYG{+w}{ }\PYG{o}{\PYGZlt{}}\PYG{o}{\PYGZlt{}}\PYG{+w}{ }\PYG{n}{match}\PYG{p}{[}\PYG{l+m+mi}{1}\PYG{p}{]}\PYG{+w}{ }\PYG{o}{\PYGZlt{}}\PYG{o}{\PYGZlt{}}\PYG{+w}{ }\PYG{l+s+sc}{\PYGZsq{}}\PYG{l+s+sc}{=}\PYG{l+s+sc}{\PYGZsq{}}\PYG{+w}{ }\PYG{o}{\PYGZlt{}}\PYG{o}{\PYGZlt{}}\PYG{+w}{ }\PYG{n}{match}\PYG{p}{[}\PYG{l+m+mi}{2}\PYG{p}{]}\PYG{+w}{ }\PYG{o}{\PYGZlt{}}\PYG{o}{\PYGZlt{}}\PYG{+w}{ }\PYG{n}{std}\PYG{o}{:}\PYG{o}{:}\PYG{n}{endl}\PYG{p}{;}
\PYG{p}{\PYGZcb{}}
\end{Verbatim}
\end{codeblockbg}


%%C020429 p

要找出给定文本中正则表达式的所有匹配，可使用迭代器


%%C020430 p

\myttbr{std::\allowbreak{}regex\_\allowbreak{}iterator}（示例2）：


%%C020431 code

\begin{codeblockbg}
\begin{Verbatim}[commandchars=\\\{\},breaklines=true,breakanywhere=true,fontsize=\codefont,xleftmargin=17bp]
\PYG{k}{auto}\PYG{+w}{ }\PYG{n}{end}\PYG{+w}{ }\PYG{o}{=}\PYG{+w}{ }\PYG{n}{std}\PYG{o}{:}\PYG{o}{:}\PYG{n}{sregex\PYGZus{}iterator}\PYG{p}{\PYGZob{}}\PYG{p}{\PYGZcb{}}\PYG{p}{;}
\PYG{k}{for}\PYG{+w}{ }\PYG{p}{(}\PYG{k}{auto}\PYG{+w}{ }\PYG{n}{it}\PYG{o}{=}\PYG{n}{std}\PYG{o}{:}\PYG{o}{:}\PYG{n}{sregex\PYGZus{}iterator}\PYG{p}{\PYGZob{}}\PYG{+w}{ }\PYG{n}{std}\PYG{o}{:}\PYG{o}{:}\PYG{n}{begin}\PYG{p}{(}\PYG{n}{text}\PYG{p}{)}\PYG{p}{,}
\PYG{+w}{                                   }\PYG{n}{std}\PYG{o}{:}\PYG{o}{:}\PYG{n}{end}\PYG{p}{(}\PYG{n}{text}\PYG{p}{)}\PYG{p}{,}\PYG{+w}{ }\PYG{n}{rx}\PYG{+w}{ }\PYG{p}{\PYGZcb{}}\PYG{p}{;}
\PYG{+w}{     }\PYG{n}{it}\PYG{+w}{ }\PYG{o}{!}\PYG{o}{=}\PYG{+w}{ }\PYG{n}{end}\PYG{p}{;}\PYG{+w}{ }\PYG{o}{+}\PYG{o}{+}\PYG{n}{it}\PYG{p}{)}
\PYG{p}{\PYGZob{}}
\PYG{+w}{  }\PYG{n}{std}\PYG{o}{:}\PYG{o}{:}\PYG{n}{cout}\PYG{+w}{ }\PYG{o}{\PYGZlt{}}\PYG{o}{\PYGZlt{}}\PYG{+w}{ }\PYG{err}{\PYGZsq{}}\PYG{err}{\PYGZsq{}}\PYG{err}{\PYGZsq{}}\PYG{+w}{ }\PYG{o}{\PYGZlt{}}\PYG{o}{\PYGZlt{}}\PYG{+w}{ }\PYG{p}{(}\PYG{o}{*}\PYG{n}{it}\PYG{p}{)}\PYG{p}{[}\PYG{l+m+mi}{1}\PYG{p}{]}\PYG{+w}{ }\PYG{o}{\PYGZlt{}}\PYG{o}{\PYGZlt{}}\PYG{+w}{ }\PYG{l+s}{\PYGZdq{}}\PYG{l+s}{\PYGZsq{}=\PYGZsq{}}\PYG{l+s}{\PYGZdq{}}
\PYG{+w}{            }\PYG{o}{\PYGZlt{}}\PYG{o}{\PYGZlt{}}\PYG{+w}{ }\PYG{p}{(}\PYG{o}{*}\PYG{n}{it}\PYG{p}{)}\PYG{p}{[}\PYG{l+m+mi}{2}\PYG{p}{]}\PYG{+w}{ }\PYG{o}{\PYGZlt{}}\PYG{o}{\PYGZlt{}}\PYG{+w}{ }\PYG{err}{\PYGZsq{}}\PYG{err}{\PYGZsq{}}\PYG{err}{\PYGZsq{}}\PYG{+w}{ }\PYG{o}{\PYGZlt{}}\PYG{o}{\PYGZlt{}}\PYG{+w}{ }\PYG{n}{std}\PYG{o}{:}\PYG{o}{:}\PYG{n}{endl}\PYG{p}{;}
\PYG{p}{\PYGZcb{}}
\end{Verbatim}
\end{codeblockbg}


%%C020432 p

要遍历某个匹配的所有子表达式，可使用迭代器


%%C020433 p

\myttbr{std::\allowbreak{}regex\_\allowbreak{}token\_\allowbreak{}iterator}（示例3）：


%%C020434 code

\begin{codeblockbg}
\begin{Verbatim}[commandchars=\\\{\},breaklines=true,breakanywhere=true,fontsize=\codefont,xleftmargin=17bp]
\PYG{k}{auto}\PYG{+w}{ }\PYG{n}{end}\PYG{+w}{ }\PYG{o}{=}\PYG{+w}{ }\PYG{n}{std}\PYG{o}{:}\PYG{o}{:}\PYG{n}{sregex\PYGZus{}token\PYGZus{}iterator}\PYG{p}{\PYGZob{}}\PYG{p}{\PYGZcb{}}\PYG{p}{;}
\PYG{k}{for}\PYG{+w}{ }\PYG{p}{(}\PYG{k}{auto}\PYG{+w}{ }\PYG{n}{it}\PYG{+w}{ }\PYG{o}{=}\PYG{+w}{ }\PYG{n}{std}\PYG{o}{:}\PYG{o}{:}\PYG{n}{sregex\PYGZus{}token\PYGZus{}iterator}\PYG{p}{\PYGZob{}}
\PYG{+w}{                  }\PYG{n}{std}\PYG{o}{:}\PYG{o}{:}\PYG{n}{begin}\PYG{p}{(}\PYG{n}{text}\PYG{p}{)}\PYG{p}{,}\PYG{+w}{  }\PYG{n}{std}\PYG{o}{:}\PYG{o}{:}\PYG{n}{end}\PYG{p}{(}\PYG{n}{text}\PYG{p}{)}\PYG{p}{,}\PYG{+w}{ }\PYG{n}{rx}\PYG{+w}{ }\PYG{p}{\PYGZcb{}}\PYG{p}{;}
\PYG{+w}{     }\PYG{n}{it}\PYG{+w}{ }\PYG{o}{!}\PYG{o}{=}\PYG{+w}{ }\PYG{n}{end}\PYG{p}{;}\PYG{+w}{ }\PYG{o}{+}\PYG{o}{+}\PYG{n}{it}\PYG{p}{)}
\PYG{p}{\PYGZob{}}
\PYG{+w}{  }\PYG{n}{std}\PYG{o}{:}\PYG{o}{:}\PYG{n}{cout}\PYG{+w}{ }\PYG{o}{\PYGZlt{}}\PYG{o}{\PYGZlt{}}\PYG{+w}{ }\PYG{o}{*}\PYG{n}{it}\PYG{+w}{ }\PYG{o}{\PYGZlt{}}\PYG{o}{\PYGZlt{}}\PYG{+w}{ }\PYG{n}{std}\PYG{o}{:}\PYG{o}{:}\PYG{n}{endl}\PYG{p}{;}
\PYG{p}{\PYGZcb{}}
\end{Verbatim}
\end{codeblockbg}


%%C020435 h3

\subsection{工作原理}


%%C020436 p

一个能解析前面所示输入文件的简单正则表达式大致如下：


%%C020437 code

\begin{codeblockbg}
\begin{Verbatim}[commandchars=\\\{\},breaklines=true,breakanywhere=true,fontsize=\codefont,xleftmargin=17bp]
\PYG{o}{\PYGZca{}}\PYG{p}{(}\PYG{o}{?}\PYG{o}{!}\PYG{err}{\PYGZsh{}}\PYG{p}{)}\PYG{p}{(}\PYG{err}{\PYGZbs{}}\PYG{n}{w}\PYG{o}{+}\PYG{p}{)}\PYG{err}{\PYGZbs{}}\PYG{n}{s}\PYG{o}{*}\PYG{o}{=}\PYG{err}{\PYGZbs{}}\PYG{n}{s}\PYG{o}{*}\PYG{p}{(}\PYG{p}{[}\PYG{err}{\PYGZbs{}}\PYG{n}{w}\PYG{err}{\PYGZbs{}}\PYG{n}{d}\PYG{p}{]}\PYG{o}{+}\PYG{p}{[}\PYG{err}{\PYGZbs{}}\PYG{n}{w}\PYG{err}{\PYGZbs{}}\PYG{n}{d}\PYG{p}{.}\PYG{n}{\PYGZus{}}\PYG{p}{,}\PYG{err}{\PYGZbs{}}\PYG{o}{\PYGZhy{}}\PYG{o}{:}\PYG{p}{]}\PYG{o}{*}\PYG{p}{)}\PYG{n}{\PYGZdl{}}
\end{Verbatim}
\end{codeblockbg}


%%C020438 p

该正则表达式应当忽略所有以\texttt{\#}开头的行；对于不以\texttt{\#}开头的行，则匹配一个名称，随后是等号，再后是一个值，该值可由字母数字字符以及其他若干字符（下划线、点、逗号等）组成。这个正则表达式的确切含义解释如下：


%%C020439 table

\myTableBegin
\begin{xltabular}{\textwidth}{|>{\hsize=0.324\hsize\raggedright\arraybackslash}X|>{\hsize=0.676\hsize\raggedright\arraybackslash}X|}
\hline
\rowcolor{tabhead} \textcolor{white}{\textbf{部分}} & \textcolor{white}{\textbf{说明}} \\
\hline
\endfirsthead
\rowcolor{tabhead} \textcolor{white}{\textbf{部分}} & \textcolor{white}{\textbf{说明}} \\
\hline
\endhead
\myttbr{\textasciicircum{}} & 行首\\
\hline
\myttbr{(\allowbreak{}?!\#)} & 负向先行断言（negative lookahead），确保无法匹配\#字符\\
\hline
\myttbr{(\textbackslash{}w)+} & 捕获组，表示至少包含一个单词字符的标识符\\
\hline
\myttbr{\textbackslash{}s*} & 任意空白字符\\
\hline
\myttbr{=} & 等号\\
\hline
\myttbr{\textbackslash{}s*} & 任意空白字符\\
\hline
\myttbr{([\textbackslash{}w\textbackslash{}d]+[\textbackslash{}w\textbackslash{}d.\_,\textbackslash{}-:]*)} & 捕获组，表示一个以字母数字字符开头的值，但也可以包含点、逗号、反斜杠、连字符、冒号或下划线。\\
\hline
\myttbr{\$} & 行尾\\
\hline
\end{xltabular}
\myTableEnd


%%C020440 table

% 原书此处是上一张表的续表（PDF分页把表切断，抽取时首行被误当表头），
% 已并入上表，本单元留空。


%%C020441 p

可以用\myttbr{std::\allowbreak{}regex\_\allowbreak{}search()}在输入文本的任意位置搜索匹配。该算法有多个重载版本，但它们的用法基本相同。你必须指定要处理的字符范围、一个用于存放匹配结果的输出对象\myttbr{std::\allowbreak{}match\_\allowbreak{}results}，以及一个表示正则表达式的\texttt{std::\allowbreak{}basic\_\allowbreak{}regex}对象和匹配标志（匹配标志决定搜索的进行方式）。函数在找到匹配时返回\texttt{true}，否则返回\texttt{false}。


%%C020442 p

在上一节的第一个示例中（见第4个列表项），\texttt{match}是\texttt{std::\allowbreak{}smatch}的实例，而\texttt{std::\allowbreak{}smatch}是\myttbr{std::\allowbreak{}match\_\allowbreak{}results}以\myttbr{string::\allowbreak{}const\_\allowbreak{}iterator}为模板类型的一个\texttt{typedef}。如果找到匹配，该对象会以值序列的形式包含所有已匹配子表达式的匹配信息。索引0处的子匹配始终是整个匹配。索引1处的子匹配是第一个被匹配的子表达式，索引2处的子匹配是第二个被匹配的子表达式，依此类推。由于我们的正则表达式中有两个捕获组（即子表达式），匹配成功时\myttbr{std::\allowbreak{}match\_\allowbreak{}results}会有三个子匹配。表示名称的标识符位于索引1，等号后面的值位于索引2。因此，这段代码只打印如下内容：


%%C020443 code

\begin{codeblockbg}
\begin{Verbatim}[commandchars=\\\{\},breaklines=true,breakanywhere=true,fontsize=\codefont,xleftmargin=17bp]
\PYG{n}{timeout}\PYG{o}{=}\PYG{l+m+mi}{120}
\end{Verbatim}
\end{codeblockbg}


%%C020444 p

\myttbr{std::\allowbreak{}regex\_\allowbreak{}search()}算法无法遍历文本中所有可能的匹配。要做到这一点，需要使用迭代器。\myttbr{std::\allowbreak{}regex\_\allowbreak{}iterator}正是为此设计的。它不仅能遍历所有匹配，还能访问某个匹配的所有子匹配。该迭代器在构造时以及每次自增时都会调用\myttbr{std::\allowbreak{}regex\_\allowbreak{}search()}，并记住这次调用得到的\myttbr{std::\allowbreak{}match\_\allowbreak{}results}结果。默认构造函数创建一个表示序列末尾的迭代器，可用它判断遍历匹配的循环应在何时结束。


%%C020445 p

在上一节的第二个示例中（见第5个列表项），我们先创建一个序列末尾迭代器，然后开始遍历所有可能的匹配。构造迭代器时会调用\myttbr{std::\allowbreak{}regex\_\allowbreak{}match()}，如果找到匹配，就可以通过当前迭代器访问其结果。这一过程会一直持续到再也找不到匹配（到达序列末尾）为止。这段代码会打印如下输出：


%%C020446 code

\begin{codeblockbg}
\begin{Verbatim}[commandchars=\\\{\},breaklines=true,breakanywhere=true,fontsize=\codefont,xleftmargin=17bp]
\PYG{err}{\PYGZsq{}}\PYG{n}{timeout}\PYG{l+s+sc}{\PYGZsq{}}\PYG{l+s+sc}{=}\PYG{l+s+sc}{\PYGZsq{}}\PYG{l+m+mi}{120}\PYG{err}{\PYGZsq{}}
\PYG{err}{\PYGZsq{}}\PYG{n}{server}\PYG{l+s+sc}{\PYGZsq{}}\PYG{l+s+sc}{=}\PYG{l+s+sc}{\PYGZsq{}}\PYG{l+m+mf}{127.0}\PYG{l+m+mf}{.0}\PYG{l+m+mf}{.1}\PYG{err}{\PYGZsq{}}
\end{Verbatim}
\end{codeblockbg}


%%C020447 p

\myttbr{std::\allowbreak{}regex\_\allowbreak{}iterator}的替代品是\myttbr{std::\allowbreak{}regex\_\allowbreak{}token\_\allowbreak{}iterator}。它的工作方式与\myttbr{std::\allowbreak{}regex\_\allowbreak{}iterator}类似，实际上它内部就包含这样一个迭代器，只是它让我们可以访问匹配中的某个特定子表达式。\textit{具体做法}一节的第三个示例（第6个列表项）展示了这一点。我们先创建一个序列末尾迭代器，然后遍历各个匹配，直到到达序列末尾。在我们使用的构造函数中，没有指定要通过迭代器访问的子表达式索引，因此使用默认值0。这意味着该程序会打印整个匹配：


%%C020448 code

\begin{codeblockbg}
\begin{Verbatim}[commandchars=\\\{\},breaklines=true,breakanywhere=true,fontsize=\codefont,xleftmargin=17bp]
\PYG{n}{timeout}\PYG{o}{=}\PYG{l+m+mi}{120}
\PYG{n}{server}\PYG{+w}{ }\PYG{o}{=}\PYG{+w}{ }\PYG{l+m+mf}{127.0}\PYG{l+m+mf}{.0}\PYG{l+m+mf}{.1}
\end{Verbatim}
\end{codeblockbg}


%%C020449 p

如果只想访问第一个子表达式（在本例中就是名称），只需在token迭代器的构造函数中指定子表达式的索引。这一次，得到的输出只有名称：


%%C020450 code

\begin{codeblockbg}
\begin{Verbatim}[commandchars=\\\{\},breaklines=true,breakanywhere=true,fontsize=\codefont,xleftmargin=17bp]
\PYG{k}{auto}\PYG{+w}{ }\PYG{n}{end}\PYG{+w}{ }\PYG{o}{=}\PYG{+w}{ }\PYG{n}{std}\PYG{o}{:}\PYG{o}{:}\PYG{n}{sregex\PYGZus{}token\PYGZus{}iterator}\PYG{p}{\PYGZob{}}\PYG{p}{\PYGZcb{}}\PYG{p}{;}
\PYG{k}{for}\PYG{+w}{ }\PYG{p}{(}\PYG{k}{auto}\PYG{+w}{ }\PYG{n}{it}\PYG{+w}{ }\PYG{o}{=}\PYG{+w}{ }\PYG{n}{std}\PYG{o}{:}\PYG{o}{:}\PYG{n}{sregex\PYGZus{}token\PYGZus{}iterator}\PYG{p}{\PYGZob{}}\PYG{+w}{ }\PYG{n}{std}\PYG{o}{:}\PYG{o}{:}\PYG{n}{begin}\PYG{p}{(}\PYG{n}{text}\PYG{p}{)}\PYG{p}{,}
\PYG{+w}{               }\PYG{n}{std}\PYG{o}{:}\PYG{o}{:}\PYG{n}{end}\PYG{p}{(}\PYG{n}{text}\PYG{p}{)}\PYG{p}{,}\PYG{+w}{ }\PYG{n}{rx}\PYG{p}{,}\PYG{+w}{ }\PYG{l+m+mi}{1}\PYG{+w}{ }\PYG{p}{\PYGZcb{}}\PYG{p}{;}
\PYG{+w}{     }\PYG{n}{it}\PYG{+w}{ }\PYG{o}{!}\PYG{o}{=}\PYG{+w}{ }\PYG{n}{end}\PYG{p}{;}\PYG{+w}{ }\PYG{o}{+}\PYG{o}{+}\PYG{n}{it}\PYG{p}{)}
\PYG{p}{\PYGZob{}}
\PYG{+w}{  }\PYG{n}{std}\PYG{o}{:}\PYG{o}{:}\PYG{n}{cout}\PYG{+w}{ }\PYG{o}{\PYGZlt{}}\PYG{o}{\PYGZlt{}}\PYG{+w}{ }\PYG{o}{*}\PYG{n}{it}\PYG{+w}{ }\PYG{o}{\PYGZlt{}}\PYG{o}{\PYGZlt{}}\PYG{+w}{ }\PYG{n}{std}\PYG{o}{:}\PYG{o}{:}\PYG{n}{endl}\PYG{p}{;}
\PYG{p}{\PYGZcb{}}
\end{Verbatim}
\end{codeblockbg}


%%C020451 p

token迭代器的一个有趣之处是：如果子表达式索引为\texttt{-\allowbreak{}1}，它可以返回字符串中未被匹配的部分，此时返回的\myttbr{std::\allowbreak{}match\_\allowbreak{}results}对象对应上一个匹配之后直到序列末尾之间的字符序列：


%%C020452 code

\begin{codeblockbg}
\begin{Verbatim}[commandchars=\\\{\},breaklines=true,breakanywhere=true,fontsize=\codefont,xleftmargin=17bp]
\PYG{k}{auto}\PYG{+w}{ }\PYG{n}{end}\PYG{+w}{ }\PYG{o}{=}\PYG{+w}{ }\PYG{n}{std}\PYG{o}{:}\PYG{o}{:}\PYG{n}{sregex\PYGZus{}token\PYGZus{}iterator}\PYG{p}{\PYGZob{}}\PYG{p}{\PYGZcb{}}\PYG{p}{;}
\PYG{k}{for}\PYG{+w}{ }\PYG{p}{(}\PYG{k}{auto}\PYG{+w}{ }\PYG{n}{it}\PYG{+w}{ }\PYG{o}{=}\PYG{+w}{ }\PYG{n}{std}\PYG{o}{:}\PYG{o}{:}\PYG{n}{sregex\PYGZus{}token\PYGZus{}iterator}\PYG{p}{\PYGZob{}}\PYG{+w}{ }\PYG{n}{std}\PYG{o}{:}\PYG{o}{:}\PYG{n}{begin}\PYG{p}{(}\PYG{n}{text}\PYG{p}{)}\PYG{p}{,}
\PYG{+w}{               }\PYG{n}{std}\PYG{o}{:}\PYG{o}{:}\PYG{n}{end}\PYG{p}{(}\PYG{n}{text}\PYG{p}{)}\PYG{p}{,}\PYG{+w}{ }\PYG{n}{rx}\PYG{p}{,}\PYG{+w}{ }\PYG{l+m+mi}{\PYGZhy{}1}\PYG{+w}{ }\PYG{p}{\PYGZcb{}}\PYG{p}{;}
\PYG{+w}{     }\PYG{n}{it}\PYG{+w}{ }\PYG{o}{!}\PYG{o}{=}\PYG{+w}{ }\PYG{n}{end}\PYG{p}{;}\PYG{+w}{ }\PYG{o}{+}\PYG{o}{+}\PYG{n}{it}\PYG{p}{)}
\PYG{p}{\PYGZob{}}
\PYG{+w}{  }\PYG{n}{std}\PYG{o}{:}\PYG{o}{:}\PYG{n}{cout}\PYG{+w}{ }\PYG{o}{\PYGZlt{}}\PYG{o}{\PYGZlt{}}\PYG{+w}{ }\PYG{o}{*}\PYG{n}{it}\PYG{+w}{ }\PYG{o}{\PYGZlt{}}\PYG{o}{\PYGZlt{}}\PYG{+w}{ }\PYG{n}{std}\PYG{o}{:}\PYG{o}{:}\PYG{n}{endl}\PYG{p}{;}
\PYG{p}{\PYGZcb{}}
\end{Verbatim}
\end{codeblockbg}


%%C020453 p

该程序会输出如下内容（注意，空行实际上也是输出的一部分）：


%%C020454 code

\begin{codeblockbg}
\begin{Verbatim}[commandchars=\\\{\},breaklines=true,breakanywhere=true,fontsize=\codefont,xleftmargin=17bp]
\PYG{c+cp}{\PYGZsh{}}\PYG{c+cp}{remove \PYGZsh{} to uncomment the following lines}
\PYG{c+cp}{\PYGZsh{}}\PYG{c+cp}{retrycount=3}
\end{Verbatim}
\end{codeblockbg}


%%C020455 h3

\subsection{延伸阅读}


%%C020456 li

\begin{itemize}
\item \textit{用正则表达式验证字符串格式}


%%C020457 li

\item \textit{用正则表达式替换字符串内容}

\end{itemize}


%%C020458 h2

\section{用正则表达式替换字符串内容}


%%C020459 p

在前两个案例中，我们了解了如何在字符串或字符串的某一部分上匹配正则表达式，以及如何遍历匹配和子匹配。正则表达式库还支持基于正则表达式的文本替换。本案例将介绍如何用\myttbr{std::\allowbreak{}regex\_\allowbreak{}replace()}完成这类文本变换。


%%C020460 h3

\subsection{准备工作}


%%C020461 p

关于C++11中正则表达式支持的一般信息，请参考\textit{用正则表达式}\textit{验证字符串格式}案例。


%%C020462 h3

\subsection{具体做法}


%%C020463 p

要用正则表达式执行文本变换，应执行以下步骤：


%%C020464 p

包含\texttt{<\allowbreak{}regex>}和\texttt{<\allowbreak{}string>}以及命名空间


%%C020465 p

\myttbr{std::\allowbreak{}string\_\allowbreak{}literals}，它提供C++14标准中用于字符串的用户定义字面量：


%%C020466 code

\begin{codeblockbg}
\begin{Verbatim}[commandchars=\\\{\},breaklines=true,breakanywhere=true,fontsize=\codefont,xleftmargin=17bp]
\PYG{c+cp}{\PYGZsh{}}\PYG{c+cp}{include}\PYG{+w}{ }\PYG{c+cpf}{\PYGZlt{}regex\PYGZgt{}}
\PYG{c+cp}{\PYGZsh{}}\PYG{c+cp}{include}\PYG{+w}{ }\PYG{c+cpf}{\PYGZlt{}string\PYGZgt{}}
\PYG{k}{using}\PYG{+w}{ }\PYG{k}{namespace}\PYG{+w}{ }\PYG{n+nn}{std}\PYG{o}{:}\PYG{o}{:}\PYG{n+nn}{string\PYGZus{}literals}\PYG{p}{;}
\end{Verbatim}
\end{codeblockbg}


%%C020467 p

使用\myttbr{std::\allowbreak{}regex\_\allowbreak{}replace()}算法，并以替换字符串作为


%%C020468 p

第三个参数。看这个例子：把全部由三个字符组成、且每个字符都是\texttt{a}、\texttt{b}或\texttt{c}的单词替换为三个连字符：


%%C020469 code

\begin{codeblockbg}
\begin{Verbatim}[commandchars=\\\{\},breaklines=true,breakanywhere=true,fontsize=\codefont,xleftmargin=17bp]
\PYG{k}{auto}\PYG{+w}{ }\PYG{n}{text}\PYG{p}{\PYGZob{}}\PYG{l+s}{\PYGZdq{}}\PYG{l+s}{abc aa bca ca bbbb}\PYG{l+s}{\PYGZdq{}}\PYG{n}{s}\PYG{p}{\PYGZcb{}}\PYG{p}{;}
\PYG{k}{auto}\PYG{+w}{ }\PYG{n}{rx}\PYG{+w}{ }\PYG{o}{=}\PYG{+w}{ }\PYG{n}{std}\PYG{o}{:}\PYG{o}{:}\PYG{n}{regex}\PYG{p}{\PYGZob{}}\PYG{+w}{ }\PYG{l+s+sa}{R}\PYG{l+s}{\PYGZdq{}}\PYG{l+s+dl}{(}\PYG{l+s}{\PYGZbs{}b[a|b|c]\PYGZob{}3\PYGZcb{}\PYGZbs{}b}\PYG{l+s+dl}{)}\PYG{l+s}{\PYGZdq{}}\PYG{n}{s}\PYG{+w}{ }\PYG{p}{\PYGZcb{}}\PYG{p}{;}
\PYG{k}{auto}\PYG{+w}{ }\PYG{n}{newtext}\PYG{+w}{ }\PYG{o}{=}\PYG{+w}{ }\PYG{n}{std}\PYG{o}{:}\PYG{o}{:}\PYG{n}{regex\PYGZus{}replace}\PYG{p}{(}\PYG{n}{text}\PYG{p}{,}\PYG{+w}{ }\PYG{n}{rx}\PYG{p}{,}\PYG{+w}{ }\PYG{l+s}{\PYGZdq{}}\PYG{l+s}{\PYGZhy{}\PYGZhy{}\PYGZhy{}}\PYG{l+s}{\PYGZdq{}}\PYG{n}{s}\PYG{p}{)}\PYG{p}{;}
\end{Verbatim}
\end{codeblockbg}


%%C020470 p

使用\myttbr{std::\allowbreak{}regex\_\allowbreak{}replace()}算法，并以带前缀的匹配标识符作为


%%C020471 p

第三个参数，该前缀为\texttt{\$}。例如，把“lastname, firstname”形式的姓名替换为“firstname lastname”格式的姓名，如下所示：


%%C020472 code

\begin{codeblockbg}
\begin{Verbatim}[commandchars=\\\{\},breaklines=true,breakanywhere=true,fontsize=\codefont,xleftmargin=17bp]
\PYG{k}{auto}\PYG{+w}{ }\PYG{n}{text}\PYG{p}{\PYGZob{}}\PYG{+w}{ }\PYG{l+s}{\PYGZdq{}}\PYG{l+s}{bancila, marius}\PYG{l+s}{\PYGZdq{}}\PYG{n}{s}\PYG{+w}{ }\PYG{p}{\PYGZcb{}}\PYG{p}{;}
\PYG{k}{auto}\PYG{+w}{ }\PYG{n}{rx}\PYG{+w}{ }\PYG{o}{=}\PYG{+w}{ }\PYG{n}{std}\PYG{o}{:}\PYG{o}{:}\PYG{n}{regex}\PYG{p}{\PYGZob{}}\PYG{+w}{ }\PYG{l+s+sa}{R}\PYG{l+s}{\PYGZdq{}}\PYG{l+s+dl}{(}\PYG{l+s}{(\PYGZbs{}w+),\PYGZbs{}s*(\PYGZbs{}w+)}\PYG{l+s+dl}{)}\PYG{l+s}{\PYGZdq{}}\PYG{n}{s}\PYG{+w}{ }\PYG{p}{\PYGZcb{}}\PYG{p}{;}
\PYG{k}{auto}\PYG{+w}{ }\PYG{n}{newtext}\PYG{+w}{ }\PYG{o}{=}\PYG{+w}{ }\PYG{n}{std}\PYG{o}{:}\PYG{o}{:}\PYG{n}{regex\PYGZus{}replace}\PYG{p}{(}\PYG{n}{text}\PYG{p}{,}\PYG{+w}{ }\PYG{n}{rx}\PYG{p}{,}\PYG{+w}{ }\PYG{l+s}{\PYGZdq{}}\PYG{l+s}{\PYGZdl{}2 \PYGZdl{}1}\PYG{l+s}{\PYGZdq{}}\PYG{n}{s}\PYG{p}{)}\PYG{p}{;}
\end{Verbatim}
\end{codeblockbg}


%%C020473 h3

\subsection{工作原理}


%%C020474 p

\myttbr{std::\allowbreak{}regex\_\allowbreak{}replace()}算法有多个参数类型不同的重载版本，但各参数的含义如下：


%%C020475 li

\begin{itemize}
\item 执行替换所针对的输入字符串。


%%C020476 li

\item 一个封装正则表达式的\texttt{std::\allowbreak{}basic\_\allowbreak{}regex}对象，该正则表达式用于找出字符串中需要替换的部分。


%%C020477 li

\item 用于替换的字符串格式。


%%C020478 li

\item 可选的匹配标志。

\end{itemize}


%%C020479 p

返回值取决于所用的重载版本，可能是一个字符串，也可能是作为参数传入的输出迭代器的一份副本。用于替换的字符串格式既可以是普通字符串，也可以是用\texttt{\$}前缀表示的匹配标识符：


%%C020480 li

\begin{itemize}
\item \texttt{\$\&}表示整个匹配。


%%C020481 li

\item \texttt{\$1}、\texttt{\$2}、\texttt{\$3}等分别表示第一个、第二个、第三个子匹配，依此类推。


%%C020482 li

\item \texttt{\$`}表示第一个匹配之前的字符串部分。


%%C020483 li

\item \texttt{\$‘}表示最后一个匹配之后的字符串部分。

\end{itemize}


%%C020484 p

在\textit{具体做法}一节所示的第一个示例中，初始文本包含两个恰好由三个\texttt{a}、\texttt{b}或\texttt{c}字符组成的单词：\texttt{abc}和\texttt{bca}。该正则表达式表示位于单词边界之间、恰好三个字符的表达式。这意味着像\texttt{bbbb}这样的子文本不会与表达式匹配。替换的结果是，字符串text变为\texttt{--- aa --- ca bbbb}。


%%C020485 p

还可以为\myttbr{std::\allowbreak{}regex\_\allowbreak{}replace()}算法指定额外的匹配标志。默认的匹配标志是\myttbr{std::\allowbreak{}regex\_\allowbreak{}constants::\allowbreak{}match\_\allowbreak{}default}，它基本上指定用ECMAScript作为构造正则表达式的文法。例如，如果只想替换第一个出现位置，可以指定\myttbr{std::\allowbreak{}regex\_\allowbreak{}constants::\allowbreak{}format\_\allowbreak{}first\_\allowbreak{}only}。在下一个示例中，结果是\texttt{---\allowbreak{}aa bca ca bbbb}，因为找到第一个匹配后替换就停止了：


%%C020486 code

\begin{codeblockbg}
\begin{Verbatim}[commandchars=\\\{\},breaklines=true,breakanywhere=true,fontsize=\codefont,xleftmargin=17bp]
\PYG{k}{auto}\PYG{+w}{ }\PYG{n}{text}\PYG{p}{\PYGZob{}}\PYG{+w}{ }\PYG{l+s}{\PYGZdq{}}\PYG{l+s}{abc aa bca ca bbbb}\PYG{l+s}{\PYGZdq{}}\PYG{n}{s}\PYG{+w}{ }\PYG{p}{\PYGZcb{}}\PYG{p}{;}
\PYG{k}{auto}\PYG{+w}{ }\PYG{n}{rx}\PYG{+w}{ }\PYG{o}{=}\PYG{+w}{ }\PYG{n}{std}\PYG{o}{:}\PYG{o}{:}\PYG{n}{regex}\PYG{p}{\PYGZob{}}\PYG{+w}{ }\PYG{l+s+sa}{R}\PYG{l+s}{\PYGZdq{}}\PYG{l+s+dl}{(}\PYG{l+s}{\PYGZbs{}b[a|b|c]\PYGZob{}3\PYGZcb{}\PYGZbs{}b}\PYG{l+s+dl}{)}\PYG{l+s}{\PYGZdq{}}\PYG{n}{s}\PYG{+w}{ }\PYG{p}{\PYGZcb{}}\PYG{p}{;}
\PYG{k}{auto}\PYG{+w}{ }\PYG{n}{newtext}\PYG{+w}{ }\PYG{o}{=}\PYG{+w}{ }\PYG{n}{std}\PYG{o}{:}\PYG{o}{:}\PYG{n}{regex\PYGZus{}replace}\PYG{p}{(}\PYG{n}{text}\PYG{p}{,}\PYG{+w}{ }\PYG{n}{rx}\PYG{p}{,}\PYG{+w}{ }\PYG{l+s}{\PYGZdq{}}\PYG{l+s}{\PYGZhy{}\PYGZhy{}\PYGZhy{}}\PYG{l+s}{\PYGZdq{}}\PYG{n}{s}\PYG{p}{,}
\PYG{+w}{                 }\PYG{n}{std}\PYG{o}{:}\PYG{o}{:}\PYG{n}{regex\PYGZus{}constants}\PYG{o}{:}\PYG{o}{:}\PYG{n}{format\PYGZus{}first\PYGZus{}only}\PYG{p}{)}\PYG{p}{;}
\end{Verbatim}
\end{codeblockbg}


%%C020487 p

不过，替换字符串可以包含特殊指示符，分别表示整个匹配、某个特定的子匹配或未被匹配的部分，如前所述。在\textit{具体做法}一节所示的第二个示例中，正则表达式识别出一个至少一个字符的单词，其后是逗号和可能存在的空白，然后又是另一个至少一个字符的单词。第一个单词应当是姓，第二个单词应当是名。替换字符串的格式为\texttt{\$2 \$1}。这条指令要求把匹配到的表达式（在本例中是整个原始字符串）替换为另一个字符串，该字符串由第二个子匹配、一个空格和第一个子匹配依次组成。


%%C020488 p

在本例中，整个字符串就是一个匹配。在下一个示例中，字符串内部会有多个匹配，它们都会被替换为指定的字符串。这个示例把出现在以元音字母开头的单词之前的不定冠词\textit{a}替换为不定冠词\textit{an}（当然，这并不涵盖以元音音素开头的单词）：


%%C020489 code

\begin{codeblockbg}
\begin{Verbatim}[commandchars=\\\{\},breaklines=true,breakanywhere=true,fontsize=\codefont,xleftmargin=17bp]
\PYG{k}{auto}\PYG{+w}{ }\PYG{n}{text}\PYG{p}{\PYGZob{}}\PYG{l+s}{\PYGZdq{}}\PYG{l+s}{this is a example with a error}\PYG{l+s}{\PYGZdq{}}\PYG{n}{s}\PYG{p}{\PYGZcb{}}\PYG{p}{;}
\PYG{k}{auto}\PYG{+w}{ }\PYG{n}{rx}\PYG{+w}{ }\PYG{o}{=}\PYG{+w}{ }\PYG{n}{std}\PYG{o}{:}\PYG{o}{:}\PYG{n}{regex}\PYG{p}{\PYGZob{}}\PYG{l+s+sa}{R}\PYG{l+s}{\PYGZdq{}}\PYG{l+s+dl}{(}\PYG{l+s}{\PYGZbs{}ba ((a|e|i|u|o)\PYGZbs{}w+)}\PYG{l+s+dl}{)}\PYG{l+s}{\PYGZdq{}}\PYG{n}{s}\PYG{p}{\PYGZcb{}}\PYG{p}{;}
\PYG{k}{auto}\PYG{+w}{ }\PYG{n}{newtext}\PYG{+w}{ }\PYG{o}{=}\PYG{+w}{ }\PYG{n}{std}\PYG{o}{:}\PYG{o}{:}\PYG{n}{regex\PYGZus{}replace}\PYG{p}{(}\PYG{n}{text}\PYG{p}{,}\PYG{+w}{ }\PYG{n}{rx}\PYG{p}{,}\PYG{+w}{ }\PYG{l+s}{\PYGZdq{}}\PYG{l+s}{an \PYGZdl{}1}\PYG{l+s}{\PYGZdq{}}\PYG{p}{)}\PYG{p}{;}
\end{Verbatim}
\end{codeblockbg}


%%C020490 p

该正则表达式把字母\textit{a}识别为单个单词（\texttt{\textbackslash{}b}表示单词边界，所以\texttt{\textbackslash{}ba}表示只含单个字母\textit{a}的单词），其后是一个空格和一个以元音字母开头、至少两个字符的单词。找到这样的匹配后，就把它替换为一个字符串，该字符串由固定字符串\textit{an}、一个空格以及匹配的第一个子表达式（即那个单词本身）组成。在本例中，\texttt{newtext}字符串将是\textit{this is an} \textit{example with an error}。


%%C020491 p

除了子表达式的标识符（\texttt{\$1}、\texttt{\$2}等）之外，还有其他标识符，分别表示整个匹配（\texttt{\$\&}）、第一个匹配之前的字符串部分（\texttt{\$`}）以及最后一个匹配之后的字符串部分（\texttt{\$‘}）。在最后一个示例中，我们把日期的格式从\texttt{dd.\allowbreak{}mm.\allowbreak{}yyyy}改为\texttt{yyyy.\allowbreak{}mm.\allowbreak{}dd}，同时还会显示匹配到的各个部分：


%%C020492 code

\begin{codeblockbg}
\begin{Verbatim}[commandchars=\\\{\},breaklines=true,breakanywhere=true,fontsize=\codefont,xleftmargin=17bp]
\PYG{k}{auto}\PYG{+w}{ }\PYG{n}{text}\PYG{p}{\PYGZob{}}\PYG{l+s}{\PYGZdq{}}\PYG{l+s}{today is 1.06.2016!!}\PYG{l+s}{\PYGZdq{}}\PYG{n}{s}\PYG{p}{\PYGZcb{}}\PYG{p}{;}
\PYG{k}{auto}\PYG{+w}{ }\PYG{n}{rx}\PYG{+w}{ }\PYG{o}{=}
\PYG{+w}{   }\PYG{n}{std}\PYG{o}{:}\PYG{o}{:}\PYG{n}{regex}\PYG{p}{\PYGZob{}}\PYG{l+s+sa}{R}\PYG{l+s}{\PYGZdq{}}\PYG{l+s+dl}{(}\PYG{l+s}{(\PYGZbs{}d\PYGZob{}1,2\PYGZcb{})(\PYGZbs{}.|\PYGZhy{}|/)(\PYGZbs{}d\PYGZob{}1,2\PYGZcb{})(\PYGZbs{}.|\PYGZhy{}|/)(\PYGZbs{}d\PYGZob{}4\PYGZcb{})}\PYG{l+s+dl}{)}\PYG{l+s}{\PYGZdq{}}\PYG{n}{s}\PYG{p}{\PYGZcb{}}\PYG{p}{;}
\PYG{c+c1}{// today is 2016.06.1!!}
\PYG{k}{auto}\PYG{+w}{ }\PYG{n}{newtext1}\PYG{+w}{ }\PYG{o}{=}\PYG{+w}{ }\PYG{n}{std}\PYG{o}{:}\PYG{o}{:}\PYG{n}{regex\PYGZus{}replace}\PYG{p}{(}\PYG{n}{text}\PYG{p}{,}\PYG{+w}{ }\PYG{n}{rx}\PYG{p}{,}\PYG{+w}{ }\PYG{l+s+sa}{R}\PYG{l+s}{\PYGZdq{}}\PYG{l+s+dl}{(}\PYG{l+s}{\PYGZdl{}5\PYGZdl{}4\PYGZdl{}3\PYGZdl{}2\PYGZdl{}1}\PYG{l+s+dl}{)}\PYG{l+s}{\PYGZdq{}}\PYG{p}{)}\PYG{p}{;}
\PYG{c+c1}{// today is [today is ][1.06.2016][!!]!!}
\PYG{k}{auto}\PYG{+w}{ }\PYG{n}{newtext2}\PYG{+w}{ }\PYG{o}{=}\PYG{+w}{ }\PYG{n}{std}\PYG{o}{:}\PYG{o}{:}\PYG{n}{regex\PYGZus{}replace}\PYG{p}{(}\PYG{n}{text}\PYG{p}{,}\PYG{+w}{ }\PYG{n}{rx}\PYG{p}{,}\PYG{+w}{ }\PYG{l+s+sa}{R}\PYG{l+s}{\PYGZdq{}}\PYG{l+s+dl}{(}\PYG{l+s}{[\PYGZdl{}`][\PYGZdl{}\PYGZam{}][\PYGZdl{}\PYGZsq{}]}\PYG{l+s+dl}{)}\PYG{l+s}{\PYGZdq{}}\PYG{p}{)}\PYG{p}{;}
\end{Verbatim}
\end{codeblockbg}


%%C020493 p

该正则表达式匹配一个一位或两位数字，随后是点、连字符或斜杠；然后又是另一个一位或两位数字；再是点、连字符或斜杠；最后是一个四位数字。


%%C020494 p

对于\texttt{newtext1}，替换字符串是\texttt{\$5\$4\$3\$2\$1}，即年、第二个分隔符、月、第一个分隔符，最后是日。因此，对于输入字符串\textit{“today is 1.06.2016!”}，结果是\textit{“today is 2016.06.1!!”}。


%%C020495 p

对于\texttt{newtext2}，替换字符串是\texttt{[\$`][\$\&][\$‘]}，即用方括号依次括起第一个匹配之前的部分、整个匹配，最后是最后一个匹配之后的部分。然而，结果并不是你乍看之下可能预期的那样\textit{“[!!][1.06.2016][today is ]”}，而是\textit{“today is [today is ][1.06.2016][!!]!!”}。原因在于被替换的是匹配到的表达式，而在这里它只是日期（\textit{“1.06.2016”}）。这个子串被替换为由初始字符串所有部分共同组成的另一个字符串。


%%C020496 h3

\subsection{延伸阅读}


%%C020497 li

\begin{itemize}
\item \textit{用正则表达式验证字符串格式}


%%C020498 li

\item \textit{用正则表达式解析字符串内容}

\end{itemize}


%%C020499 h2

\section{用string\_view代替常量字符串引用}

%%C020500 p

处理字符串时，临时对象随时都在产生，你可能根本没有意识到这一点。很多时候这些临时对象毫无用处，只是把数据从一个地方复制到另一个地方（例如从函数复制给调用者）。这会带来性能问题，因为它们需要分配内存并复制数据，而这些正是应当避免的。为此，C++17标准提供了一个新的字符串类模板\myttbr{std::\allowbreak{}basic\_\allowbreak{}string\_\allowbreak{}view}，它表示对字符串（即字符序列）的非拥有常量引用。本案例将介绍应当在何时以及如何使用这个类。


%%C020501 h3

\subsection{准备工作}


%%C020502 p

\texttt{string\_\allowbreak{}view}类位于\texttt{std}命名空间中，定义在\texttt{string\_\allowbreak{}view}头文件里。


%%C020503 h3

\subsection{具体做法}


%%C020504 p

应当用\texttt{std::\allowbreak{}string\_\allowbreak{}view}向函数传递参数（或从函数返回值），而不是\myttbr{std::\allowbreak{}string const \&}，除非你的代码需要调用其他接受\texttt{std::\allowbreak{}string}参数的函数（这种情况下就需要进行转换）：


%%C020505 code

\begin{codeblockbg}
\begin{Verbatim}[commandchars=\\\{\},breaklines=true,breakanywhere=true,fontsize=\codefont,xleftmargin=17bp]
\PYG{n}{std}\PYG{o}{:}\PYG{o}{:}\PYG{n}{string\PYGZus{}view}\PYG{+w}{ }\PYG{n+nf}{get\PYGZus{}filename}\PYG{p}{(}\PYG{n}{std}\PYG{o}{:}\PYG{o}{:}\PYG{n}{string\PYGZus{}view}\PYG{+w}{ }\PYG{n}{str}\PYG{p}{)}
\PYG{p}{\PYGZob{}}
\PYG{+w}{  }\PYG{k}{auto}\PYG{+w}{ }\PYG{k}{const}\PYG{+w}{ }\PYG{n}{pos1}\PYG{+w}{ }\PYG{p}{\PYGZob{}}\PYG{n}{str}\PYG{p}{.}\PYG{n}{find\PYGZus{}last\PYGZus{}of}\PYG{p}{(}\PYG{err}{\PYGZsq{}}\PYG{err}{\PYGZsq{}}\PYG{p}{)}\PYG{p}{\PYGZcb{}}\PYG{p}{;}
\PYG{+w}{  }\PYG{k}{auto}\PYG{+w}{ }\PYG{k}{const}\PYG{+w}{ }\PYG{n}{pos2}\PYG{+w}{ }\PYG{p}{\PYGZob{}}\PYG{n}{str}\PYG{p}{.}\PYG{n}{find\PYGZus{}last\PYGZus{}of}\PYG{p}{(}\PYG{l+s+sc}{\PYGZsq{}}\PYG{l+s+sc}{.}\PYG{l+s+sc}{\PYGZsq{}}\PYG{p}{)}\PYG{p}{\PYGZcb{}}\PYG{p}{;}
\PYG{+w}{  }\PYG{k}{return}\PYG{+w}{ }\PYG{n}{str}\PYG{p}{.}\PYG{n}{substr}\PYG{p}{(}\PYG{n}{pos1}\PYG{+w}{ }\PYG{o}{+}\PYG{+w}{ }\PYG{l+m+mi}{1}\PYG{p}{,}\PYG{+w}{ }\PYG{n}{pos2}\PYG{+w}{ }\PYG{o}{\PYGZhy{}}\PYG{+w}{ }\PYG{n}{pos1}\PYG{+w}{ }\PYG{o}{\PYGZhy{}}\PYG{+w}{ }\PYG{l+m+mi}{1}\PYG{p}{)}\PYG{p}{;}
\PYG{p}{\PYGZcb{}}
\PYG{k+kt}{char}\PYG{+w}{ }\PYG{k}{const}\PYG{+w}{ }\PYG{n}{file1}\PYG{p}{[}\PYG{p}{]}\PYG{+w}{ }\PYG{p}{\PYGZob{}}\PYG{l+s+sa}{R}\PYG{l+s}{\PYGZdq{}}\PYG{l+s+dl}{(}\PYG{l+s}{c:\PYGZbs{}test\PYGZbs{}example1.doc}\PYG{l+s+dl}{)}\PYG{l+s}{\PYGZdq{}}\PYG{p}{\PYGZcb{}}\PYG{p}{;}
\PYG{k}{auto}\PYG{+w}{ }\PYG{n}{name1}\PYG{+w}{ }\PYG{o}{=}\PYG{+w}{ }\PYG{n}{get\PYGZus{}filename}\PYG{p}{(}\PYG{n}{file1}\PYG{p}{)}\PYG{p}{;}
\PYG{n}{std}\PYG{o}{:}\PYG{o}{:}\PYG{n}{string}\PYG{+w}{ }\PYG{n}{file2}\PYG{+w}{ }\PYG{p}{\PYGZob{}}\PYG{l+s+sa}{R}\PYG{l+s}{\PYGZdq{}}\PYG{l+s+dl}{(}\PYG{l+s}{c:\PYGZbs{}test\PYGZbs{}example2}\PYG{l+s+dl}{)}\PYG{l+s}{\PYGZdq{}}\PYG{p}{\PYGZcb{}}\PYG{p}{;}
\PYG{k}{auto}\PYG{+w}{ }\PYG{n}{name2}\PYG{+w}{ }\PYG{o}{=}\PYG{+w}{ }\PYG{n}{get\PYGZus{}filename}\PYG{p}{(}\PYG{n}{file2}\PYG{p}{)}\PYG{p}{;}
\PYG{k}{auto}\PYG{+w}{ }\PYG{n}{name3}\PYG{+w}{ }\PYG{o}{=}\PYG{+w}{ }\PYG{n}{get\PYGZus{}filename}\PYG{p}{(}\PYG{n}{std}\PYG{o}{:}\PYG{o}{:}\PYG{n}{string\PYGZus{}view}\PYG{p}{\PYGZob{}}\PYG{n}{file1}\PYG{p}{,}\PYG{+w}{ }\PYG{l+m+mi}{16}\PYG{p}{\PYGZcb{}}\PYG{p}{)}\PYG{p}{;}
\end{Verbatim}
\end{codeblockbg}


%%C020506 h3

\subsection{工作原理}


%%C020507 p

在考察这个新的字符串类型如何工作之前，先看看下面这个函数示例，它的作用是在去掉扩展名的情况下提取文件名。这基本上就是你在C++17之前会写出的上一节那个函数。


%%C020508 note

\begin{myNote}
注意，本示例中的文件分隔符是\texttt{\textbackslash{}}（反斜杠），与Windows中一致。对于基于Linux的系统，必须把它改成\texttt{/}（斜杠）。
\end{myNote}


%%C020509 code

\begin{codeblockbg}
\begin{Verbatim}[commandchars=\\\{\},breaklines=true,breakanywhere=true,fontsize=\codefont,xleftmargin=17bp]
\PYG{n}{std}\PYG{o}{:}\PYG{o}{:}\PYG{n}{string}\PYG{+w}{ }\PYG{n+nf}{get\PYGZus{}filename}\PYG{p}{(}\PYG{n}{std}\PYG{o}{:}\PYG{o}{:}\PYG{n}{string}\PYG{+w}{ }\PYG{k}{const}\PYG{+w}{ }\PYG{o}{\PYGZam{}}\PYG{+w}{ }\PYG{n}{str}\PYG{p}{)}
\PYG{p}{\PYGZob{}}
\PYG{+w}{  }\PYG{k}{auto}\PYG{+w}{ }\PYG{k}{const}\PYG{+w}{ }\PYG{n}{pos1}\PYG{+w}{ }\PYG{p}{\PYGZob{}}\PYG{n}{str}\PYG{p}{.}\PYG{n}{find\PYGZus{}last\PYGZus{}of}\PYG{p}{(}\PYG{err}{\PYGZsq{}}\PYG{err}{\PYGZsq{}}\PYG{p}{)}\PYG{p}{\PYGZcb{}}\PYG{p}{;}
\PYG{+w}{  }\PYG{k}{auto}\PYG{+w}{ }\PYG{k}{const}\PYG{+w}{ }\PYG{n}{pos2}\PYG{+w}{ }\PYG{p}{\PYGZob{}}\PYG{n}{str}\PYG{p}{.}\PYG{n}{find\PYGZus{}last\PYGZus{}of}\PYG{p}{(}\PYG{l+s+sc}{\PYGZsq{}}\PYG{l+s+sc}{.}\PYG{l+s+sc}{\PYGZsq{}}\PYG{p}{)}\PYG{p}{\PYGZcb{}}\PYG{p}{;}
\PYG{+w}{  }\PYG{k}{return}\PYG{+w}{ }\PYG{n}{str}\PYG{p}{.}\PYG{n}{substr}\PYG{p}{(}\PYG{n}{pos1}\PYG{+w}{ }\PYG{o}{+}\PYG{+w}{ }\PYG{l+m+mi}{1}\PYG{p}{,}\PYG{+w}{ }\PYG{n}{pos2}\PYG{+w}{ }\PYG{o}{\PYGZhy{}}\PYG{+w}{ }\PYG{n}{pos1}\PYG{+w}{ }\PYG{o}{\PYGZhy{}}\PYG{+w}{ }\PYG{l+m+mi}{1}\PYG{p}{)}\PYG{p}{;}
\PYG{p}{\PYGZcb{}}
\PYG{k}{auto}\PYG{+w}{ }\PYG{n}{name1}\PYG{+w}{ }\PYG{o}{=}\PYG{+w}{ }\PYG{n}{get\PYGZus{}filename}\PYG{p}{(}\PYG{l+s+sa}{R}\PYG{l+s}{\PYGZdq{}}\PYG{l+s+dl}{(}\PYG{l+s}{c:\PYGZbs{}test\PYGZbs{}example1.doc}\PYG{l+s+dl}{)}\PYG{l+s}{\PYGZdq{}}\PYG{p}{)}\PYG{p}{;}\PYG{+w}{ }\PYG{c+c1}{// example1}
\PYG{k}{auto}\PYG{+w}{ }\PYG{n}{name2}\PYG{+w}{ }\PYG{o}{=}\PYG{+w}{ }\PYG{n}{get\PYGZus{}filename}\PYG{p}{(}\PYG{l+s+sa}{R}\PYG{l+s}{\PYGZdq{}}\PYG{l+s+dl}{(}\PYG{l+s}{c:\PYGZbs{}test\PYGZbs{}example2}\PYG{l+s+dl}{)}\PYG{l+s}{\PYGZdq{}}\PYG{p}{)}\PYG{p}{;}\PYG{+w}{     }\PYG{c+c1}{// example2}
\PYG{k}{if}\PYG{p}{(}\PYG{n}{get\PYGZus{}filename}\PYG{p}{(}\PYG{l+s+sa}{R}\PYG{l+s}{\PYGZdq{}}\PYG{l+s+dl}{(}\PYG{l+s}{c:\PYGZbs{}test\PYGZbs{}\PYGZus{}sample\PYGZus{}.tmp}\PYG{l+s+dl}{)}\PYG{l+s}{\PYGZdq{}}\PYG{p}{)}\PYG{p}{.}\PYG{n}{front}\PYG{p}{(}\PYG{p}{)}\PYG{+w}{ }\PYG{o}{=}\PYG{o}{=}\PYG{+w}{ }\PYG{l+s+sc}{\PYGZsq{}}\PYG{l+s+sc}{\PYGZus{}}\PYG{l+s+sc}{\PYGZsq{}}\PYG{p}{)}\PYG{+w}{ }\PYG{p}{\PYGZob{}}\PYG{p}{\PYGZcb{}}
\end{Verbatim}
\end{codeblockbg}


%%C020510 p

这是一个相当简单的函数。它接受一个\texttt{std::\allowbreak{}string}的常量引用，找出位于最后一个文件分隔符与最后一个点之间的子串，这基本上就是不带扩展名（也不含目录名）的文件名。


%%C020511 p

然而这段代码的问题在于，它会创建一到两个、甚至更多的临时对象，具体取决于编译器的优化。函数参数是\texttt{std::\allowbreak{}string}的常量引用，但调用该函数时传入的是字符串字面量，这意味着需要由字面量构造出\texttt{std::\allowbreak{}string}。这些临时对象需要分配和复制数据，既耗时又耗费资源。在最后一个示例中，我们只想检查文件名的第一个字符是不是下划线，却为此至少创建了两个临时字符串对象。


%%C020512 p

\myttbr{std::\allowbreak{}basic\_\allowbreak{}string\_\allowbreak{}view}类模板正是为解决这个问题而提供的。这个类模板与\texttt{std::\allowbreak{}basic\_\allowbreak{}string}非常相似，两者的接口几乎相同。原因在于，\myttbr{std::\allowbreak{}basic\_\allowbreak{}string\_\allowbreak{}view}的定位就是替代对\texttt{std::\allowbreak{}basic\_\allowbreak{}string}的常量引用，而不需要改动其他代码。与\texttt{std::\allowbreak{}basic\_\allowbreak{}string}一样，它为所有标准字符类型都提供了特化：


%%C020513 code

\begin{codeblockbg}
\begin{Verbatim}[commandchars=\\\{\},breaklines=true,breakanywhere=true,fontsize=\codefont,xleftmargin=17bp]
\PYG{k}{typedef}\PYG{+w}{ }\PYG{n}{basic\PYGZus{}string\PYGZus{}view}\PYG{o}{\PYGZlt{}}\PYG{k+kt}{char}\PYG{o}{\PYGZgt{}}\PYG{+w}{     }\PYG{n}{string\PYGZus{}view}\PYG{p}{;}
\PYG{k}{typedef}\PYG{+w}{ }\PYG{n}{basic\PYGZus{}string\PYGZus{}view}\PYG{o}{\PYGZlt{}}\PYG{k+kt}{wchar\PYGZus{}t}\PYG{o}{\PYGZgt{}}\PYG{+w}{  }\PYG{n}{wstring\PYGZus{}view}\PYG{p}{;}
\PYG{k}{typedef}\PYG{+w}{ }\PYG{n}{basic\PYGZus{}string\PYGZus{}view}\PYG{o}{\PYGZlt{}}\PYG{k+kt}{char16\PYGZus{}t}\PYG{o}{\PYGZgt{}}\PYG{+w}{ }\PYG{n}{u16string\PYGZus{}view}\PYG{p}{;}
\PYG{k}{typedef}\PYG{+w}{ }\PYG{n}{basic\PYGZus{}string\PYGZus{}view}\PYG{o}{\PYGZlt{}}\PYG{k+kt}{char32\PYGZus{}t}\PYG{o}{\PYGZgt{}}\PYG{+w}{ }\PYG{n}{u32string\PYGZus{}view}\PYG{p}{;}
\end{Verbatim}
\end{codeblockbg}


%%C020514 p

\myttbr{std::\allowbreak{}basic\_\allowbreak{}string\_\allowbreak{}view}类模板定义了对一段连续常量字符序列的引用。顾名思义，它表示一个视图，不能用来修改所引用的字符序列。\myttbr{std::\allowbreak{}basic\_\allowbreak{}string\_\allowbreak{}view}对象占用的空间相当小，因为它只需要一个指向序列首字符的指针和长度。它不仅可以用\texttt{std::\allowbreak{}basic\_\allowbreak{}string}对象构造，还可以用指针加长度构造，或者用以空字符结尾的字符序列构造（后一种情况需要先遍历字符串以求出长度）。因此，\myttbr{std::\allowbreak{}basic\_\allowbreak{}string\_\allowbreak{}view}类模板也可以作为多种字符串类型的公共接口（只要数据只需读取即可）。另一方面，从\myttbr{std::\allowbreak{}basic\_\allowbreak{}string\_\allowbreak{}view}转换成\texttt{std::\allowbreak{}basic\_\allowbreak{}string}很容易，因为前者既有\texttt{to\_\allowbreak{}string()}，也有到\texttt{std::\allowbreak{}basic\_\allowbreak{}string}的转换\texttt{operator}，可以创建新的\texttt{std::\allowbreak{}basic\_\allowbreak{}string}对象。


%%C020515 p

向函数传递\myttbr{std::\allowbreak{}basic\_\allowbreak{}string\_\allowbreak{}view}、以及从函数返回\myttbr{std::\allowbreak{}basic\_\allowbreak{}string\_\allowbreak{}view}，仍然会创建该类型的临时对象，但这些是栈上的小对象（在64位平台上，一个指针加一个长度可能是16字节）；因此，它们的性能开销应当比分配堆空间并复制数据更小。


%%C020516 note

\begin{myNote}
注意，微软对\texttt{std::\allowbreak{}basic\_\allowbreak{}string}的实现针对短字符串提供了一项优化：静态分配一个16字符的缓冲区，不涉及堆操作，只有当字符串长度超过16个字符时才需要堆操作。
\end{myNote}


%%C020517 p

除了与\texttt{std::\allowbreak{}basic\_\allowbreak{}string}中完全相同的方法之外，\myttbr{std::\allowbreak{}basic\_\allowbreak{}string\_\allowbreak{}view}还多了两个方法：


%%C020518 li

\begin{itemize}
\item \texttt{remove\_\allowbreak{}prefix()}：把起点后移\textit{N}个字符、长度减少\textit{N}个字符，从而缩小视图。


%%C020519 li

\item \texttt{remove\_\allowbreak{}suffix()}：把长度减少\textit{N}个字符，从而缩小视图。

\end{itemize}


%%C020520 p

下面的示例用这两个成员函数去掉\texttt{std::\allowbreak{}string\_\allowbreak{}view}开头和结尾的空格。函数的实现先找出第一个不是空格的元素，再找出最后一个不是空格的元素。然后从尾部删除最后一个非空格字符之后的所有内容，从头部删除第一个非空格字符之前的所有内容。函数返回两端都已去除空格的新视图：


%%C020521 code

\begin{codeblockbg}
\begin{Verbatim}[commandchars=\\\{\},breaklines=true,breakanywhere=true,fontsize=\codefont,xleftmargin=17bp]
\PYG{n}{std}\PYG{o}{:}\PYG{o}{:}\PYG{n}{string\PYGZus{}view}\PYG{+w}{ }\PYG{n+nf}{trim\PYGZus{}view}\PYG{p}{(}\PYG{n}{std}\PYG{o}{:}\PYG{o}{:}\PYG{n}{string\PYGZus{}view}\PYG{+w}{ }\PYG{n}{str}\PYG{p}{)}
\PYG{p}{\PYGZob{}}
\PYG{+w}{  }\PYG{k}{auto}\PYG{+w}{ }\PYG{k}{const}\PYG{+w}{ }\PYG{n}{pos1}\PYG{p}{\PYGZob{}}\PYG{+w}{ }\PYG{n}{str}\PYG{p}{.}\PYG{n}{find\PYGZus{}first\PYGZus{}not\PYGZus{}of}\PYG{p}{(}\PYG{l+s}{\PYGZdq{}}\PYG{l+s}{ }\PYG{l+s}{\PYGZdq{}}\PYG{p}{)}\PYG{+w}{ }\PYG{p}{\PYGZcb{}}\PYG{p}{;}
\PYG{+w}{  }\PYG{k}{auto}\PYG{+w}{ }\PYG{k}{const}\PYG{+w}{ }\PYG{n}{pos2}\PYG{p}{\PYGZob{}}\PYG{+w}{ }\PYG{n}{str}\PYG{p}{.}\PYG{n}{find\PYGZus{}last\PYGZus{}not\PYGZus{}of}\PYG{p}{(}\PYG{l+s}{\PYGZdq{}}\PYG{l+s}{ }\PYG{l+s}{\PYGZdq{}}\PYG{p}{)}\PYG{+w}{ }\PYG{p}{\PYGZcb{}}\PYG{p}{;}
\PYG{+w}{  }\PYG{n}{str}\PYG{p}{.}\PYG{n}{remove\PYGZus{}suffix}\PYG{p}{(}\PYG{n}{str}\PYG{p}{.}\PYG{n}{length}\PYG{p}{(}\PYG{p}{)}\PYG{+w}{ }\PYG{o}{\PYGZhy{}}\PYG{+w}{ }\PYG{n}{pos2}\PYG{+w}{ }\PYG{o}{\PYGZhy{}}\PYG{+w}{ }\PYG{l+m+mi}{1}\PYG{p}{)}\PYG{p}{;}
\PYG{+w}{  }\PYG{n}{str}\PYG{p}{.}\PYG{n}{remove\PYGZus{}prefix}\PYG{p}{(}\PYG{n}{pos1}\PYG{p}{)}\PYG{p}{;}
\PYG{+w}{  }\PYG{k}{return}\PYG{+w}{ }\PYG{n}{str}\PYG{p}{;}
\PYG{p}{\PYGZcb{}}
\PYG{k}{auto}\PYG{+w}{ }\PYG{n}{sv1}\PYG{p}{\PYGZob{}}\PYG{+w}{ }\PYG{n}{trim\PYGZus{}view}\PYG{p}{(}\PYG{l+s}{\PYGZdq{}}\PYG{l+s}{sample}\PYG{l+s}{\PYGZdq{}}\PYG{p}{)}\PYG{+w}{ }\PYG{p}{\PYGZcb{}}\PYG{p}{;}
\PYG{k}{auto}\PYG{+w}{ }\PYG{n}{sv2}\PYG{p}{\PYGZob{}}\PYG{+w}{ }\PYG{n}{trim\PYGZus{}view}\PYG{p}{(}\PYG{l+s}{\PYGZdq{}}\PYG{l+s}{  sample}\PYG{l+s}{\PYGZdq{}}\PYG{p}{)}\PYG{+w}{ }\PYG{p}{\PYGZcb{}}\PYG{p}{;}
\PYG{k}{auto}\PYG{+w}{ }\PYG{n}{sv3}\PYG{p}{\PYGZob{}}\PYG{+w}{ }\PYG{n}{trim\PYGZus{}view}\PYG{p}{(}\PYG{l+s}{\PYGZdq{}}\PYG{l+s}{sample  }\PYG{l+s}{\PYGZdq{}}\PYG{p}{)}\PYG{+w}{ }\PYG{p}{\PYGZcb{}}\PYG{p}{;}
\PYG{k}{auto}\PYG{+w}{ }\PYG{n}{sv4}\PYG{p}{\PYGZob{}}\PYG{+w}{ }\PYG{n}{trim\PYGZus{}view}\PYG{p}{(}\PYG{l+s}{\PYGZdq{}}\PYG{l+s}{  sample  }\PYG{l+s}{\PYGZdq{}}\PYG{p}{)}\PYG{+w}{ }\PYG{p}{\PYGZcb{}}\PYG{p}{;}
\PYG{k}{auto}\PYG{+w}{ }\PYG{n}{s1}\PYG{p}{\PYGZob{}}\PYG{+w}{ }\PYG{n}{sv1}\PYG{p}{.}\PYG{n}{to\PYGZus{}string}\PYG{p}{(}\PYG{p}{)}\PYG{+w}{ }\PYG{p}{\PYGZcb{}}\PYG{p}{;}
\PYG{k}{auto}\PYG{+w}{ }\PYG{n}{s2}\PYG{p}{\PYGZob{}}\PYG{+w}{ }\PYG{n}{sv2}\PYG{p}{.}\PYG{n}{to\PYGZus{}string}\PYG{p}{(}\PYG{p}{)}\PYG{+w}{ }\PYG{p}{\PYGZcb{}}\PYG{p}{;}
\PYG{k}{auto}\PYG{+w}{ }\PYG{n}{s3}\PYG{p}{\PYGZob{}}\PYG{+w}{ }\PYG{n}{sv3}\PYG{p}{.}\PYG{n}{to\PYGZus{}string}\PYG{p}{(}\PYG{p}{)}\PYG{+w}{ }\PYG{p}{\PYGZcb{}}\PYG{p}{;}
\PYG{k}{auto}\PYG{+w}{ }\PYG{n}{s4}\PYG{p}{\PYGZob{}}\PYG{+w}{ }\PYG{n}{sv4}\PYG{p}{.}\PYG{n}{to\PYGZus{}string}\PYG{p}{(}\PYG{p}{)}\PYG{+w}{ }\PYG{p}{\PYGZcb{}}\PYG{p}{;}
\end{Verbatim}
\end{codeblockbg}


%%C020522 note

\begin{myNote}
使用\myttbr{std::\allowbreak{}basic\_\allowbreak{}string\_\allowbreak{}view}时，你必须注意两件事：不能修改视图所引用的底层数据；并且由于视图是非拥有引用，你必须自行管理数据的生存期。
\end{myNote}


%%C020523 h3

\subsection{延伸阅读}


%%C020524 li

\begin{itemize}
\item \textit{构建字符串工具库}

\end{itemize}
